% V4.2 新增操作化章节；V4.1正文保留为历史基线。

\part{V4.2 网页数据可用化与赛题逐条落地}

\chapter{V4.2 当前真实状态与权威顺序}

\begin{decisionbox}[当前状态替代规则]
本章替代前文 V4.1 时点状态，前文章节保留用于证明方案演进。当前权威顺序为：\texttt{00\_}当前项目交接\texttt{\_V4.2/README\_}先读我\texttt{.md}，V4.2 计划书，调研资料归档\texttt{\_V1.1/00\_}交接入口\texttt{/source\_action\_registry\_V1.1.jsonl}，V1.0 原始归档，V4.1 及更早历史文件。V1.1 是增量可用化层，不重写 V1.0 的原件、终态或哈希。
\end{decisionbox}

截至 2026 年 8 月 16 日，本轮已经完成 688 项来源的四态动作闭包。清单中共有 196 条扩展名为 HTML 的来源记录，其中 159 条存在本地 HTML 文件，37 条从未取得本地网页。本轮只对有合法本地原件的 159 条执行内容提取：154 条完成正文或表格提取，1 条因内容过少标记为部分提取，3 条为拦截页，1 条为空壳页面。\Field{WEB-0335} 与 \Field{WEB-0336} 不再作为可用论文页面，统一由已恢复的官方开放 PDF \Field{WEB-0115} 替代；\Field{WEB-0345} 是无正文的访问阻断页；\Field{LCL-0156} 是动态空壳。

补取阶段从 1,909 个页面链接中按“结构化数据优先、替代来源优先、每个父资产最多三个”收敛到 36 个公开 GET 请求；26 个实际取得并通过文件头检查，10 个明确失败。取得资源展开后形成 155 个 XLSX 工作表、8 个 XML 对象、6 个 CSV、7 个 ZIP 归档和 2 个文本对象；XLSX 共转出 124,178 行，CSV 原件共识别 1,090 行。另有 2 个旧式 XLS 原件已保存并校验，但为避免引入带高危漏洞的解析依赖，暂列人工转换。

\begin{longtable}{p{4.2cm}p{2.5cm}p{7.5cm}}
\caption{V4.2 数据可用化结果}\label{tab:v42-usability-result}\\
\toprule
\textbf{验收对象} & \textbf{实测结果} & \textbf{准确含义}\\
\midrule
\endfirsthead
\toprule \textbf{验收对象} & \textbf{实测结果} & \textbf{准确含义}\\ \midrule
\endhead
688 项来源 & 688 个唯一 \Field{asset_id} & 每项都有获取、内容、提取、正式用途和下一动作；不代表全部全文可合法下载。\\
本地 HTML & 159 项 & 90 项混合内容、49 项叙述内容、8 项结构化表格、7 项下载入口、1 项低内容、3 项拦截页、1 项空壳。\\
HTML 自动提取 & 154 完成、1 部分、4 失败 & 自动结果是人工复核候选；关键数字仍需页码、单位、期间和双人复核。\\
底层资源补取 & 26 成功、10 失败 & 失败包括 HTTP 403、请求结构化文件却返回 HTML、SEC 文件超过 50 MB 上限；均保留 URL、状态和失败原因。\\
结构化展开 & 180 条对象清单 & 155 个 XLSX 工作表、8 个 XML、7 个 ZIP、6 个 CSV、2 个文本和2个待人工转换 XLS。\\
正式项目数据 & 尚未建设 & 130 国、40 国、20 家企业、模型、Smartbi 写入、AIChat 与 XML 恢复仍为后续阶段。\\
\bottomrule
\end{longtable}

\section{四态接口}

以后任何来源不得只使用“已下载”一个状态。\Field{source_registry} 和动作清单必须同时保存：

\begin{longtable}{p{3.5cm}p{4.1cm}p{6.6cm}}
\toprule
\textbf{状态轴} & \textbf{允许值} & \textbf{判定规则}\\
\midrule
获取状态 & \Field{ACQUIRED / NOT_ACQUIRED / NOT_APPLICABLE} & 取得合法原件或底层响应才是 ACQUIRED；许可禁止、需认证或仅元数据时为 NOT\_APPLICABLE。\\
内容状态 & \Field{STRUCTURED_DATA / NARRATIVE_CONTENT / DOWNLOAD_LANDING / BLOCK_PAGE / EMPTY_SHELL} 等 & HTTP 200、非空HTML和文件存在均不能代替内容判断。\\
提取状态 & \Field{EXTRACTED / PARTIAL / FAILED / NOT_REQUIRED} & 表格、正文或数据对象必须实际落到提取路径；失败保留原因。\\
正式用途 & CORE、SUPPLEMENT、REPLACED、EXCLUDED & A级可用原件可进入核心候选；B/C级通常为补充；受限、拦截、重复和秘密来源不得直接进入正式数据。\\
\bottomrule
\end{longtable}

机器权威字段除上述四态外还包括：\Field{required_fields}、\Field{data_usage}、\Field{preferred_endpoint}、\Field{request_parameters}、\Field{extraction_method}、\Field{raw_local_path}、\Field{extracted_local_path}、\Field{target_table}、\Field{smartbi_page}、\Field{acceptance_rule}、\Field{replacement_source}、\Field{failure_reason}、\Field{owner} 和 \Field{next_action}。

\chapter{赛题评选标准原文与执行解释}

\begin{warningbox}[原文保留说明]
下列“原文”根据项目本地赛题文件 \texttt{XH-202612\_Smartbi AI}驱动的数据创新平台研究\texttt{.pdf} 第 6—8 页转录，不改变原意。每段原文后另列执行解释；执行解释不是对原文的替换。机器追踪表为 \texttt{00\_}当前项目交接\texttt{\_V4.2/SCORING\_TRACEABILITY\_V4.2.csv}。
\end{warningbox}

\section{一、问题价值与现实意义（20 分）}

\subsection*{S1.1 原文}
（1）选题是否紧扣社会、经济、民生或行业中的真实痛点，需引用至少 1 份权威资料（如国家统计局报告、行业白皮书、主流媒体报道等）证明问题存在，以及明确界定问题影响范围（如“涉及全国超 3000 万灵活就业人员”或“导致某省年医保支出增加 5 亿元”）；

\begin{decisionbox}[S1.1 怎么做]
从商务部公报提取境外企业覆盖、投资流量和存量，从国别事件表提取货币危机、通胀、资本限制和冻结事件，从企业年报提取海外收入与汇兑损失。问题规模页必须同时展示权威数字、来源ID、表号或页码、发布日期和当前数据完整度。任何“130国、40国、20家企业”在正式表未完成前只能写作计划范围。
\end{decisionbox}

\subsection*{S1.2 原文}
（2）是否具有明确的应用场景和潜在推广价值，明确指出目标用户或应用主体（如“适用于县级疾控中心”“可部署于连锁商超运营系统”），提出至少 1 条可行的落地路径（如试点合作机制、政策对接建议、SaaS 化方案）；

\begin{decisionbox}[S1.2 怎么做]
目标用户固定为企业总部财资部门、区域财资中心、境外子公司财务负责人和风险管理人员。试点路径固定为：一家公司或一个行业填写企业敞口模板，导入 Smartbi 小样，完成国别预警、周期判断、配置实验和管理层简报，再扩展到区域模板。验收不是“写出应用场景”，而是陌生成员在十分钟内完成一次指定风险决策任务。
\end{decisionbox}

\subsection*{S1.3 原文}
（3）是否体现数据驱动决策的必要性与可行性，对比说明传统经验决策与数据驱动方式的差异（如预测准确率提升$\geq 10\%$、响应效率提高$\geq 20\%$），同时数据来源合法、稳定且可获取（提供数据接口、公开链接或授权证明）。

\begin{decisionbox}[S1.3 怎么做]
建立30项固定任务：人工查找汇率、通胀、资本限制、企业披露和历史类比，与平台输出比较完成时长、数值正确率、来源追溯率和预警提前量。提升比例必须实测，不能预填。合法性由 \Field{source_registry} 的许可、原URL、获取方式、哈希和再分发范围证明；稳定性由发布时间、抓取时间、替代源和重跑记录证明。
\end{decisionbox}

\section{二、数据分析方法与技术应用（25 分）}

\subsection*{S2.1 原文}
（1）数据处理流程是否规范（包括数据获取、清洗、整合等），提供完整的数据处理流程说明（含 ETL 流程图或步骤清单）、清洗规则（如：异常值剔除比例$\leq 10\%$）、多源数据整合逻辑（如：主外键关联说明），创建分析主题的数据模型；

\begin{decisionbox}[S2.1 怎么做]
V1.1先完成来源层四态闭包；正式Python管线再按 raw、staging、curated、smartbi、qa 五层运行。国家以ISO3、月份以月末、企业以稳定企业ID、事件以 \Field{event_id} 关联。极端贬值和通胀不删除，只加异常标记；清洗剔除比例单列，超过10\%必须失败并解释。事实表之间不得直接多对多连接。
\end{decisionbox}

\subsection*{S2.2 原文}
（2）分析方法是否科学合理（如统计分析、归因分析、RAG知识增强等），提供每类分析方法的使用背景与结果说明；

\begin{decisionbox}[S2.2 怎么做]
国别层使用等权主模型、熵权和PCA挑战模型；事件层使用事件窗口；跨国层使用国家和时间固定效应及按国家聚类稳健误差；资产层使用周期状态、压力测试和美元—黄金对称剔除；证据问答使用带 \Field{asset_id}、\Field{source_locator} 和证据等级的RAG。所有结果限定为历史相关性和情景结果。
\end{decisionbox}

\subsection*{S2.3 原文}
（3）是否有效运用 Smartbi 平台的核心功能（如可视化仪表板、智能问数、智能报告、数据洞察分析等），提供每类功能实现的效果数量（如：通过可视化仪表板开发了 3 张决策分析看板）；

\begin{decisionbox}[S2.3 怎么做]
数量合同固定为六页决策看板、25个AIChat测试问题、1套管理层智能报告和数据洞察结果。每项功能保存资源名、截图、操作日志、输入表、指标口径、正确答案和验收状态；仅有页面设计不计作已实现。
\end{decisionbox}

\subsection*{S2.4 原文}
（4）技术实现是否体现创新性或对现有方法的优化，在性能优化$\geq 5\%$，数据获取由 T+1 提升至小时或分钟级。

\begin{decisionbox}[S2.4 怎么做]
把“数据发布频率”和“发布后入库延迟”分开。月度、季度和年度指标不得伪造成实时数据；支持增量获取的来源以官方发布到 raw 层落盘不超过60分钟为目标。性能比较使用同一机器、同一输入，分别运行全量和增量管线；增量总耗时至少减少5\%，并保存命令、起止时间、CPU/内存和行数。
\end{decisionbox}

\section{三、可视化呈现与交互体验（25 分）}

\subsection*{S3.1 原文}
（1）仪表板/报告是否逻辑清晰、重点突出（标题和结论加粗）、美观易读，使用对比色/放大字体等方式突出核心结论（如KPI、预警信号）；

\begin{decisionbox}[S3.1 怎么做]
六页统一使用深蓝主色、金色风险提示和红色禁止状态；每页只设置一个主结论区、三至六个KPI和一个“数据可信度/限制”区。视觉QA逐页检查标题、数字层级、色盲可读性、中文换行和来源脚注。
\end{decisionbox}

\subsection*{S3.2 原文}
（2）可视化形式是否契合数据特征与分析目标，如：类别数据用柱状图/条形图，趋势用折线图，占比用环形图/堆叠图，地理分布用热力地图；

\begin{decisionbox}[S3.2 怎么做]
国别分布用地图和四象限，时间变化用折线和周期状态带，行业和企业比较用条形图，资产权重用堆叠图，危机用事件时间轴，风险与投资关系用散点图。每张图必须在图表绑定表中登记来源表、字段、公式、筛选、下钻、预期结论和核对值。
\end{decisionbox}

\subsection*{S3.3 原文}
（3）是否支持良好的用户交互（如下钻、联动、筛选、动态更新等），交互响应时间$\leq 10$ 秒；

\begin{decisionbox}[S3.3 怎么做]
固定测试地区、国家、行业、企业、风险类型、周期、事件、资产、期限和计价口径筛选。使用冷缓存、热缓存和连续联动三组测试，记录至少30次操作；P95响应不超过10秒，且联动前后行数与Python基准一致。
\end{decisionbox}

\subsection*{S3.4 原文}
（4）是否通过 Smartbi 实现“讲好数据故事”的能力，包含完整叙事逻辑：问题背景→数据洞察→根因分析→行动建议，让听众可在 10 分钟内掌握核心结论。

\begin{decisionbox}[S3.4 怎么做]
六页顺序固定为：出海版图与痛点，国别危机，企业敞口，多币种与黄金实验，决策与约束，全球周期类比。由未参与制作的人观看一次后复述问题、洞察、原因、边界和行动；缺一项即重写故事线。视频按用户冻结口径制作3—5分钟，目标4分钟。
\end{decisionbox}

\section{四、创新性与原创性（20 分）}

\subsection*{S4.1 原文}
（1）在选题、方法、模型或应用场景上是否有独到见解，如选题在近 3 年主流数据分析竞赛或核心期刊中出现频率$\leq 2$次或将成熟方法创新应用于新领域（如用 NLP 情感分析优化政务服务满意度评估）；

\begin{decisionbox}[S4.1 怎么做]
创新性不得只写“黄金配置”。需分别证明全球周期—国别风险—企业披露三层联动、周期协同、制度可比性和Smartbi决策闭环相对已有研究新增了什么。每项创新登记检索范围、已有方法、项目新增接口和可核验产物。
\end{decisionbox}

\subsection*{S4.2 原文}
（2）是否提出新颖的分析视角或解决方案，提供可验证的假设并用数据验证；

\begin{decisionbox}[S4.2 怎么做]
固定检验H1—H5、0\%黄金反事实、2.5\%条件底仓、去黄金和去美元对称实验。报告必须保留至少三个黄金不占优、三个美元短债不占优和两个两者都不能完全覆盖损失的窗口，禁止只展示支持性案例。
\end{decisionbox}

\subsection*{S4.3 原文}
（3）作品是否为团队独立完成，无抄袭或不当引用，提供含团队成员基于 Smartbi 平台的操作日志或账号信息。

\begin{decisionbox}[S4.3 怎么做]
所有成员在Smartbi使用独立可追踪操作记录，代码、配置、图表和文字均进入提交清单。引用必须回到 \Field{asset_id} 和原始定位；账号通过私有交付单提供，密码、Cookie和令牌不得进入PDF、ZIP、截图或版本文件。
\end{decisionbox}

\section{五、完整性与可复现性（10 分）}

\subsection*{S5.1 原文}
（1）提交材料是否齐全（报告、基于 Smartbi 实现的资源、视频、示例数据及数据说明等）。缺失任一材料扣 2 分：分析报告、Smartbi 实现的资源（xml 文件或导出包）、演示视频（3–5 分钟，包含关键功能操作与核心结论展示）、示例数据语文档、数据字典/字段说明；

\begin{decisionbox}[S5.1 怎么做]
提交清单逐项登记分析报告、XML或导出包、4分钟视频、示例数据、数据说明、字段字典、来源说明、合规声明、操作指南和哈希。赛题其他位置存在“3分钟以内”表述，当前按用户冻结的3—5分钟口径执行，并在提交前再次核对主办方书面通知。
\end{decisionbox}

\subsection*{S5.2 原文}
（2）分析过程是否可追溯、结果是否可复现，基于 Smartbi 导入 Smartbi 实现的资源（xml 文件或导出包）可完整成功恢复资源；

\begin{decisionbox}[S5.2 怎么做]
在无项目缓存、无原资源的干净环境中导入XML，核对数据源、主题模型、六页看板、25问配置、行数、主外键和关键指标。Python QA与Smartbi关键结果误差不超过1\%；恢复全过程录屏并保存日志。
\end{decisionbox}

\subsection*{S5.3 原文}
（3）是否提供清晰的操作说明或部署指南，便于评审验证，提供 Smartbi 登录地址和账号信息，同时提供启动成功窗口界面和登录成功界面图。

\begin{decisionbox}[S5.3 怎么做]
操作指南从登录、导入、模型、筛选、下钻、AIChat到XML恢复逐步列出菜单、输入、预期画面和失败处理。共享材料只放登录地址和脱敏截图；账号密码通过私有渠道交付。由非原制作成员按指南独立完成一次验收。
\end{decisionbox}

\chapter{D0 第一天操作表与D2数据可用化步骤卡}

\section{D0 第一天逐时操作}

\begin{longtable}{p{2.0cm}p{5.3cm}p{6.7cm}}
\toprule
\textbf{时间} & \textbf{必须执行} & \textbf{必须留下的文件和门槛}\\
\midrule
09:00 & 逐条登记赛题原文，分配 S1.1—S5.3 & \File{SCORING_TRACEABILITY_V4.2.csv}；原文和解释分栏。\\
09:30 & 冻结问题、用户、应用场景和禁止误述 & \File{D0_scope_signoff.xlsx}；A、B、C签字。\\
10:30 & 冻结传统人工决策基线和30项对照任务 & \File{traditional_vs_data_benchmark.xlsx}；不得预填提升比例。\\
11:30 & 冻结11表、6维表、桥表、主键、频率和计价口径 & \File{data_contract_V4.2.xlsx}；年度数据禁止复制成月度观测。\\
13:30 & 导入688项动作清单，按哈希、规范URL和许可复用 & \File{archive_to_source_registry_crosswalk.xlsx}；只补缺口。\\
14:30 & 冻结指标—来源—字段合同 & 每个字段有主源、替代源、提取位置、单位和验收阈值。\\
15:30 & 冻结危机、周期、收益和组合公式 & \File{formula_contract_V4.2.xlsx}；美元与黄金对称实验。\\
16:30 & 映射六页、25问和17条评分项 & 每项评分可定位到数据、脚本、Smartbi资源和验收证据。\\
17:30 & G0签署与配置哈希 & 未签署不得进入130国和40国正式采集。\\
\bottomrule
\end{longtable}

\section{D2A：688项内容审计}

\StepTitle{D2A}{来源取得状态与内容可用性审计}{负责人：B；复核：D}
\CardRow{输入}{V1.0 研究资料总清单\texttt{.csv}、本地原件和许可矩阵。}
\CardRow{命令}{进入调研资料归档\texttt{\_V1.1/05\_QA}与运行记录\texttt{/scripts}，运行 \Field{npm ci}，再运行 \Field{node 02_audit_web_assets.mjs}。}
\CardRow{操作}{验证本地路径和哈希；解析标题、可见字符、表格数、下载链接、iframe和表单；同时检测拦截页、动态空壳和低内容页。}
\CardRow{输出}{\File{web_asset_audit_V1.1.csv}、审计摘要和两个HTML动作矩阵。}
\CardRow{验收}{688项清单主键唯一；159个本地HTML全部得到内容状态；拦截和空壳不得标为可用。}
\CardRow{失败处理}{读取错误保留错误信息；原件不修改；编码无法识别时保留字节并进入人工复核。}
\CardRow{下一接收人}{D2B使用候选链接和许可状态决定补取。}

\section{D2B：公开原件和底层数据补取}

\StepTitle{D2B}{合法开放底层资源补取}{负责人：B；许可复核：A/D}
\CardRow{输入}{页面下载链接、V1.0替代来源、许可和再分发范围。}
\CardRow{命令}{\Field{node 03_discover_official_resources.mjs}；复核 \File{selected_resource_fetch_plan_V1.1.csv} 后运行 \Field{node 04_fetch_registered_resources.mjs}。}
\CardRow{操作}{替代源优先；CSV、XLSX、JSON、XML、ZIP优先；每父资产最多3个；同域限速；30秒超时；单文件上限50MB；只使用无认证公开GET。}
\CardRow{输出}{新增原件、替代原件、请求时间、最终URL、安全响应头、MIME、大小、SHA-256和失败原因。}
\CardRow{验收}{文件扩展名、MIME和文件头一致；请求结构化资源却返回HTML必须失败；不保存Cookie、令牌或认证头。}
\CardRow{失败处理}{403、登录、付费、验证码、超限或条款限制全部转为人工任务或权威替代；不得绕过。}
\CardRow{下一接收人}{D2C展开和提取实际内容。}

\section{D2C：HTML、PDF和表格内容提取}

\StepTitle{D2C}{正文、表格和结构化对象提取}{负责人：B；关键数字复核：D}
\CardRow{命令}{\Field{node 05_extract_html_tables.mjs}、\Field{node 06_extract_narrative_content.mjs}、\Field{node 10_inspect_structured_resources.mjs}。}
\CardRow{HTML表格}{每张表输出独立CSV；记录表序号、行号、字段、原始值、输出路径。}
\CardRow{HTML正文}{Readability提取主文，输出JSON和Markdown；每段保存 \Field{html:paragraph[n]} 定位。}
\CardRow{ZIP/XLSX}{ZIP先检查路径并展开；XLSX按工作表转CSV并保留Excel行号；XML、CSV和JSON生成结构清单。}
\CardRow{PDF}{正式数据工程使用 \Field{pdftotext -layout}；无文本层则生成OCR候选；关键数字必须双人复核页码。}
\CardRow{验收}{提取结果包含 \Field{source_id}、\Field{source_locator}、原值、解析值、单位、期间、地区、发布时间、方法、置信度和复核状态。}
\CardRow{失败处理}{旧XLS不使用带高危漏洞的库，保存原件后采用受控转换；无法读取时不得猜测。}
\CardRow{下一接收人}{D2D建立正式来源表映射。}

\section{D2D：来源到业务表、看板和评分的映射}

\StepTitle{D2D}{来源用途闭包}{负责人：A/B；模型复核：C}
\CardRow{命令}{\Field{node 08_build_source_action_registry.mjs} 和 \Field{node 11_build_v42_plan_artifacts.mjs}。}
\CardRow{输出}{688项CSV/JSONL动作清单、159项HTML矩阵、评分追踪矩阵、受限待办、PDF全量附录。}
\CardRow{验收}{每个 \Field{asset_id} 有唯一四态、目标字段、目标表、Smartbi页面、验收、替代和下一动作；总数严格为688。}
\CardRow{失败处理}{目标字段缺少合法可用来源时，先启用登记替代源；仍失败则删除字段承诺和相关评分目标，不用零值填补。}
\CardRow{下一接收人}{D1/D3—D6只从正式 \Field{source_registry} 读取已批准来源。}

\chapter{按来源类别提取哪些数据、如何处理、输出到哪里}

\clearpage
\begin{landscape}
\begingroup
\setlength{\tabcolsep}{2pt}
\scriptsize
\begin{longtable}{p{3.2cm}p{5.2cm}p{5.4cm}p{4.5cm}p{4.8cm}}
\toprule
\textbf{来源} & \textbf{必须提取字段} & \textbf{具体操作} & \textbf{进入何处及用途} & \textbf{验收与失败处理}\\
\midrule
\endfirsthead
\toprule \textbf{来源} & \textbf{必须提取字段} & \textbf{具体操作} & \textbf{进入何处及用途} & \textbf{验收与失败处理}\\ \midrule
\endhead
商务部公报 & ISO3、国家、年份、投资流量、存量、企业存在、表号、页码、发布日期 & 下载PDF/XLSX；PDF按页转文本；定位“分国家/地区”附表；国家名映射ISO3；缺失保持空值 & \Field{country_exposure}；DB01问题规模、130国版图和40国筛选 & 唯一ISO3；数字须有表号页码；通道经济体单列；\Field{WEB-0344}现为入口资料，正式附表仍需复核。\\
世界银行WDI/GEM & 汇率、CPI、储备、进口、贸易、GDP、频率、版本、发布日期 & WDI使用官方API并保存查询URL与JSON；GEM使用已取得批量包 \Field{WEB-0048}；解压后按序列元数据和国家代码整理 & \Field{country_monthly_risk}、\Field{global_cycle_month}；DB01/DB02/DB06 & 汇率和CPI月度完整率分别至少80\%；WDI年度值不得冒充月度；失败用央行/IMF替代。\\
IMF WEO/IFS/AREAER & 年度宏观、月度金融序列、汇率制度、兑换与汇出限制、原文编码 & WEO工作簿已补取并按工作表转CSV；IFS使用登记API；AREAER保留年度频率和原文定位 & \Field{country_policy_year}、\Field{country_event}、\Field{global_cycle_month} & 年度数据用有效期关联，不复制12次；登录或条款限制进入人工任务。\\
BIS、央行和财政部门 & 政策利率、期限利差、信用、房地产、储备、原始频率与发布日期 & 优先SDMX/CSV/XLSX；季度房地产记录真实观察期和可用期；已补取表格按工作表转CSV & \Field{global_cycle_month}；DB06周期识别 & 核心序列完整率至少90\%；季度值不得伪造月度观测；无覆盖时为空。\\
美联储与美国财政部 & 工业生产、H.10汇率、3月短债、实际利率、发布时间 & 已取得美联储G.17与H.10 XML ZIP并展开；财政部CSV请求本轮403，保留失败，后续使用公开XML接口或人工下载 & \Field{global_cycle_month}、\Field{asset_monthly_return}；DB04/DB06 & XML根节点、日期、数值和单位可解析；财政部未取得前不得声称短债正式序列完成。\\
上海黄金交易所 & 日期、合约、开高低收、加权均价、成交量、金额、持仓、币种 & 159项HTML中的SGE静态表格已转CSV；后续登记公开查询参数，按日去重并聚合月度 & \Field{asset_monthly_return}；DB04 & 日期与合约复合主键唯一；价格单位统一；无可靠接口时使用人工导出而非页面猜取。\\
企业年报与交易所 & 企业、财年、海外收入、地区、汇兑损益、外币资产负债、套保、现金、授信、页码 & PDF文本+页码候选；表格和XBRL优先；人工双录关键数字；仅地区披露不得推断国家 & \Field{company_overseas_exposure}；DB03 & 20家基础画像；量化样本需同时满足收入与外汇风险披露；SEC超50MB文件使用官方分件或人工下载。\\
案例、法规和历史资料 & 事件、机制、起止期、制度环境、政策响应、企业损失、证据等级、页码 & 45案例链接导入；按 \Field{event_cluster_id} 去重；17宏观案例字段化；A级或双B级复核 & \Field{case_evidence}、\Field{historical_crisis_event}；DB05/DB06 & 不把45文档当45独立事件；广场协议为政策节点；关键数字不得只引机构主页。\\
Smartbi官方资料和实例 & 版本、许可、菜单、导入、ETL、模型、AIChat、XML、限制、操作证据 & 只读盘点复用；写入、导入和恢复另经P0授权；截图脱敏；不记录凭据 & 平台操作手册、评分证据和D10—D12 & 能力盘点不等于功能已实现；六页、25问和XML须分别验收。\\
\bottomrule
\end{longtable}
\endgroup
\end{landscape}

\chapter{V1.1到正式Python数据工程的交接接口}

V1.1使用当前可运行且依赖审计为零漏洞的Node.js 24工具链。正式数据工程继续使用Python；双方只通过原始文件、CSV、JSONL和哈希衔接。Node侧不得直接写入11张正式业务表，Python侧不得重新盲目抓取已取得原件。

\begin{verbatim}
cd 调研资料归档_V1.1/05_QA与运行记录/scripts
npm ci
node run_archive_v11.mjs

python scripts/run_pipeline.py \
  --as-of 2025-12-31 \
  --run-id 20260831_v42_formal
\end{verbatim}

正式 \Field{source_registry} 新增并强制保留：\Field{archive_asset_id}、\Field{archive_run_id}、\Field{reuse_status}、\Field{acquisition_state}、\Field{content_state}、\Field{extraction_state}、\Field{formal_use_state}、\Field{source_locator}、\Field{redistribution_scope} 和 \Field{revalidation_required}。V1.1自动提取段落与表格在人工复核前只能进入 staging，不能直接进入 curated 或Smartbi。

\chapter{赛题得分证据、Smartbi产品和交接验收}

\section{评分证据闭环}

17条评分子项全部进入 \File{SCORING_TRACEABILITY_V4.2.csv}。每条必须同时存在原文、项目响应、所需数据、来源ID、输出表或交付物、处理脚本、Smartbi功能、页面、量化验收、负责人和当前状态。评分状态只有 \Field{COMPLETE_VERIFIED}、\Field{PARTIAL_REUSABLE}、\Field{NOT_STARTED} 和 \Field{EXTERNAL_ACTION}；“写入计划书”不得提升实现分数。

\section{六页和25问保持不变}

六页仍为亚非拉风险全景、国别与危机下钻、中国企业海外风险、多币种—黄金周期实验室、决策约束与证据中心、全球周期与历史危机。25问必须覆盖事实查询、资产比较、周期识别、历史类比、来源追溯和约束路由。企业仅披露地区时，AIChat不得回答到具体国家；历史制度问题必须展示可比性限制。

\section{V4.2交接验收}

\begin{enumerate}
  \item 688项动作清单行数和唯一ID均为688；本地HTML矩阵严格为159项；
  \item 26个新取得资源逐一校验MIME、文件头、SHA-256和结构；10个失败请求有原因和影响；
  \item 155个XLSX工作表、8个XML、6个CSV和7个ZIP均能从资源ID定位；
  \item 旧V4.1、V1.0及旧流程图前后哈希相同；
  \item V4.2 PDF包含赛题原文、操作卡、688项总表和159项HTML矩阵；
  \item 交接包不含密码、Cookie、令牌、认证头、Office锁文件或 \File{node_modules}；
  \item 在新目录解压后重新计算清单哈希，并以相对路径完成不少于40项离线定位测试；
  \item PDF逐页渲染检查，无乱码、裁切、重叠、异常空白页，最后一页显示V4.2。
\end{enumerate}

% 本文件由11_build_v42_plan_artifacts.mjs生成，请勿手工维护。
\part{V4.2 全量来源动作与提取附录}
\chapter{688 项来源动作总表}
机器权威清单位于调研资料归档 V1.1 的交接入口，文件为 \Field{source_action_registry_V1.1.csv} 和 \Field{source_action_registry_V1.1.jsonl}。本表保证每个 \Field{asset_id} 至少出现一次；详细URL、请求参数、输出路径和失败原因以机器清单为准。
\clearpage
\begin{landscape}\begingroup\setlength{\tabcolsep}{2pt}\scriptsize
\begin{longtable}{p{1.8cm}p{2.6cm}p{2.6cm}p{2.6cm}p{4.8cm}p{7.2cm}}
\toprule
ID & 证据/类别 & 获取与内容 & 提取/用途 & 目标表 & 下一动作\\
\midrule\endfirsthead
\toprule ID & 证据/类别 & 获取与内容 & 提取/用途 & 目标表 & 下一动作\\\midrule\endhead
LCL-0001 & 内部/\allowbreak{}01\_\allowbreak{}比赛与项目基线 & NOT\_\allowbreak{}APPLICABLE/\allowbreak{}NOT\_\allowbreak{}ACQUIRED & FAILED/\allowbreak{}EXCLUDED & source\_\allowbreak{}registry | case\_\allowbreak{}evidence & 仅保留元数据；正式分析改用许可允许的权威替代源。\\
LCL-0002 & A/\allowbreak{}01\_\allowbreak{}比赛与项目基线 & ACQUIRED/\allowbreak{}DOCUMENT\_\allowbreak{}LOCAL & NOT\_\allowbreak{}REQUIRED/\allowbreak{}CORE & source\_\allowbreak{}registry | case\_\allowbreak{}evidence & 人工复核关键字段与source\_\allowbreak{}locator；通过后映射至正式source\_\allowbreak{}registry。\\
LCL-0003 & 内部/\allowbreak{}03\_\allowbreak{}案例库与证据 & ACQUIRED/\allowbreak{}LOCAL\_\allowbreak{}ASSET & NOT\_\allowbreak{}REQUIRED/\allowbreak{}SUPPLEMENT & company\_\allowbreak{}overseas\_\allowbreak{}exposure | country\_\allowbreak{}monthly\_\allowbreak{}risk | country\_\allowbreak{}policy\_\allowbreak{}year | c… & 按项目用途复核；不足以单独支持核心结论。\\
LCL-0004 & 内部/\allowbreak{}03\_\allowbreak{}案例库与证据 & ACQUIRED/\allowbreak{}LOCAL\_\allowbreak{}ASSET & NOT\_\allowbreak{}REQUIRED/\allowbreak{}SUPPLEMENT & company\_\allowbreak{}overseas\_\allowbreak{}exposure | country\_\allowbreak{}monthly\_\allowbreak{}risk | country\_\allowbreak{}policy\_\allowbreak{}year | c… & 按项目用途复核；不足以单独支持核心结论。\\
LCL-0005 & 内部/\allowbreak{}03\_\allowbreak{}案例库与证据 & ACQUIRED/\allowbreak{}LOCAL\_\allowbreak{}ASSET & NOT\_\allowbreak{}REQUIRED/\allowbreak{}SUPPLEMENT & company\_\allowbreak{}overseas\_\allowbreak{}exposure | country\_\allowbreak{}monthly\_\allowbreak{}risk | country\_\allowbreak{}policy\_\allowbreak{}year | c… & 按项目用途复核；不足以单独支持核心结论。\\
LCL-0006 & 内部/\allowbreak{}03\_\allowbreak{}案例库与证据 & ACQUIRED/\allowbreak{}LOCAL\_\allowbreak{}ASSET & NOT\_\allowbreak{}REQUIRED/\allowbreak{}SUPPLEMENT & company\_\allowbreak{}overseas\_\allowbreak{}exposure | country\_\allowbreak{}monthly\_\allowbreak{}risk | country\_\allowbreak{}policy\_\allowbreak{}year | c… & 按项目用途复核；不足以单独支持核心结论。\\
LCL-0007 & 内部/\allowbreak{}03\_\allowbreak{}案例库与证据 & ACQUIRED/\allowbreak{}LOCAL\_\allowbreak{}ASSET & NOT\_\allowbreak{}REQUIRED/\allowbreak{}SUPPLEMENT & company\_\allowbreak{}overseas\_\allowbreak{}exposure | country\_\allowbreak{}monthly\_\allowbreak{}risk | country\_\allowbreak{}policy\_\allowbreak{}year | c… & 按项目用途复核；不足以单独支持核心结论。\\
LCL-0008 & 内部/\allowbreak{}03\_\allowbreak{}案例库与证据 & ACQUIRED/\allowbreak{}LOCAL\_\allowbreak{}ASSET & NOT\_\allowbreak{}REQUIRED/\allowbreak{}SUPPLEMENT & company\_\allowbreak{}overseas\_\allowbreak{}exposure | country\_\allowbreak{}monthly\_\allowbreak{}risk | country\_\allowbreak{}policy\_\allowbreak{}year | c… & 按项目用途复核；不足以单独支持核心结论。\\
LCL-0009 & 内部/\allowbreak{}03\_\allowbreak{}案例库与证据 & ACQUIRED/\allowbreak{}LOCAL\_\allowbreak{}ASSET & NOT\_\allowbreak{}REQUIRED/\allowbreak{}SUPPLEMENT & company\_\allowbreak{}overseas\_\allowbreak{}exposure | country\_\allowbreak{}monthly\_\allowbreak{}risk | country\_\allowbreak{}policy\_\allowbreak{}year | c… & 按项目用途复核；不足以单独支持核心结论。\\
LCL-0010 & 内部/\allowbreak{}03\_\allowbreak{}案例库与证据 & ACQUIRED/\allowbreak{}LOCAL\_\allowbreak{}ASSET & NOT\_\allowbreak{}REQUIRED/\allowbreak{}SUPPLEMENT & company\_\allowbreak{}overseas\_\allowbreak{}exposure | country\_\allowbreak{}monthly\_\allowbreak{}risk | country\_\allowbreak{}policy\_\allowbreak{}year | c… & 按项目用途复核；不足以单独支持核心结论。\\
LCL-0011 & 内部/\allowbreak{}03\_\allowbreak{}案例库与证据 & ACQUIRED/\allowbreak{}LOCAL\_\allowbreak{}ASSET & NOT\_\allowbreak{}REQUIRED/\allowbreak{}SUPPLEMENT & company\_\allowbreak{}overseas\_\allowbreak{}exposure | country\_\allowbreak{}monthly\_\allowbreak{}risk | country\_\allowbreak{}policy\_\allowbreak{}year | c… & 按项目用途复核；不足以单独支持核心结论。\\
LCL-0012 & 内部/\allowbreak{}03\_\allowbreak{}案例库与证据 & ACQUIRED/\allowbreak{}LOCAL\_\allowbreak{}ASSET & NOT\_\allowbreak{}REQUIRED/\allowbreak{}SUPPLEMENT & company\_\allowbreak{}overseas\_\allowbreak{}exposure | country\_\allowbreak{}monthly\_\allowbreak{}risk | country\_\allowbreak{}policy\_\allowbreak{}year | c… & 按项目用途复核；不足以单独支持核心结论。\\
LCL-0013 & 内部/\allowbreak{}03\_\allowbreak{}案例库与证据 & ACQUIRED/\allowbreak{}LOCAL\_\allowbreak{}ASSET & NOT\_\allowbreak{}REQUIRED/\allowbreak{}SUPPLEMENT & company\_\allowbreak{}overseas\_\allowbreak{}exposure | country\_\allowbreak{}monthly\_\allowbreak{}risk | country\_\allowbreak{}policy\_\allowbreak{}year | c… & 按项目用途复核；不足以单独支持核心结论。\\
LCL-0014 & 内部/\allowbreak{}03\_\allowbreak{}案例库与证据 & ACQUIRED/\allowbreak{}LOCAL\_\allowbreak{}ASSET & NOT\_\allowbreak{}REQUIRED/\allowbreak{}SUPPLEMENT & company\_\allowbreak{}overseas\_\allowbreak{}exposure | country\_\allowbreak{}monthly\_\allowbreak{}risk | country\_\allowbreak{}policy\_\allowbreak{}year | c… & 按项目用途复核；不足以单独支持核心结论。\\
LCL-0015 & 内部/\allowbreak{}03\_\allowbreak{}案例库与证据 & ACQUIRED/\allowbreak{}LOCAL\_\allowbreak{}ASSET & NOT\_\allowbreak{}REQUIRED/\allowbreak{}SUPPLEMENT & company\_\allowbreak{}overseas\_\allowbreak{}exposure | country\_\allowbreak{}monthly\_\allowbreak{}risk | country\_\allowbreak{}policy\_\allowbreak{}year | c… & 按项目用途复核；不足以单独支持核心结论。\\
LCL-0016 & 内部/\allowbreak{}03\_\allowbreak{}案例库与证据 & ACQUIRED/\allowbreak{}LOCAL\_\allowbreak{}ASSET & NOT\_\allowbreak{}REQUIRED/\allowbreak{}SUPPLEMENT & company\_\allowbreak{}overseas\_\allowbreak{}exposure | country\_\allowbreak{}monthly\_\allowbreak{}risk | country\_\allowbreak{}policy\_\allowbreak{}year | c… & 按项目用途复核；不足以单独支持核心结论。\\
LCL-0017 & 内部/\allowbreak{}03\_\allowbreak{}案例库与证据 & ACQUIRED/\allowbreak{}LOCAL\_\allowbreak{}ASSET & NOT\_\allowbreak{}REQUIRED/\allowbreak{}SUPPLEMENT & company\_\allowbreak{}overseas\_\allowbreak{}exposure | country\_\allowbreak{}monthly\_\allowbreak{}risk | country\_\allowbreak{}policy\_\allowbreak{}year | c… & 按项目用途复核；不足以单独支持核心结论。\\
LCL-0018 & 内部/\allowbreak{}03\_\allowbreak{}案例库与证据 & ACQUIRED/\allowbreak{}LOCAL\_\allowbreak{}ASSET & NOT\_\allowbreak{}REQUIRED/\allowbreak{}SUPPLEMENT & company\_\allowbreak{}overseas\_\allowbreak{}exposure | country\_\allowbreak{}monthly\_\allowbreak{}risk | country\_\allowbreak{}policy\_\allowbreak{}year | c… & 按项目用途复核；不足以单独支持核心结论。\\
LCL-0019 & 内部/\allowbreak{}03\_\allowbreak{}案例库与证据 & ACQUIRED/\allowbreak{}LOCAL\_\allowbreak{}ASSET & NOT\_\allowbreak{}REQUIRED/\allowbreak{}SUPPLEMENT & company\_\allowbreak{}overseas\_\allowbreak{}exposure | country\_\allowbreak{}monthly\_\allowbreak{}risk | country\_\allowbreak{}policy\_\allowbreak{}year | c… & 按项目用途复核；不足以单独支持核心结论。\\
LCL-0020 & 内部/\allowbreak{}03\_\allowbreak{}案例库与证据 & ACQUIRED/\allowbreak{}LOCAL\_\allowbreak{}ASSET & NOT\_\allowbreak{}REQUIRED/\allowbreak{}SUPPLEMENT & company\_\allowbreak{}overseas\_\allowbreak{}exposure | country\_\allowbreak{}monthly\_\allowbreak{}risk | country\_\allowbreak{}policy\_\allowbreak{}year | c… & 按项目用途复核；不足以单独支持核心结论。\\
LCL-0021 & 内部/\allowbreak{}03\_\allowbreak{}案例库与证据 & ACQUIRED/\allowbreak{}LOCAL\_\allowbreak{}ASSET & NOT\_\allowbreak{}REQUIRED/\allowbreak{}SUPPLEMENT & company\_\allowbreak{}overseas\_\allowbreak{}exposure | country\_\allowbreak{}monthly\_\allowbreak{}risk | country\_\allowbreak{}policy\_\allowbreak{}year | c… & 按项目用途复核；不足以单独支持核心结论。\\
LCL-0022 & 内部/\allowbreak{}03\_\allowbreak{}案例库与证据 & ACQUIRED/\allowbreak{}LOCAL\_\allowbreak{}ASSET & NOT\_\allowbreak{}REQUIRED/\allowbreak{}SUPPLEMENT & company\_\allowbreak{}overseas\_\allowbreak{}exposure | country\_\allowbreak{}monthly\_\allowbreak{}risk | country\_\allowbreak{}policy\_\allowbreak{}year | c… & 按项目用途复核；不足以单独支持核心结论。\\
LCL-0023 & 内部/\allowbreak{}03\_\allowbreak{}案例库与证据 & ACQUIRED/\allowbreak{}LOCAL\_\allowbreak{}ASSET & NOT\_\allowbreak{}REQUIRED/\allowbreak{}SUPPLEMENT & company\_\allowbreak{}overseas\_\allowbreak{}exposure | country\_\allowbreak{}monthly\_\allowbreak{}risk | country\_\allowbreak{}policy\_\allowbreak{}year | c… & 按项目用途复核；不足以单独支持核心结论。\\
LCL-0024 & 内部/\allowbreak{}03\_\allowbreak{}案例库与证据 & ACQUIRED/\allowbreak{}LOCAL\_\allowbreak{}ASSET & NOT\_\allowbreak{}REQUIRED/\allowbreak{}SUPPLEMENT & company\_\allowbreak{}overseas\_\allowbreak{}exposure | country\_\allowbreak{}monthly\_\allowbreak{}risk | country\_\allowbreak{}policy\_\allowbreak{}year | c… & 按项目用途复核；不足以单独支持核心结论。\\
LCL-0025 & 内部/\allowbreak{}03\_\allowbreak{}案例库与证据 & ACQUIRED/\allowbreak{}LOCAL\_\allowbreak{}ASSET & NOT\_\allowbreak{}REQUIRED/\allowbreak{}SUPPLEMENT & company\_\allowbreak{}overseas\_\allowbreak{}exposure | country\_\allowbreak{}monthly\_\allowbreak{}risk | country\_\allowbreak{}policy\_\allowbreak{}year | c… & 按项目用途复核；不足以单独支持核心结论。\\
LCL-0026 & 内部/\allowbreak{}03\_\allowbreak{}案例库与证据 & ACQUIRED/\allowbreak{}LOCAL\_\allowbreak{}ASSET & NOT\_\allowbreak{}REQUIRED/\allowbreak{}SUPPLEMENT & company\_\allowbreak{}overseas\_\allowbreak{}exposure | country\_\allowbreak{}monthly\_\allowbreak{}risk | country\_\allowbreak{}policy\_\allowbreak{}year | c… & 按项目用途复核；不足以单独支持核心结论。\\
LCL-0027 & 内部/\allowbreak{}03\_\allowbreak{}案例库与证据 & ACQUIRED/\allowbreak{}LOCAL\_\allowbreak{}ASSET & NOT\_\allowbreak{}REQUIRED/\allowbreak{}SUPPLEMENT & company\_\allowbreak{}overseas\_\allowbreak{}exposure | country\_\allowbreak{}monthly\_\allowbreak{}risk | country\_\allowbreak{}policy\_\allowbreak{}year | c… & 按项目用途复核；不足以单独支持核心结论。\\
LCL-0028 & 内部/\allowbreak{}03\_\allowbreak{}案例库与证据 & ACQUIRED/\allowbreak{}LOCAL\_\allowbreak{}ASSET & NOT\_\allowbreak{}REQUIRED/\allowbreak{}SUPPLEMENT & company\_\allowbreak{}overseas\_\allowbreak{}exposure | country\_\allowbreak{}monthly\_\allowbreak{}risk | country\_\allowbreak{}policy\_\allowbreak{}year | c… & 按项目用途复核；不足以单独支持核心结论。\\
LCL-0029 & 内部/\allowbreak{}03\_\allowbreak{}案例库与证据 & ACQUIRED/\allowbreak{}LOCAL\_\allowbreak{}ASSET & NOT\_\allowbreak{}REQUIRED/\allowbreak{}SUPPLEMENT & company\_\allowbreak{}overseas\_\allowbreak{}exposure | country\_\allowbreak{}monthly\_\allowbreak{}risk | country\_\allowbreak{}policy\_\allowbreak{}year | c… & 按项目用途复核；不足以单独支持核心结论。\\
LCL-0030 & 内部/\allowbreak{}03\_\allowbreak{}案例库与证据 & ACQUIRED/\allowbreak{}LOCAL\_\allowbreak{}ASSET & NOT\_\allowbreak{}REQUIRED/\allowbreak{}SUPPLEMENT & company\_\allowbreak{}overseas\_\allowbreak{}exposure | country\_\allowbreak{}monthly\_\allowbreak{}risk | country\_\allowbreak{}policy\_\allowbreak{}year | c… & 按项目用途复核；不足以单独支持核心结论。\\
LCL-0031 & 内部/\allowbreak{}03\_\allowbreak{}案例库与证据 & ACQUIRED/\allowbreak{}LOCAL\_\allowbreak{}ASSET & NOT\_\allowbreak{}REQUIRED/\allowbreak{}SUPPLEMENT & company\_\allowbreak{}overseas\_\allowbreak{}exposure | country\_\allowbreak{}monthly\_\allowbreak{}risk | country\_\allowbreak{}policy\_\allowbreak{}year | c… & 按项目用途复核；不足以单独支持核心结论。\\
LCL-0032 & 内部/\allowbreak{}03\_\allowbreak{}案例库与证据 & ACQUIRED/\allowbreak{}LOCAL\_\allowbreak{}ASSET & NOT\_\allowbreak{}REQUIRED/\allowbreak{}SUPPLEMENT & company\_\allowbreak{}overseas\_\allowbreak{}exposure | country\_\allowbreak{}monthly\_\allowbreak{}risk | country\_\allowbreak{}policy\_\allowbreak{}year | c… & 按项目用途复核；不足以单独支持核心结论。\\
LCL-0033 & 内部/\allowbreak{}03\_\allowbreak{}案例库与证据 & ACQUIRED/\allowbreak{}LOCAL\_\allowbreak{}ASSET & NOT\_\allowbreak{}REQUIRED/\allowbreak{}SUPPLEMENT & company\_\allowbreak{}overseas\_\allowbreak{}exposure | country\_\allowbreak{}monthly\_\allowbreak{}risk | country\_\allowbreak{}policy\_\allowbreak{}year | c… & 按项目用途复核；不足以单独支持核心结论。\\
LCL-0034 & 内部/\allowbreak{}03\_\allowbreak{}案例库与证据 & ACQUIRED/\allowbreak{}LOCAL\_\allowbreak{}ASSET & NOT\_\allowbreak{}REQUIRED/\allowbreak{}SUPPLEMENT & company\_\allowbreak{}overseas\_\allowbreak{}exposure | country\_\allowbreak{}monthly\_\allowbreak{}risk | country\_\allowbreak{}policy\_\allowbreak{}year | c… & 按项目用途复核；不足以单独支持核心结论。\\
LCL-0035 & 内部/\allowbreak{}03\_\allowbreak{}案例库与证据 & ACQUIRED/\allowbreak{}LOCAL\_\allowbreak{}ASSET & NOT\_\allowbreak{}REQUIRED/\allowbreak{}SUPPLEMENT & company\_\allowbreak{}overseas\_\allowbreak{}exposure | country\_\allowbreak{}monthly\_\allowbreak{}risk | country\_\allowbreak{}policy\_\allowbreak{}year | c… & 按项目用途复核；不足以单独支持核心结论。\\
LCL-0036 & 内部/\allowbreak{}03\_\allowbreak{}案例库与证据 & ACQUIRED/\allowbreak{}LOCAL\_\allowbreak{}ASSET & NOT\_\allowbreak{}REQUIRED/\allowbreak{}SUPPLEMENT & company\_\allowbreak{}overseas\_\allowbreak{}exposure | country\_\allowbreak{}monthly\_\allowbreak{}risk | country\_\allowbreak{}policy\_\allowbreak{}year | c… & 按项目用途复核；不足以单独支持核心结论。\\
LCL-0037 & 内部/\allowbreak{}03\_\allowbreak{}案例库与证据 & ACQUIRED/\allowbreak{}LOCAL\_\allowbreak{}ASSET & NOT\_\allowbreak{}REQUIRED/\allowbreak{}SUPPLEMENT & company\_\allowbreak{}overseas\_\allowbreak{}exposure | country\_\allowbreak{}monthly\_\allowbreak{}risk | country\_\allowbreak{}policy\_\allowbreak{}year | c… & 按项目用途复核；不足以单独支持核心结论。\\
LCL-0038 & 内部/\allowbreak{}03\_\allowbreak{}案例库与证据 & ACQUIRED/\allowbreak{}LOCAL\_\allowbreak{}ASSET & NOT\_\allowbreak{}REQUIRED/\allowbreak{}SUPPLEMENT & company\_\allowbreak{}overseas\_\allowbreak{}exposure | country\_\allowbreak{}monthly\_\allowbreak{}risk | country\_\allowbreak{}policy\_\allowbreak{}year | c… & 按项目用途复核；不足以单独支持核心结论。\\
LCL-0039 & 内部/\allowbreak{}03\_\allowbreak{}案例库与证据 & ACQUIRED/\allowbreak{}LOCAL\_\allowbreak{}ASSET & NOT\_\allowbreak{}REQUIRED/\allowbreak{}SUPPLEMENT & company\_\allowbreak{}overseas\_\allowbreak{}exposure | country\_\allowbreak{}monthly\_\allowbreak{}risk | country\_\allowbreak{}policy\_\allowbreak{}year | c… & 按项目用途复核；不足以单独支持核心结论。\\
LCL-0040 & 内部/\allowbreak{}03\_\allowbreak{}案例库与证据 & ACQUIRED/\allowbreak{}LOCAL\_\allowbreak{}ASSET & NOT\_\allowbreak{}REQUIRED/\allowbreak{}SUPPLEMENT & company\_\allowbreak{}overseas\_\allowbreak{}exposure | country\_\allowbreak{}monthly\_\allowbreak{}risk | country\_\allowbreak{}policy\_\allowbreak{}year | c… & 按项目用途复核；不足以单独支持核心结论。\\
LCL-0041 & 内部/\allowbreak{}03\_\allowbreak{}案例库与证据 & ACQUIRED/\allowbreak{}LOCAL\_\allowbreak{}ASSET & NOT\_\allowbreak{}REQUIRED/\allowbreak{}SUPPLEMENT & company\_\allowbreak{}overseas\_\allowbreak{}exposure | country\_\allowbreak{}monthly\_\allowbreak{}risk | country\_\allowbreak{}policy\_\allowbreak{}year | c… & 按项目用途复核；不足以单独支持核心结论。\\
LCL-0042 & 内部/\allowbreak{}03\_\allowbreak{}案例库与证据 & ACQUIRED/\allowbreak{}LOCAL\_\allowbreak{}ASSET & NOT\_\allowbreak{}REQUIRED/\allowbreak{}SUPPLEMENT & company\_\allowbreak{}overseas\_\allowbreak{}exposure | country\_\allowbreak{}monthly\_\allowbreak{}risk | country\_\allowbreak{}policy\_\allowbreak{}year | c… & 按项目用途复核；不足以单独支持核心结论。\\
LCL-0043 & 内部/\allowbreak{}03\_\allowbreak{}案例库与证据 & ACQUIRED/\allowbreak{}LOCAL\_\allowbreak{}ASSET & NOT\_\allowbreak{}REQUIRED/\allowbreak{}SUPPLEMENT & company\_\allowbreak{}overseas\_\allowbreak{}exposure | country\_\allowbreak{}monthly\_\allowbreak{}risk | country\_\allowbreak{}policy\_\allowbreak{}year | c… & 按项目用途复核；不足以单独支持核心结论。\\
LCL-0044 & 内部/\allowbreak{}03\_\allowbreak{}案例库与证据 & ACQUIRED/\allowbreak{}LOCAL\_\allowbreak{}ASSET & NOT\_\allowbreak{}REQUIRED/\allowbreak{}SUPPLEMENT & company\_\allowbreak{}overseas\_\allowbreak{}exposure | country\_\allowbreak{}monthly\_\allowbreak{}risk | country\_\allowbreak{}policy\_\allowbreak{}year | c… & 按项目用途复核；不足以单独支持核心结论。\\
LCL-0045 & 内部/\allowbreak{}03\_\allowbreak{}案例库与证据 & ACQUIRED/\allowbreak{}LOCAL\_\allowbreak{}ASSET & NOT\_\allowbreak{}REQUIRED/\allowbreak{}SUPPLEMENT & company\_\allowbreak{}overseas\_\allowbreak{}exposure | country\_\allowbreak{}monthly\_\allowbreak{}risk | country\_\allowbreak{}policy\_\allowbreak{}year | c… & 按项目用途复核；不足以单独支持核心结论。\\
LCL-0046 & 内部/\allowbreak{}03\_\allowbreak{}案例库与证据 & ACQUIRED/\allowbreak{}LOCAL\_\allowbreak{}ASSET & NOT\_\allowbreak{}REQUIRED/\allowbreak{}SUPPLEMENT & company\_\allowbreak{}overseas\_\allowbreak{}exposure | country\_\allowbreak{}monthly\_\allowbreak{}risk | country\_\allowbreak{}policy\_\allowbreak{}year | c… & 按项目用途复核；不足以单独支持核心结论。\\
LCL-0047 & 内部/\allowbreak{}03\_\allowbreak{}案例库与证据 & ACQUIRED/\allowbreak{}LOCAL\_\allowbreak{}ASSET & NOT\_\allowbreak{}REQUIRED/\allowbreak{}SUPPLEMENT & company\_\allowbreak{}overseas\_\allowbreak{}exposure | country\_\allowbreak{}monthly\_\allowbreak{}risk | country\_\allowbreak{}policy\_\allowbreak{}year | c… & 按项目用途复核；不足以单独支持核心结论。\\
LCL-0048 & 内部/\allowbreak{}03\_\allowbreak{}案例库与证据 & ACQUIRED/\allowbreak{}LOCAL\_\allowbreak{}ASSET & NOT\_\allowbreak{}REQUIRED/\allowbreak{}SUPPLEMENT & company\_\allowbreak{}overseas\_\allowbreak{}exposure | country\_\allowbreak{}monthly\_\allowbreak{}risk | country\_\allowbreak{}policy\_\allowbreak{}year | c… & 按项目用途复核；不足以单独支持核心结论。\\
LCL-0049 & 内部/\allowbreak{}03\_\allowbreak{}案例库与证据 & ACQUIRED/\allowbreak{}LOCAL\_\allowbreak{}ASSET & NOT\_\allowbreak{}REQUIRED/\allowbreak{}SUPPLEMENT & company\_\allowbreak{}overseas\_\allowbreak{}exposure | country\_\allowbreak{}monthly\_\allowbreak{}risk | country\_\allowbreak{}policy\_\allowbreak{}year | c… & 按项目用途复核；不足以单独支持核心结论。\\
LCL-0050 & 内部/\allowbreak{}02\_\allowbreak{}方案演进与调研报告 & ACQUIRED/\allowbreak{}DOCUMENT\_\allowbreak{}LOCAL & NOT\_\allowbreak{}REQUIRED/\allowbreak{}SUPPLEMENT & asset\_\allowbreak{}monthly\_\allowbreak{}return | case\_\allowbreak{}evidence | global\_\allowbreak{}cycle\_\allowbreak{}month & 按项目用途复核；不足以单独支持核心结论。\\
LCL-0051 & 内部/\allowbreak{}02\_\allowbreak{}方案演进与调研报告 & ACQUIRED/\allowbreak{}LOCAL\_\allowbreak{}ASSET & NOT\_\allowbreak{}REQUIRED/\allowbreak{}SUPPLEMENT & asset\_\allowbreak{}monthly\_\allowbreak{}return | case\_\allowbreak{}evidence | global\_\allowbreak{}cycle\_\allowbreak{}month & 按项目用途复核；不足以单独支持核心结论。\\
LCL-0052 & 内部/\allowbreak{}02\_\allowbreak{}方案演进与调研报告 & ACQUIRED/\allowbreak{}DOCUMENT\_\allowbreak{}LOCAL & NOT\_\allowbreak{}REQUIRED/\allowbreak{}SUPPLEMENT & asset\_\allowbreak{}monthly\_\allowbreak{}return | case\_\allowbreak{}evidence | global\_\allowbreak{}cycle\_\allowbreak{}month & 按项目用途复核；不足以单独支持核心结论。\\
LCL-0053 & 内部/\allowbreak{}02\_\allowbreak{}方案演进与调研报告 & ACQUIRED/\allowbreak{}DOCUMENT\_\allowbreak{}LOCAL & NOT\_\allowbreak{}REQUIRED/\allowbreak{}SUPPLEMENT & asset\_\allowbreak{}monthly\_\allowbreak{}return | case\_\allowbreak{}evidence | global\_\allowbreak{}cycle\_\allowbreak{}month & 按项目用途复核；不足以单独支持核心结论。\\
LCL-0054 & 内部/\allowbreak{}02\_\allowbreak{}方案演进与调研报告 & ACQUIRED/\allowbreak{}LOCAL\_\allowbreak{}ASSET & NOT\_\allowbreak{}REQUIRED/\allowbreak{}SUPPLEMENT & asset\_\allowbreak{}monthly\_\allowbreak{}return | case\_\allowbreak{}evidence | global\_\allowbreak{}cycle\_\allowbreak{}month & 按项目用途复核；不足以单独支持核心结论。\\
LCL-0055 & 内部/\allowbreak{}02\_\allowbreak{}方案演进与调研报告 & ACQUIRED/\allowbreak{}DOCUMENT\_\allowbreak{}LOCAL & NOT\_\allowbreak{}REQUIRED/\allowbreak{}SUPPLEMENT & asset\_\allowbreak{}monthly\_\allowbreak{}return | case\_\allowbreak{}evidence | global\_\allowbreak{}cycle\_\allowbreak{}month & 按项目用途复核；不足以单独支持核心结论。\\
LCL-0056 & 内部/\allowbreak{}02\_\allowbreak{}方案演进与调研报告 & ACQUIRED/\allowbreak{}LOCAL\_\allowbreak{}ASSET & NOT\_\allowbreak{}REQUIRED/\allowbreak{}SUPPLEMENT & asset\_\allowbreak{}monthly\_\allowbreak{}return | case\_\allowbreak{}evidence | global\_\allowbreak{}cycle\_\allowbreak{}month & 按项目用途复核；不足以单独支持核心结论。\\
LCL-0057 & 内部/\allowbreak{}02\_\allowbreak{}方案演进与调研报告 & ACQUIRED/\allowbreak{}DOCUMENT\_\allowbreak{}LOCAL & NOT\_\allowbreak{}REQUIRED/\allowbreak{}SUPPLEMENT & asset\_\allowbreak{}monthly\_\allowbreak{}return | case\_\allowbreak{}evidence | global\_\allowbreak{}cycle\_\allowbreak{}month & 按项目用途复核；不足以单独支持核心结论。\\
LCL-0058 & 内部/\allowbreak{}02\_\allowbreak{}方案演进与调研报告 & ACQUIRED/\allowbreak{}LOCAL\_\allowbreak{}ASSET & NOT\_\allowbreak{}REQUIRED/\allowbreak{}SUPPLEMENT & asset\_\allowbreak{}monthly\_\allowbreak{}return | case\_\allowbreak{}evidence | global\_\allowbreak{}cycle\_\allowbreak{}month & 按项目用途复核；不足以单独支持核心结论。\\
LCL-0059 & 内部/\allowbreak{}02\_\allowbreak{}方案演进与调研报告 & ACQUIRED/\allowbreak{}DOCUMENT\_\allowbreak{}LOCAL & NOT\_\allowbreak{}REQUIRED/\allowbreak{}SUPPLEMENT & source\_\allowbreak{}registry | asset\_\allowbreak{}monthly\_\allowbreak{}return | case\_\allowbreak{}evidence | global\_\allowbreak{}cycle\_\allowbreak{}month & 按项目用途复核；不足以单独支持核心结论。\\
LCL-0060 & 内部/\allowbreak{}02\_\allowbreak{}方案演进与调研报告 & ACQUIRED/\allowbreak{}LOCAL\_\allowbreak{}ASSET & NOT\_\allowbreak{}REQUIRED/\allowbreak{}SUPPLEMENT & source\_\allowbreak{}registry | asset\_\allowbreak{}monthly\_\allowbreak{}return | case\_\allowbreak{}evidence | global\_\allowbreak{}cycle\_\allowbreak{}month & 按项目用途复核；不足以单独支持核心结论。\\
LCL-0061 & 内部/\allowbreak{}02\_\allowbreak{}方案演进与调研报告 & ACQUIRED/\allowbreak{}DOCUMENT\_\allowbreak{}LOCAL & NOT\_\allowbreak{}REQUIRED/\allowbreak{}SUPPLEMENT & source\_\allowbreak{}registry | asset\_\allowbreak{}monthly\_\allowbreak{}return | case\_\allowbreak{}evidence | global\_\allowbreak{}cycle\_\allowbreak{}month & 按项目用途复核；不足以单独支持核心结论。\\
LCL-0062 & 内部/\allowbreak{}02\_\allowbreak{}方案演进与调研报告 & ACQUIRED/\allowbreak{}LOCAL\_\allowbreak{}ASSET & NOT\_\allowbreak{}REQUIRED/\allowbreak{}SUPPLEMENT & source\_\allowbreak{}registry | asset\_\allowbreak{}monthly\_\allowbreak{}return | case\_\allowbreak{}evidence | global\_\allowbreak{}cycle\_\allowbreak{}month & 按项目用途复核；不足以单独支持核心结论。\\
LCL-0063 & 内部/\allowbreak{}02\_\allowbreak{}方案演进与调研报告 & ACQUIRED/\allowbreak{}LOCAL\_\allowbreak{}ASSET & NOT\_\allowbreak{}REQUIRED/\allowbreak{}SUPPLEMENT & asset\_\allowbreak{}monthly\_\allowbreak{}return | case\_\allowbreak{}evidence | global\_\allowbreak{}cycle\_\allowbreak{}month & 按项目用途复核；不足以单独支持核心结论。\\
LCL-0064 & 内部/\allowbreak{}02\_\allowbreak{}方案演进与调研报告 & ACQUIRED/\allowbreak{}LOCAL\_\allowbreak{}ASSET & NOT\_\allowbreak{}REQUIRED/\allowbreak{}SUPPLEMENT & asset\_\allowbreak{}monthly\_\allowbreak{}return | case\_\allowbreak{}evidence | global\_\allowbreak{}cycle\_\allowbreak{}month & 按项目用途复核；不足以单独支持核心结论。\\
LCL-0065 & 内部/\allowbreak{}02\_\allowbreak{}方案演进与调研报告 & ACQUIRED/\allowbreak{}LOCAL\_\allowbreak{}ASSET & NOT\_\allowbreak{}REQUIRED/\allowbreak{}SUPPLEMENT & asset\_\allowbreak{}monthly\_\allowbreak{}return | case\_\allowbreak{}evidence | global\_\allowbreak{}cycle\_\allowbreak{}month & 按项目用途复核；不足以单独支持核心结论。\\
LCL-0066 & 内部/\allowbreak{}02\_\allowbreak{}方案演进与调研报告 & ACQUIRED/\allowbreak{}LOCAL\_\allowbreak{}ASSET & NOT\_\allowbreak{}REQUIRED/\allowbreak{}SUPPLEMENT & source\_\allowbreak{}registry | asset\_\allowbreak{}monthly\_\allowbreak{}return | case\_\allowbreak{}evidence | global\_\allowbreak{}cycle\_\allowbreak{}month & 按项目用途复核；不足以单独支持核心结论。\\
LCL-0067 & 内部/\allowbreak{}02\_\allowbreak{}方案演进与调研报告 & ACQUIRED/\allowbreak{}LOCAL\_\allowbreak{}ASSET & NOT\_\allowbreak{}REQUIRED/\allowbreak{}SUPPLEMENT & source\_\allowbreak{}registry | asset\_\allowbreak{}monthly\_\allowbreak{}return | case\_\allowbreak{}evidence | global\_\allowbreak{}cycle\_\allowbreak{}month & 按项目用途复核；不足以单独支持核心结论。\\
LCL-0068 & 内部/\allowbreak{}02\_\allowbreak{}方案演进与调研报告 & ACQUIRED/\allowbreak{}LOCAL\_\allowbreak{}ASSET & NOT\_\allowbreak{}REQUIRED/\allowbreak{}SUPPLEMENT & source\_\allowbreak{}registry | asset\_\allowbreak{}monthly\_\allowbreak{}return | case\_\allowbreak{}evidence | global\_\allowbreak{}cycle\_\allowbreak{}month & 按项目用途复核；不足以单独支持核心结论。\\
LCL-0069 & 内部/\allowbreak{}02\_\allowbreak{}方案演进与调研报告 & ACQUIRED/\allowbreak{}LOCAL\_\allowbreak{}ASSET & NOT\_\allowbreak{}REQUIRED/\allowbreak{}SUPPLEMENT & asset\_\allowbreak{}monthly\_\allowbreak{}return | case\_\allowbreak{}evidence | global\_\allowbreak{}cycle\_\allowbreak{}month & 按项目用途复核；不足以单独支持核心结论。\\
LCL-0070 & 内部/\allowbreak{}02\_\allowbreak{}方案演进与调研报告 & ACQUIRED/\allowbreak{}LOCAL\_\allowbreak{}ASSET & NOT\_\allowbreak{}REQUIRED/\allowbreak{}SUPPLEMENT & asset\_\allowbreak{}monthly\_\allowbreak{}return | case\_\allowbreak{}evidence | global\_\allowbreak{}cycle\_\allowbreak{}month & 按项目用途复核；不足以单独支持核心结论。\\
LCL-0071 & 内部/\allowbreak{}02\_\allowbreak{}方案演进与调研报告 & ACQUIRED/\allowbreak{}LOCAL\_\allowbreak{}ASSET & NOT\_\allowbreak{}REQUIRED/\allowbreak{}SUPPLEMENT & asset\_\allowbreak{}monthly\_\allowbreak{}return | case\_\allowbreak{}evidence | global\_\allowbreak{}cycle\_\allowbreak{}month & 按项目用途复核；不足以单独支持核心结论。\\
LCL-0072 & 内部/\allowbreak{}02\_\allowbreak{}方案演进与调研报告 & ACQUIRED/\allowbreak{}LOCAL\_\allowbreak{}ASSET & NOT\_\allowbreak{}REQUIRED/\allowbreak{}SUPPLEMENT & asset\_\allowbreak{}monthly\_\allowbreak{}return | case\_\allowbreak{}evidence | global\_\allowbreak{}cycle\_\allowbreak{}month & 按项目用途复核；不足以单独支持核心结论。\\
LCL-0073 & 内部/\allowbreak{}02\_\allowbreak{}方案演进与调研报告 & ACQUIRED/\allowbreak{}LOCAL\_\allowbreak{}ASSET & NOT\_\allowbreak{}REQUIRED/\allowbreak{}SUPPLEMENT & asset\_\allowbreak{}monthly\_\allowbreak{}return | case\_\allowbreak{}evidence | global\_\allowbreak{}cycle\_\allowbreak{}month & 按项目用途复核；不足以单独支持核心结论。\\
LCL-0074 & 内部/\allowbreak{}02\_\allowbreak{}方案演进与调研报告 & ACQUIRED/\allowbreak{}LOCAL\_\allowbreak{}ASSET & NOT\_\allowbreak{}REQUIRED/\allowbreak{}SUPPLEMENT & asset\_\allowbreak{}monthly\_\allowbreak{}return | case\_\allowbreak{}evidence | global\_\allowbreak{}cycle\_\allowbreak{}month & 按项目用途复核；不足以单独支持核心结论。\\
LCL-0075 & 内部/\allowbreak{}02\_\allowbreak{}方案演进与调研报告 & ACQUIRED/\allowbreak{}LOCAL\_\allowbreak{}ASSET & NOT\_\allowbreak{}REQUIRED/\allowbreak{}SUPPLEMENT & asset\_\allowbreak{}monthly\_\allowbreak{}return | case\_\allowbreak{}evidence | global\_\allowbreak{}cycle\_\allowbreak{}month & 按项目用途复核；不足以单独支持核心结论。\\
LCL-0076 & 内部/\allowbreak{}02\_\allowbreak{}方案演进与调研报告 & ACQUIRED/\allowbreak{}LOCAL\_\allowbreak{}ASSET & NOT\_\allowbreak{}REQUIRED/\allowbreak{}SUPPLEMENT & company\_\allowbreak{}overseas\_\allowbreak{}exposure | asset\_\allowbreak{}monthly\_\allowbreak{}return | case\_\allowbreak{}evidence | global\_\allowbreak{}… & 按项目用途复核；不足以单独支持核心结论。\\
LCL-0077 & 内部/\allowbreak{}01\_\allowbreak{}比赛与项目基线 & ACQUIRED/\allowbreak{}LOCAL\_\allowbreak{}ASSET & NOT\_\allowbreak{}REQUIRED/\allowbreak{}SUPPLEMENT & source\_\allowbreak{}registry | case\_\allowbreak{}evidence & 按项目用途复核；不足以单独支持核心结论。\\
LCL-0078 & 内部/\allowbreak{}02\_\allowbreak{}方案演进与调研报告 & ACQUIRED/\allowbreak{}LOCAL\_\allowbreak{}ASSET & NOT\_\allowbreak{}REQUIRED/\allowbreak{}SUPPLEMENT & asset\_\allowbreak{}monthly\_\allowbreak{}return | case\_\allowbreak{}evidence | global\_\allowbreak{}cycle\_\allowbreak{}month & 按项目用途复核；不足以单独支持核心结论。\\
LCL-0079 & 内部/\allowbreak{}01\_\allowbreak{}比赛与项目基线 & ACQUIRED/\allowbreak{}LOCAL\_\allowbreak{}ASSET & NOT\_\allowbreak{}REQUIRED/\allowbreak{}SUPPLEMENT & source\_\allowbreak{}registry | case\_\allowbreak{}evidence & 按项目用途复核；不足以单独支持核心结论。\\
LCL-0080 & 内部/\allowbreak{}01\_\allowbreak{}比赛与项目基线 & ACQUIRED/\allowbreak{}LOCAL\_\allowbreak{}ASSET & NOT\_\allowbreak{}REQUIRED/\allowbreak{}SUPPLEMENT & source\_\allowbreak{}registry | case\_\allowbreak{}evidence & 按项目用途复核；不足以单独支持核心结论。\\
LCL-0081 & 内部/\allowbreak{}01\_\allowbreak{}比赛与项目基线 & NOT\_\allowbreak{}ACQUIRED/\allowbreak{}NOT\_\allowbreak{}ACQUIRED & FAILED/\allowbreak{}EXCLUDED & source\_\allowbreak{}registry | case\_\allowbreak{}evidence & 仅保留元数据；正式分析改用许可允许的权威替代源。\\
LCL-0082 & 内部/\allowbreak{}10\_\allowbreak{}数据集原件与说明 & ACQUIRED/\allowbreak{}STRUCTURED\_\allowbreak{}LOCAL & NOT\_\allowbreak{}REQUIRED/\allowbreak{}SUPPLEMENT & asset\_\allowbreak{}monthly\_\allowbreak{}return | case\_\allowbreak{}evidence & 按项目用途复核；不足以单独支持核心结论。\\
LCL-0083 & 内部/\allowbreak{}10\_\allowbreak{}数据集原件与说明 & ACQUIRED/\allowbreak{}LOCAL\_\allowbreak{}ASSET & NOT\_\allowbreak{}REQUIRED/\allowbreak{}SUPPLEMENT & asset\_\allowbreak{}monthly\_\allowbreak{}return | case\_\allowbreak{}evidence & 按项目用途复核；不足以单独支持核心结论。\\
LCL-0084 & 内部/\allowbreak{}10\_\allowbreak{}数据集原件与说明 & ACQUIRED/\allowbreak{}STRUCTURED\_\allowbreak{}LOCAL & NOT\_\allowbreak{}REQUIRED/\allowbreak{}SUPPLEMENT & case\_\allowbreak{}evidence & 按项目用途复核；不足以单独支持核心结论。\\
LCL-0085 & 内部/\allowbreak{}10\_\allowbreak{}数据集原件与说明 & ACQUIRED/\allowbreak{}LOCAL\_\allowbreak{}ASSET & NOT\_\allowbreak{}REQUIRED/\allowbreak{}SUPPLEMENT & case\_\allowbreak{}evidence & 按项目用途复核；不足以单独支持核心结论。\\
LCL-0086 & 内部/\allowbreak{}10\_\allowbreak{}数据集原件与说明 & ACQUIRED/\allowbreak{}STRUCTURED\_\allowbreak{}LOCAL & NOT\_\allowbreak{}REQUIRED/\allowbreak{}SUPPLEMENT & case\_\allowbreak{}evidence & 按项目用途复核；不足以单独支持核心结论。\\
LCL-0087 & 内部/\allowbreak{}10\_\allowbreak{}数据集原件与说明 & ACQUIRED/\allowbreak{}LOCAL\_\allowbreak{}ASSET & NOT\_\allowbreak{}REQUIRED/\allowbreak{}SUPPLEMENT & case\_\allowbreak{}evidence & 按项目用途复核；不足以单独支持核心结论。\\
LCL-0088 & 内部/\allowbreak{}10\_\allowbreak{}数据集原件与说明 & ACQUIRED/\allowbreak{}STRUCTURED\_\allowbreak{}LOCAL & NOT\_\allowbreak{}REQUIRED/\allowbreak{}SUPPLEMENT & case\_\allowbreak{}evidence & 按项目用途复核；不足以单独支持核心结论。\\
LCL-0089 & 内部/\allowbreak{}10\_\allowbreak{}数据集原件与说明 & ACQUIRED/\allowbreak{}LOCAL\_\allowbreak{}ASSET & NOT\_\allowbreak{}REQUIRED/\allowbreak{}SUPPLEMENT & case\_\allowbreak{}evidence & 按项目用途复核；不足以单独支持核心结论。\\
LCL-0090 & 内部/\allowbreak{}10\_\allowbreak{}数据集原件与说明 & ACQUIRED/\allowbreak{}LOCAL\_\allowbreak{}ASSET & NOT\_\allowbreak{}REQUIRED/\allowbreak{}SUPPLEMENT & case\_\allowbreak{}evidence & 按项目用途复核；不足以单独支持核心结论。\\
LCL-0091 & 内部/\allowbreak{}10\_\allowbreak{}数据集原件与说明 & ACQUIRED/\allowbreak{}LOCAL\_\allowbreak{}ASSET & NOT\_\allowbreak{}REQUIRED/\allowbreak{}SUPPLEMENT & case\_\allowbreak{}evidence & 按项目用途复核；不足以单独支持核心结论。\\
LCL-0092 & 内部/\allowbreak{}10\_\allowbreak{}数据集原件与说明 & ACQUIRED/\allowbreak{}LOCAL\_\allowbreak{}ASSET & NOT\_\allowbreak{}REQUIRED/\allowbreak{}SUPPLEMENT & case\_\allowbreak{}evidence & 按项目用途复核；不足以单独支持核心结论。\\
LCL-0093 & 内部/\allowbreak{}10\_\allowbreak{}数据集原件与说明 & ACQUIRED/\allowbreak{}LOCAL\_\allowbreak{}ASSET & NOT\_\allowbreak{}REQUIRED/\allowbreak{}SUPPLEMENT & case\_\allowbreak{}evidence & 按项目用途复核；不足以单独支持核心结论。\\
LCL-0094 & 内部/\allowbreak{}10\_\allowbreak{}数据集原件与说明 & ACQUIRED/\allowbreak{}LOCAL\_\allowbreak{}ASSET & NOT\_\allowbreak{}REQUIRED/\allowbreak{}SUPPLEMENT & case\_\allowbreak{}evidence & 按项目用途复核；不足以单独支持核心结论。\\
LCL-0095 & 内部/\allowbreak{}10\_\allowbreak{}数据集原件与说明 & ACQUIRED/\allowbreak{}LOCAL\_\allowbreak{}ASSET & NOT\_\allowbreak{}REQUIRED/\allowbreak{}SUPPLEMENT & asset\_\allowbreak{}monthly\_\allowbreak{}return | case\_\allowbreak{}evidence & 按项目用途复核；不足以单独支持核心结论。\\
LCL-0096 & 内部/\allowbreak{}10\_\allowbreak{}数据集原件与说明 & ACQUIRED/\allowbreak{}LOCAL\_\allowbreak{}ASSET & NOT\_\allowbreak{}REQUIRED/\allowbreak{}SUPPLEMENT & asset\_\allowbreak{}monthly\_\allowbreak{}return | case\_\allowbreak{}evidence & 按项目用途复核；不足以单独支持核心结论。\\
LCL-0097 & 内部/\allowbreak{}10\_\allowbreak{}数据集原件与说明 & ACQUIRED/\allowbreak{}LOCAL\_\allowbreak{}ASSET & NOT\_\allowbreak{}REQUIRED/\allowbreak{}SUPPLEMENT & case\_\allowbreak{}evidence & 按项目用途复核；不足以单独支持核心结论。\\
LCL-0098 & 内部/\allowbreak{}10\_\allowbreak{}数据集原件与说明 & ACQUIRED/\allowbreak{}LOCAL\_\allowbreak{}ASSET & NOT\_\allowbreak{}REQUIRED/\allowbreak{}SUPPLEMENT & case\_\allowbreak{}evidence & 按项目用途复核；不足以单独支持核心结论。\\
LCL-0099 & 内部/\allowbreak{}10\_\allowbreak{}数据集原件与说明 & ACQUIRED/\allowbreak{}LOCAL\_\allowbreak{}ASSET & NOT\_\allowbreak{}REQUIRED/\allowbreak{}SUPPLEMENT & case\_\allowbreak{}evidence & 按项目用途复核；不足以单独支持核心结论。\\
LCL-0100 & 内部/\allowbreak{}10\_\allowbreak{}数据集原件与说明 & ACQUIRED/\allowbreak{}LOCAL\_\allowbreak{}ASSET & NOT\_\allowbreak{}REQUIRED/\allowbreak{}SUPPLEMENT & case\_\allowbreak{}evidence & 按项目用途复核；不足以单独支持核心结论。\\
LCL-0101 & 内部/\allowbreak{}10\_\allowbreak{}数据集原件与说明 & ACQUIRED/\allowbreak{}LOCAL\_\allowbreak{}ASSET & NOT\_\allowbreak{}REQUIRED/\allowbreak{}SUPPLEMENT & case\_\allowbreak{}evidence & 按项目用途复核；不足以单独支持核心结论。\\
LCL-0102 & 内部/\allowbreak{}10\_\allowbreak{}数据集原件与说明 & ACQUIRED/\allowbreak{}LOCAL\_\allowbreak{}ASSET & NOT\_\allowbreak{}REQUIRED/\allowbreak{}SUPPLEMENT & case\_\allowbreak{}evidence & 按项目用途复核；不足以单独支持核心结论。\\
LCL-0103 & 内部/\allowbreak{}10\_\allowbreak{}数据集原件与说明 & ACQUIRED/\allowbreak{}LOCAL\_\allowbreak{}ASSET & NOT\_\allowbreak{}REQUIRED/\allowbreak{}SUPPLEMENT & case\_\allowbreak{}evidence & 按项目用途复核；不足以单独支持核心结论。\\
LCL-0104 & 内部/\allowbreak{}10\_\allowbreak{}数据集原件与说明 & ACQUIRED/\allowbreak{}LOCAL\_\allowbreak{}ASSET & NOT\_\allowbreak{}REQUIRED/\allowbreak{}SUPPLEMENT & case\_\allowbreak{}evidence & 按项目用途复核；不足以单独支持核心结论。\\
LCL-0105 & 内部/\allowbreak{}10\_\allowbreak{}数据集原件与说明 & ACQUIRED/\allowbreak{}LOCAL\_\allowbreak{}ASSET & NOT\_\allowbreak{}REQUIRED/\allowbreak{}SUPPLEMENT & case\_\allowbreak{}evidence & 按项目用途复核；不足以单独支持核心结论。\\
LCL-0106 & 内部/\allowbreak{}10\_\allowbreak{}数据集原件与说明 & ACQUIRED/\allowbreak{}LOCAL\_\allowbreak{}ASSET & NOT\_\allowbreak{}REQUIRED/\allowbreak{}SUPPLEMENT & case\_\allowbreak{}evidence & 按项目用途复核；不足以单独支持核心结论。\\
LCL-0107 & 内部/\allowbreak{}10\_\allowbreak{}数据集原件与说明 & ACQUIRED/\allowbreak{}LOCAL\_\allowbreak{}ASSET & NOT\_\allowbreak{}REQUIRED/\allowbreak{}SUPPLEMENT & case\_\allowbreak{}evidence & 按项目用途复核；不足以单独支持核心结论。\\
LCL-0108 & 内部/\allowbreak{}10\_\allowbreak{}数据集原件与说明 & ACQUIRED/\allowbreak{}LOCAL\_\allowbreak{}ASSET & NOT\_\allowbreak{}REQUIRED/\allowbreak{}SUPPLEMENT & case\_\allowbreak{}evidence & 按项目用途复核；不足以单独支持核心结论。\\
LCL-0109 & 内部/\allowbreak{}10\_\allowbreak{}数据集原件与说明 & ACQUIRED/\allowbreak{}LOCAL\_\allowbreak{}ASSET & NOT\_\allowbreak{}REQUIRED/\allowbreak{}SUPPLEMENT & case\_\allowbreak{}evidence & 按项目用途复核；不足以单独支持核心结论。\\
LCL-0110 & 内部/\allowbreak{}10\_\allowbreak{}数据集原件与说明 & ACQUIRED/\allowbreak{}LOCAL\_\allowbreak{}ASSET & NOT\_\allowbreak{}REQUIRED/\allowbreak{}SUPPLEMENT & asset\_\allowbreak{}monthly\_\allowbreak{}return | case\_\allowbreak{}evidence & 按项目用途复核；不足以单独支持核心结论。\\
LCL-0111 & 内部/\allowbreak{}10\_\allowbreak{}数据集原件与说明 & ACQUIRED/\allowbreak{}LOCAL\_\allowbreak{}ASSET & NOT\_\allowbreak{}REQUIRED/\allowbreak{}SUPPLEMENT & asset\_\allowbreak{}monthly\_\allowbreak{}return | case\_\allowbreak{}evidence & 按项目用途复核；不足以单独支持核心结论。\\
LCL-0112 & 内部/\allowbreak{}10\_\allowbreak{}数据集原件与说明 & ACQUIRED/\allowbreak{}LOCAL\_\allowbreak{}ASSET & NOT\_\allowbreak{}REQUIRED/\allowbreak{}SUPPLEMENT & case\_\allowbreak{}evidence & 按项目用途复核；不足以单独支持核心结论。\\
LCL-0113 & 内部/\allowbreak{}10\_\allowbreak{}数据集原件与说明 & ACQUIRED/\allowbreak{}LOCAL\_\allowbreak{}ASSET & NOT\_\allowbreak{}REQUIRED/\allowbreak{}SUPPLEMENT & case\_\allowbreak{}evidence & 按项目用途复核；不足以单独支持核心结论。\\
LCL-0114 & 内部/\allowbreak{}10\_\allowbreak{}数据集原件与说明 & ACQUIRED/\allowbreak{}LOCAL\_\allowbreak{}ASSET & NOT\_\allowbreak{}REQUIRED/\allowbreak{}SUPPLEMENT & case\_\allowbreak{}evidence & 按项目用途复核；不足以单独支持核心结论。\\
LCL-0115 & 内部/\allowbreak{}10\_\allowbreak{}数据集原件与说明 & ACQUIRED/\allowbreak{}LOCAL\_\allowbreak{}ASSET & NOT\_\allowbreak{}REQUIRED/\allowbreak{}SUPPLEMENT & case\_\allowbreak{}evidence & 按项目用途复核；不足以单独支持核心结论。\\
LCL-0116 & 内部/\allowbreak{}10\_\allowbreak{}数据集原件与说明 & ACQUIRED/\allowbreak{}LOCAL\_\allowbreak{}ASSET & NOT\_\allowbreak{}REQUIRED/\allowbreak{}SUPPLEMENT & case\_\allowbreak{}evidence & 按项目用途复核；不足以单独支持核心结论。\\
LCL-0117 & 内部/\allowbreak{}10\_\allowbreak{}数据集原件与说明 & ACQUIRED/\allowbreak{}LOCAL\_\allowbreak{}ASSET & NOT\_\allowbreak{}REQUIRED/\allowbreak{}SUPPLEMENT & case\_\allowbreak{}evidence & 按项目用途复核；不足以单独支持核心结论。\\
LCL-0118 & 内部/\allowbreak{}10\_\allowbreak{}数据集原件与说明 & ACQUIRED/\allowbreak{}LOCAL\_\allowbreak{}ASSET & NOT\_\allowbreak{}REQUIRED/\allowbreak{}SUPPLEMENT & case\_\allowbreak{}evidence & 按项目用途复核；不足以单独支持核心结论。\\
LCL-0119 & 内部/\allowbreak{}10\_\allowbreak{}数据集原件与说明 & ACQUIRED/\allowbreak{}LOCAL\_\allowbreak{}ASSET & NOT\_\allowbreak{}REQUIRED/\allowbreak{}REPLACED & case\_\allowbreak{}evidence & 使用替代来源：LCL-0105\\
LCL-0120 & 内部/\allowbreak{}10\_\allowbreak{}数据集原件与说明 & ACQUIRED/\allowbreak{}LOCAL\_\allowbreak{}ASSET & NOT\_\allowbreak{}REQUIRED/\allowbreak{}REPLACED & case\_\allowbreak{}evidence & 使用替代来源：LCL-0106\\
LCL-0121 & 内部/\allowbreak{}10\_\allowbreak{}数据集原件与说明 & ACQUIRED/\allowbreak{}LOCAL\_\allowbreak{}ASSET & NOT\_\allowbreak{}REQUIRED/\allowbreak{}SUPPLEMENT & case\_\allowbreak{}evidence & 按项目用途复核；不足以单独支持核心结论。\\
LCL-0122 & 内部/\allowbreak{}10\_\allowbreak{}数据集原件与说明 & ACQUIRED/\allowbreak{}LOCAL\_\allowbreak{}ASSET & NOT\_\allowbreak{}REQUIRED/\allowbreak{}SUPPLEMENT & case\_\allowbreak{}evidence & 按项目用途复核；不足以单独支持核心结论。\\
LCL-0123 & 内部/\allowbreak{}10\_\allowbreak{}数据集原件与说明 & ACQUIRED/\allowbreak{}LOCAL\_\allowbreak{}ASSET & NOT\_\allowbreak{}REQUIRED/\allowbreak{}REPLACED & asset\_\allowbreak{}monthly\_\allowbreak{}return | case\_\allowbreak{}evidence & 使用替代来源：LCL-0095\\
LCL-0124 & 内部/\allowbreak{}10\_\allowbreak{}数据集原件与说明 & ACQUIRED/\allowbreak{}LOCAL\_\allowbreak{}ASSET & NOT\_\allowbreak{}REQUIRED/\allowbreak{}REPLACED & asset\_\allowbreak{}monthly\_\allowbreak{}return | case\_\allowbreak{}evidence & 使用替代来源：LCL-0096\\
LCL-0125 & 内部/\allowbreak{}10\_\allowbreak{}数据集原件与说明 & ACQUIRED/\allowbreak{}LOCAL\_\allowbreak{}ASSET & NOT\_\allowbreak{}REQUIRED/\allowbreak{}REPLACED & case\_\allowbreak{}evidence & 使用替代来源：LCL-0097\\
LCL-0126 & 内部/\allowbreak{}10\_\allowbreak{}数据集原件与说明 & ACQUIRED/\allowbreak{}LOCAL\_\allowbreak{}ASSET & NOT\_\allowbreak{}REQUIRED/\allowbreak{}REPLACED & case\_\allowbreak{}evidence & 使用替代来源：LCL-0098\\
LCL-0127 & 内部/\allowbreak{}10\_\allowbreak{}数据集原件与说明 & ACQUIRED/\allowbreak{}LOCAL\_\allowbreak{}ASSET & NOT\_\allowbreak{}REQUIRED/\allowbreak{}REPLACED & case\_\allowbreak{}evidence & 使用替代来源：LCL-0099\\
LCL-0128 & 内部/\allowbreak{}10\_\allowbreak{}数据集原件与说明 & ACQUIRED/\allowbreak{}LOCAL\_\allowbreak{}ASSET & NOT\_\allowbreak{}REQUIRED/\allowbreak{}SUPPLEMENT & case\_\allowbreak{}evidence & 按项目用途复核；不足以单独支持核心结论。\\
LCL-0129 & 内部/\allowbreak{}10\_\allowbreak{}数据集原件与说明 & ACQUIRED/\allowbreak{}LOCAL\_\allowbreak{}ASSET & NOT\_\allowbreak{}REQUIRED/\allowbreak{}REPLACED & case\_\allowbreak{}evidence & 使用替代来源：LCL-0116\\
LCL-0130 & 内部/\allowbreak{}10\_\allowbreak{}数据集原件与说明 & ACQUIRED/\allowbreak{}LOCAL\_\allowbreak{}ASSET & NOT\_\allowbreak{}REQUIRED/\allowbreak{}SUPPLEMENT & case\_\allowbreak{}evidence & 按项目用途复核；不足以单独支持核心结论。\\
LCL-0131 & 内部/\allowbreak{}10\_\allowbreak{}数据集原件与说明 & ACQUIRED/\allowbreak{}LOCAL\_\allowbreak{}ASSET & NOT\_\allowbreak{}REQUIRED/\allowbreak{}SUPPLEMENT & case\_\allowbreak{}evidence & 按项目用途复核；不足以单独支持核心结论。\\
LCL-0132 & 内部/\allowbreak{}10\_\allowbreak{}数据集原件与说明 & ACQUIRED/\allowbreak{}LOCAL\_\allowbreak{}ASSET & NOT\_\allowbreak{}REQUIRED/\allowbreak{}SUPPLEMENT & case\_\allowbreak{}evidence & 按项目用途复核；不足以单独支持核心结论。\\
LCL-0133 & 内部/\allowbreak{}10\_\allowbreak{}数据集原件与说明 & ACQUIRED/\allowbreak{}LOCAL\_\allowbreak{}ASSET & NOT\_\allowbreak{}REQUIRED/\allowbreak{}SUPPLEMENT & case\_\allowbreak{}evidence & 按项目用途复核；不足以单独支持核心结论。\\
LCL-0134 & 内部/\allowbreak{}10\_\allowbreak{}数据集原件与说明 & ACQUIRED/\allowbreak{}LOCAL\_\allowbreak{}ASSET & NOT\_\allowbreak{}REQUIRED/\allowbreak{}SUPPLEMENT & case\_\allowbreak{}evidence & 按项目用途复核；不足以单独支持核心结论。\\
LCL-0135 & 内部/\allowbreak{}10\_\allowbreak{}数据集原件与说明 & ACQUIRED/\allowbreak{}LOCAL\_\allowbreak{}ASSET & NOT\_\allowbreak{}REQUIRED/\allowbreak{}SUPPLEMENT & case\_\allowbreak{}evidence & 按项目用途复核；不足以单独支持核心结论。\\
LCL-0136 & 内部/\allowbreak{}10\_\allowbreak{}数据集原件与说明 & ACQUIRED/\allowbreak{}LOCAL\_\allowbreak{}ASSET & NOT\_\allowbreak{}REQUIRED/\allowbreak{}SUPPLEMENT & case\_\allowbreak{}evidence & 按项目用途复核；不足以单独支持核心结论。\\
LCL-0137 & 内部/\allowbreak{}10\_\allowbreak{}数据集原件与说明 & ACQUIRED/\allowbreak{}LOCAL\_\allowbreak{}ASSET & NOT\_\allowbreak{}REQUIRED/\allowbreak{}SUPPLEMENT & case\_\allowbreak{}evidence & 按项目用途复核；不足以单独支持核心结论。\\
LCL-0138 & 内部/\allowbreak{}10\_\allowbreak{}数据集原件与说明 & ACQUIRED/\allowbreak{}LOCAL\_\allowbreak{}ASSET & NOT\_\allowbreak{}REQUIRED/\allowbreak{}SUPPLEMENT & case\_\allowbreak{}evidence & 按项目用途复核；不足以单独支持核心结论。\\
LCL-0139 & 内部/\allowbreak{}10\_\allowbreak{}数据集原件与说明 & ACQUIRED/\allowbreak{}LOCAL\_\allowbreak{}ASSET & NOT\_\allowbreak{}REQUIRED/\allowbreak{}SUPPLEMENT & asset\_\allowbreak{}monthly\_\allowbreak{}return | case\_\allowbreak{}evidence & 按项目用途复核；不足以单独支持核心结论。\\
LCL-0140 & 内部/\allowbreak{}10\_\allowbreak{}数据集原件与说明 & ACQUIRED/\allowbreak{}LOCAL\_\allowbreak{}ASSET & NOT\_\allowbreak{}REQUIRED/\allowbreak{}SUPPLEMENT & case\_\allowbreak{}evidence & 按项目用途复核；不足以单独支持核心结论。\\
LCL-0141 & 内部/\allowbreak{}10\_\allowbreak{}数据集原件与说明 & ACQUIRED/\allowbreak{}LOCAL\_\allowbreak{}ASSET & NOT\_\allowbreak{}REQUIRED/\allowbreak{}SUPPLEMENT & case\_\allowbreak{}evidence & 按项目用途复核；不足以单独支持核心结论。\\
LCL-0142 & 内部/\allowbreak{}10\_\allowbreak{}数据集原件与说明 & ACQUIRED/\allowbreak{}LOCAL\_\allowbreak{}ASSET & NOT\_\allowbreak{}REQUIRED/\allowbreak{}SUPPLEMENT & case\_\allowbreak{}evidence & 按项目用途复核；不足以单独支持核心结论。\\
LCL-0143 & 内部/\allowbreak{}10\_\allowbreak{}数据集原件与说明 & ACQUIRED/\allowbreak{}STRUCTURED\_\allowbreak{}LOCAL & NOT\_\allowbreak{}REQUIRED/\allowbreak{}SUPPLEMENT & case\_\allowbreak{}evidence & 按项目用途复核；不足以单独支持核心结论。\\
LCL-0144 & 内部/\allowbreak{}10\_\allowbreak{}数据集原件与说明 & ACQUIRED/\allowbreak{}STRUCTURED\_\allowbreak{}LOCAL & NOT\_\allowbreak{}REQUIRED/\allowbreak{}SUPPLEMENT & case\_\allowbreak{}evidence & 按项目用途复核；不足以单独支持核心结论。\\
LCL-0145 & 内部/\allowbreak{}10\_\allowbreak{}数据集原件与说明 & ACQUIRED/\allowbreak{}STRUCTURED\_\allowbreak{}LOCAL & NOT\_\allowbreak{}REQUIRED/\allowbreak{}SUPPLEMENT & case\_\allowbreak{}evidence & 按项目用途复核；不足以单独支持核心结论。\\
LCL-0146 & 内部/\allowbreak{}10\_\allowbreak{}数据集原件与说明 & ACQUIRED/\allowbreak{}STRUCTURED\_\allowbreak{}LOCAL & NOT\_\allowbreak{}REQUIRED/\allowbreak{}SUPPLEMENT & asset\_\allowbreak{}monthly\_\allowbreak{}return | case\_\allowbreak{}evidence & 按项目用途复核；不足以单独支持核心结论。\\
LCL-0147 & 内部/\allowbreak{}10\_\allowbreak{}数据集原件与说明 & ACQUIRED/\allowbreak{}STRUCTURED\_\allowbreak{}LOCAL & NOT\_\allowbreak{}REQUIRED/\allowbreak{}SUPPLEMENT & country\_\allowbreak{}exposure | case\_\allowbreak{}evidence & 按项目用途复核；不足以单独支持核心结论。\\
LCL-0148 & 内部/\allowbreak{}10\_\allowbreak{}数据集原件与说明 & NOT\_\allowbreak{}APPLICABLE/\allowbreak{}NOT\_\allowbreak{}ACQUIRED & FAILED/\allowbreak{}EXCLUDED & case\_\allowbreak{}evidence & 仅保留元数据；正式分析改用许可允许的权威替代源。\\
LCL-0149 & 内部/\allowbreak{}10\_\allowbreak{}数据集原件与说明 & ACQUIRED/\allowbreak{}STRUCTURED\_\allowbreak{}LOCAL & NOT\_\allowbreak{}REQUIRED/\allowbreak{}SUPPLEMENT & case\_\allowbreak{}evidence & 按项目用途复核；不足以单独支持核心结论。\\
LCL-0150 & 内部/\allowbreak{}10\_\allowbreak{}数据集原件与说明 & ACQUIRED/\allowbreak{}LOCAL\_\allowbreak{}ASSET & NOT\_\allowbreak{}REQUIRED/\allowbreak{}SUPPLEMENT & case\_\allowbreak{}evidence & 按项目用途复核；不足以单独支持核心结论。\\
LCL-0151 & 内部/\allowbreak{}10\_\allowbreak{}数据集原件与说明 & ACQUIRED/\allowbreak{}LOCAL\_\allowbreak{}ASSET & NOT\_\allowbreak{}REQUIRED/\allowbreak{}SUPPLEMENT & case\_\allowbreak{}evidence & 按项目用途复核；不足以单独支持核心结论。\\
LCL-0152 & 内部/\allowbreak{}10\_\allowbreak{}数据集原件与说明 & ACQUIRED/\allowbreak{}LOCAL\_\allowbreak{}ASSET & NOT\_\allowbreak{}REQUIRED/\allowbreak{}SUPPLEMENT & case\_\allowbreak{}evidence & 按项目用途复核；不足以单独支持核心结论。\\
LCL-0153 & 内部/\allowbreak{}10\_\allowbreak{}数据集原件与说明 & ACQUIRED/\allowbreak{}STRUCTURED\_\allowbreak{}LOCAL & NOT\_\allowbreak{}REQUIRED/\allowbreak{}SUPPLEMENT & case\_\allowbreak{}evidence & 按项目用途复核；不足以单独支持核心结论。\\
LCL-0154 & 内部/\allowbreak{}10\_\allowbreak{}数据集原件与说明 & ACQUIRED/\allowbreak{}LOCAL\_\allowbreak{}ASSET & NOT\_\allowbreak{}REQUIRED/\allowbreak{}SUPPLEMENT & case\_\allowbreak{}evidence & 按项目用途复核；不足以单独支持核心结论。\\
LCL-0155 & 内部/\allowbreak{}10\_\allowbreak{}数据集原件与说明 & ACQUIRED/\allowbreak{}LOCAL\_\allowbreak{}ASSET & NOT\_\allowbreak{}REQUIRED/\allowbreak{}SUPPLEMENT & case\_\allowbreak{}evidence & 按项目用途复核；不足以单独支持核心结论。\\
LCL-0156 & 内部/\allowbreak{}10\_\allowbreak{}数据集原件与说明 & ACQUIRED/\allowbreak{}EMPTY\_\allowbreak{}SHELL & FAILED/\allowbreak{}EXCLUDED & case\_\allowbreak{}evidence & 补充权威替代来源；无法取得时删除相应字段和结论承诺。\\
LCL-0157 & 内部/\allowbreak{}10\_\allowbreak{}数据集原件与说明 & ACQUIRED/\allowbreak{}NARRATIVE\_\allowbreak{}CONTENT & EXTRACTED/\allowbreak{}SUPPLEMENT & asset\_\allowbreak{}monthly\_\allowbreak{}return | case\_\allowbreak{}evidence & 按项目用途复核；不足以单独支持核心结论。\\
LCL-0158 & 内部/\allowbreak{}10\_\allowbreak{}数据集原件与说明 & ACQUIRED/\allowbreak{}LOCAL\_\allowbreak{}ASSET & NOT\_\allowbreak{}REQUIRED/\allowbreak{}SUPPLEMENT & asset\_\allowbreak{}monthly\_\allowbreak{}return | case\_\allowbreak{}evidence & 按项目用途复核；不足以单独支持核心结论。\\
LCL-0159 & 内部/\allowbreak{}10\_\allowbreak{}数据集原件与说明 & ACQUIRED/\allowbreak{}LOCAL\_\allowbreak{}ASSET & NOT\_\allowbreak{}REQUIRED/\allowbreak{}SUPPLEMENT & asset\_\allowbreak{}monthly\_\allowbreak{}return | case\_\allowbreak{}evidence & 按项目用途复核；不足以单独支持核心结论。\\
LCL-0160 & 内部/\allowbreak{}10\_\allowbreak{}数据集原件与说明 & ACQUIRED/\allowbreak{}NARRATIVE\_\allowbreak{}CONTENT & EXTRACTED/\allowbreak{}SUPPLEMENT & asset\_\allowbreak{}monthly\_\allowbreak{}return | case\_\allowbreak{}evidence & 按项目用途复核；不足以单独支持核心结论。\\
LCL-0161 & 内部/\allowbreak{}10\_\allowbreak{}数据集原件与说明 & ACQUIRED/\allowbreak{}LOCAL\_\allowbreak{}ASSET & NOT\_\allowbreak{}REQUIRED/\allowbreak{}SUPPLEMENT & asset\_\allowbreak{}monthly\_\allowbreak{}return | case\_\allowbreak{}evidence & 按项目用途复核；不足以单独支持核心结论。\\
LCL-0162 & 内部/\allowbreak{}10\_\allowbreak{}数据集原件与说明 & NOT\_\allowbreak{}APPLICABLE/\allowbreak{}NOT\_\allowbreak{}ACQUIRED & FAILED/\allowbreak{}EXCLUDED & case\_\allowbreak{}evidence & 仅保留元数据；正式分析改用许可允许的权威替代源。\\
LCL-0163 & 内部/\allowbreak{}10\_\allowbreak{}数据集原件与说明 & NOT\_\allowbreak{}APPLICABLE/\allowbreak{}NOT\_\allowbreak{}ACQUIRED & FAILED/\allowbreak{}EXCLUDED & case\_\allowbreak{}evidence & 仅保留元数据；正式分析改用许可允许的权威替代源。\\
LCL-0164 & 内部/\allowbreak{}10\_\allowbreak{}数据集原件与说明 & ACQUIRED/\allowbreak{}STRUCTURED\_\allowbreak{}LOCAL & NOT\_\allowbreak{}REQUIRED/\allowbreak{}SUPPLEMENT & case\_\allowbreak{}evidence & 按项目用途复核；不足以单独支持核心结论。\\
LCL-0165 & 内部/\allowbreak{}10\_\allowbreak{}数据集原件与说明 & ACQUIRED/\allowbreak{}LOCAL\_\allowbreak{}ASSET & NOT\_\allowbreak{}REQUIRED/\allowbreak{}SUPPLEMENT & case\_\allowbreak{}evidence & 按项目用途复核；不足以单独支持核心结论。\\
LCL-0166 & 内部/\allowbreak{}10\_\allowbreak{}数据集原件与说明 & ACQUIRED/\allowbreak{}LOCAL\_\allowbreak{}ASSET & NOT\_\allowbreak{}REQUIRED/\allowbreak{}SUPPLEMENT & case\_\allowbreak{}evidence & 按项目用途复核；不足以单独支持核心结论。\\
LCL-0167 & 内部/\allowbreak{}10\_\allowbreak{}数据集原件与说明 & ACQUIRED/\allowbreak{}LOCAL\_\allowbreak{}ASSET & NOT\_\allowbreak{}REQUIRED/\allowbreak{}SUPPLEMENT & case\_\allowbreak{}evidence & 按项目用途复核；不足以单独支持核心结论。\\
LCL-0168 & 内部/\allowbreak{}10\_\allowbreak{}数据集原件与说明 & ACQUIRED/\allowbreak{}LOCAL\_\allowbreak{}ASSET & NOT\_\allowbreak{}REQUIRED/\allowbreak{}SUPPLEMENT & case\_\allowbreak{}evidence & 按项目用途复核；不足以单独支持核心结论。\\
LCL-0169 & 内部/\allowbreak{}10\_\allowbreak{}数据集原件与说明 & ACQUIRED/\allowbreak{}LOCAL\_\allowbreak{}ASSET & NOT\_\allowbreak{}REQUIRED/\allowbreak{}SUPPLEMENT & asset\_\allowbreak{}monthly\_\allowbreak{}return | case\_\allowbreak{}evidence & 按项目用途复核；不足以单独支持核心结论。\\
LCL-0170 & 内部/\allowbreak{}10\_\allowbreak{}数据集原件与说明 & ACQUIRED/\allowbreak{}LOCAL\_\allowbreak{}ASSET & NOT\_\allowbreak{}REQUIRED/\allowbreak{}SUPPLEMENT & case\_\allowbreak{}evidence & 按项目用途复核；不足以单独支持核心结论。\\
LCL-0171 & 内部/\allowbreak{}10\_\allowbreak{}数据集原件与说明 & ACQUIRED/\allowbreak{}LOCAL\_\allowbreak{}ASSET & NOT\_\allowbreak{}REQUIRED/\allowbreak{}SUPPLEMENT & asset\_\allowbreak{}monthly\_\allowbreak{}return | case\_\allowbreak{}evidence & 按项目用途复核；不足以单独支持核心结论。\\
LCL-0172 & 内部/\allowbreak{}10\_\allowbreak{}数据集原件与说明 & ACQUIRED/\allowbreak{}LOCAL\_\allowbreak{}ASSET & NOT\_\allowbreak{}REQUIRED/\allowbreak{}SUPPLEMENT & case\_\allowbreak{}evidence & 按项目用途复核；不足以单独支持核心结论。\\
LCL-0173 & 内部/\allowbreak{}10\_\allowbreak{}数据集原件与说明 & ACQUIRED/\allowbreak{}LOCAL\_\allowbreak{}ASSET & NOT\_\allowbreak{}REQUIRED/\allowbreak{}SUPPLEMENT & case\_\allowbreak{}evidence & 按项目用途复核；不足以单独支持核心结论。\\
LCL-0174 & 内部/\allowbreak{}10\_\allowbreak{}数据集原件与说明 & ACQUIRED/\allowbreak{}LOCAL\_\allowbreak{}ASSET & NOT\_\allowbreak{}REQUIRED/\allowbreak{}SUPPLEMENT & case\_\allowbreak{}evidence & 按项目用途复核；不足以单独支持核心结论。\\
LCL-0175 & 内部/\allowbreak{}10\_\allowbreak{}数据集原件与说明 & ACQUIRED/\allowbreak{}LOCAL\_\allowbreak{}ASSET & NOT\_\allowbreak{}REQUIRED/\allowbreak{}SUPPLEMENT & case\_\allowbreak{}evidence & 按项目用途复核；不足以单独支持核心结论。\\
LCL-0176 & 内部/\allowbreak{}10\_\allowbreak{}数据集原件与说明 & ACQUIRED/\allowbreak{}LOCAL\_\allowbreak{}ASSET & NOT\_\allowbreak{}REQUIRED/\allowbreak{}SUPPLEMENT & case\_\allowbreak{}evidence & 按项目用途复核；不足以单独支持核心结论。\\
LCL-0177 & 内部/\allowbreak{}10\_\allowbreak{}数据集原件与说明 & ACQUIRED/\allowbreak{}LOCAL\_\allowbreak{}ASSET & NOT\_\allowbreak{}REQUIRED/\allowbreak{}SUPPLEMENT & case\_\allowbreak{}evidence & 按项目用途复核；不足以单独支持核心结论。\\
LCL-0178 & 内部/\allowbreak{}10\_\allowbreak{}数据集原件与说明 & ACQUIRED/\allowbreak{}LOCAL\_\allowbreak{}ASSET & NOT\_\allowbreak{}REQUIRED/\allowbreak{}SUPPLEMENT & case\_\allowbreak{}evidence & 按项目用途复核；不足以单独支持核心结论。\\
LCL-0179 & A/\allowbreak{}01\_\allowbreak{}比赛与项目基线 & ACQUIRED/\allowbreak{}DOCUMENT\_\allowbreak{}LOCAL & NOT\_\allowbreak{}REQUIRED/\allowbreak{}CORE & source\_\allowbreak{}registry | case\_\allowbreak{}evidence & 人工复核关键字段与source\_\allowbreak{}locator；通过后映射至正式source\_\allowbreak{}registry。\\
LCL-0180 & 内部/\allowbreak{}01\_\allowbreak{}比赛与项目基线 & ACQUIRED/\allowbreak{}STRUCTURED\_\allowbreak{}LOCAL & NOT\_\allowbreak{}REQUIRED/\allowbreak{}SUPPLEMENT & source\_\allowbreak{}registry | case\_\allowbreak{}evidence & 按项目用途复核；不足以单独支持核心结论。\\
WEB-0001 & D/\allowbreak{}05\_\allowbreak{}全球周期与历史危机 & NOT\_\allowbreak{}APPLICABLE/\allowbreak{}NOT\_\allowbreak{}ACQUIRED & FAILED/\allowbreak{}EXCLUDED & asset\_\allowbreak{}monthly\_\allowbreak{}return | case\_\allowbreak{}evidence | historical\_\allowbreak{}crisis\_\allowbreak{}event | global\_\allowbreak{}cy… & 仅保留元数据；正式分析改用许可允许的权威替代源。\\
WEB-0002 & B/\allowbreak{}06\_\allowbreak{}黄金美元与储备配置 & ACQUIRED/\allowbreak{}MIXED\_\allowbreak{}CONTENT & EXTRACTED/\allowbreak{}SUPPLEMENT & asset\_\allowbreak{}monthly\_\allowbreak{}return | portfolio\_\allowbreak{}scenario & 按项目用途复核；不足以单独支持核心结论。\\
WEB-0003 & D/\allowbreak{}05\_\allowbreak{}全球周期与历史危机 & NOT\_\allowbreak{}APPLICABLE/\allowbreak{}NOT\_\allowbreak{}ACQUIRED & FAILED/\allowbreak{}EXCLUDED & asset\_\allowbreak{}monthly\_\allowbreak{}return | case\_\allowbreak{}evidence | historical\_\allowbreak{}crisis\_\allowbreak{}event | global\_\allowbreak{}cy… & 仅保留元数据；正式分析改用许可允许的权威替代源。\\
WEB-0004 & D/\allowbreak{}08\_\allowbreak{}企业案例与公开披露 & NOT\_\allowbreak{}APPLICABLE/\allowbreak{}NOT\_\allowbreak{}ACQUIRED & FAILED/\allowbreak{}EXCLUDED & company\_\allowbreak{}overseas\_\allowbreak{}exposure | country\_\allowbreak{}event | asset\_\allowbreak{}monthly\_\allowbreak{}return | portfol… & 仅保留元数据；正式分析改用许可允许的权威替代源。\\
WEB-0005 & D/\allowbreak{}06\_\allowbreak{}黄金美元与储备配置 & NOT\_\allowbreak{}APPLICABLE/\allowbreak{}NOT\_\allowbreak{}ACQUIRED & FAILED/\allowbreak{}EXCLUDED & asset\_\allowbreak{}monthly\_\allowbreak{}return | portfolio\_\allowbreak{}scenario & 仅保留元数据；正式分析改用许可允许的权威替代源。\\
WEB-0006 & D/\allowbreak{}06\_\allowbreak{}黄金美元与储备配置 & NOT\_\allowbreak{}APPLICABLE/\allowbreak{}NOT\_\allowbreak{}ACQUIRED & FAILED/\allowbreak{}EXCLUDED & asset\_\allowbreak{}monthly\_\allowbreak{}return | portfolio\_\allowbreak{}scenario & 仅保留元数据；正式分析改用许可允许的权威替代源。\\
WEB-0007 & D/\allowbreak{}06\_\allowbreak{}黄金美元与储备配置 & NOT\_\allowbreak{}APPLICABLE/\allowbreak{}NOT\_\allowbreak{}ACQUIRED & FAILED/\allowbreak{}EXCLUDED & asset\_\allowbreak{}monthly\_\allowbreak{}return | portfolio\_\allowbreak{}scenario & 仅保留元数据；正式分析改用许可允许的权威替代源。\\
WEB-0008 & D/\allowbreak{}08\_\allowbreak{}企业案例与公开披露 & ACQUIRED/\allowbreak{}DOCUMENT\_\allowbreak{}LOCAL & NOT\_\allowbreak{}REQUIRED/\allowbreak{}EXCLUDED & company\_\allowbreak{}overseas\_\allowbreak{}exposure | country\_\allowbreak{}event | asset\_\allowbreak{}monthly\_\allowbreak{}return | portfol… & 补充权威替代来源；无法取得时删除相应字段和结论承诺。\\
WEB-0009 & D/\allowbreak{}08\_\allowbreak{}企业案例与公开披露 & NOT\_\allowbreak{}APPLICABLE/\allowbreak{}NOT\_\allowbreak{}ACQUIRED & FAILED/\allowbreak{}EXCLUDED & company\_\allowbreak{}overseas\_\allowbreak{}exposure | country\_\allowbreak{}event | asset\_\allowbreak{}monthly\_\allowbreak{}return | portfol… & 仅保留元数据；正式分析改用许可允许的权威替代源。\\
WEB-0010 & D/\allowbreak{}08\_\allowbreak{}企业案例与公开披露 & NOT\_\allowbreak{}APPLICABLE/\allowbreak{}NOT\_\allowbreak{}ACQUIRED & FAILED/\allowbreak{}EXCLUDED & company\_\allowbreak{}overseas\_\allowbreak{}exposure | country\_\allowbreak{}event | asset\_\allowbreak{}monthly\_\allowbreak{}return | portfol… & 仅保留元数据；正式分析改用许可允许的权威替代源。\\
WEB-0011 & A/\allowbreak{}08\_\allowbreak{}企业案例与公开披露 & ACQUIRED/\allowbreak{}MIXED\_\allowbreak{}CONTENT & EXTRACTED/\allowbreak{}CORE & company\_\allowbreak{}overseas\_\allowbreak{}exposure | country\_\allowbreak{}event | asset\_\allowbreak{}monthly\_\allowbreak{}return | portfol… & 人工复核关键字段与source\_\allowbreak{}locator；通过后映射至正式source\_\allowbreak{}registry。\\
WEB-0012 & D/\allowbreak{}08\_\allowbreak{}企业案例与公开披露 & NOT\_\allowbreak{}APPLICABLE/\allowbreak{}NOT\_\allowbreak{}ACQUIRED & FAILED/\allowbreak{}EXCLUDED & company\_\allowbreak{}overseas\_\allowbreak{}exposure | country\_\allowbreak{}event | asset\_\allowbreak{}monthly\_\allowbreak{}return | portfol… & 仅保留元数据；正式分析改用许可允许的权威替代源。\\
WEB-0013 & B/\allowbreak{}06\_\allowbreak{}黄金美元与储备配置 & ACQUIRED/\allowbreak{}STRUCTURED\_\allowbreak{}DATA & EXTRACTED/\allowbreak{}SUPPLEMENT & asset\_\allowbreak{}monthly\_\allowbreak{}return | portfolio\_\allowbreak{}scenario & 按项目用途复核；不足以单独支持核心结论。\\
WEB-0014 & D/\allowbreak{}06\_\allowbreak{}黄金美元与储备配置 & NOT\_\allowbreak{}APPLICABLE/\allowbreak{}NOT\_\allowbreak{}ACQUIRED & FAILED/\allowbreak{}EXCLUDED & asset\_\allowbreak{}monthly\_\allowbreak{}return | portfolio\_\allowbreak{}scenario & 仅保留元数据；正式分析改用许可允许的权威替代源。\\
WEB-0015 & D/\allowbreak{}07\_\allowbreak{}会计合规与制裁 & NOT\_\allowbreak{}APPLICABLE/\allowbreak{}NOT\_\allowbreak{}ACQUIRED & FAILED/\allowbreak{}EXCLUDED & country\_\allowbreak{}policy\_\allowbreak{}year & 按人工任务卡接受条款或完成机构授权后重新验收。\\
WEB-0016 & D/\allowbreak{}08\_\allowbreak{}企业案例与公开披露 & ACQUIRED/\allowbreak{}DOCUMENT\_\allowbreak{}LOCAL & NOT\_\allowbreak{}REQUIRED/\allowbreak{}EXCLUDED & company\_\allowbreak{}overseas\_\allowbreak{}exposure | country\_\allowbreak{}event | asset\_\allowbreak{}monthly\_\allowbreak{}return | portfol… & 补充权威替代来源；无法取得时删除相应字段和结论承诺。\\
WEB-0017 & D/\allowbreak{}08\_\allowbreak{}企业案例与公开披露 & NOT\_\allowbreak{}APPLICABLE/\allowbreak{}NOT\_\allowbreak{}ACQUIRED & FAILED/\allowbreak{}EXCLUDED & company\_\allowbreak{}overseas\_\allowbreak{}exposure | country\_\allowbreak{}event | asset\_\allowbreak{}monthly\_\allowbreak{}return | portfol… & 仅保留元数据；正式分析改用许可允许的权威替代源。\\
WEB-0018 & D/\allowbreak{}08\_\allowbreak{}企业案例与公开披露 & NOT\_\allowbreak{}APPLICABLE/\allowbreak{}NOT\_\allowbreak{}ACQUIRED & FAILED/\allowbreak{}EXCLUDED & company\_\allowbreak{}overseas\_\allowbreak{}exposure | country\_\allowbreak{}event | asset\_\allowbreak{}monthly\_\allowbreak{}return | portfol… & 仅保留元数据；正式分析改用许可允许的权威替代源。\\
WEB-0019 & D/\allowbreak{}06\_\allowbreak{}黄金美元与储备配置 & NOT\_\allowbreak{}APPLICABLE/\allowbreak{}NOT\_\allowbreak{}ACQUIRED & FAILED/\allowbreak{}EXCLUDED & asset\_\allowbreak{}monthly\_\allowbreak{}return | portfolio\_\allowbreak{}scenario & 按人工任务卡接受条款或完成机构授权后重新验收。\\
WEB-0020 & D/\allowbreak{}05\_\allowbreak{}全球周期与历史危机 & ACQUIRED/\allowbreak{}DOCUMENT\_\allowbreak{}LOCAL & NOT\_\allowbreak{}REQUIRED/\allowbreak{}EXCLUDED & asset\_\allowbreak{}monthly\_\allowbreak{}return | case\_\allowbreak{}evidence | historical\_\allowbreak{}crisis\_\allowbreak{}event | global\_\allowbreak{}cy… & 补充权威替代来源；无法取得时删除相应字段和结论承诺。\\
WEB-0021 & D/\allowbreak{}06\_\allowbreak{}黄金美元与储备配置 & NOT\_\allowbreak{}APPLICABLE/\allowbreak{}NOT\_\allowbreak{}ACQUIRED & FAILED/\allowbreak{}EXCLUDED & country\_\allowbreak{}monthly\_\allowbreak{}risk | country\_\allowbreak{}event | asset\_\allowbreak{}monthly\_\allowbreak{}return | portfolio\_\allowbreak{}sc… & 仅保留元数据；正式分析改用许可允许的权威替代源。\\
WEB-0022 & D/\allowbreak{}08\_\allowbreak{}企业案例与公开披露 & NOT\_\allowbreak{}APPLICABLE/\allowbreak{}NOT\_\allowbreak{}ACQUIRED & FAILED/\allowbreak{}EXCLUDED & company\_\allowbreak{}overseas\_\allowbreak{}exposure | country\_\allowbreak{}event | asset\_\allowbreak{}monthly\_\allowbreak{}return | portfol… & 仅保留元数据；正式分析改用许可允许的权威替代源。\\
WEB-0023 & D/\allowbreak{}08\_\allowbreak{}企业案例与公开披露 & NOT\_\allowbreak{}APPLICABLE/\allowbreak{}NOT\_\allowbreak{}ACQUIRED & FAILED/\allowbreak{}EXCLUDED & company\_\allowbreak{}overseas\_\allowbreak{}exposure | country\_\allowbreak{}event | asset\_\allowbreak{}monthly\_\allowbreak{}return | portfol… & 仅保留元数据；正式分析改用许可允许的权威替代源。\\
WEB-0024 & D/\allowbreak{}05\_\allowbreak{}全球周期与历史危机 & NOT\_\allowbreak{}APPLICABLE/\allowbreak{}NOT\_\allowbreak{}ACQUIRED & FAILED/\allowbreak{}EXCLUDED & case\_\allowbreak{}evidence | historical\_\allowbreak{}crisis\_\allowbreak{}event | global\_\allowbreak{}cycle\_\allowbreak{}month & 仅保留元数据；正式分析改用许可允许的权威替代源。\\
WEB-0025 & D/\allowbreak{}08\_\allowbreak{}企业案例与公开披露 & NOT\_\allowbreak{}APPLICABLE/\allowbreak{}NOT\_\allowbreak{}ACQUIRED & FAILED/\allowbreak{}EXCLUDED & company\_\allowbreak{}overseas\_\allowbreak{}exposure | country\_\allowbreak{}event | asset\_\allowbreak{}monthly\_\allowbreak{}return | portfol… & 仅保留元数据；正式分析改用许可允许的权威替代源。\\
WEB-0026 & B/\allowbreak{}06\_\allowbreak{}黄金美元与储备配置 & ACQUIRED/\allowbreak{}NARRATIVE\_\allowbreak{}CONTENT & EXTRACTED/\allowbreak{}SUPPLEMENT & asset\_\allowbreak{}monthly\_\allowbreak{}return | portfolio\_\allowbreak{}scenario & 按项目用途复核；不足以单独支持核心结论。\\
WEB-0027 & D/\allowbreak{}05\_\allowbreak{}全球周期与历史危机 & NOT\_\allowbreak{}APPLICABLE/\allowbreak{}NOT\_\allowbreak{}ACQUIRED & FAILED/\allowbreak{}EXCLUDED & case\_\allowbreak{}evidence | historical\_\allowbreak{}crisis\_\allowbreak{}event | global\_\allowbreak{}cycle\_\allowbreak{}month & 仅保留元数据；正式分析改用许可允许的权威替代源。\\
WEB-0028 & D/\allowbreak{}08\_\allowbreak{}企业案例与公开披露 & NOT\_\allowbreak{}APPLICABLE/\allowbreak{}NOT\_\allowbreak{}ACQUIRED & FAILED/\allowbreak{}EXCLUDED & company\_\allowbreak{}overseas\_\allowbreak{}exposure | country\_\allowbreak{}event | asset\_\allowbreak{}monthly\_\allowbreak{}return | portfol… & 仅保留元数据；正式分析改用许可允许的权威替代源。\\
WEB-0029 & D/\allowbreak{}08\_\allowbreak{}企业案例与公开披露 & NOT\_\allowbreak{}APPLICABLE/\allowbreak{}NOT\_\allowbreak{}ACQUIRED & FAILED/\allowbreak{}EXCLUDED & company\_\allowbreak{}overseas\_\allowbreak{}exposure | country\_\allowbreak{}event | asset\_\allowbreak{}monthly\_\allowbreak{}return | portfol… & 仅保留元数据；正式分析改用许可允许的权威替代源。\\
WEB-0030 & D/\allowbreak{}06\_\allowbreak{}黄金美元与储备配置 & NOT\_\allowbreak{}APPLICABLE/\allowbreak{}NOT\_\allowbreak{}ACQUIRED & FAILED/\allowbreak{}EXCLUDED & asset\_\allowbreak{}monthly\_\allowbreak{}return | portfolio\_\allowbreak{}scenario & 仅保留元数据；正式分析改用许可允许的权威替代源。\\
WEB-0031 & D/\allowbreak{}06\_\allowbreak{}黄金美元与储备配置 & NOT\_\allowbreak{}APPLICABLE/\allowbreak{}NOT\_\allowbreak{}ACQUIRED & FAILED/\allowbreak{}EXCLUDED & country\_\allowbreak{}monthly\_\allowbreak{}risk | country\_\allowbreak{}event | asset\_\allowbreak{}monthly\_\allowbreak{}return | portfolio\_\allowbreak{}sc… & 仅保留元数据；正式分析改用许可允许的权威替代源。\\
WEB-0032 & D/\allowbreak{}08\_\allowbreak{}企业案例与公开披露 & NOT\_\allowbreak{}APPLICABLE/\allowbreak{}NOT\_\allowbreak{}ACQUIRED & FAILED/\allowbreak{}EXCLUDED & company\_\allowbreak{}overseas\_\allowbreak{}exposure | country\_\allowbreak{}event | asset\_\allowbreak{}monthly\_\allowbreak{}return | portfol… & 仅保留元数据；正式分析改用许可允许的权威替代源。\\
WEB-0033 & D/\allowbreak{}06\_\allowbreak{}黄金美元与储备配置 & NOT\_\allowbreak{}APPLICABLE/\allowbreak{}NOT\_\allowbreak{}ACQUIRED & FAILED/\allowbreak{}EXCLUDED & asset\_\allowbreak{}monthly\_\allowbreak{}return | portfolio\_\allowbreak{}scenario & 仅保留元数据；正式分析改用许可允许的权威替代源。\\
WEB-0034 & D/\allowbreak{}05\_\allowbreak{}全球周期与历史危机 & NOT\_\allowbreak{}APPLICABLE/\allowbreak{}NOT\_\allowbreak{}ACQUIRED & FAILED/\allowbreak{}EXCLUDED & case\_\allowbreak{}evidence | historical\_\allowbreak{}crisis\_\allowbreak{}event | global\_\allowbreak{}cycle\_\allowbreak{}month & 仅保留元数据；正式分析改用许可允许的权威替代源。\\
WEB-0035 & D/\allowbreak{}05\_\allowbreak{}全球周期与历史危机 & NOT\_\allowbreak{}APPLICABLE/\allowbreak{}NOT\_\allowbreak{}ACQUIRED & FAILED/\allowbreak{}EXCLUDED & case\_\allowbreak{}evidence | historical\_\allowbreak{}crisis\_\allowbreak{}event | global\_\allowbreak{}cycle\_\allowbreak{}month & 仅保留元数据；正式分析改用许可允许的权威替代源。\\
WEB-0036 & D/\allowbreak{}07\_\allowbreak{}会计合规与制裁 & NOT\_\allowbreak{}APPLICABLE/\allowbreak{}NOT\_\allowbreak{}ACQUIRED & FAILED/\allowbreak{}EXCLUDED & country\_\allowbreak{}policy\_\allowbreak{}year | asset\_\allowbreak{}monthly\_\allowbreak{}return | case\_\allowbreak{}evidence & 仅保留元数据；正式分析改用许可允许的权威替代源。\\
WEB-0037 & D/\allowbreak{}06\_\allowbreak{}黄金美元与储备配置 & NOT\_\allowbreak{}APPLICABLE/\allowbreak{}NOT\_\allowbreak{}ACQUIRED & FAILED/\allowbreak{}EXCLUDED & asset\_\allowbreak{}monthly\_\allowbreak{}return | portfolio\_\allowbreak{}scenario & 仅保留元数据；正式分析改用许可允许的权威替代源。\\
WEB-0038 & D/\allowbreak{}06\_\allowbreak{}黄金美元与储备配置 & NOT\_\allowbreak{}APPLICABLE/\allowbreak{}NOT\_\allowbreak{}ACQUIRED & FAILED/\allowbreak{}EXCLUDED & asset\_\allowbreak{}monthly\_\allowbreak{}return | portfolio\_\allowbreak{}scenario & 仅保留元数据；正式分析改用许可允许的权威替代源。\\
WEB-0039 & B/\allowbreak{}04\_\allowbreak{}国别货币与中国出海 & ACQUIRED/\allowbreak{}DOWNLOAD\_\allowbreak{}LANDING & EXTRACTED/\allowbreak{}SUPPLEMENT & source\_\allowbreak{}registry | country\_\allowbreak{}monthly\_\allowbreak{}risk | country\_\allowbreak{}policy\_\allowbreak{}year & 按项目用途复核；不足以单独支持核心结论。\\
WEB-0040 & B/\allowbreak{}06\_\allowbreak{}黄金美元与储备配置 & ACQUIRED/\allowbreak{}NARRATIVE\_\allowbreak{}CONTENT & EXTRACTED/\allowbreak{}SUPPLEMENT & source\_\allowbreak{}registry | asset\_\allowbreak{}monthly\_\allowbreak{}return | portfolio\_\allowbreak{}scenario & 按项目用途复核；不足以单独支持核心结论。\\
WEB-0041 & D/\allowbreak{}04\_\allowbreak{}国别货币与中国出海 & NOT\_\allowbreak{}APPLICABLE/\allowbreak{}NOT\_\allowbreak{}ACQUIRED & FAILED/\allowbreak{}EXCLUDED & country\_\allowbreak{}monthly\_\allowbreak{}risk | country\_\allowbreak{}policy\_\allowbreak{}year | country\_\allowbreak{}event | case\_\allowbreak{}evidence & 仅保留元数据；正式分析改用许可允许的权威替代源。\\
WEB-0042 & B/\allowbreak{}04\_\allowbreak{}国别货币与中国出海 & ACQUIRED/\allowbreak{}DOWNLOAD\_\allowbreak{}LANDING & EXTRACTED/\allowbreak{}SUPPLEMENT & country\_\allowbreak{}monthly\_\allowbreak{}risk | country\_\allowbreak{}policy\_\allowbreak{}year & 按项目用途复核；不足以单独支持核心结论。\\
WEB-0043 & B/\allowbreak{}04\_\allowbreak{}国别货币与中国出海 & ACQUIRED/\allowbreak{}DOWNLOAD\_\allowbreak{}LANDING & EXTRACTED/\allowbreak{}SUPPLEMENT & country\_\allowbreak{}monthly\_\allowbreak{}risk | country\_\allowbreak{}policy\_\allowbreak{}year & 按项目用途复核；不足以单独支持核心结论。\\
WEB-0044 & B/\allowbreak{}04\_\allowbreak{}国别货币与中国出海 & ACQUIRED/\allowbreak{}NARRATIVE\_\allowbreak{}CONTENT & EXTRACTED/\allowbreak{}SUPPLEMENT & country\_\allowbreak{}monthly\_\allowbreak{}risk | country\_\allowbreak{}policy\_\allowbreak{}year & 按项目用途复核；不足以单独支持核心结论。\\
WEB-0045 & B/\allowbreak{}07\_\allowbreak{}会计合规与制裁 & ACQUIRED/\allowbreak{}MIXED\_\allowbreak{}CONTENT & EXTRACTED/\allowbreak{}SUPPLEMENT & country\_\allowbreak{}monthly\_\allowbreak{}risk | country\_\allowbreak{}policy\_\allowbreak{}year | country\_\allowbreak{}event | case\_\allowbreak{}evidence… & 按项目用途复核；不足以单独支持核心结论。\\
WEB-0046 & B/\allowbreak{}06\_\allowbreak{}黄金美元与储备配置 & ACQUIRED/\allowbreak{}NARRATIVE\_\allowbreak{}CONTENT & EXTRACTED/\allowbreak{}SUPPLEMENT & source\_\allowbreak{}registry | asset\_\allowbreak{}monthly\_\allowbreak{}return | portfolio\_\allowbreak{}scenario & 按项目用途复核；不足以单独支持核心结论。\\
WEB-0047 & B/\allowbreak{}04\_\allowbreak{}国别货币与中国出海 & NOT\_\allowbreak{}ACQUIRED/\allowbreak{}NOT\_\allowbreak{}ACQUIRED & FAILED/\allowbreak{}SUPPLEMENT & country\_\allowbreak{}monthly\_\allowbreak{}risk | country\_\allowbreak{}policy\_\allowbreak{}year & 补充权威替代来源；无法取得时删除相应字段和结论承诺。\\
WEB-0048 & B/\allowbreak{}06\_\allowbreak{}黄金美元与储备配置 & ACQUIRED/\allowbreak{}DOWNLOAD\_\allowbreak{}LANDING & EXTRACTED/\allowbreak{}SUPPLEMENT & source\_\allowbreak{}registry | asset\_\allowbreak{}monthly\_\allowbreak{}return | portfolio\_\allowbreak{}scenario & 按项目用途复核；不足以单独支持核心结论。\\
WEB-0049 & B/\allowbreak{}04\_\allowbreak{}国别货币与中国出海 & ACQUIRED/\allowbreak{}NARRATIVE\_\allowbreak{}CONTENT & EXTRACTED/\allowbreak{}SUPPLEMENT & source\_\allowbreak{}registry | country\_\allowbreak{}monthly\_\allowbreak{}risk | country\_\allowbreak{}policy\_\allowbreak{}year & 按项目用途复核；不足以单独支持核心结论。\\
WEB-0050 & B/\allowbreak{}06\_\allowbreak{}黄金美元与储备配置 & ACQUIRED/\allowbreak{}NARRATIVE\_\allowbreak{}CONTENT & EXTRACTED/\allowbreak{}SUPPLEMENT & source\_\allowbreak{}registry | asset\_\allowbreak{}monthly\_\allowbreak{}return | portfolio\_\allowbreak{}scenario & 按项目用途复核；不足以单独支持核心结论。\\
WEB-0051 & D/\allowbreak{}04\_\allowbreak{}国别货币与中国出海 & NOT\_\allowbreak{}APPLICABLE/\allowbreak{}NOT\_\allowbreak{}ACQUIRED & FAILED/\allowbreak{}EXCLUDED & country\_\allowbreak{}monthly\_\allowbreak{}risk | country\_\allowbreak{}policy\_\allowbreak{}year | country\_\allowbreak{}event | case\_\allowbreak{}evidence & 仅保留元数据；正式分析改用许可允许的权威替代源。\\
WEB-0052 & D/\allowbreak{}08\_\allowbreak{}企业案例与公开披露 & NOT\_\allowbreak{}APPLICABLE/\allowbreak{}NOT\_\allowbreak{}ACQUIRED & FAILED/\allowbreak{}EXCLUDED & company\_\allowbreak{}overseas\_\allowbreak{}exposure | country\_\allowbreak{}event | asset\_\allowbreak{}monthly\_\allowbreak{}return | portfol… & 仅保留元数据；正式分析改用许可允许的权威替代源。\\
WEB-0053 & D/\allowbreak{}05\_\allowbreak{}全球周期与历史危机 & NOT\_\allowbreak{}APPLICABLE/\allowbreak{}NOT\_\allowbreak{}ACQUIRED & FAILED/\allowbreak{}EXCLUDED & country\_\allowbreak{}monthly\_\allowbreak{}risk | case\_\allowbreak{}evidence | historical\_\allowbreak{}crisis\_\allowbreak{}event | global\_\allowbreak{}cy… & 仅保留元数据；正式分析改用许可允许的权威替代源。\\
WEB-0054 & D/\allowbreak{}08\_\allowbreak{}企业案例与公开披露 & NOT\_\allowbreak{}ACQUIRED/\allowbreak{}NOT\_\allowbreak{}ACQUIRED & FAILED/\allowbreak{}EXCLUDED & company\_\allowbreak{}overseas\_\allowbreak{}exposure | country\_\allowbreak{}event | asset\_\allowbreak{}monthly\_\allowbreak{}return | portfol… & 补充权威替代来源；无法取得时删除相应字段和结论承诺。\\
WEB-0055 & D/\allowbreak{}04\_\allowbreak{}国别货币与中国出海 & NOT\_\allowbreak{}APPLICABLE/\allowbreak{}NOT\_\allowbreak{}ACQUIRED & FAILED/\allowbreak{}EXCLUDED & country\_\allowbreak{}monthly\_\allowbreak{}risk | country\_\allowbreak{}policy\_\allowbreak{}year & 仅保留元数据；正式分析改用许可允许的权威替代源。\\
WEB-0056 & D/\allowbreak{}06\_\allowbreak{}黄金美元与储备配置 & NOT\_\allowbreak{}APPLICABLE/\allowbreak{}NOT\_\allowbreak{}ACQUIRED & FAILED/\allowbreak{}EXCLUDED & asset\_\allowbreak{}monthly\_\allowbreak{}return | portfolio\_\allowbreak{}scenario & 仅保留元数据；正式分析改用许可允许的权威替代源。\\
WEB-0057 & D/\allowbreak{}06\_\allowbreak{}黄金美元与储备配置 & NOT\_\allowbreak{}APPLICABLE/\allowbreak{}NOT\_\allowbreak{}ACQUIRED & FAILED/\allowbreak{}EXCLUDED & country\_\allowbreak{}monthly\_\allowbreak{}risk | country\_\allowbreak{}event | asset\_\allowbreak{}monthly\_\allowbreak{}return | portfolio\_\allowbreak{}sc… & 仅保留元数据；正式分析改用许可允许的权威替代源。\\
WEB-0058 & D/\allowbreak{}05\_\allowbreak{}全球周期与历史危机 & NOT\_\allowbreak{}APPLICABLE/\allowbreak{}NOT\_\allowbreak{}ACQUIRED & FAILED/\allowbreak{}EXCLUDED & case\_\allowbreak{}evidence | historical\_\allowbreak{}crisis\_\allowbreak{}event | global\_\allowbreak{}cycle\_\allowbreak{}month & 仅保留元数据；正式分析改用许可允许的权威替代源。\\
WEB-0059 & D/\allowbreak{}05\_\allowbreak{}全球周期与历史危机 & NOT\_\allowbreak{}APPLICABLE/\allowbreak{}NOT\_\allowbreak{}ACQUIRED & FAILED/\allowbreak{}EXCLUDED & case\_\allowbreak{}evidence | historical\_\allowbreak{}crisis\_\allowbreak{}event | global\_\allowbreak{}cycle\_\allowbreak{}month & 仅保留元数据；正式分析改用许可允许的权威替代源。\\
WEB-0060 & B/\allowbreak{}06\_\allowbreak{}黄金美元与储备配置 & ACQUIRED/\allowbreak{}STRUCTURED\_\allowbreak{}DATA & EXTRACTED/\allowbreak{}SUPPLEMENT & asset\_\allowbreak{}monthly\_\allowbreak{}return | portfolio\_\allowbreak{}scenario & 按项目用途复核；不足以单独支持核心结论。\\
WEB-0061 & B/\allowbreak{}06\_\allowbreak{}黄金美元与储备配置 & NOT\_\allowbreak{}APPLICABLE/\allowbreak{}NOT\_\allowbreak{}ACQUIRED & FAILED/\allowbreak{}EXCLUDED & asset\_\allowbreak{}monthly\_\allowbreak{}return | portfolio\_\allowbreak{}scenario & 仅保留元数据；正式分析改用许可允许的权威替代源。\\
WEB-0062 & D/\allowbreak{}08\_\allowbreak{}企业案例与公开披露 & NOT\_\allowbreak{}APPLICABLE/\allowbreak{}NOT\_\allowbreak{}ACQUIRED & FAILED/\allowbreak{}EXCLUDED & company\_\allowbreak{}overseas\_\allowbreak{}exposure | country\_\allowbreak{}event | asset\_\allowbreak{}monthly\_\allowbreak{}return | portfol… & 仅保留元数据；正式分析改用许可允许的权威替代源。\\
WEB-0063 & C/\allowbreak{}08\_\allowbreak{}企业案例与公开披露 & NOT\_\allowbreak{}APPLICABLE/\allowbreak{}NOT\_\allowbreak{}ACQUIRED & FAILED/\allowbreak{}EXCLUDED & company\_\allowbreak{}overseas\_\allowbreak{}exposure | country\_\allowbreak{}event | asset\_\allowbreak{}monthly\_\allowbreak{}return | portfol… & 仅保留元数据；正式分析改用许可允许的权威替代源。\\
WEB-0064 & D/\allowbreak{}08\_\allowbreak{}企业案例与公开披露 & NOT\_\allowbreak{}APPLICABLE/\allowbreak{}NOT\_\allowbreak{}ACQUIRED & FAILED/\allowbreak{}EXCLUDED & company\_\allowbreak{}overseas\_\allowbreak{}exposure | country\_\allowbreak{}event | asset\_\allowbreak{}monthly\_\allowbreak{}return | portfol… & 仅保留元数据；正式分析改用许可允许的权威替代源。\\
WEB-0065 & D/\allowbreak{}08\_\allowbreak{}企业案例与公开披露 & NOT\_\allowbreak{}APPLICABLE/\allowbreak{}NOT\_\allowbreak{}ACQUIRED & FAILED/\allowbreak{}EXCLUDED & company\_\allowbreak{}overseas\_\allowbreak{}exposure | country\_\allowbreak{}event | asset\_\allowbreak{}monthly\_\allowbreak{}return | portfol… & 仅保留元数据；正式分析改用许可允许的权威替代源。\\
WEB-0066 & C/\allowbreak{}08\_\allowbreak{}企业案例与公开披露 & NOT\_\allowbreak{}APPLICABLE/\allowbreak{}NOT\_\allowbreak{}ACQUIRED & FAILED/\allowbreak{}EXCLUDED & company\_\allowbreak{}overseas\_\allowbreak{}exposure | country\_\allowbreak{}event | asset\_\allowbreak{}monthly\_\allowbreak{}return | portfol… & 仅保留元数据；正式分析改用许可允许的权威替代源。\\
WEB-0067 & C/\allowbreak{}05\_\allowbreak{}全球周期与历史危机 & NOT\_\allowbreak{}APPLICABLE/\allowbreak{}NOT\_\allowbreak{}ACQUIRED & FAILED/\allowbreak{}EXCLUDED & asset\_\allowbreak{}monthly\_\allowbreak{}return | case\_\allowbreak{}evidence | historical\_\allowbreak{}crisis\_\allowbreak{}event | global\_\allowbreak{}cy… & 仅保留元数据；正式分析改用许可允许的权威替代源。\\
WEB-0068 & C/\allowbreak{}08\_\allowbreak{}企业案例与公开披露 & NOT\_\allowbreak{}APPLICABLE/\allowbreak{}NOT\_\allowbreak{}ACQUIRED & FAILED/\allowbreak{}EXCLUDED & company\_\allowbreak{}overseas\_\allowbreak{}exposure | country\_\allowbreak{}event | asset\_\allowbreak{}monthly\_\allowbreak{}return | portfol… & 仅保留元数据；正式分析改用许可允许的权威替代源。\\
WEB-0069 & C/\allowbreak{}08\_\allowbreak{}企业案例与公开披露 & NOT\_\allowbreak{}APPLICABLE/\allowbreak{}NOT\_\allowbreak{}ACQUIRED & FAILED/\allowbreak{}EXCLUDED & company\_\allowbreak{}overseas\_\allowbreak{}exposure | country\_\allowbreak{}event | asset\_\allowbreak{}monthly\_\allowbreak{}return | portfol… & 仅保留元数据；正式分析改用许可允许的权威替代源。\\
WEB-0070 & C/\allowbreak{}08\_\allowbreak{}企业案例与公开披露 & NOT\_\allowbreak{}APPLICABLE/\allowbreak{}NOT\_\allowbreak{}ACQUIRED & FAILED/\allowbreak{}EXCLUDED & company\_\allowbreak{}overseas\_\allowbreak{}exposure | country\_\allowbreak{}event | asset\_\allowbreak{}monthly\_\allowbreak{}return | portfol… & 仅保留元数据；正式分析改用许可允许的权威替代源。\\
WEB-0071 & D/\allowbreak{}05\_\allowbreak{}全球周期与历史危机 & NOT\_\allowbreak{}APPLICABLE/\allowbreak{}NOT\_\allowbreak{}ACQUIRED & FAILED/\allowbreak{}EXCLUDED & case\_\allowbreak{}evidence | historical\_\allowbreak{}crisis\_\allowbreak{}event | global\_\allowbreak{}cycle\_\allowbreak{}month & 仅保留元数据；正式分析改用许可允许的权威替代源。\\
WEB-0072 & D/\allowbreak{}05\_\allowbreak{}全球周期与历史危机 & NOT\_\allowbreak{}APPLICABLE/\allowbreak{}NOT\_\allowbreak{}ACQUIRED & FAILED/\allowbreak{}EXCLUDED & portfolio\_\allowbreak{}scenario | case\_\allowbreak{}evidence | historical\_\allowbreak{}crisis\_\allowbreak{}event | global\_\allowbreak{}cycl… & 仅保留元数据；正式分析改用许可允许的权威替代源。\\
WEB-0073 & D/\allowbreak{}08\_\allowbreak{}企业案例与公开披露 & ACQUIRED/\allowbreak{}DOCUMENT\_\allowbreak{}LOCAL & NOT\_\allowbreak{}REQUIRED/\allowbreak{}EXCLUDED & company\_\allowbreak{}overseas\_\allowbreak{}exposure | country\_\allowbreak{}event | portfolio\_\allowbreak{}scenario | case\_\allowbreak{}evid… & 补充权威替代来源；无法取得时删除相应字段和结论承诺。\\
WEB-0074 & D/\allowbreak{}06\_\allowbreak{}黄金美元与储备配置 & NOT\_\allowbreak{}APPLICABLE/\allowbreak{}NOT\_\allowbreak{}ACQUIRED & FAILED/\allowbreak{}EXCLUDED & source\_\allowbreak{}registry | asset\_\allowbreak{}monthly\_\allowbreak{}return | portfolio\_\allowbreak{}scenario & 仅保留元数据；正式分析改用许可允许的权威替代源。\\
WEB-0075 & D/\allowbreak{}06\_\allowbreak{}黄金美元与储备配置 & NOT\_\allowbreak{}APPLICABLE/\allowbreak{}NOT\_\allowbreak{}ACQUIRED & FAILED/\allowbreak{}EXCLUDED & country\_\allowbreak{}policy\_\allowbreak{}year | asset\_\allowbreak{}monthly\_\allowbreak{}return | portfolio\_\allowbreak{}scenario & 仅保留元数据；正式分析改用许可允许的权威替代源。\\
WEB-0076 & D/\allowbreak{}06\_\allowbreak{}黄金美元与储备配置 & NOT\_\allowbreak{}APPLICABLE/\allowbreak{}NOT\_\allowbreak{}ACQUIRED & FAILED/\allowbreak{}EXCLUDED & country\_\allowbreak{}monthly\_\allowbreak{}risk | country\_\allowbreak{}event | asset\_\allowbreak{}monthly\_\allowbreak{}return | portfolio\_\allowbreak{}sc… & 仅保留元数据；正式分析改用许可允许的权威替代源。\\
WEB-0077 & D/\allowbreak{}06\_\allowbreak{}黄金美元与储备配置 & NOT\_\allowbreak{}APPLICABLE/\allowbreak{}NOT\_\allowbreak{}ACQUIRED & FAILED/\allowbreak{}EXCLUDED & country\_\allowbreak{}monthly\_\allowbreak{}risk | country\_\allowbreak{}event | asset\_\allowbreak{}monthly\_\allowbreak{}return | portfolio\_\allowbreak{}sc… & 仅保留元数据；正式分析改用许可允许的权威替代源。\\
WEB-0078 & D/\allowbreak{}08\_\allowbreak{}企业案例与公开披露 & NOT\_\allowbreak{}APPLICABLE/\allowbreak{}NOT\_\allowbreak{}ACQUIRED & FAILED/\allowbreak{}EXCLUDED & company\_\allowbreak{}overseas\_\allowbreak{}exposure | country\_\allowbreak{}monthly\_\allowbreak{}risk | country\_\allowbreak{}event | asset\_\allowbreak{}m… & 仅保留元数据；正式分析改用许可允许的权威替代源。\\
WEB-0079 & D/\allowbreak{}06\_\allowbreak{}黄金美元与储备配置 & NOT\_\allowbreak{}APPLICABLE/\allowbreak{}NOT\_\allowbreak{}ACQUIRED & FAILED/\allowbreak{}EXCLUDED & source\_\allowbreak{}registry | asset\_\allowbreak{}monthly\_\allowbreak{}return | portfolio\_\allowbreak{}scenario & 仅保留元数据；正式分析改用许可允许的权威替代源。\\
WEB-0080 & D/\allowbreak{}07\_\allowbreak{}会计合规与制裁 & NOT\_\allowbreak{}APPLICABLE/\allowbreak{}NOT\_\allowbreak{}ACQUIRED & FAILED/\allowbreak{}EXCLUDED & country\_\allowbreak{}monthly\_\allowbreak{}risk | country\_\allowbreak{}policy\_\allowbreak{}year | country\_\allowbreak{}event | asset\_\allowbreak{}monthly… & 仅保留元数据；正式分析改用许可允许的权威替代源。\\
WEB-0081 & D/\allowbreak{}05\_\allowbreak{}全球周期与历史危机 & NOT\_\allowbreak{}APPLICABLE/\allowbreak{}NOT\_\allowbreak{}ACQUIRED & FAILED/\allowbreak{}EXCLUDED & country\_\allowbreak{}monthly\_\allowbreak{}risk | asset\_\allowbreak{}monthly\_\allowbreak{}return | case\_\allowbreak{}evidence | historical\_\allowbreak{}c… & 仅保留元数据；正式分析改用许可允许的权威替代源。\\
WEB-0082 & D/\allowbreak{}08\_\allowbreak{}企业案例与公开披露 & NOT\_\allowbreak{}APPLICABLE/\allowbreak{}NOT\_\allowbreak{}ACQUIRED & FAILED/\allowbreak{}EXCLUDED & company\_\allowbreak{}overseas\_\allowbreak{}exposure | country\_\allowbreak{}monthly\_\allowbreak{}risk | country\_\allowbreak{}event | asset\_\allowbreak{}m… & 仅保留元数据；正式分析改用许可允许的权威替代源。\\
WEB-0083 & D/\allowbreak{}06\_\allowbreak{}黄金美元与储备配置 & NOT\_\allowbreak{}APPLICABLE/\allowbreak{}NOT\_\allowbreak{}ACQUIRED & FAILED/\allowbreak{}EXCLUDED & country\_\allowbreak{}monthly\_\allowbreak{}risk | country\_\allowbreak{}event | asset\_\allowbreak{}monthly\_\allowbreak{}return | portfolio\_\allowbreak{}sc… & 仅保留元数据；正式分析改用许可允许的权威替代源。\\
WEB-0084 & D/\allowbreak{}05\_\allowbreak{}全球周期与历史危机 & NOT\_\allowbreak{}APPLICABLE/\allowbreak{}NOT\_\allowbreak{}ACQUIRED & FAILED/\allowbreak{}EXCLUDED & asset\_\allowbreak{}monthly\_\allowbreak{}return | case\_\allowbreak{}evidence | historical\_\allowbreak{}crisis\_\allowbreak{}event | global\_\allowbreak{}cy… & 仅保留元数据；正式分析改用许可允许的权威替代源。\\
WEB-0085 & D/\allowbreak{}07\_\allowbreak{}会计合规与制裁 & NOT\_\allowbreak{}APPLICABLE/\allowbreak{}NOT\_\allowbreak{}ACQUIRED & FAILED/\allowbreak{}EXCLUDED & country\_\allowbreak{}monthly\_\allowbreak{}risk | country\_\allowbreak{}policy\_\allowbreak{}year | asset\_\allowbreak{}monthly\_\allowbreak{}return | case\_\allowbreak{}e… & 仅保留元数据；正式分析改用许可允许的权威替代源。\\
WEB-0086 & D/\allowbreak{}05\_\allowbreak{}全球周期与历史危机 & NOT\_\allowbreak{}APPLICABLE/\allowbreak{}NOT\_\allowbreak{}ACQUIRED & FAILED/\allowbreak{}EXCLUDED & asset\_\allowbreak{}monthly\_\allowbreak{}return | case\_\allowbreak{}evidence | historical\_\allowbreak{}crisis\_\allowbreak{}event | global\_\allowbreak{}cy… & 仅保留元数据；正式分析改用许可允许的权威替代源。\\
WEB-0087 & D/\allowbreak{}06\_\allowbreak{}黄金美元与储备配置 & NOT\_\allowbreak{}APPLICABLE/\allowbreak{}NOT\_\allowbreak{}ACQUIRED & FAILED/\allowbreak{}EXCLUDED & source\_\allowbreak{}registry | asset\_\allowbreak{}monthly\_\allowbreak{}return | portfolio\_\allowbreak{}scenario & 仅保留元数据；正式分析改用许可允许的权威替代源。\\
WEB-0088 & D/\allowbreak{}05\_\allowbreak{}全球周期与历史危机 & NOT\_\allowbreak{}APPLICABLE/\allowbreak{}NOT\_\allowbreak{}ACQUIRED & FAILED/\allowbreak{}EXCLUDED & country\_\allowbreak{}monthly\_\allowbreak{}risk | asset\_\allowbreak{}monthly\_\allowbreak{}return | case\_\allowbreak{}evidence | historical\_\allowbreak{}c… & 仅保留元数据；正式分析改用许可允许的权威替代源。\\
WEB-0089 & D/\allowbreak{}06\_\allowbreak{}黄金美元与储备配置 & NOT\_\allowbreak{}APPLICABLE/\allowbreak{}NOT\_\allowbreak{}ACQUIRED & FAILED/\allowbreak{}EXCLUDED & source\_\allowbreak{}registry | asset\_\allowbreak{}monthly\_\allowbreak{}return | portfolio\_\allowbreak{}scenario & 仅保留元数据；正式分析改用许可允许的权威替代源。\\
WEB-0090 & D/\allowbreak{}06\_\allowbreak{}黄金美元与储备配置 & NOT\_\allowbreak{}APPLICABLE/\allowbreak{}NOT\_\allowbreak{}ACQUIRED & FAILED/\allowbreak{}EXCLUDED & source\_\allowbreak{}registry | asset\_\allowbreak{}monthly\_\allowbreak{}return | portfolio\_\allowbreak{}scenario & 仅保留元数据；正式分析改用许可允许的权威替代源。\\
WEB-0091 & D/\allowbreak{}08\_\allowbreak{}企业案例与公开披露 & NOT\_\allowbreak{}APPLICABLE/\allowbreak{}NOT\_\allowbreak{}ACQUIRED & FAILED/\allowbreak{}EXCLUDED & company\_\allowbreak{}overseas\_\allowbreak{}exposure | country\_\allowbreak{}event | asset\_\allowbreak{}monthly\_\allowbreak{}return | portfol… & 仅保留元数据；正式分析改用许可允许的权威替代源。\\
WEB-0092 & D/\allowbreak{}06\_\allowbreak{}黄金美元与储备配置 & NOT\_\allowbreak{}APPLICABLE/\allowbreak{}NOT\_\allowbreak{}ACQUIRED & FAILED/\allowbreak{}EXCLUDED & asset\_\allowbreak{}monthly\_\allowbreak{}return | portfolio\_\allowbreak{}scenario & 由资料负责人合法登录、采购或申请授权；不得自动绕过。\\
WEB-0093 & D/\allowbreak{}08\_\allowbreak{}企业案例与公开披露 & NOT\_\allowbreak{}APPLICABLE/\allowbreak{}NOT\_\allowbreak{}ACQUIRED & FAILED/\allowbreak{}EXCLUDED & company\_\allowbreak{}overseas\_\allowbreak{}exposure | country\_\allowbreak{}event | asset\_\allowbreak{}monthly\_\allowbreak{}return | portfol… & 仅保留元数据；正式分析改用许可允许的权威替代源。\\
WEB-0094 & D/\allowbreak{}05\_\allowbreak{}全球周期与历史危机 & NOT\_\allowbreak{}APPLICABLE/\allowbreak{}NOT\_\allowbreak{}ACQUIRED & FAILED/\allowbreak{}EXCLUDED & asset\_\allowbreak{}monthly\_\allowbreak{}return | case\_\allowbreak{}evidence | historical\_\allowbreak{}crisis\_\allowbreak{}event | global\_\allowbreak{}cy… & 仅保留元数据；正式分析改用许可允许的权威替代源。\\
WEB-0095 & D/\allowbreak{}05\_\allowbreak{}全球周期与历史危机 & NOT\_\allowbreak{}APPLICABLE/\allowbreak{}NOT\_\allowbreak{}ACQUIRED & FAILED/\allowbreak{}EXCLUDED & case\_\allowbreak{}evidence | historical\_\allowbreak{}crisis\_\allowbreak{}event | global\_\allowbreak{}cycle\_\allowbreak{}month & 仅保留元数据；正式分析改用许可允许的权威替代源。\\
WEB-0096 & D/\allowbreak{}06\_\allowbreak{}黄金美元与储备配置 & NOT\_\allowbreak{}APPLICABLE/\allowbreak{}NOT\_\allowbreak{}ACQUIRED & FAILED/\allowbreak{}EXCLUDED & asset\_\allowbreak{}monthly\_\allowbreak{}return | portfolio\_\allowbreak{}scenario & 仅保留元数据；正式分析改用许可允许的权威替代源。\\
WEB-0097 & D/\allowbreak{}06\_\allowbreak{}黄金美元与储备配置 & NOT\_\allowbreak{}APPLICABLE/\allowbreak{}NOT\_\allowbreak{}ACQUIRED & FAILED/\allowbreak{}EXCLUDED & asset\_\allowbreak{}monthly\_\allowbreak{}return | portfolio\_\allowbreak{}scenario & 仅保留元数据；正式分析改用许可允许的权威替代源。\\
WEB-0098 & A/\allowbreak{}06\_\allowbreak{}黄金美元与储备配置 & ACQUIRED/\allowbreak{}NARRATIVE\_\allowbreak{}CONTENT & EXTRACTED/\allowbreak{}CORE & asset\_\allowbreak{}monthly\_\allowbreak{}return | portfolio\_\allowbreak{}scenario & 人工复核关键字段与source\_\allowbreak{}locator；通过后映射至正式source\_\allowbreak{}registry。\\
WEB-0099 & A/\allowbreak{}05\_\allowbreak{}全球周期与历史危机 & ACQUIRED/\allowbreak{}NARRATIVE\_\allowbreak{}CONTENT & EXTRACTED/\allowbreak{}CORE & asset\_\allowbreak{}monthly\_\allowbreak{}return | case\_\allowbreak{}evidence | historical\_\allowbreak{}crisis\_\allowbreak{}event | global\_\allowbreak{}cy… & 人工复核关键字段与source\_\allowbreak{}locator；通过后映射至正式source\_\allowbreak{}registry。\\
WEB-0100 & A/\allowbreak{}05\_\allowbreak{}全球周期与历史危机 & NOT\_\allowbreak{}ACQUIRED/\allowbreak{}NOT\_\allowbreak{}ACQUIRED & FAILED/\allowbreak{}CORE & asset\_\allowbreak{}monthly\_\allowbreak{}return | case\_\allowbreak{}evidence | historical\_\allowbreak{}crisis\_\allowbreak{}event | global\_\allowbreak{}cy… & 补充权威替代来源；无法取得时删除相应字段和结论承诺。\\
WEB-0101 & A/\allowbreak{}06\_\allowbreak{}黄金美元与储备配置 & ACQUIRED/\allowbreak{}MIXED\_\allowbreak{}CONTENT & EXTRACTED/\allowbreak{}CORE & asset\_\allowbreak{}monthly\_\allowbreak{}return | portfolio\_\allowbreak{}scenario & 人工复核关键字段与source\_\allowbreak{}locator；通过后映射至正式source\_\allowbreak{}registry。\\
WEB-0102 & D/\allowbreak{}05\_\allowbreak{}全球周期与历史危机 & NOT\_\allowbreak{}APPLICABLE/\allowbreak{}NOT\_\allowbreak{}ACQUIRED & FAILED/\allowbreak{}EXCLUDED & case\_\allowbreak{}evidence | historical\_\allowbreak{}crisis\_\allowbreak{}event | global\_\allowbreak{}cycle\_\allowbreak{}month & 仅保留元数据；正式分析改用许可允许的权威替代源。\\
WEB-0103 & D/\allowbreak{}08\_\allowbreak{}企业案例与公开披露 & ACQUIRED/\allowbreak{}DOCUMENT\_\allowbreak{}LOCAL & NOT\_\allowbreak{}REQUIRED/\allowbreak{}EXCLUDED & company\_\allowbreak{}overseas\_\allowbreak{}exposure | country\_\allowbreak{}event | asset\_\allowbreak{}monthly\_\allowbreak{}return | portfol… & 补充权威替代来源；无法取得时删除相应字段和结论承诺。\\
WEB-0104 & D/\allowbreak{}05\_\allowbreak{}全球周期与历史危机 & NOT\_\allowbreak{}APPLICABLE/\allowbreak{}NOT\_\allowbreak{}ACQUIRED & FAILED/\allowbreak{}EXCLUDED & case\_\allowbreak{}evidence | historical\_\allowbreak{}crisis\_\allowbreak{}event | global\_\allowbreak{}cycle\_\allowbreak{}month & 仅保留元数据；正式分析改用许可允许的权威替代源。\\
WEB-0105 & D/\allowbreak{}08\_\allowbreak{}企业案例与公开披露 & NOT\_\allowbreak{}APPLICABLE/\allowbreak{}NOT\_\allowbreak{}ACQUIRED & FAILED/\allowbreak{}EXCLUDED & company\_\allowbreak{}overseas\_\allowbreak{}exposure | country\_\allowbreak{}event | asset\_\allowbreak{}monthly\_\allowbreak{}return | portfol… & 仅保留元数据；正式分析改用许可允许的权威替代源。\\
WEB-0106 & D/\allowbreak{}05\_\allowbreak{}全球周期与历史危机 & NOT\_\allowbreak{}APPLICABLE/\allowbreak{}NOT\_\allowbreak{}ACQUIRED & FAILED/\allowbreak{}EXCLUDED & asset\_\allowbreak{}monthly\_\allowbreak{}return | case\_\allowbreak{}evidence | historical\_\allowbreak{}crisis\_\allowbreak{}event | global\_\allowbreak{}cy… & 仅保留元数据；正式分析改用许可允许的权威替代源。\\
WEB-0107 & D/\allowbreak{}08\_\allowbreak{}企业案例与公开披露 & NOT\_\allowbreak{}APPLICABLE/\allowbreak{}NOT\_\allowbreak{}ACQUIRED & FAILED/\allowbreak{}EXCLUDED & company\_\allowbreak{}overseas\_\allowbreak{}exposure | country\_\allowbreak{}event | asset\_\allowbreak{}monthly\_\allowbreak{}return | portfol… & 仅保留元数据；正式分析改用许可允许的权威替代源。\\
WEB-0108 & A/\allowbreak{}08\_\allowbreak{}企业案例与公开披露 & ACQUIRED/\allowbreak{}DOWNLOAD\_\allowbreak{}LANDING & EXTRACTED/\allowbreak{}CORE & company\_\allowbreak{}overseas\_\allowbreak{}exposure | country\_\allowbreak{}event | asset\_\allowbreak{}monthly\_\allowbreak{}return | portfol… & 人工复核关键字段与source\_\allowbreak{}locator；通过后映射至正式source\_\allowbreak{}registry。\\
WEB-0109 & A/\allowbreak{}08\_\allowbreak{}企业案例与公开披露 & ACQUIRED/\allowbreak{}NARRATIVE\_\allowbreak{}CONTENT & EXTRACTED/\allowbreak{}CORE & company\_\allowbreak{}overseas\_\allowbreak{}exposure | country\_\allowbreak{}event | asset\_\allowbreak{}monthly\_\allowbreak{}return | portfol… & 人工复核关键字段与source\_\allowbreak{}locator；通过后映射至正式source\_\allowbreak{}registry。\\
WEB-0110 & D/\allowbreak{}07\_\allowbreak{}会计合规与制裁 & NOT\_\allowbreak{}APPLICABLE/\allowbreak{}NOT\_\allowbreak{}ACQUIRED & FAILED/\allowbreak{}EXCLUDED & company\_\allowbreak{}overseas\_\allowbreak{}exposure | country\_\allowbreak{}policy\_\allowbreak{}year | country\_\allowbreak{}event | case\_\allowbreak{}evi… & 仅保留元数据；正式分析改用许可允许的权威替代源。\\
WEB-0111 & C/\allowbreak{}08\_\allowbreak{}企业案例与公开披露 & NOT\_\allowbreak{}APPLICABLE/\allowbreak{}NOT\_\allowbreak{}ACQUIRED & FAILED/\allowbreak{}EXCLUDED & company\_\allowbreak{}overseas\_\allowbreak{}exposure | country\_\allowbreak{}event | asset\_\allowbreak{}monthly\_\allowbreak{}return | portfol… & 仅保留元数据；正式分析改用许可允许的权威替代源。\\
WEB-0112 & C/\allowbreak{}08\_\allowbreak{}企业案例与公开披露 & NOT\_\allowbreak{}APPLICABLE/\allowbreak{}NOT\_\allowbreak{}ACQUIRED & FAILED/\allowbreak{}EXCLUDED & company\_\allowbreak{}overseas\_\allowbreak{}exposure | country\_\allowbreak{}event | asset\_\allowbreak{}monthly\_\allowbreak{}return | portfol… & 仅保留元数据；正式分析改用许可允许的权威替代源。\\
WEB-0113 & D/\allowbreak{}08\_\allowbreak{}企业案例与公开披露 & NOT\_\allowbreak{}APPLICABLE/\allowbreak{}NOT\_\allowbreak{}ACQUIRED & FAILED/\allowbreak{}EXCLUDED & company\_\allowbreak{}overseas\_\allowbreak{}exposure | country\_\allowbreak{}event | asset\_\allowbreak{}monthly\_\allowbreak{}return | portfol… & 仅保留元数据；正式分析改用许可允许的权威替代源。\\
WEB-0114 & D/\allowbreak{}05\_\allowbreak{}全球周期与历史危机 & NOT\_\allowbreak{}APPLICABLE/\allowbreak{}NOT\_\allowbreak{}ACQUIRED & FAILED/\allowbreak{}EXCLUDED & asset\_\allowbreak{}monthly\_\allowbreak{}return | case\_\allowbreak{}evidence | historical\_\allowbreak{}crisis\_\allowbreak{}event | global\_\allowbreak{}cy… & 仅保留元数据；正式分析改用许可允许的权威替代源。\\
WEB-0115 & B/\allowbreak{}06\_\allowbreak{}黄金美元与储备配置 & ACQUIRED/\allowbreak{}DOCUMENT\_\allowbreak{}LOCAL & NOT\_\allowbreak{}REQUIRED/\allowbreak{}SUPPLEMENT & asset\_\allowbreak{}monthly\_\allowbreak{}return | portfolio\_\allowbreak{}scenario & 按项目用途复核；不足以单独支持核心结论。\\
WEB-0116 & D/\allowbreak{}05\_\allowbreak{}全球周期与历史危机 & NOT\_\allowbreak{}APPLICABLE/\allowbreak{}NOT\_\allowbreak{}ACQUIRED & FAILED/\allowbreak{}EXCLUDED & asset\_\allowbreak{}monthly\_\allowbreak{}return | case\_\allowbreak{}evidence | historical\_\allowbreak{}crisis\_\allowbreak{}event | global\_\allowbreak{}cy… & 仅保留元数据；正式分析改用许可允许的权威替代源。\\
WEB-0117 & D/\allowbreak{}04\_\allowbreak{}国别货币与中国出海 & NOT\_\allowbreak{}APPLICABLE/\allowbreak{}NOT\_\allowbreak{}ACQUIRED & FAILED/\allowbreak{}EXCLUDED & country\_\allowbreak{}monthly\_\allowbreak{}risk | country\_\allowbreak{}policy\_\allowbreak{}year | country\_\allowbreak{}event | case\_\allowbreak{}evidence & 仅保留元数据；正式分析改用许可允许的权威替代源。\\
WEB-0118 & D/\allowbreak{}06\_\allowbreak{}黄金美元与储备配置 & NOT\_\allowbreak{}APPLICABLE/\allowbreak{}NOT\_\allowbreak{}ACQUIRED & FAILED/\allowbreak{}EXCLUDED & asset\_\allowbreak{}monthly\_\allowbreak{}return | portfolio\_\allowbreak{}scenario & 仅保留元数据；正式分析改用许可允许的权威替代源。\\
WEB-0119 & C/\allowbreak{}08\_\allowbreak{}企业案例与公开披露 & NOT\_\allowbreak{}APPLICABLE/\allowbreak{}NOT\_\allowbreak{}ACQUIRED & FAILED/\allowbreak{}EXCLUDED & company\_\allowbreak{}overseas\_\allowbreak{}exposure | country\_\allowbreak{}event | asset\_\allowbreak{}monthly\_\allowbreak{}return | portfol… & 仅保留元数据；正式分析改用许可允许的权威替代源。\\
WEB-0120 & D/\allowbreak{}06\_\allowbreak{}黄金美元与储备配置 & NOT\_\allowbreak{}APPLICABLE/\allowbreak{}NOT\_\allowbreak{}ACQUIRED & FAILED/\allowbreak{}EXCLUDED & asset\_\allowbreak{}monthly\_\allowbreak{}return | portfolio\_\allowbreak{}scenario & 仅保留元数据；正式分析改用许可允许的权威替代源。\\
WEB-0121 & D/\allowbreak{}06\_\allowbreak{}黄金美元与储备配置 & NOT\_\allowbreak{}APPLICABLE/\allowbreak{}NOT\_\allowbreak{}ACQUIRED & FAILED/\allowbreak{}EXCLUDED & asset\_\allowbreak{}monthly\_\allowbreak{}return | portfolio\_\allowbreak{}scenario & 仅保留元数据；正式分析改用许可允许的权威替代源。\\
WEB-0122 & D/\allowbreak{}06\_\allowbreak{}黄金美元与储备配置 & NOT\_\allowbreak{}APPLICABLE/\allowbreak{}NOT\_\allowbreak{}ACQUIRED & FAILED/\allowbreak{}EXCLUDED & asset\_\allowbreak{}monthly\_\allowbreak{}return | portfolio\_\allowbreak{}scenario & 仅保留元数据；正式分析改用许可允许的权威替代源。\\
WEB-0123 & D/\allowbreak{}05\_\allowbreak{}全球周期与历史危机 & NOT\_\allowbreak{}APPLICABLE/\allowbreak{}NOT\_\allowbreak{}ACQUIRED & FAILED/\allowbreak{}EXCLUDED & asset\_\allowbreak{}monthly\_\allowbreak{}return | case\_\allowbreak{}evidence | historical\_\allowbreak{}crisis\_\allowbreak{}event | global\_\allowbreak{}cy… & 仅保留元数据；正式分析改用许可允许的权威替代源。\\
WEB-0124 & D/\allowbreak{}08\_\allowbreak{}企业案例与公开披露 & NOT\_\allowbreak{}APPLICABLE/\allowbreak{}NOT\_\allowbreak{}ACQUIRED & FAILED/\allowbreak{}EXCLUDED & company\_\allowbreak{}overseas\_\allowbreak{}exposure | country\_\allowbreak{}event | asset\_\allowbreak{}monthly\_\allowbreak{}return | portfol… & 仅保留元数据；正式分析改用许可允许的权威替代源。\\
WEB-0125 & D/\allowbreak{}08\_\allowbreak{}企业案例与公开披露 & NOT\_\allowbreak{}APPLICABLE/\allowbreak{}NOT\_\allowbreak{}ACQUIRED & FAILED/\allowbreak{}EXCLUDED & company\_\allowbreak{}overseas\_\allowbreak{}exposure | country\_\allowbreak{}event | asset\_\allowbreak{}monthly\_\allowbreak{}return | portfol… & 仅保留元数据；正式分析改用许可允许的权威替代源。\\
WEB-0126 & C/\allowbreak{}08\_\allowbreak{}企业案例与公开披露 & NOT\_\allowbreak{}APPLICABLE/\allowbreak{}NOT\_\allowbreak{}ACQUIRED & FAILED/\allowbreak{}EXCLUDED & company\_\allowbreak{}overseas\_\allowbreak{}exposure | country\_\allowbreak{}event | asset\_\allowbreak{}monthly\_\allowbreak{}return | portfol… & 仅保留元数据；正式分析改用许可允许的权威替代源。\\
WEB-0127 & D/\allowbreak{}05\_\allowbreak{}全球周期与历史危机 & NOT\_\allowbreak{}APPLICABLE/\allowbreak{}NOT\_\allowbreak{}ACQUIRED & FAILED/\allowbreak{}EXCLUDED & asset\_\allowbreak{}monthly\_\allowbreak{}return | case\_\allowbreak{}evidence | historical\_\allowbreak{}crisis\_\allowbreak{}event | global\_\allowbreak{}cy… & 仅保留元数据；正式分析改用许可允许的权威替代源。\\
WEB-0128 & A/\allowbreak{}07\_\allowbreak{}会计合规与制裁 & ACQUIRED/\allowbreak{}NARRATIVE\_\allowbreak{}CONTENT & EXTRACTED/\allowbreak{}CORE & source\_\allowbreak{}registry | country\_\allowbreak{}policy\_\allowbreak{}year | case\_\allowbreak{}evidence & 人工复核关键字段与source\_\allowbreak{}locator；通过后映射至正式source\_\allowbreak{}registry。\\
WEB-0129 & A/\allowbreak{}07\_\allowbreak{}会计合规与制裁 & ACQUIRED/\allowbreak{}NARRATIVE\_\allowbreak{}CONTENT & EXTRACTED/\allowbreak{}CORE & country\_\allowbreak{}policy\_\allowbreak{}year | asset\_\allowbreak{}monthly\_\allowbreak{}return & 人工复核关键字段与source\_\allowbreak{}locator；通过后映射至正式source\_\allowbreak{}registry。\\
WEB-0130 & A/\allowbreak{}07\_\allowbreak{}会计合规与制裁 & NOT\_\allowbreak{}APPLICABLE/\allowbreak{}NOT\_\allowbreak{}ACQUIRED & FAILED/\allowbreak{}EXCLUDED & country\_\allowbreak{}policy\_\allowbreak{}year | asset\_\allowbreak{}monthly\_\allowbreak{}return & 仅保留元数据；正式分析改用许可允许的权威替代源。\\
WEB-0131 & D/\allowbreak{}05\_\allowbreak{}全球周期与历史危机 & NOT\_\allowbreak{}APPLICABLE/\allowbreak{}NOT\_\allowbreak{}ACQUIRED & FAILED/\allowbreak{}EXCLUDED & asset\_\allowbreak{}monthly\_\allowbreak{}return | case\_\allowbreak{}evidence | historical\_\allowbreak{}crisis\_\allowbreak{}event | global\_\allowbreak{}cy… & 仅保留元数据；正式分析改用许可允许的权威替代源。\\
WEB-0132 & B/\allowbreak{}05\_\allowbreak{}全球周期与历史危机 & ACQUIRED/\allowbreak{}DOCUMENT\_\allowbreak{}LOCAL & NOT\_\allowbreak{}REQUIRED/\allowbreak{}SUPPLEMENT & country\_\allowbreak{}monthly\_\allowbreak{}risk | asset\_\allowbreak{}monthly\_\allowbreak{}return | case\_\allowbreak{}evidence | historical\_\allowbreak{}c… & 按项目用途复核；不足以单独支持核心结论。\\
WEB-0133 & C/\allowbreak{}04\_\allowbreak{}国别货币与中国出海 & NOT\_\allowbreak{}APPLICABLE/\allowbreak{}NOT\_\allowbreak{}ACQUIRED & FAILED/\allowbreak{}EXCLUDED & country\_\allowbreak{}monthly\_\allowbreak{}risk | country\_\allowbreak{}policy\_\allowbreak{}year | country\_\allowbreak{}event | case\_\allowbreak{}evidence & 仅保留元数据；正式分析改用许可允许的权威替代源。\\
WEB-0134 & C/\allowbreak{}08\_\allowbreak{}企业案例与公开披露 & NOT\_\allowbreak{}APPLICABLE/\allowbreak{}NOT\_\allowbreak{}ACQUIRED & FAILED/\allowbreak{}EXCLUDED & company\_\allowbreak{}overseas\_\allowbreak{}exposure | country\_\allowbreak{}event | asset\_\allowbreak{}monthly\_\allowbreak{}return | portfol… & 仅保留元数据；正式分析改用许可允许的权威替代源。\\
WEB-0135 & D/\allowbreak{}08\_\allowbreak{}企业案例与公开披露 & ACQUIRED/\allowbreak{}DOCUMENT\_\allowbreak{}LOCAL & NOT\_\allowbreak{}REQUIRED/\allowbreak{}EXCLUDED & company\_\allowbreak{}overseas\_\allowbreak{}exposure | country\_\allowbreak{}event | asset\_\allowbreak{}monthly\_\allowbreak{}return | portfol… & 补充权威替代来源；无法取得时删除相应字段和结论承诺。\\
WEB-0136 & D/\allowbreak{}08\_\allowbreak{}企业案例与公开披露 & ACQUIRED/\allowbreak{}DOCUMENT\_\allowbreak{}LOCAL & NOT\_\allowbreak{}REQUIRED/\allowbreak{}EXCLUDED & company\_\allowbreak{}overseas\_\allowbreak{}exposure | country\_\allowbreak{}event | asset\_\allowbreak{}monthly\_\allowbreak{}return | portfol… & 补充权威替代来源；无法取得时删除相应字段和结论承诺。\\
WEB-0137 & D/\allowbreak{}08\_\allowbreak{}企业案例与公开披露 & ACQUIRED/\allowbreak{}DOCUMENT\_\allowbreak{}LOCAL & NOT\_\allowbreak{}REQUIRED/\allowbreak{}EXCLUDED & company\_\allowbreak{}overseas\_\allowbreak{}exposure | country\_\allowbreak{}event | asset\_\allowbreak{}monthly\_\allowbreak{}return | portfol… & 补充权威替代来源；无法取得时删除相应字段和结论承诺。\\
WEB-0138 & D/\allowbreak{}08\_\allowbreak{}企业案例与公开披露 & ACQUIRED/\allowbreak{}DOCUMENT\_\allowbreak{}LOCAL & NOT\_\allowbreak{}REQUIRED/\allowbreak{}EXCLUDED & company\_\allowbreak{}overseas\_\allowbreak{}exposure | country\_\allowbreak{}event | asset\_\allowbreak{}monthly\_\allowbreak{}return | portfol… & 补充权威替代来源；无法取得时删除相应字段和结论承诺。\\
WEB-0139 & A/\allowbreak{}06\_\allowbreak{}黄金美元与储备配置 & ACQUIRED/\allowbreak{}MIXED\_\allowbreak{}CONTENT & EXTRACTED/\allowbreak{}CORE & asset\_\allowbreak{}monthly\_\allowbreak{}return | portfolio\_\allowbreak{}scenario & 人工复核关键字段与source\_\allowbreak{}locator；通过后映射至正式source\_\allowbreak{}registry。\\
WEB-0140 & A/\allowbreak{}06\_\allowbreak{}黄金美元与储备配置 & NOT\_\allowbreak{}ACQUIRED/\allowbreak{}NOT\_\allowbreak{}ACQUIRED & FAILED/\allowbreak{}CORE & asset\_\allowbreak{}monthly\_\allowbreak{}return | portfolio\_\allowbreak{}scenario & 补充权威替代来源；无法取得时删除相应字段和结论承诺。\\
WEB-0141 & D/\allowbreak{}06\_\allowbreak{}黄金美元与储备配置 & NOT\_\allowbreak{}APPLICABLE/\allowbreak{}NOT\_\allowbreak{}ACQUIRED & FAILED/\allowbreak{}EXCLUDED & asset\_\allowbreak{}monthly\_\allowbreak{}return | portfolio\_\allowbreak{}scenario & 仅保留元数据；正式分析改用许可允许的权威替代源。\\
WEB-0142 & D/\allowbreak{}06\_\allowbreak{}黄金美元与储备配置 & NOT\_\allowbreak{}APPLICABLE/\allowbreak{}NOT\_\allowbreak{}ACQUIRED & FAILED/\allowbreak{}EXCLUDED & asset\_\allowbreak{}monthly\_\allowbreak{}return | portfolio\_\allowbreak{}scenario & 仅保留元数据；正式分析改用许可允许的权威替代源。\\
WEB-0143 & D/\allowbreak{}08\_\allowbreak{}企业案例与公开披露 & ACQUIRED/\allowbreak{}DOCUMENT\_\allowbreak{}LOCAL & NOT\_\allowbreak{}REQUIRED/\allowbreak{}EXCLUDED & company\_\allowbreak{}overseas\_\allowbreak{}exposure | country\_\allowbreak{}event | asset\_\allowbreak{}monthly\_\allowbreak{}return | portfol… & 补充权威替代来源；无法取得时删除相应字段和结论承诺。\\
WEB-0144 & A/\allowbreak{}04\_\allowbreak{}国别货币与中国出海 & NOT\_\allowbreak{}APPLICABLE/\allowbreak{}NOT\_\allowbreak{}ACQUIRED & FAILED/\allowbreak{}EXCLUDED & country\_\allowbreak{}monthly\_\allowbreak{}risk | country\_\allowbreak{}policy\_\allowbreak{}year | country\_\allowbreak{}event | case\_\allowbreak{}evidence & 仅保留元数据；正式分析改用许可允许的权威替代源。\\
WEB-0145 & B/\allowbreak{}06\_\allowbreak{}黄金美元与储备配置 & NOT\_\allowbreak{}ACQUIRED/\allowbreak{}NOT\_\allowbreak{}ACQUIRED & FAILED/\allowbreak{}SUPPLEMENT & asset\_\allowbreak{}monthly\_\allowbreak{}return | portfolio\_\allowbreak{}scenario & 补充权威替代来源；无法取得时删除相应字段和结论承诺。\\
WEB-0146 & D/\allowbreak{}07\_\allowbreak{}会计合规与制裁 & ACQUIRED/\allowbreak{}DOCUMENT\_\allowbreak{}LOCAL & NOT\_\allowbreak{}REQUIRED/\allowbreak{}EXCLUDED & country\_\allowbreak{}policy\_\allowbreak{}year | asset\_\allowbreak{}monthly\_\allowbreak{}return & 补充权威替代来源；无法取得时删除相应字段和结论承诺。\\
WEB-0147 & D/\allowbreak{}06\_\allowbreak{}黄金美元与储备配置 & NOT\_\allowbreak{}APPLICABLE/\allowbreak{}NOT\_\allowbreak{}ACQUIRED & FAILED/\allowbreak{}EXCLUDED & asset\_\allowbreak{}monthly\_\allowbreak{}return | portfolio\_\allowbreak{}scenario & 仅保留元数据；正式分析改用许可允许的权威替代源。\\
WEB-0148 & D/\allowbreak{}06\_\allowbreak{}黄金美元与储备配置 & NOT\_\allowbreak{}APPLICABLE/\allowbreak{}NOT\_\allowbreak{}ACQUIRED & FAILED/\allowbreak{}EXCLUDED & asset\_\allowbreak{}monthly\_\allowbreak{}return | portfolio\_\allowbreak{}scenario & 仅保留元数据；正式分析改用许可允许的权威替代源。\\
WEB-0149 & D/\allowbreak{}05\_\allowbreak{}全球周期与历史危机 & ACQUIRED/\allowbreak{}DOCUMENT\_\allowbreak{}LOCAL & NOT\_\allowbreak{}REQUIRED/\allowbreak{}EXCLUDED & asset\_\allowbreak{}monthly\_\allowbreak{}return | case\_\allowbreak{}evidence | historical\_\allowbreak{}crisis\_\allowbreak{}event | global\_\allowbreak{}cy… & 补充权威替代来源；无法取得时删除相应字段和结论承诺。\\
WEB-0150 & D/\allowbreak{}05\_\allowbreak{}全球周期与历史危机 & NOT\_\allowbreak{}APPLICABLE/\allowbreak{}NOT\_\allowbreak{}ACQUIRED & FAILED/\allowbreak{}EXCLUDED & asset\_\allowbreak{}monthly\_\allowbreak{}return | case\_\allowbreak{}evidence | historical\_\allowbreak{}crisis\_\allowbreak{}event | global\_\allowbreak{}cy… & 仅保留元数据；正式分析改用许可允许的权威替代源。\\
WEB-0151 & D/\allowbreak{}06\_\allowbreak{}黄金美元与储备配置 & NOT\_\allowbreak{}APPLICABLE/\allowbreak{}NOT\_\allowbreak{}ACQUIRED & FAILED/\allowbreak{}EXCLUDED & asset\_\allowbreak{}monthly\_\allowbreak{}return | portfolio\_\allowbreak{}scenario & 按人工任务卡接受条款或完成机构授权后重新验收。\\
WEB-0152 & D/\allowbreak{}06\_\allowbreak{}黄金美元与储备配置 & NOT\_\allowbreak{}APPLICABLE/\allowbreak{}NOT\_\allowbreak{}ACQUIRED & FAILED/\allowbreak{}EXCLUDED & asset\_\allowbreak{}monthly\_\allowbreak{}return | portfolio\_\allowbreak{}scenario & 按人工任务卡接受条款或完成机构授权后重新验收。\\
WEB-0153 & D/\allowbreak{}06\_\allowbreak{}黄金美元与储备配置 & NOT\_\allowbreak{}APPLICABLE/\allowbreak{}NOT\_\allowbreak{}ACQUIRED & FAILED/\allowbreak{}EXCLUDED & asset\_\allowbreak{}monthly\_\allowbreak{}return | portfolio\_\allowbreak{}scenario & 按人工任务卡接受条款或完成机构授权后重新验收。\\
WEB-0154 & D/\allowbreak{}06\_\allowbreak{}黄金美元与储备配置 & NOT\_\allowbreak{}APPLICABLE/\allowbreak{}NOT\_\allowbreak{}ACQUIRED & FAILED/\allowbreak{}EXCLUDED & asset\_\allowbreak{}monthly\_\allowbreak{}return | portfolio\_\allowbreak{}scenario & 按人工任务卡接受条款或完成机构授权后重新验收。\\
WEB-0155 & D/\allowbreak{}06\_\allowbreak{}黄金美元与储备配置 & NOT\_\allowbreak{}APPLICABLE/\allowbreak{}NOT\_\allowbreak{}ACQUIRED & FAILED/\allowbreak{}EXCLUDED & asset\_\allowbreak{}monthly\_\allowbreak{}return | portfolio\_\allowbreak{}scenario & 按人工任务卡接受条款或完成机构授权后重新验收。\\
WEB-0156 & D/\allowbreak{}08\_\allowbreak{}企业案例与公开披露 & ACQUIRED/\allowbreak{}DOCUMENT\_\allowbreak{}LOCAL & NOT\_\allowbreak{}REQUIRED/\allowbreak{}EXCLUDED & company\_\allowbreak{}overseas\_\allowbreak{}exposure | country\_\allowbreak{}event | asset\_\allowbreak{}monthly\_\allowbreak{}return | portfol… & 补充权威替代来源；无法取得时删除相应字段和结论承诺。\\
WEB-0157 & D/\allowbreak{}08\_\allowbreak{}企业案例与公开披露 & NOT\_\allowbreak{}APPLICABLE/\allowbreak{}NOT\_\allowbreak{}ACQUIRED & FAILED/\allowbreak{}EXCLUDED & company\_\allowbreak{}overseas\_\allowbreak{}exposure | country\_\allowbreak{}event | asset\_\allowbreak{}monthly\_\allowbreak{}return | portfol… & 仅保留元数据；正式分析改用许可允许的权威替代源。\\
WEB-0158 & B/\allowbreak{}06\_\allowbreak{}黄金美元与储备配置 & ACQUIRED/\allowbreak{}DOCUMENT\_\allowbreak{}LOCAL & NOT\_\allowbreak{}REQUIRED/\allowbreak{}SUPPLEMENT & asset\_\allowbreak{}monthly\_\allowbreak{}return | portfolio\_\allowbreak{}scenario & 按项目用途复核；不足以单独支持核心结论。\\
WEB-0159 & B/\allowbreak{}06\_\allowbreak{}黄金美元与储备配置 & ACQUIRED/\allowbreak{}NARRATIVE\_\allowbreak{}CONTENT & EXTRACTED/\allowbreak{}SUPPLEMENT & country\_\allowbreak{}exposure | asset\_\allowbreak{}monthly\_\allowbreak{}return | portfolio\_\allowbreak{}scenario & 按项目用途复核；不足以单独支持核心结论。\\
WEB-0160 & D/\allowbreak{}09\_\allowbreak{}Smartbi与比赛平台 & NOT\_\allowbreak{}APPLICABLE/\allowbreak{}NOT\_\allowbreak{}ACQUIRED & FAILED/\allowbreak{}EXCLUDED & source\_\allowbreak{}registry & 由资料负责人合法登录、采购或申请授权；不得自动绕过。\\
WEB-0161 & D/\allowbreak{}09\_\allowbreak{}Smartbi与比赛平台 & NOT\_\allowbreak{}APPLICABLE/\allowbreak{}NOT\_\allowbreak{}ACQUIRED & FAILED/\allowbreak{}EXCLUDED & source\_\allowbreak{}registry & 由资料负责人合法登录、采购或申请授权；不得自动绕过。\\
WEB-0162 & A/\allowbreak{}05\_\allowbreak{}全球周期与历史危机 & ACQUIRED/\allowbreak{}NARRATIVE\_\allowbreak{}CONTENT & EXTRACTED/\allowbreak{}CORE & asset\_\allowbreak{}monthly\_\allowbreak{}return | case\_\allowbreak{}evidence | historical\_\allowbreak{}crisis\_\allowbreak{}event | global\_\allowbreak{}cy… & 人工复核关键字段与source\_\allowbreak{}locator；通过后映射至正式source\_\allowbreak{}registry。\\
WEB-0163 & D/\allowbreak{}05\_\allowbreak{}全球周期与历史危机 & NOT\_\allowbreak{}APPLICABLE/\allowbreak{}NOT\_\allowbreak{}ACQUIRED & FAILED/\allowbreak{}EXCLUDED & case\_\allowbreak{}evidence | historical\_\allowbreak{}crisis\_\allowbreak{}event | global\_\allowbreak{}cycle\_\allowbreak{}month & 仅保留元数据；正式分析改用许可允许的权威替代源。\\
WEB-0164 & D/\allowbreak{}08\_\allowbreak{}企业案例与公开披露 & NOT\_\allowbreak{}APPLICABLE/\allowbreak{}NOT\_\allowbreak{}ACQUIRED & FAILED/\allowbreak{}EXCLUDED & company\_\allowbreak{}overseas\_\allowbreak{}exposure | country\_\allowbreak{}event | asset\_\allowbreak{}monthly\_\allowbreak{}return | portfol… & 仅保留元数据；正式分析改用许可允许的权威替代源。\\
WEB-0165 & D/\allowbreak{}06\_\allowbreak{}黄金美元与储备配置 & ACQUIRED/\allowbreak{}DOCUMENT\_\allowbreak{}LOCAL & NOT\_\allowbreak{}REQUIRED/\allowbreak{}EXCLUDED & source\_\allowbreak{}registry | asset\_\allowbreak{}monthly\_\allowbreak{}return | portfolio\_\allowbreak{}scenario & 补充权威替代来源；无法取得时删除相应字段和结论承诺。\\
WEB-0166 & D/\allowbreak{}06\_\allowbreak{}黄金美元与储备配置 & ACQUIRED/\allowbreak{}DOCUMENT\_\allowbreak{}LOCAL & NOT\_\allowbreak{}REQUIRED/\allowbreak{}EXCLUDED & asset\_\allowbreak{}monthly\_\allowbreak{}return | portfolio\_\allowbreak{}scenario & 补充权威替代来源；无法取得时删除相应字段和结论承诺。\\
WEB-0167 & D/\allowbreak{}06\_\allowbreak{}黄金美元与储备配置 & NOT\_\allowbreak{}APPLICABLE/\allowbreak{}NOT\_\allowbreak{}ACQUIRED & FAILED/\allowbreak{}EXCLUDED & asset\_\allowbreak{}monthly\_\allowbreak{}return | portfolio\_\allowbreak{}scenario & 仅保留元数据；正式分析改用许可允许的权威替代源。\\
WEB-0168 & D/\allowbreak{}06\_\allowbreak{}黄金美元与储备配置 & NOT\_\allowbreak{}APPLICABLE/\allowbreak{}NOT\_\allowbreak{}ACQUIRED & FAILED/\allowbreak{}EXCLUDED & country\_\allowbreak{}monthly\_\allowbreak{}risk | country\_\allowbreak{}event | asset\_\allowbreak{}monthly\_\allowbreak{}return | portfolio\_\allowbreak{}sc… & 仅保留元数据；正式分析改用许可允许的权威替代源。\\
WEB-0169 & D/\allowbreak{}05\_\allowbreak{}全球周期与历史危机 & NOT\_\allowbreak{}APPLICABLE/\allowbreak{}NOT\_\allowbreak{}ACQUIRED & FAILED/\allowbreak{}EXCLUDED & asset\_\allowbreak{}monthly\_\allowbreak{}return | case\_\allowbreak{}evidence | historical\_\allowbreak{}crisis\_\allowbreak{}event | global\_\allowbreak{}cy… & 由资料负责人合法登录、采购或申请授权；不得自动绕过。\\
WEB-0170 & D/\allowbreak{}08\_\allowbreak{}企业案例与公开披露 & NOT\_\allowbreak{}APPLICABLE/\allowbreak{}NOT\_\allowbreak{}ACQUIRED & FAILED/\allowbreak{}EXCLUDED & company\_\allowbreak{}overseas\_\allowbreak{}exposure | country\_\allowbreak{}event | asset\_\allowbreak{}monthly\_\allowbreak{}return | portfol… & 仅保留元数据；正式分析改用许可允许的权威替代源。\\
WEB-0171 & D/\allowbreak{}04\_\allowbreak{}国别货币与中国出海 & NOT\_\allowbreak{}APPLICABLE/\allowbreak{}NOT\_\allowbreak{}ACQUIRED & FAILED/\allowbreak{}EXCLUDED & country\_\allowbreak{}monthly\_\allowbreak{}risk | country\_\allowbreak{}policy\_\allowbreak{}year | country\_\allowbreak{}event | case\_\allowbreak{}evidence & 仅保留元数据；正式分析改用许可允许的权威替代源。\\
WEB-0172 & D/\allowbreak{}05\_\allowbreak{}全球周期与历史危机 & NOT\_\allowbreak{}APPLICABLE/\allowbreak{}NOT\_\allowbreak{}ACQUIRED & FAILED/\allowbreak{}EXCLUDED & case\_\allowbreak{}evidence | historical\_\allowbreak{}crisis\_\allowbreak{}event | global\_\allowbreak{}cycle\_\allowbreak{}month & 仅保留元数据；正式分析改用许可允许的权威替代源。\\
WEB-0173 & D/\allowbreak{}07\_\allowbreak{}会计合规与制裁 & NOT\_\allowbreak{}APPLICABLE/\allowbreak{}NOT\_\allowbreak{}ACQUIRED & FAILED/\allowbreak{}EXCLUDED & country\_\allowbreak{}policy\_\allowbreak{}year | asset\_\allowbreak{}monthly\_\allowbreak{}return & 仅保留元数据；正式分析改用许可允许的权威替代源。\\
WEB-0174 & D/\allowbreak{}05\_\allowbreak{}全球周期与历史危机 & NOT\_\allowbreak{}APPLICABLE/\allowbreak{}NOT\_\allowbreak{}ACQUIRED & FAILED/\allowbreak{}EXCLUDED & country\_\allowbreak{}monthly\_\allowbreak{}risk | asset\_\allowbreak{}monthly\_\allowbreak{}return | case\_\allowbreak{}evidence | historical\_\allowbreak{}c… & 仅保留元数据；正式分析改用许可允许的权威替代源。\\
WEB-0175 & C/\allowbreak{}07\_\allowbreak{}会计合规与制裁 & NOT\_\allowbreak{}APPLICABLE/\allowbreak{}NOT\_\allowbreak{}ACQUIRED & FAILED/\allowbreak{}EXCLUDED & country\_\allowbreak{}monthly\_\allowbreak{}risk | country\_\allowbreak{}policy\_\allowbreak{}year | country\_\allowbreak{}event | case\_\allowbreak{}evidence & 仅保留元数据；正式分析改用许可允许的权威替代源。\\
WEB-0176 & C/\allowbreak{}04\_\allowbreak{}国别货币与中国出海 & NOT\_\allowbreak{}APPLICABLE/\allowbreak{}NOT\_\allowbreak{}ACQUIRED & FAILED/\allowbreak{}EXCLUDED & country\_\allowbreak{}monthly\_\allowbreak{}risk | country\_\allowbreak{}policy\_\allowbreak{}year | country\_\allowbreak{}event | asset\_\allowbreak{}monthly… & 仅保留元数据；正式分析改用许可允许的权威替代源。\\
WEB-0177 & C/\allowbreak{}08\_\allowbreak{}企业案例与公开披露 & NOT\_\allowbreak{}APPLICABLE/\allowbreak{}NOT\_\allowbreak{}ACQUIRED & FAILED/\allowbreak{}EXCLUDED & company\_\allowbreak{}overseas\_\allowbreak{}exposure | country\_\allowbreak{}monthly\_\allowbreak{}risk | country\_\allowbreak{}event | asset\_\allowbreak{}m… & 仅保留元数据；正式分析改用许可允许的权威替代源。\\
WEB-0178 & C/\allowbreak{}07\_\allowbreak{}会计合规与制裁 & NOT\_\allowbreak{}APPLICABLE/\allowbreak{}NOT\_\allowbreak{}ACQUIRED & FAILED/\allowbreak{}EXCLUDED & company\_\allowbreak{}overseas\_\allowbreak{}exposure | country\_\allowbreak{}monthly\_\allowbreak{}risk | country\_\allowbreak{}policy\_\allowbreak{}year | c… & 仅保留元数据；正式分析改用许可允许的权威替代源。\\
WEB-0179 & C/\allowbreak{}06\_\allowbreak{}黄金美元与储备配置 & NOT\_\allowbreak{}APPLICABLE/\allowbreak{}NOT\_\allowbreak{}ACQUIRED & FAILED/\allowbreak{}EXCLUDED & country\_\allowbreak{}monthly\_\allowbreak{}risk | country\_\allowbreak{}event | asset\_\allowbreak{}monthly\_\allowbreak{}return | portfolio\_\allowbreak{}sc… & 仅保留元数据；正式分析改用许可允许的权威替代源。\\
WEB-0180 & C/\allowbreak{}06\_\allowbreak{}黄金美元与储备配置 & NOT\_\allowbreak{}APPLICABLE/\allowbreak{}NOT\_\allowbreak{}ACQUIRED & FAILED/\allowbreak{}EXCLUDED & country\_\allowbreak{}monthly\_\allowbreak{}risk | country\_\allowbreak{}event | asset\_\allowbreak{}monthly\_\allowbreak{}return | portfolio\_\allowbreak{}sc… & 仅保留元数据；正式分析改用许可允许的权威替代源。\\
WEB-0181 & A/\allowbreak{}05\_\allowbreak{}全球周期与历史危机 & NOT\_\allowbreak{}APPLICABLE/\allowbreak{}NOT\_\allowbreak{}ACQUIRED & FAILED/\allowbreak{}EXCLUDED & country\_\allowbreak{}monthly\_\allowbreak{}risk | country\_\allowbreak{}event | case\_\allowbreak{}evidence | historical\_\allowbreak{}crisis\_\allowbreak{}e… & 仅保留元数据；正式分析改用许可允许的权威替代源。\\
WEB-0182 & D/\allowbreak{}04\_\allowbreak{}国别货币与中国出海 & NOT\_\allowbreak{}APPLICABLE/\allowbreak{}NOT\_\allowbreak{}ACQUIRED & FAILED/\allowbreak{}EXCLUDED & country\_\allowbreak{}monthly\_\allowbreak{}risk | country\_\allowbreak{}policy\_\allowbreak{}year & 仅保留元数据；正式分析改用许可允许的权威替代源。\\
WEB-0183 & D/\allowbreak{}06\_\allowbreak{}黄金美元与储备配置 & NOT\_\allowbreak{}APPLICABLE/\allowbreak{}NOT\_\allowbreak{}ACQUIRED & FAILED/\allowbreak{}EXCLUDED & source\_\allowbreak{}registry | asset\_\allowbreak{}monthly\_\allowbreak{}return | portfolio\_\allowbreak{}scenario & 仅保留元数据；正式分析改用许可允许的权威替代源。\\
WEB-0184 & D/\allowbreak{}07\_\allowbreak{}会计合规与制裁 & NOT\_\allowbreak{}APPLICABLE/\allowbreak{}NOT\_\allowbreak{}ACQUIRED & FAILED/\allowbreak{}EXCLUDED & country\_\allowbreak{}policy\_\allowbreak{}year | asset\_\allowbreak{}monthly\_\allowbreak{}return & 仅保留元数据；正式分析改用许可允许的权威替代源。\\
WEB-0185 & D/\allowbreak{}08\_\allowbreak{}企业案例与公开披露 & NOT\_\allowbreak{}APPLICABLE/\allowbreak{}NOT\_\allowbreak{}ACQUIRED & FAILED/\allowbreak{}EXCLUDED & company\_\allowbreak{}overseas\_\allowbreak{}exposure | country\_\allowbreak{}event | portfolio\_\allowbreak{}scenario | case\_\allowbreak{}evid… & 仅保留元数据；正式分析改用许可允许的权威替代源。\\
WEB-0186 & D/\allowbreak{}08\_\allowbreak{}企业案例与公开披露 & ACQUIRED/\allowbreak{}DOCUMENT\_\allowbreak{}LOCAL & NOT\_\allowbreak{}REQUIRED/\allowbreak{}EXCLUDED & company\_\allowbreak{}overseas\_\allowbreak{}exposure | country\_\allowbreak{}event | asset\_\allowbreak{}monthly\_\allowbreak{}return | portfol… & 补充权威替代来源；无法取得时删除相应字段和结论承诺。\\
WEB-0187 & B/\allowbreak{}06\_\allowbreak{}黄金美元与储备配置 & ACQUIRED/\allowbreak{}NARRATIVE\_\allowbreak{}CONTENT & EXTRACTED/\allowbreak{}SUPPLEMENT & asset\_\allowbreak{}monthly\_\allowbreak{}return | portfolio\_\allowbreak{}scenario & 按项目用途复核；不足以单独支持核心结论。\\
WEB-0188 & B/\allowbreak{}08\_\allowbreak{}企业案例与公开披露 & ACQUIRED/\allowbreak{}DOCUMENT\_\allowbreak{}LOCAL & NOT\_\allowbreak{}REQUIRED/\allowbreak{}SUPPLEMENT & company\_\allowbreak{}overseas\_\allowbreak{}exposure | country\_\allowbreak{}event | portfolio\_\allowbreak{}scenario | case\_\allowbreak{}evid… & 按项目用途复核；不足以单独支持核心结论。\\
WEB-0189 & B/\allowbreak{}04\_\allowbreak{}国别货币与中国出海 & ACQUIRED/\allowbreak{}DOCUMENT\_\allowbreak{}LOCAL & NOT\_\allowbreak{}REQUIRED/\allowbreak{}SUPPLEMENT & country\_\allowbreak{}monthly\_\allowbreak{}risk | country\_\allowbreak{}policy\_\allowbreak{}year & 按项目用途复核；不足以单独支持核心结论。\\
WEB-0190 & B/\allowbreak{}06\_\allowbreak{}黄金美元与储备配置 & ACQUIRED/\allowbreak{}DOCUMENT\_\allowbreak{}LOCAL & NOT\_\allowbreak{}REQUIRED/\allowbreak{}SUPPLEMENT & asset\_\allowbreak{}monthly\_\allowbreak{}return | portfolio\_\allowbreak{}scenario & 按项目用途复核；不足以单独支持核心结论。\\
WEB-0191 & B/\allowbreak{}06\_\allowbreak{}黄金美元与储备配置 & ACQUIRED/\allowbreak{}NARRATIVE\_\allowbreak{}CONTENT & EXTRACTED/\allowbreak{}SUPPLEMENT & asset\_\allowbreak{}monthly\_\allowbreak{}return | portfolio\_\allowbreak{}scenario & 按项目用途复核；不足以单独支持核心结论。\\
WEB-0192 & B/\allowbreak{}04\_\allowbreak{}国别货币与中国出海 & ACQUIRED/\allowbreak{}DOCUMENT\_\allowbreak{}LOCAL & NOT\_\allowbreak{}REQUIRED/\allowbreak{}SUPPLEMENT & source\_\allowbreak{}registry | country\_\allowbreak{}monthly\_\allowbreak{}risk | country\_\allowbreak{}policy\_\allowbreak{}year & 按项目用途复核；不足以单独支持核心结论。\\
WEB-0193 & D/\allowbreak{}06\_\allowbreak{}黄金美元与储备配置 & NOT\_\allowbreak{}APPLICABLE/\allowbreak{}NOT\_\allowbreak{}ACQUIRED & FAILED/\allowbreak{}EXCLUDED & asset\_\allowbreak{}monthly\_\allowbreak{}return | portfolio\_\allowbreak{}scenario & 仅保留元数据；正式分析改用许可允许的权威替代源。\\
WEB-0194 & A/\allowbreak{}05\_\allowbreak{}全球周期与历史危机 & ACQUIRED/\allowbreak{}NARRATIVE\_\allowbreak{}CONTENT & EXTRACTED/\allowbreak{}CORE & country\_\allowbreak{}monthly\_\allowbreak{}risk | asset\_\allowbreak{}monthly\_\allowbreak{}return | case\_\allowbreak{}evidence | historical\_\allowbreak{}c… & 人工复核关键字段与source\_\allowbreak{}locator；通过后映射至正式source\_\allowbreak{}registry。\\
WEB-0195 & A/\allowbreak{}06\_\allowbreak{}黄金美元与储备配置 & ACQUIRED/\allowbreak{}MIXED\_\allowbreak{}CONTENT & EXTRACTED/\allowbreak{}CORE & asset\_\allowbreak{}monthly\_\allowbreak{}return | portfolio\_\allowbreak{}scenario & 人工复核关键字段与source\_\allowbreak{}locator；通过后映射至正式source\_\allowbreak{}registry。\\
WEB-0196 & D/\allowbreak{}05\_\allowbreak{}全球周期与历史危机 & NOT\_\allowbreak{}APPLICABLE/\allowbreak{}NOT\_\allowbreak{}ACQUIRED & FAILED/\allowbreak{}EXCLUDED & case\_\allowbreak{}evidence | historical\_\allowbreak{}crisis\_\allowbreak{}event | global\_\allowbreak{}cycle\_\allowbreak{}month & 仅保留元数据；正式分析改用许可允许的权威替代源。\\
WEB-0197 & B/\allowbreak{}05\_\allowbreak{}全球周期与历史危机 & ACQUIRED/\allowbreak{}NARRATIVE\_\allowbreak{}CONTENT & EXTRACTED/\allowbreak{}SUPPLEMENT & source\_\allowbreak{}registry | historical\_\allowbreak{}crisis\_\allowbreak{}event | global\_\allowbreak{}cycle\_\allowbreak{}month & 按项目用途复核；不足以单独支持核心结论。\\
WEB-0198 & B/\allowbreak{}05\_\allowbreak{}全球周期与历史危机 & ACQUIRED/\allowbreak{}NARRATIVE\_\allowbreak{}CONTENT & EXTRACTED/\allowbreak{}SUPPLEMENT & source\_\allowbreak{}registry | historical\_\allowbreak{}crisis\_\allowbreak{}event | global\_\allowbreak{}cycle\_\allowbreak{}month & 按项目用途复核；不足以单独支持核心结论。\\
WEB-0199 & B/\allowbreak{}05\_\allowbreak{}全球周期与历史危机 & ACQUIRED/\allowbreak{}NARRATIVE\_\allowbreak{}CONTENT & EXTRACTED/\allowbreak{}SUPPLEMENT & source\_\allowbreak{}registry | historical\_\allowbreak{}crisis\_\allowbreak{}event | global\_\allowbreak{}cycle\_\allowbreak{}month & 按项目用途复核；不足以单独支持核心结论。\\
WEB-0200 & B/\allowbreak{}05\_\allowbreak{}全球周期与历史危机 & ACQUIRED/\allowbreak{}MIXED\_\allowbreak{}CONTENT & EXTRACTED/\allowbreak{}SUPPLEMENT & case\_\allowbreak{}evidence | historical\_\allowbreak{}crisis\_\allowbreak{}event | global\_\allowbreak{}cycle\_\allowbreak{}month & 按项目用途复核；不足以单独支持核心结论。\\
WEB-0201 & D/\allowbreak{}04\_\allowbreak{}国别货币与中国出海 & NOT\_\allowbreak{}APPLICABLE/\allowbreak{}NOT\_\allowbreak{}ACQUIRED & FAILED/\allowbreak{}EXCLUDED & country\_\allowbreak{}monthly\_\allowbreak{}risk | country\_\allowbreak{}policy\_\allowbreak{}year | country\_\allowbreak{}event | case\_\allowbreak{}evidence & 仅保留元数据；正式分析改用许可允许的权威替代源。\\
WEB-0202 & D/\allowbreak{}05\_\allowbreak{}全球周期与历史危机 & NOT\_\allowbreak{}APPLICABLE/\allowbreak{}NOT\_\allowbreak{}ACQUIRED & FAILED/\allowbreak{}EXCLUDED & case\_\allowbreak{}evidence | historical\_\allowbreak{}crisis\_\allowbreak{}event | global\_\allowbreak{}cycle\_\allowbreak{}month & 仅保留元数据；正式分析改用许可允许的权威替代源。\\
WEB-0203 & D/\allowbreak{}05\_\allowbreak{}全球周期与历史危机 & NOT\_\allowbreak{}APPLICABLE/\allowbreak{}NOT\_\allowbreak{}ACQUIRED & FAILED/\allowbreak{}EXCLUDED & case\_\allowbreak{}evidence | historical\_\allowbreak{}crisis\_\allowbreak{}event | global\_\allowbreak{}cycle\_\allowbreak{}month & 仅保留元数据；正式分析改用许可允许的权威替代源。\\
WEB-0204 & D/\allowbreak{}05\_\allowbreak{}全球周期与历史危机 & NOT\_\allowbreak{}APPLICABLE/\allowbreak{}NOT\_\allowbreak{}ACQUIRED & FAILED/\allowbreak{}EXCLUDED & asset\_\allowbreak{}monthly\_\allowbreak{}return | case\_\allowbreak{}evidence | historical\_\allowbreak{}crisis\_\allowbreak{}event | global\_\allowbreak{}cy… & 仅保留元数据；正式分析改用许可允许的权威替代源。\\
WEB-0205 & D/\allowbreak{}05\_\allowbreak{}全球周期与历史危机 & NOT\_\allowbreak{}APPLICABLE/\allowbreak{}NOT\_\allowbreak{}ACQUIRED & FAILED/\allowbreak{}EXCLUDED & case\_\allowbreak{}evidence | historical\_\allowbreak{}crisis\_\allowbreak{}event | global\_\allowbreak{}cycle\_\allowbreak{}month & 仅保留元数据；正式分析改用许可允许的权威替代源。\\
WEB-0206 & D/\allowbreak{}05\_\allowbreak{}全球周期与历史危机 & NOT\_\allowbreak{}APPLICABLE/\allowbreak{}NOT\_\allowbreak{}ACQUIRED & FAILED/\allowbreak{}EXCLUDED & case\_\allowbreak{}evidence | historical\_\allowbreak{}crisis\_\allowbreak{}event | global\_\allowbreak{}cycle\_\allowbreak{}month & 仅保留元数据；正式分析改用许可允许的权威替代源。\\
WEB-0207 & D/\allowbreak{}08\_\allowbreak{}企业案例与公开披露 & ACQUIRED/\allowbreak{}DOCUMENT\_\allowbreak{}LOCAL & NOT\_\allowbreak{}REQUIRED/\allowbreak{}EXCLUDED & company\_\allowbreak{}overseas\_\allowbreak{}exposure | country\_\allowbreak{}event | asset\_\allowbreak{}monthly\_\allowbreak{}return | portfol… & 补充权威替代来源；无法取得时删除相应字段和结论承诺。\\
WEB-0208 & A/\allowbreak{}04\_\allowbreak{}国别货币与中国出海 & NOT\_\allowbreak{}APPLICABLE/\allowbreak{}NOT\_\allowbreak{}ACQUIRED & FAILED/\allowbreak{}EXCLUDED & country\_\allowbreak{}monthly\_\allowbreak{}risk | country\_\allowbreak{}policy\_\allowbreak{}year | country\_\allowbreak{}event | case\_\allowbreak{}evidence & 仅保留元数据；正式分析改用许可允许的权威替代源。\\
WEB-0209 & D/\allowbreak{}04\_\allowbreak{}国别货币与中国出海 & NOT\_\allowbreak{}APPLICABLE/\allowbreak{}NOT\_\allowbreak{}ACQUIRED & FAILED/\allowbreak{}EXCLUDED & country\_\allowbreak{}monthly\_\allowbreak{}risk | country\_\allowbreak{}policy\_\allowbreak{}year | country\_\allowbreak{}event | case\_\allowbreak{}evidence & 由资料负责人合法登录、采购或申请授权；不得自动绕过。\\
WEB-0210 & D/\allowbreak{}04\_\allowbreak{}国别货币与中国出海 & NOT\_\allowbreak{}APPLICABLE/\allowbreak{}NOT\_\allowbreak{}ACQUIRED & FAILED/\allowbreak{}EXCLUDED & country\_\allowbreak{}monthly\_\allowbreak{}risk | country\_\allowbreak{}policy\_\allowbreak{}year | country\_\allowbreak{}event | case\_\allowbreak{}evidence & 仅保留元数据；正式分析改用许可允许的权威替代源。\\
WEB-0211 & A/\allowbreak{}08\_\allowbreak{}企业案例与公开披露 & NOT\_\allowbreak{}APPLICABLE/\allowbreak{}NOT\_\allowbreak{}ACQUIRED & FAILED/\allowbreak{}EXCLUDED & company\_\allowbreak{}overseas\_\allowbreak{}exposure | country\_\allowbreak{}monthly\_\allowbreak{}risk | country\_\allowbreak{}event | portfol… & 仅保留元数据；正式分析改用许可允许的权威替代源。\\
WEB-0212 & D/\allowbreak{}06\_\allowbreak{}黄金美元与储备配置 & NOT\_\allowbreak{}APPLICABLE/\allowbreak{}NOT\_\allowbreak{}ACQUIRED & FAILED/\allowbreak{}EXCLUDED & asset\_\allowbreak{}monthly\_\allowbreak{}return | portfolio\_\allowbreak{}scenario & 仅保留元数据；正式分析改用许可允许的权威替代源。\\
WEB-0213 & D/\allowbreak{}07\_\allowbreak{}会计合规与制裁 & NOT\_\allowbreak{}APPLICABLE/\allowbreak{}NOT\_\allowbreak{}ACQUIRED & FAILED/\allowbreak{}EXCLUDED & country\_\allowbreak{}policy\_\allowbreak{}year | asset\_\allowbreak{}monthly\_\allowbreak{}return | case\_\allowbreak{}evidence & 仅保留元数据；正式分析改用许可允许的权威替代源。\\
WEB-0214 & A/\allowbreak{}06\_\allowbreak{}黄金美元与储备配置 & NOT\_\allowbreak{}APPLICABLE/\allowbreak{}NOT\_\allowbreak{}ACQUIRED & FAILED/\allowbreak{}EXCLUDED & asset\_\allowbreak{}monthly\_\allowbreak{}return | portfolio\_\allowbreak{}scenario & 仅保留元数据；正式分析改用许可允许的权威替代源。\\
WEB-0215 & D/\allowbreak{}08\_\allowbreak{}企业案例与公开披露 & NOT\_\allowbreak{}APPLICABLE/\allowbreak{}NOT\_\allowbreak{}ACQUIRED & FAILED/\allowbreak{}EXCLUDED & company\_\allowbreak{}overseas\_\allowbreak{}exposure | country\_\allowbreak{}event | asset\_\allowbreak{}monthly\_\allowbreak{}return | portfol… & 仅保留元数据；正式分析改用许可允许的权威替代源。\\
WEB-0216 & B/\allowbreak{}04\_\allowbreak{}国别货币与中国出海 & ACQUIRED/\allowbreak{}NARRATIVE\_\allowbreak{}CONTENT & EXTRACTED/\allowbreak{}SUPPLEMENT & source\_\allowbreak{}registry | country\_\allowbreak{}monthly\_\allowbreak{}risk | country\_\allowbreak{}policy\_\allowbreak{}year & 按项目用途复核；不足以单独支持核心结论。\\
WEB-0217 & D/\allowbreak{}07\_\allowbreak{}会计合规与制裁 & NOT\_\allowbreak{}APPLICABLE/\allowbreak{}NOT\_\allowbreak{}ACQUIRED & FAILED/\allowbreak{}EXCLUDED & company\_\allowbreak{}overseas\_\allowbreak{}exposure | country\_\allowbreak{}policy\_\allowbreak{}year | asset\_\allowbreak{}monthly\_\allowbreak{}return & 仅保留元数据；正式分析改用许可允许的权威替代源。\\
WEB-0218 & D/\allowbreak{}07\_\allowbreak{}会计合规与制裁 & NOT\_\allowbreak{}APPLICABLE/\allowbreak{}NOT\_\allowbreak{}ACQUIRED & FAILED/\allowbreak{}EXCLUDED & company\_\allowbreak{}overseas\_\allowbreak{}exposure | country\_\allowbreak{}policy\_\allowbreak{}year | asset\_\allowbreak{}monthly\_\allowbreak{}return | c… & 仅保留元数据；正式分析改用许可允许的权威替代源。\\
WEB-0219 & D/\allowbreak{}08\_\allowbreak{}企业案例与公开披露 & ACQUIRED/\allowbreak{}DOCUMENT\_\allowbreak{}LOCAL & NOT\_\allowbreak{}REQUIRED/\allowbreak{}EXCLUDED & company\_\allowbreak{}overseas\_\allowbreak{}exposure | country\_\allowbreak{}event | asset\_\allowbreak{}monthly\_\allowbreak{}return | portfol… & 补充权威替代来源；无法取得时删除相应字段和结论承诺。\\
WEB-0220 & C/\allowbreak{}08\_\allowbreak{}企业案例与公开披露 & NOT\_\allowbreak{}APPLICABLE/\allowbreak{}NOT\_\allowbreak{}ACQUIRED & FAILED/\allowbreak{}EXCLUDED & company\_\allowbreak{}overseas\_\allowbreak{}exposure | country\_\allowbreak{}event | portfolio\_\allowbreak{}scenario | case\_\allowbreak{}evid… & 仅保留元数据；正式分析改用许可允许的权威替代源。\\
WEB-0221 & C/\allowbreak{}08\_\allowbreak{}企业案例与公开披露 & NOT\_\allowbreak{}APPLICABLE/\allowbreak{}NOT\_\allowbreak{}ACQUIRED & FAILED/\allowbreak{}EXCLUDED & company\_\allowbreak{}overseas\_\allowbreak{}exposure | country\_\allowbreak{}event | asset\_\allowbreak{}monthly\_\allowbreak{}return | portfol… & 仅保留元数据；正式分析改用许可允许的权威替代源。\\
WEB-0222 & C/\allowbreak{}08\_\allowbreak{}企业案例与公开披露 & NOT\_\allowbreak{}APPLICABLE/\allowbreak{}NOT\_\allowbreak{}ACQUIRED & FAILED/\allowbreak{}EXCLUDED & company\_\allowbreak{}overseas\_\allowbreak{}exposure | country\_\allowbreak{}event | asset\_\allowbreak{}monthly\_\allowbreak{}return | portfol… & 仅保留元数据；正式分析改用许可允许的权威替代源。\\
WEB-0223 & C/\allowbreak{}05\_\allowbreak{}全球周期与历史危机 & NOT\_\allowbreak{}APPLICABLE/\allowbreak{}NOT\_\allowbreak{}ACQUIRED & FAILED/\allowbreak{}EXCLUDED & asset\_\allowbreak{}monthly\_\allowbreak{}return | case\_\allowbreak{}evidence | historical\_\allowbreak{}crisis\_\allowbreak{}event | global\_\allowbreak{}cy… & 仅保留元数据；正式分析改用许可允许的权威替代源。\\
WEB-0224 & D/\allowbreak{}05\_\allowbreak{}全球周期与历史危机 & NOT\_\allowbreak{}APPLICABLE/\allowbreak{}NOT\_\allowbreak{}ACQUIRED & FAILED/\allowbreak{}EXCLUDED & case\_\allowbreak{}evidence | historical\_\allowbreak{}crisis\_\allowbreak{}event | global\_\allowbreak{}cycle\_\allowbreak{}month & 仅保留元数据；正式分析改用许可允许的权威替代源。\\
WEB-0225 & D/\allowbreak{}06\_\allowbreak{}黄金美元与储备配置 & NOT\_\allowbreak{}APPLICABLE/\allowbreak{}NOT\_\allowbreak{}ACQUIRED & FAILED/\allowbreak{}EXCLUDED & asset\_\allowbreak{}monthly\_\allowbreak{}return | portfolio\_\allowbreak{}scenario & 仅保留元数据；正式分析改用许可允许的权威替代源。\\
WEB-0226 & D/\allowbreak{}05\_\allowbreak{}全球周期与历史危机 & NOT\_\allowbreak{}APPLICABLE/\allowbreak{}NOT\_\allowbreak{}ACQUIRED & FAILED/\allowbreak{}EXCLUDED & country\_\allowbreak{}monthly\_\allowbreak{}risk | case\_\allowbreak{}evidence | historical\_\allowbreak{}crisis\_\allowbreak{}event | global\_\allowbreak{}cy… & 仅保留元数据；正式分析改用许可允许的权威替代源。\\
WEB-0227 & D/\allowbreak{}06\_\allowbreak{}黄金美元与储备配置 & NOT\_\allowbreak{}APPLICABLE/\allowbreak{}NOT\_\allowbreak{}ACQUIRED & FAILED/\allowbreak{}EXCLUDED & country\_\allowbreak{}monthly\_\allowbreak{}risk | country\_\allowbreak{}event | asset\_\allowbreak{}monthly\_\allowbreak{}return | portfolio\_\allowbreak{}sc… & 仅保留元数据；正式分析改用许可允许的权威替代源。\\
WEB-0228 & D/\allowbreak{}06\_\allowbreak{}黄金美元与储备配置 & NOT\_\allowbreak{}APPLICABLE/\allowbreak{}NOT\_\allowbreak{}ACQUIRED & FAILED/\allowbreak{}EXCLUDED & asset\_\allowbreak{}monthly\_\allowbreak{}return | portfolio\_\allowbreak{}scenario & 仅保留元数据；正式分析改用许可允许的权威替代源。\\
WEB-0229 & A/\allowbreak{}07\_\allowbreak{}会计合规与制裁 & NOT\_\allowbreak{}APPLICABLE/\allowbreak{}NOT\_\allowbreak{}ACQUIRED & FAILED/\allowbreak{}EXCLUDED & country\_\allowbreak{}policy\_\allowbreak{}year | asset\_\allowbreak{}monthly\_\allowbreak{}return | case\_\allowbreak{}evidence & 仅保留元数据；正式分析改用许可允许的权威替代源。\\
WEB-0230 & D/\allowbreak{}08\_\allowbreak{}企业案例与公开披露 & NOT\_\allowbreak{}APPLICABLE/\allowbreak{}NOT\_\allowbreak{}ACQUIRED & FAILED/\allowbreak{}EXCLUDED & company\_\allowbreak{}overseas\_\allowbreak{}exposure | country\_\allowbreak{}event | asset\_\allowbreak{}monthly\_\allowbreak{}return | portfol… & 仅保留元数据；正式分析改用许可允许的权威替代源。\\
WEB-0231 & D/\allowbreak{}07\_\allowbreak{}会计合规与制裁 & NOT\_\allowbreak{}APPLICABLE/\allowbreak{}NOT\_\allowbreak{}ACQUIRED & FAILED/\allowbreak{}EXCLUDED & country\_\allowbreak{}policy\_\allowbreak{}year | asset\_\allowbreak{}monthly\_\allowbreak{}return & 仅保留元数据；正式分析改用许可允许的权威替代源。\\
WEB-0232 & D/\allowbreak{}05\_\allowbreak{}全球周期与历史危机 & NOT\_\allowbreak{}APPLICABLE/\allowbreak{}NOT\_\allowbreak{}ACQUIRED & FAILED/\allowbreak{}EXCLUDED & asset\_\allowbreak{}monthly\_\allowbreak{}return | case\_\allowbreak{}evidence | historical\_\allowbreak{}crisis\_\allowbreak{}event | global\_\allowbreak{}cy… & 仅保留元数据；正式分析改用许可允许的权威替代源。\\
WEB-0233 & D/\allowbreak{}05\_\allowbreak{}全球周期与历史危机 & NOT\_\allowbreak{}APPLICABLE/\allowbreak{}NOT\_\allowbreak{}ACQUIRED & FAILED/\allowbreak{}EXCLUDED & case\_\allowbreak{}evidence | historical\_\allowbreak{}crisis\_\allowbreak{}event | global\_\allowbreak{}cycle\_\allowbreak{}month & 仅保留元数据；正式分析改用许可允许的权威替代源。\\
WEB-0234 & D/\allowbreak{}07\_\allowbreak{}会计合规与制裁 & NOT\_\allowbreak{}APPLICABLE/\allowbreak{}NOT\_\allowbreak{}ACQUIRED & FAILED/\allowbreak{}EXCLUDED & country\_\allowbreak{}policy\_\allowbreak{}year | asset\_\allowbreak{}monthly\_\allowbreak{}return | case\_\allowbreak{}evidence & 仅保留元数据；正式分析改用许可允许的权威替代源。\\
WEB-0235 & A/\allowbreak{}05\_\allowbreak{}全球周期与历史危机 & ACQUIRED/\allowbreak{}MIXED\_\allowbreak{}CONTENT & EXTRACTED/\allowbreak{}CORE & case\_\allowbreak{}evidence | historical\_\allowbreak{}crisis\_\allowbreak{}event | global\_\allowbreak{}cycle\_\allowbreak{}month & 人工复核关键字段与source\_\allowbreak{}locator；通过后映射至正式source\_\allowbreak{}registry。\\
WEB-0236 & D/\allowbreak{}06\_\allowbreak{}黄金美元与储备配置 & ACQUIRED/\allowbreak{}LOCAL\_\allowbreak{}ASSET & NOT\_\allowbreak{}REQUIRED/\allowbreak{}REPLACED & asset\_\allowbreak{}monthly\_\allowbreak{}return | portfolio\_\allowbreak{}scenario & 使用替代来源：LCL-0154\\
WEB-0237 & B/\allowbreak{}05\_\allowbreak{}全球周期与历史危机 & ACQUIRED/\allowbreak{}MIXED\_\allowbreak{}CONTENT & EXTRACTED/\allowbreak{}SUPPLEMENT & asset\_\allowbreak{}monthly\_\allowbreak{}return | case\_\allowbreak{}evidence | historical\_\allowbreak{}crisis\_\allowbreak{}event | global\_\allowbreak{}cy… & 按项目用途复核；不足以单独支持核心结论。\\
WEB-0238 & B/\allowbreak{}05\_\allowbreak{}全球周期与历史危机 & ACQUIRED/\allowbreak{}NARRATIVE\_\allowbreak{}CONTENT & EXTRACTED/\allowbreak{}SUPPLEMENT & case\_\allowbreak{}evidence | historical\_\allowbreak{}crisis\_\allowbreak{}event | global\_\allowbreak{}cycle\_\allowbreak{}month & 按项目用途复核；不足以单独支持核心结论。\\
WEB-0239 & D/\allowbreak{}05\_\allowbreak{}全球周期与历史危机 & ACQUIRED/\allowbreak{}DOCUMENT\_\allowbreak{}LOCAL & NOT\_\allowbreak{}REQUIRED/\allowbreak{}EXCLUDED & case\_\allowbreak{}evidence | historical\_\allowbreak{}crisis\_\allowbreak{}event | global\_\allowbreak{}cycle\_\allowbreak{}month & 补充权威替代来源；无法取得时删除相应字段和结论承诺。\\
WEB-0240 & A/\allowbreak{}07\_\allowbreak{}会计合规与制裁 & NOT\_\allowbreak{}APPLICABLE/\allowbreak{}NOT\_\allowbreak{}ACQUIRED & FAILED/\allowbreak{}EXCLUDED & country\_\allowbreak{}policy\_\allowbreak{}year | asset\_\allowbreak{}monthly\_\allowbreak{}return | case\_\allowbreak{}evidence & 仅保留元数据；正式分析改用许可允许的权威替代源。\\
WEB-0241 & A/\allowbreak{}07\_\allowbreak{}会计合规与制裁 & NOT\_\allowbreak{}APPLICABLE/\allowbreak{}NOT\_\allowbreak{}ACQUIRED & FAILED/\allowbreak{}EXCLUDED & country\_\allowbreak{}monthly\_\allowbreak{}risk | country\_\allowbreak{}policy\_\allowbreak{}year | country\_\allowbreak{}event | case\_\allowbreak{}evidence & 仅保留元数据；正式分析改用许可允许的权威替代源。\\
WEB-0242 & B/\allowbreak{}08\_\allowbreak{}企业案例与公开披露 & NOT\_\allowbreak{}APPLICABLE/\allowbreak{}NOT\_\allowbreak{}ACQUIRED & FAILED/\allowbreak{}EXCLUDED & source\_\allowbreak{}registry | company\_\allowbreak{}overseas\_\allowbreak{}exposure | country\_\allowbreak{}event | asset\_\allowbreak{}monthl… & 按人工任务卡接受条款或完成机构授权后重新验收。\\
WEB-0243 & B/\allowbreak{}08\_\allowbreak{}企业案例与公开披露 & NOT\_\allowbreak{}APPLICABLE/\allowbreak{}NOT\_\allowbreak{}ACQUIRED & FAILED/\allowbreak{}EXCLUDED & company\_\allowbreak{}overseas\_\allowbreak{}exposure | country\_\allowbreak{}event | asset\_\allowbreak{}monthly\_\allowbreak{}return | portfol… & 按人工任务卡接受条款或完成机构授权后重新验收。\\
WEB-0244 & B/\allowbreak{}06\_\allowbreak{}黄金美元与储备配置 & NOT\_\allowbreak{}APPLICABLE/\allowbreak{}NOT\_\allowbreak{}ACQUIRED & FAILED/\allowbreak{}EXCLUDED & country\_\allowbreak{}monthly\_\allowbreak{}risk | country\_\allowbreak{}event | asset\_\allowbreak{}monthly\_\allowbreak{}return | portfolio\_\allowbreak{}sc… & 按人工任务卡接受条款或完成机构授权后重新验收。\\
WEB-0245 & B/\allowbreak{}06\_\allowbreak{}黄金美元与储备配置 & NOT\_\allowbreak{}APPLICABLE/\allowbreak{}NOT\_\allowbreak{}ACQUIRED & FAILED/\allowbreak{}EXCLUDED & asset\_\allowbreak{}monthly\_\allowbreak{}return | portfolio\_\allowbreak{}scenario & 按人工任务卡接受条款或完成机构授权后重新验收。\\
WEB-0246 & B/\allowbreak{}06\_\allowbreak{}黄金美元与储备配置 & NOT\_\allowbreak{}APPLICABLE/\allowbreak{}NOT\_\allowbreak{}ACQUIRED & FAILED/\allowbreak{}EXCLUDED & asset\_\allowbreak{}monthly\_\allowbreak{}return | portfolio\_\allowbreak{}scenario & 按人工任务卡接受条款或完成机构授权后重新验收。\\
WEB-0247 & D/\allowbreak{}05\_\allowbreak{}全球周期与历史危机 & NOT\_\allowbreak{}APPLICABLE/\allowbreak{}NOT\_\allowbreak{}ACQUIRED & FAILED/\allowbreak{}EXCLUDED & case\_\allowbreak{}evidence | historical\_\allowbreak{}crisis\_\allowbreak{}event | global\_\allowbreak{}cycle\_\allowbreak{}month & 仅保留元数据；正式分析改用许可允许的权威替代源。\\
WEB-0248 & D/\allowbreak{}08\_\allowbreak{}企业案例与公开披露 & NOT\_\allowbreak{}APPLICABLE/\allowbreak{}NOT\_\allowbreak{}ACQUIRED & FAILED/\allowbreak{}EXCLUDED & company\_\allowbreak{}overseas\_\allowbreak{}exposure | country\_\allowbreak{}event | asset\_\allowbreak{}monthly\_\allowbreak{}return | portfol… & 仅保留元数据；正式分析改用许可允许的权威替代源。\\
WEB-0249 & D/\allowbreak{}08\_\allowbreak{}企业案例与公开披露 & ACQUIRED/\allowbreak{}DOCUMENT\_\allowbreak{}LOCAL & NOT\_\allowbreak{}REQUIRED/\allowbreak{}EXCLUDED & company\_\allowbreak{}overseas\_\allowbreak{}exposure | country\_\allowbreak{}event | asset\_\allowbreak{}monthly\_\allowbreak{}return | portfol… & 补充权威替代来源；无法取得时删除相应字段和结论承诺。\\
WEB-0250 & D/\allowbreak{}08\_\allowbreak{}企业案例与公开披露 & NOT\_\allowbreak{}APPLICABLE/\allowbreak{}NOT\_\allowbreak{}ACQUIRED & FAILED/\allowbreak{}EXCLUDED & company\_\allowbreak{}overseas\_\allowbreak{}exposure | country\_\allowbreak{}event | asset\_\allowbreak{}monthly\_\allowbreak{}return | portfol… & 仅保留元数据；正式分析改用许可允许的权威替代源。\\
WEB-0251 & D/\allowbreak{}08\_\allowbreak{}企业案例与公开披露 & NOT\_\allowbreak{}APPLICABLE/\allowbreak{}NOT\_\allowbreak{}ACQUIRED & FAILED/\allowbreak{}EXCLUDED & company\_\allowbreak{}overseas\_\allowbreak{}exposure | country\_\allowbreak{}event | asset\_\allowbreak{}monthly\_\allowbreak{}return | portfol… & 仅保留元数据；正式分析改用许可允许的权威替代源。\\
WEB-0252 & A/\allowbreak{}06\_\allowbreak{}黄金美元与储备配置 & NOT\_\allowbreak{}APPLICABLE/\allowbreak{}NOT\_\allowbreak{}ACQUIRED & FAILED/\allowbreak{}EXCLUDED & asset\_\allowbreak{}monthly\_\allowbreak{}return | portfolio\_\allowbreak{}scenario & 仅保留元数据；正式分析改用许可允许的权威替代源。\\
WEB-0253 & A/\allowbreak{}06\_\allowbreak{}黄金美元与储备配置 & ACQUIRED/\allowbreak{}MIXED\_\allowbreak{}CONTENT & EXTRACTED/\allowbreak{}CORE & asset\_\allowbreak{}monthly\_\allowbreak{}return | portfolio\_\allowbreak{}scenario & 人工复核关键字段与source\_\allowbreak{}locator；通过后映射至正式source\_\allowbreak{}registry。\\
WEB-0254 & A/\allowbreak{}05\_\allowbreak{}全球周期与历史危机 & ACQUIRED/\allowbreak{}NARRATIVE\_\allowbreak{}CONTENT & EXTRACTED/\allowbreak{}CORE & case\_\allowbreak{}evidence | historical\_\allowbreak{}crisis\_\allowbreak{}event | global\_\allowbreak{}cycle\_\allowbreak{}month & 人工复核关键字段与source\_\allowbreak{}locator；通过后映射至正式source\_\allowbreak{}registry。\\
WEB-0255 & A/\allowbreak{}05\_\allowbreak{}全球周期与历史危机 & NOT\_\allowbreak{}ACQUIRED/\allowbreak{}NOT\_\allowbreak{}ACQUIRED & FAILED/\allowbreak{}CORE & case\_\allowbreak{}evidence | historical\_\allowbreak{}crisis\_\allowbreak{}event | global\_\allowbreak{}cycle\_\allowbreak{}month & 补充权威替代来源；无法取得时删除相应字段和结论承诺。\\
WEB-0256 & A/\allowbreak{}04\_\allowbreak{}国别货币与中国出海 & ACQUIRED/\allowbreak{}NARRATIVE\_\allowbreak{}CONTENT & EXTRACTED/\allowbreak{}CORE & country\_\allowbreak{}monthly\_\allowbreak{}risk | country\_\allowbreak{}policy\_\allowbreak{}year & 人工复核关键字段与source\_\allowbreak{}locator；通过后映射至正式source\_\allowbreak{}registry。\\
WEB-0257 & A/\allowbreak{}06\_\allowbreak{}黄金美元与储备配置 & ACQUIRED/\allowbreak{}NARRATIVE\_\allowbreak{}CONTENT & EXTRACTED/\allowbreak{}CORE & asset\_\allowbreak{}monthly\_\allowbreak{}return | portfolio\_\allowbreak{}scenario & 人工复核关键字段与source\_\allowbreak{}locator；通过后映射至正式source\_\allowbreak{}registry。\\
WEB-0258 & A/\allowbreak{}06\_\allowbreak{}黄金美元与储备配置 & ACQUIRED/\allowbreak{}NARRATIVE\_\allowbreak{}CONTENT & EXTRACTED/\allowbreak{}CORE & asset\_\allowbreak{}monthly\_\allowbreak{}return | portfolio\_\allowbreak{}scenario & 人工复核关键字段与source\_\allowbreak{}locator；通过后映射至正式source\_\allowbreak{}registry。\\
WEB-0259 & A/\allowbreak{}06\_\allowbreak{}黄金美元与储备配置 & ACQUIRED/\allowbreak{}NARRATIVE\_\allowbreak{}CONTENT & EXTRACTED/\allowbreak{}CORE & asset\_\allowbreak{}monthly\_\allowbreak{}return | portfolio\_\allowbreak{}scenario & 人工复核关键字段与source\_\allowbreak{}locator；通过后映射至正式source\_\allowbreak{}registry。\\
WEB-0260 & A/\allowbreak{}06\_\allowbreak{}黄金美元与储备配置 & ACQUIRED/\allowbreak{}NARRATIVE\_\allowbreak{}CONTENT & EXTRACTED/\allowbreak{}CORE & asset\_\allowbreak{}monthly\_\allowbreak{}return | portfolio\_\allowbreak{}scenario & 人工复核关键字段与source\_\allowbreak{}locator；通过后映射至正式source\_\allowbreak{}registry。\\
WEB-0261 & A/\allowbreak{}06\_\allowbreak{}黄金美元与储备配置 & ACQUIRED/\allowbreak{}NARRATIVE\_\allowbreak{}CONTENT & EXTRACTED/\allowbreak{}CORE & asset\_\allowbreak{}monthly\_\allowbreak{}return | portfolio\_\allowbreak{}scenario & 人工复核关键字段与source\_\allowbreak{}locator；通过后映射至正式source\_\allowbreak{}registry。\\
WEB-0262 & A/\allowbreak{}06\_\allowbreak{}黄金美元与储备配置 & ACQUIRED/\allowbreak{}NARRATIVE\_\allowbreak{}CONTENT & EXTRACTED/\allowbreak{}CORE & asset\_\allowbreak{}monthly\_\allowbreak{}return | portfolio\_\allowbreak{}scenario & 人工复核关键字段与source\_\allowbreak{}locator；通过后映射至正式source\_\allowbreak{}registry。\\
WEB-0263 & A/\allowbreak{}06\_\allowbreak{}黄金美元与储备配置 & ACQUIRED/\allowbreak{}NARRATIVE\_\allowbreak{}CONTENT & EXTRACTED/\allowbreak{}CORE & asset\_\allowbreak{}monthly\_\allowbreak{}return | portfolio\_\allowbreak{}scenario & 人工复核关键字段与source\_\allowbreak{}locator；通过后映射至正式source\_\allowbreak{}registry。\\
WEB-0264 & A/\allowbreak{}06\_\allowbreak{}黄金美元与储备配置 & ACQUIRED/\allowbreak{}MIXED\_\allowbreak{}CONTENT & EXTRACTED/\allowbreak{}CORE & asset\_\allowbreak{}monthly\_\allowbreak{}return | portfolio\_\allowbreak{}scenario & 人工复核关键字段与source\_\allowbreak{}locator；通过后映射至正式source\_\allowbreak{}registry。\\
WEB-0265 & A/\allowbreak{}06\_\allowbreak{}黄金美元与储备配置 & ACQUIRED/\allowbreak{}MIXED\_\allowbreak{}CONTENT & EXTRACTED/\allowbreak{}CORE & asset\_\allowbreak{}monthly\_\allowbreak{}return | portfolio\_\allowbreak{}scenario & 人工复核关键字段与source\_\allowbreak{}locator；通过后映射至正式source\_\allowbreak{}registry。\\
WEB-0266 & A/\allowbreak{}06\_\allowbreak{}黄金美元与储备配置 & ACQUIRED/\allowbreak{}MIXED\_\allowbreak{}CONTENT & EXTRACTED/\allowbreak{}CORE & asset\_\allowbreak{}monthly\_\allowbreak{}return | portfolio\_\allowbreak{}scenario & 人工复核关键字段与source\_\allowbreak{}locator；通过后映射至正式source\_\allowbreak{}registry。\\
WEB-0267 & B/\allowbreak{}05\_\allowbreak{}全球周期与历史危机 & NOT\_\allowbreak{}APPLICABLE/\allowbreak{}NOT\_\allowbreak{}ACQUIRED & FAILED/\allowbreak{}EXCLUDED & case\_\allowbreak{}evidence | historical\_\allowbreak{}crisis\_\allowbreak{}event | global\_\allowbreak{}cycle\_\allowbreak{}month & 仅保留元数据；正式分析改用许可允许的权威替代源。\\
WEB-0268 & B/\allowbreak{}05\_\allowbreak{}全球周期与历史危机 & NOT\_\allowbreak{}APPLICABLE/\allowbreak{}NOT\_\allowbreak{}ACQUIRED & FAILED/\allowbreak{}EXCLUDED & asset\_\allowbreak{}monthly\_\allowbreak{}return | case\_\allowbreak{}evidence | historical\_\allowbreak{}crisis\_\allowbreak{}event | global\_\allowbreak{}cy… & 仅保留元数据；正式分析改用许可允许的权威替代源。\\
WEB-0269 & D/\allowbreak{}08\_\allowbreak{}企业案例与公开披露 & NOT\_\allowbreak{}APPLICABLE/\allowbreak{}NOT\_\allowbreak{}ACQUIRED & FAILED/\allowbreak{}EXCLUDED & company\_\allowbreak{}overseas\_\allowbreak{}exposure | country\_\allowbreak{}event | asset\_\allowbreak{}monthly\_\allowbreak{}return | portfol… & 仅保留元数据；正式分析改用许可允许的权威替代源。\\
WEB-0270 & D/\allowbreak{}08\_\allowbreak{}企业案例与公开披露 & NOT\_\allowbreak{}APPLICABLE/\allowbreak{}NOT\_\allowbreak{}ACQUIRED & FAILED/\allowbreak{}EXCLUDED & company\_\allowbreak{}overseas\_\allowbreak{}exposure | country\_\allowbreak{}event | asset\_\allowbreak{}monthly\_\allowbreak{}return | portfol… & 仅保留元数据；正式分析改用许可允许的权威替代源。\\
WEB-0271 & D/\allowbreak{}08\_\allowbreak{}企业案例与公开披露 & NOT\_\allowbreak{}APPLICABLE/\allowbreak{}NOT\_\allowbreak{}ACQUIRED & FAILED/\allowbreak{}EXCLUDED & company\_\allowbreak{}overseas\_\allowbreak{}exposure | country\_\allowbreak{}event | asset\_\allowbreak{}monthly\_\allowbreak{}return | portfol… & 仅保留元数据；正式分析改用许可允许的权威替代源。\\
WEB-0272 & D/\allowbreak{}08\_\allowbreak{}企业案例与公开披露 & NOT\_\allowbreak{}APPLICABLE/\allowbreak{}NOT\_\allowbreak{}ACQUIRED & FAILED/\allowbreak{}EXCLUDED & company\_\allowbreak{}overseas\_\allowbreak{}exposure | country\_\allowbreak{}event | asset\_\allowbreak{}monthly\_\allowbreak{}return | portfol… & 仅保留元数据；正式分析改用许可允许的权威替代源。\\
WEB-0273 & D/\allowbreak{}05\_\allowbreak{}全球周期与历史危机 & NOT\_\allowbreak{}APPLICABLE/\allowbreak{}NOT\_\allowbreak{}ACQUIRED & FAILED/\allowbreak{}EXCLUDED & case\_\allowbreak{}evidence | historical\_\allowbreak{}crisis\_\allowbreak{}event | global\_\allowbreak{}cycle\_\allowbreak{}month & 仅保留元数据；正式分析改用许可允许的权威替代源。\\
WEB-0274 & D/\allowbreak{}06\_\allowbreak{}黄金美元与储备配置 & NOT\_\allowbreak{}APPLICABLE/\allowbreak{}NOT\_\allowbreak{}ACQUIRED & FAILED/\allowbreak{}EXCLUDED & asset\_\allowbreak{}monthly\_\allowbreak{}return | portfolio\_\allowbreak{}scenario & 仅保留元数据；正式分析改用许可允许的权威替代源。\\
WEB-0275 & D/\allowbreak{}07\_\allowbreak{}会计合规与制裁 & NOT\_\allowbreak{}APPLICABLE/\allowbreak{}NOT\_\allowbreak{}ACQUIRED & FAILED/\allowbreak{}EXCLUDED & company\_\allowbreak{}overseas\_\allowbreak{}exposure | country\_\allowbreak{}policy\_\allowbreak{}year | asset\_\allowbreak{}monthly\_\allowbreak{}return | c… & 仅保留元数据；正式分析改用许可允许的权威替代源。\\
WEB-0276 & D/\allowbreak{}06\_\allowbreak{}黄金美元与储备配置 & NOT\_\allowbreak{}APPLICABLE/\allowbreak{}NOT\_\allowbreak{}ACQUIRED & FAILED/\allowbreak{}EXCLUDED & asset\_\allowbreak{}monthly\_\allowbreak{}return | portfolio\_\allowbreak{}scenario & 仅保留元数据；正式分析改用许可允许的权威替代源。\\
WEB-0277 & D/\allowbreak{}07\_\allowbreak{}会计合规与制裁 & NOT\_\allowbreak{}APPLICABLE/\allowbreak{}NOT\_\allowbreak{}ACQUIRED & FAILED/\allowbreak{}EXCLUDED & country\_\allowbreak{}policy\_\allowbreak{}year | country\_\allowbreak{}event | asset\_\allowbreak{}monthly\_\allowbreak{}return | case\_\allowbreak{}evidence & 仅保留元数据；正式分析改用许可允许的权威替代源。\\
WEB-0278 & D/\allowbreak{}07\_\allowbreak{}会计合规与制裁 & NOT\_\allowbreak{}APPLICABLE/\allowbreak{}NOT\_\allowbreak{}ACQUIRED & FAILED/\allowbreak{}EXCLUDED & country\_\allowbreak{}monthly\_\allowbreak{}risk | country\_\allowbreak{}policy\_\allowbreak{}year | country\_\allowbreak{}event | asset\_\allowbreak{}monthly… & 仅保留元数据；正式分析改用许可允许的权威替代源。\\
WEB-0279 & D/\allowbreak{}07\_\allowbreak{}会计合规与制裁 & NOT\_\allowbreak{}APPLICABLE/\allowbreak{}NOT\_\allowbreak{}ACQUIRED & FAILED/\allowbreak{}EXCLUDED & company\_\allowbreak{}overseas\_\allowbreak{}exposure | country\_\allowbreak{}monthly\_\allowbreak{}risk | country\_\allowbreak{}policy\_\allowbreak{}year | c… & 仅保留元数据；正式分析改用许可允许的权威替代源。\\
WEB-0280 & D/\allowbreak{}07\_\allowbreak{}会计合规与制裁 & NOT\_\allowbreak{}APPLICABLE/\allowbreak{}NOT\_\allowbreak{}ACQUIRED & FAILED/\allowbreak{}EXCLUDED & company\_\allowbreak{}overseas\_\allowbreak{}exposure | country\_\allowbreak{}policy\_\allowbreak{}year | country\_\allowbreak{}event | asset\_\allowbreak{}mo… & 仅保留元数据；正式分析改用许可允许的权威替代源。\\
WEB-0281 & D/\allowbreak{}06\_\allowbreak{}黄金美元与储备配置 & NOT\_\allowbreak{}APPLICABLE/\allowbreak{}NOT\_\allowbreak{}ACQUIRED & FAILED/\allowbreak{}EXCLUDED & asset\_\allowbreak{}monthly\_\allowbreak{}return | portfolio\_\allowbreak{}scenario & 仅保留元数据；正式分析改用许可允许的权威替代源。\\
WEB-0282 & D/\allowbreak{}05\_\allowbreak{}全球周期与历史危机 & NOT\_\allowbreak{}APPLICABLE/\allowbreak{}NOT\_\allowbreak{}ACQUIRED & FAILED/\allowbreak{}EXCLUDED & asset\_\allowbreak{}monthly\_\allowbreak{}return | case\_\allowbreak{}evidence | historical\_\allowbreak{}crisis\_\allowbreak{}event | global\_\allowbreak{}cy… & 仅保留元数据；正式分析改用许可允许的权威替代源。\\
WEB-0283 & D/\allowbreak{}05\_\allowbreak{}全球周期与历史危机 & NOT\_\allowbreak{}APPLICABLE/\allowbreak{}NOT\_\allowbreak{}ACQUIRED & FAILED/\allowbreak{}EXCLUDED & asset\_\allowbreak{}monthly\_\allowbreak{}return | case\_\allowbreak{}evidence | historical\_\allowbreak{}crisis\_\allowbreak{}event | global\_\allowbreak{}cy… & 仅保留元数据；正式分析改用许可允许的权威替代源。\\
WEB-0284 & D/\allowbreak{}05\_\allowbreak{}全球周期与历史危机 & NOT\_\allowbreak{}APPLICABLE/\allowbreak{}NOT\_\allowbreak{}ACQUIRED & FAILED/\allowbreak{}EXCLUDED & asset\_\allowbreak{}monthly\_\allowbreak{}return | portfolio\_\allowbreak{}scenario | case\_\allowbreak{}evidence | historical\_\allowbreak{}cri… & 仅保留元数据；正式分析改用许可允许的权威替代源。\\
WEB-0285 & D/\allowbreak{}06\_\allowbreak{}黄金美元与储备配置 & NOT\_\allowbreak{}APPLICABLE/\allowbreak{}NOT\_\allowbreak{}ACQUIRED & FAILED/\allowbreak{}EXCLUDED & asset\_\allowbreak{}monthly\_\allowbreak{}return | portfolio\_\allowbreak{}scenario & 仅保留元数据；正式分析改用许可允许的权威替代源。\\
WEB-0286 & D/\allowbreak{}06\_\allowbreak{}黄金美元与储备配置 & NOT\_\allowbreak{}APPLICABLE/\allowbreak{}NOT\_\allowbreak{}ACQUIRED & FAILED/\allowbreak{}EXCLUDED & asset\_\allowbreak{}monthly\_\allowbreak{}return | portfolio\_\allowbreak{}scenario & 仅保留元数据；正式分析改用许可允许的权威替代源。\\
WEB-0287 & D/\allowbreak{}05\_\allowbreak{}全球周期与历史危机 & NOT\_\allowbreak{}APPLICABLE/\allowbreak{}NOT\_\allowbreak{}ACQUIRED & FAILED/\allowbreak{}EXCLUDED & case\_\allowbreak{}evidence | historical\_\allowbreak{}crisis\_\allowbreak{}event | global\_\allowbreak{}cycle\_\allowbreak{}month & 仅保留元数据；正式分析改用许可允许的权威替代源。\\
WEB-0288 & D/\allowbreak{}05\_\allowbreak{}全球周期与历史危机 & NOT\_\allowbreak{}APPLICABLE/\allowbreak{}NOT\_\allowbreak{}ACQUIRED & FAILED/\allowbreak{}EXCLUDED & asset\_\allowbreak{}monthly\_\allowbreak{}return | case\_\allowbreak{}evidence | historical\_\allowbreak{}crisis\_\allowbreak{}event | global\_\allowbreak{}cy… & 仅保留元数据；正式分析改用许可允许的权威替代源。\\
WEB-0289 & A/\allowbreak{}08\_\allowbreak{}企业案例与公开披露 & ACQUIRED/\allowbreak{}NARRATIVE\_\allowbreak{}CONTENT & EXTRACTED/\allowbreak{}CORE & company\_\allowbreak{}overseas\_\allowbreak{}exposure | country\_\allowbreak{}event | asset\_\allowbreak{}monthly\_\allowbreak{}return | portfol… & 人工复核关键字段与source\_\allowbreak{}locator；通过后映射至正式source\_\allowbreak{}registry。\\
WEB-0290 & A/\allowbreak{}05\_\allowbreak{}全球周期与历史危机 & NOT\_\allowbreak{}APPLICABLE/\allowbreak{}NOT\_\allowbreak{}ACQUIRED & FAILED/\allowbreak{}EXCLUDED & asset\_\allowbreak{}monthly\_\allowbreak{}return | case\_\allowbreak{}evidence | historical\_\allowbreak{}crisis\_\allowbreak{}event | global\_\allowbreak{}cy… & 仅保留元数据；正式分析改用许可允许的权威替代源。\\
WEB-0291 & D/\allowbreak{}08\_\allowbreak{}企业案例与公开披露 & NOT\_\allowbreak{}APPLICABLE/\allowbreak{}NOT\_\allowbreak{}ACQUIRED & FAILED/\allowbreak{}EXCLUDED & company\_\allowbreak{}overseas\_\allowbreak{}exposure | country\_\allowbreak{}event | asset\_\allowbreak{}monthly\_\allowbreak{}return | portfol… & 仅保留元数据；正式分析改用许可允许的权威替代源。\\
WEB-0292 & D/\allowbreak{}08\_\allowbreak{}企业案例与公开披露 & NOT\_\allowbreak{}APPLICABLE/\allowbreak{}NOT\_\allowbreak{}ACQUIRED & FAILED/\allowbreak{}EXCLUDED & company\_\allowbreak{}overseas\_\allowbreak{}exposure | country\_\allowbreak{}event | asset\_\allowbreak{}monthly\_\allowbreak{}return | portfol… & 仅保留元数据；正式分析改用许可允许的权威替代源。\\
WEB-0293 & D/\allowbreak{}06\_\allowbreak{}黄金美元与储备配置 & NOT\_\allowbreak{}APPLICABLE/\allowbreak{}NOT\_\allowbreak{}ACQUIRED & FAILED/\allowbreak{}EXCLUDED & asset\_\allowbreak{}monthly\_\allowbreak{}return | portfolio\_\allowbreak{}scenario & 按人工任务卡接受条款或完成机构授权后重新验收。\\
WEB-0294 & D/\allowbreak{}06\_\allowbreak{}黄金美元与储备配置 & NOT\_\allowbreak{}APPLICABLE/\allowbreak{}NOT\_\allowbreak{}ACQUIRED & FAILED/\allowbreak{}EXCLUDED & asset\_\allowbreak{}monthly\_\allowbreak{}return | portfolio\_\allowbreak{}scenario & 按人工任务卡接受条款或完成机构授权后重新验收。\\
WEB-0295 & D/\allowbreak{}06\_\allowbreak{}黄金美元与储备配置 & NOT\_\allowbreak{}APPLICABLE/\allowbreak{}NOT\_\allowbreak{}ACQUIRED & FAILED/\allowbreak{}EXCLUDED & asset\_\allowbreak{}monthly\_\allowbreak{}return | portfolio\_\allowbreak{}scenario & 按人工任务卡接受条款或完成机构授权后重新验收。\\
WEB-0296 & D/\allowbreak{}08\_\allowbreak{}企业案例与公开披露 & NOT\_\allowbreak{}APPLICABLE/\allowbreak{}NOT\_\allowbreak{}ACQUIRED & FAILED/\allowbreak{}EXCLUDED & company\_\allowbreak{}overseas\_\allowbreak{}exposure | country\_\allowbreak{}event | portfolio\_\allowbreak{}scenario | case\_\allowbreak{}evid… & 仅保留元数据；正式分析改用许可允许的权威替代源。\\
WEB-0297 & D/\allowbreak{}06\_\allowbreak{}黄金美元与储备配置 & ACQUIRED/\allowbreak{}DOCUMENT\_\allowbreak{}LOCAL & NOT\_\allowbreak{}REQUIRED/\allowbreak{}EXCLUDED & asset\_\allowbreak{}monthly\_\allowbreak{}return | portfolio\_\allowbreak{}scenario & 补充权威替代来源；无法取得时删除相应字段和结论承诺。\\
WEB-0298 & B/\allowbreak{}08\_\allowbreak{}企业案例与公开披露 & ACQUIRED/\allowbreak{}NARRATIVE\_\allowbreak{}CONTENT & EXTRACTED/\allowbreak{}SUPPLEMENT & company\_\allowbreak{}overseas\_\allowbreak{}exposure | country\_\allowbreak{}event | portfolio\_\allowbreak{}scenario | case\_\allowbreak{}evid… & 按项目用途复核；不足以单独支持核心结论。\\
WEB-0299 & B/\allowbreak{}04\_\allowbreak{}国别货币与中国出海 & ACQUIRED/\allowbreak{}MIXED\_\allowbreak{}CONTENT & EXTRACTED/\allowbreak{}SUPPLEMENT & country\_\allowbreak{}monthly\_\allowbreak{}risk | country\_\allowbreak{}policy\_\allowbreak{}year & 按项目用途复核；不足以单独支持核心结论。\\
WEB-0300 & B/\allowbreak{}08\_\allowbreak{}企业案例与公开披露 & ACQUIRED/\allowbreak{}NARRATIVE\_\allowbreak{}CONTENT & EXTRACTED/\allowbreak{}SUPPLEMENT & company\_\allowbreak{}overseas\_\allowbreak{}exposure | country\_\allowbreak{}event | portfolio\_\allowbreak{}scenario | case\_\allowbreak{}evid… & 按项目用途复核；不足以单独支持核心结论。\\
WEB-0301 & B/\allowbreak{}04\_\allowbreak{}国别货币与中国出海 & ACQUIRED/\allowbreak{}MIXED\_\allowbreak{}CONTENT & EXTRACTED/\allowbreak{}SUPPLEMENT & country\_\allowbreak{}monthly\_\allowbreak{}risk | country\_\allowbreak{}policy\_\allowbreak{}year & 按项目用途复核；不足以单独支持核心结论。\\
WEB-0302 & B/\allowbreak{}04\_\allowbreak{}国别货币与中国出海 & NOT\_\allowbreak{}APPLICABLE/\allowbreak{}NOT\_\allowbreak{}ACQUIRED & FAILED/\allowbreak{}EXCLUDED & country\_\allowbreak{}monthly\_\allowbreak{}risk | country\_\allowbreak{}policy\_\allowbreak{}year & 仅保留元数据；正式分析改用许可允许的权威替代源。\\
WEB-0303 & B/\allowbreak{}07\_\allowbreak{}会计合规与制裁 & ACQUIRED/\allowbreak{}NARRATIVE\_\allowbreak{}CONTENT & EXTRACTED/\allowbreak{}SUPPLEMENT & company\_\allowbreak{}overseas\_\allowbreak{}exposure | country\_\allowbreak{}policy\_\allowbreak{}year | asset\_\allowbreak{}monthly\_\allowbreak{}return | c… & 按项目用途复核；不足以单独支持核心结论。\\
WEB-0304 & B/\allowbreak{}07\_\allowbreak{}会计合规与制裁 & NOT\_\allowbreak{}ACQUIRED/\allowbreak{}NOT\_\allowbreak{}ACQUIRED & FAILED/\allowbreak{}SUPPLEMENT & country\_\allowbreak{}policy\_\allowbreak{}year | asset\_\allowbreak{}monthly\_\allowbreak{}return & 补充权威替代来源；无法取得时删除相应字段和结论承诺。\\
WEB-0305 & B/\allowbreak{}07\_\allowbreak{}会计合规与制裁 & NOT\_\allowbreak{}ACQUIRED/\allowbreak{}NOT\_\allowbreak{}ACQUIRED & FAILED/\allowbreak{}SUPPLEMENT & country\_\allowbreak{}policy\_\allowbreak{}year & 补充权威替代来源；无法取得时删除相应字段和结论承诺。\\
WEB-0306 & B/\allowbreak{}07\_\allowbreak{}会计合规与制裁 & ACQUIRED/\allowbreak{}DOCUMENT\_\allowbreak{}LOCAL & NOT\_\allowbreak{}REQUIRED/\allowbreak{}SUPPLEMENT & country\_\allowbreak{}policy\_\allowbreak{}year | asset\_\allowbreak{}monthly\_\allowbreak{}return & 按项目用途复核；不足以单独支持核心结论。\\
WEB-0307 & B/\allowbreak{}07\_\allowbreak{}会计合规与制裁 & ACQUIRED/\allowbreak{}DOCUMENT\_\allowbreak{}LOCAL & NOT\_\allowbreak{}REQUIRED/\allowbreak{}SUPPLEMENT & country\_\allowbreak{}policy\_\allowbreak{}year & 按项目用途复核；不足以单独支持核心结论。\\
WEB-0308 & B/\allowbreak{}07\_\allowbreak{}会计合规与制裁 & NOT\_\allowbreak{}ACQUIRED/\allowbreak{}NOT\_\allowbreak{}ACQUIRED & FAILED/\allowbreak{}SUPPLEMENT & country\_\allowbreak{}policy\_\allowbreak{}year & 补充权威替代来源；无法取得时删除相应字段和结论承诺。\\
WEB-0309 & B/\allowbreak{}07\_\allowbreak{}会计合规与制裁 & ACQUIRED/\allowbreak{}NARRATIVE\_\allowbreak{}CONTENT & EXTRACTED/\allowbreak{}SUPPLEMENT & country\_\allowbreak{}policy\_\allowbreak{}year | asset\_\allowbreak{}monthly\_\allowbreak{}return & 按项目用途复核；不足以单独支持核心结论。\\
WEB-0310 & B/\allowbreak{}07\_\allowbreak{}会计合规与制裁 & ACQUIRED/\allowbreak{}NARRATIVE\_\allowbreak{}CONTENT & EXTRACTED/\allowbreak{}SUPPLEMENT & country\_\allowbreak{}policy\_\allowbreak{}year | asset\_\allowbreak{}monthly\_\allowbreak{}return & 按项目用途复核；不足以单独支持核心结论。\\
WEB-0311 & B/\allowbreak{}08\_\allowbreak{}企业案例与公开披露 & NOT\_\allowbreak{}APPLICABLE/\allowbreak{}NOT\_\allowbreak{}ACQUIRED & FAILED/\allowbreak{}SUPPLEMENT & company\_\allowbreak{}overseas\_\allowbreak{}exposure | country\_\allowbreak{}event | portfolio\_\allowbreak{}scenario | case\_\allowbreak{}evid… & 由资料负责人合法登录、采购或申请授权；不得自动绕过。\\
WEB-0312 & B/\allowbreak{}05\_\allowbreak{}全球周期与历史危机 & NOT\_\allowbreak{}APPLICABLE/\allowbreak{}NOT\_\allowbreak{}ACQUIRED & FAILED/\allowbreak{}SUPPLEMENT & country\_\allowbreak{}monthly\_\allowbreak{}risk | country\_\allowbreak{}event | case\_\allowbreak{}evidence | historical\_\allowbreak{}crisis\_\allowbreak{}e… & 由资料负责人合法登录、采购或申请授权；不得自动绕过。\\
WEB-0313 & B/\allowbreak{}06\_\allowbreak{}黄金美元与储备配置 & NOT\_\allowbreak{}APPLICABLE/\allowbreak{}NOT\_\allowbreak{}ACQUIRED & FAILED/\allowbreak{}SUPPLEMENT & country\_\allowbreak{}monthly\_\allowbreak{}risk | country\_\allowbreak{}event | asset\_\allowbreak{}monthly\_\allowbreak{}return | portfolio\_\allowbreak{}sc… & 由资料负责人合法登录、采购或申请授权；不得自动绕过。\\
WEB-0314 & B/\allowbreak{}05\_\allowbreak{}全球周期与历史危机 & NOT\_\allowbreak{}APPLICABLE/\allowbreak{}NOT\_\allowbreak{}ACQUIRED & FAILED/\allowbreak{}SUPPLEMENT & country\_\allowbreak{}monthly\_\allowbreak{}risk | country\_\allowbreak{}event | case\_\allowbreak{}evidence | historical\_\allowbreak{}crisis\_\allowbreak{}e… & 由资料负责人合法登录、采购或申请授权；不得自动绕过。\\
WEB-0315 & B/\allowbreak{}07\_\allowbreak{}会计合规与制裁 & NOT\_\allowbreak{}APPLICABLE/\allowbreak{}NOT\_\allowbreak{}ACQUIRED & FAILED/\allowbreak{}SUPPLEMENT & country\_\allowbreak{}policy\_\allowbreak{}year | asset\_\allowbreak{}monthly\_\allowbreak{}return | case\_\allowbreak{}evidence & 由资料负责人合法登录、采购或申请授权；不得自动绕过。\\
WEB-0316 & B/\allowbreak{}06\_\allowbreak{}黄金美元与储备配置 & NOT\_\allowbreak{}APPLICABLE/\allowbreak{}NOT\_\allowbreak{}ACQUIRED & FAILED/\allowbreak{}SUPPLEMENT & asset\_\allowbreak{}monthly\_\allowbreak{}return | portfolio\_\allowbreak{}scenario & 由资料负责人合法登录、采购或申请授权；不得自动绕过。\\
WEB-0317 & B/\allowbreak{}06\_\allowbreak{}黄金美元与储备配置 & NOT\_\allowbreak{}APPLICABLE/\allowbreak{}NOT\_\allowbreak{}ACQUIRED & FAILED/\allowbreak{}SUPPLEMENT & asset\_\allowbreak{}monthly\_\allowbreak{}return | portfolio\_\allowbreak{}scenario & 由资料负责人合法登录、采购或申请授权；不得自动绕过。\\
WEB-0318 & B/\allowbreak{}05\_\allowbreak{}全球周期与历史危机 & ACQUIRED/\allowbreak{}STRUCTURED\_\allowbreak{}DATA & EXTRACTED/\allowbreak{}SUPPLEMENT & case\_\allowbreak{}evidence | historical\_\allowbreak{}crisis\_\allowbreak{}event | global\_\allowbreak{}cycle\_\allowbreak{}month & 按项目用途复核；不足以单独支持核心结论。\\
WEB-0319 & A/\allowbreak{}08\_\allowbreak{}企业案例与公开披露 & NOT\_\allowbreak{}APPLICABLE/\allowbreak{}NOT\_\allowbreak{}ACQUIRED & FAILED/\allowbreak{}EXCLUDED & company\_\allowbreak{}overseas\_\allowbreak{}exposure | country\_\allowbreak{}monthly\_\allowbreak{}risk | country\_\allowbreak{}event | portfol… & 仅保留元数据；正式分析改用许可允许的权威替代源。\\
WEB-0320 & D/\allowbreak{}05\_\allowbreak{}全球周期与历史危机 & NOT\_\allowbreak{}APPLICABLE/\allowbreak{}NOT\_\allowbreak{}ACQUIRED & FAILED/\allowbreak{}EXCLUDED & case\_\allowbreak{}evidence | historical\_\allowbreak{}crisis\_\allowbreak{}event | global\_\allowbreak{}cycle\_\allowbreak{}month & 仅保留元数据；正式分析改用许可允许的权威替代源。\\
WEB-0321 & D/\allowbreak{}08\_\allowbreak{}企业案例与公开披露 & NOT\_\allowbreak{}APPLICABLE/\allowbreak{}NOT\_\allowbreak{}ACQUIRED & FAILED/\allowbreak{}EXCLUDED & company\_\allowbreak{}overseas\_\allowbreak{}exposure | country\_\allowbreak{}monthly\_\allowbreak{}risk | country\_\allowbreak{}event | portfol… & 仅保留元数据；正式分析改用许可允许的权威替代源。\\
WEB-0322 & D/\allowbreak{}08\_\allowbreak{}企业案例与公开披露 & NOT\_\allowbreak{}APPLICABLE/\allowbreak{}NOT\_\allowbreak{}ACQUIRED & FAILED/\allowbreak{}EXCLUDED & company\_\allowbreak{}overseas\_\allowbreak{}exposure | country\_\allowbreak{}event | asset\_\allowbreak{}monthly\_\allowbreak{}return | portfol… & 仅保留元数据；正式分析改用许可允许的权威替代源。\\
WEB-0323 & D/\allowbreak{}05\_\allowbreak{}全球周期与历史危机 & NOT\_\allowbreak{}APPLICABLE/\allowbreak{}NOT\_\allowbreak{}ACQUIRED & FAILED/\allowbreak{}EXCLUDED & asset\_\allowbreak{}monthly\_\allowbreak{}return | case\_\allowbreak{}evidence | historical\_\allowbreak{}crisis\_\allowbreak{}event | global\_\allowbreak{}cy… & 仅保留元数据；正式分析改用许可允许的权威替代源。\\
WEB-0324 & D/\allowbreak{}04\_\allowbreak{}国别货币与中国出海 & NOT\_\allowbreak{}APPLICABLE/\allowbreak{}NOT\_\allowbreak{}ACQUIRED & FAILED/\allowbreak{}EXCLUDED & country\_\allowbreak{}monthly\_\allowbreak{}risk | country\_\allowbreak{}policy\_\allowbreak{}year | country\_\allowbreak{}event | case\_\allowbreak{}evidence & 仅保留元数据；正式分析改用许可允许的权威替代源。\\
WEB-0325 & D/\allowbreak{}08\_\allowbreak{}企业案例与公开披露 & NOT\_\allowbreak{}APPLICABLE/\allowbreak{}NOT\_\allowbreak{}ACQUIRED & FAILED/\allowbreak{}EXCLUDED & company\_\allowbreak{}overseas\_\allowbreak{}exposure | country\_\allowbreak{}event | asset\_\allowbreak{}monthly\_\allowbreak{}return | portfol… & 仅保留元数据；正式分析改用许可允许的权威替代源。\\
WEB-0326 & D/\allowbreak{}05\_\allowbreak{}全球周期与历史危机 & NOT\_\allowbreak{}APPLICABLE/\allowbreak{}NOT\_\allowbreak{}ACQUIRED & FAILED/\allowbreak{}EXCLUDED & asset\_\allowbreak{}monthly\_\allowbreak{}return | case\_\allowbreak{}evidence | historical\_\allowbreak{}crisis\_\allowbreak{}event | global\_\allowbreak{}cy… & 仅保留元数据；正式分析改用许可允许的权威替代源。\\
WEB-0327 & D/\allowbreak{}07\_\allowbreak{}会计合规与制裁 & NOT\_\allowbreak{}APPLICABLE/\allowbreak{}NOT\_\allowbreak{}ACQUIRED & FAILED/\allowbreak{}EXCLUDED & country\_\allowbreak{}policy\_\allowbreak{}year | asset\_\allowbreak{}monthly\_\allowbreak{}return | case\_\allowbreak{}evidence & 仅保留元数据；正式分析改用许可允许的权威替代源。\\
WEB-0328 & D/\allowbreak{}05\_\allowbreak{}全球周期与历史危机 & NOT\_\allowbreak{}APPLICABLE/\allowbreak{}NOT\_\allowbreak{}ACQUIRED & FAILED/\allowbreak{}EXCLUDED & asset\_\allowbreak{}monthly\_\allowbreak{}return | case\_\allowbreak{}evidence | historical\_\allowbreak{}crisis\_\allowbreak{}event | global\_\allowbreak{}cy… & 仅保留元数据；正式分析改用许可允许的权威替代源。\\
WEB-0329 & D/\allowbreak{}06\_\allowbreak{}黄金美元与储备配置 & NOT\_\allowbreak{}APPLICABLE/\allowbreak{}NOT\_\allowbreak{}ACQUIRED & FAILED/\allowbreak{}EXCLUDED & asset\_\allowbreak{}monthly\_\allowbreak{}return | portfolio\_\allowbreak{}scenario & 仅保留元数据；正式分析改用许可允许的权威替代源。\\
WEB-0330 & D/\allowbreak{}06\_\allowbreak{}黄金美元与储备配置 & NOT\_\allowbreak{}APPLICABLE/\allowbreak{}NOT\_\allowbreak{}ACQUIRED & FAILED/\allowbreak{}EXCLUDED & asset\_\allowbreak{}monthly\_\allowbreak{}return | portfolio\_\allowbreak{}scenario & 仅保留元数据；正式分析改用许可允许的权威替代源。\\
WEB-0331 & D/\allowbreak{}06\_\allowbreak{}黄金美元与储备配置 & NOT\_\allowbreak{}APPLICABLE/\allowbreak{}NOT\_\allowbreak{}ACQUIRED & FAILED/\allowbreak{}EXCLUDED & asset\_\allowbreak{}monthly\_\allowbreak{}return | portfolio\_\allowbreak{}scenario & 仅保留元数据；正式分析改用许可允许的权威替代源。\\
WEB-0332 & D/\allowbreak{}05\_\allowbreak{}全球周期与历史危机 & NOT\_\allowbreak{}APPLICABLE/\allowbreak{}NOT\_\allowbreak{}ACQUIRED & FAILED/\allowbreak{}EXCLUDED & country\_\allowbreak{}monthly\_\allowbreak{}risk | country\_\allowbreak{}event | case\_\allowbreak{}evidence | historical\_\allowbreak{}crisis\_\allowbreak{}e… & 仅保留元数据；正式分析改用许可允许的权威替代源。\\
WEB-0333 & D/\allowbreak{}05\_\allowbreak{}全球周期与历史危机 & NOT\_\allowbreak{}APPLICABLE/\allowbreak{}NOT\_\allowbreak{}ACQUIRED & FAILED/\allowbreak{}EXCLUDED & case\_\allowbreak{}evidence | historical\_\allowbreak{}crisis\_\allowbreak{}event | global\_\allowbreak{}cycle\_\allowbreak{}month & 仅保留元数据；正式分析改用许可允许的权威替代源。\\
WEB-0334 & D/\allowbreak{}06\_\allowbreak{}黄金美元与储备配置 & NOT\_\allowbreak{}APPLICABLE/\allowbreak{}NOT\_\allowbreak{}ACQUIRED & FAILED/\allowbreak{}EXCLUDED & asset\_\allowbreak{}monthly\_\allowbreak{}return | portfolio\_\allowbreak{}scenario & 仅保留元数据；正式分析改用许可允许的权威替代源。\\
WEB-0335 & B/\allowbreak{}06\_\allowbreak{}黄金美元与储备配置 & ACQUIRED/\allowbreak{}BLOCK\_\allowbreak{}PAGE & FAILED/\allowbreak{}REPLACED & asset\_\allowbreak{}monthly\_\allowbreak{}return | portfolio\_\allowbreak{}scenario & 使用WEB-0115官方开放PDF，不再使用拦截页。\\
WEB-0336 & B/\allowbreak{}04\_\allowbreak{}国别货币与中国出海 & ACQUIRED/\allowbreak{}BLOCK\_\allowbreak{}PAGE & FAILED/\allowbreak{}REPLACED & country\_\allowbreak{}monthly\_\allowbreak{}risk | country\_\allowbreak{}policy\_\allowbreak{}year & 使用WEB-0115官方开放PDF，不再使用拦截页。\\
WEB-0337 & D/\allowbreak{}07\_\allowbreak{}会计合规与制裁 & NOT\_\allowbreak{}APPLICABLE/\allowbreak{}NOT\_\allowbreak{}ACQUIRED & FAILED/\allowbreak{}EXCLUDED & country\_\allowbreak{}policy\_\allowbreak{}year | asset\_\allowbreak{}monthly\_\allowbreak{}return & 仅保留元数据；正式分析改用许可允许的权威替代源。\\
WEB-0338 & D/\allowbreak{}08\_\allowbreak{}企业案例与公开披露 & NOT\_\allowbreak{}APPLICABLE/\allowbreak{}NOT\_\allowbreak{}ACQUIRED & FAILED/\allowbreak{}EXCLUDED & company\_\allowbreak{}overseas\_\allowbreak{}exposure | country\_\allowbreak{}event | asset\_\allowbreak{}monthly\_\allowbreak{}return | portfol… & 仅保留元数据；正式分析改用许可允许的权威替代源。\\
WEB-0339 & D/\allowbreak{}06\_\allowbreak{}黄金美元与储备配置 & NOT\_\allowbreak{}APPLICABLE/\allowbreak{}NOT\_\allowbreak{}ACQUIRED & FAILED/\allowbreak{}EXCLUDED & asset\_\allowbreak{}monthly\_\allowbreak{}return | portfolio\_\allowbreak{}scenario & 仅保留元数据；正式分析改用许可允许的权威替代源。\\
WEB-0340 & D/\allowbreak{}06\_\allowbreak{}黄金美元与储备配置 & NOT\_\allowbreak{}APPLICABLE/\allowbreak{}NOT\_\allowbreak{}ACQUIRED & FAILED/\allowbreak{}EXCLUDED & asset\_\allowbreak{}monthly\_\allowbreak{}return | portfolio\_\allowbreak{}scenario & 仅保留元数据；正式分析改用许可允许的权威替代源。\\
WEB-0341 & D/\allowbreak{}05\_\allowbreak{}全球周期与历史危机 & NOT\_\allowbreak{}APPLICABLE/\allowbreak{}NOT\_\allowbreak{}ACQUIRED & FAILED/\allowbreak{}EXCLUDED & case\_\allowbreak{}evidence | historical\_\allowbreak{}crisis\_\allowbreak{}event | global\_\allowbreak{}cycle\_\allowbreak{}month & 仅保留元数据；正式分析改用许可允许的权威替代源。\\
WEB-0342 & A/\allowbreak{}04\_\allowbreak{}国别货币与中国出海 & NOT\_\allowbreak{}ACQUIRED/\allowbreak{}NOT\_\allowbreak{}ACQUIRED & FAILED/\allowbreak{}CORE & country\_\allowbreak{}monthly\_\allowbreak{}risk | country\_\allowbreak{}policy\_\allowbreak{}year & 补充权威替代来源；无法取得时删除相应字段和结论承诺。\\
WEB-0343 & A/\allowbreak{}06\_\allowbreak{}黄金美元与储备配置 & NOT\_\allowbreak{}ACQUIRED/\allowbreak{}NOT\_\allowbreak{}ACQUIRED & FAILED/\allowbreak{}CORE & asset\_\allowbreak{}monthly\_\allowbreak{}return | portfolio\_\allowbreak{}scenario & 补充权威替代来源；无法取得时删除相应字段和结论承诺。\\
WEB-0344 & A/\allowbreak{}06\_\allowbreak{}黄金美元与储备配置 & ACQUIRED/\allowbreak{}NARRATIVE\_\allowbreak{}CONTENT & EXTRACTED/\allowbreak{}CORE & source\_\allowbreak{}registry | country\_\allowbreak{}exposure | asset\_\allowbreak{}monthly\_\allowbreak{}return | portfolio\_\allowbreak{}scen… & 人工复核关键字段与source\_\allowbreak{}locator；通过后映射至正式source\_\allowbreak{}registry。\\
WEB-0345 & A/\allowbreak{}08\_\allowbreak{}企业案例与公开披露 & ACQUIRED/\allowbreak{}BLOCK\_\allowbreak{}PAGE & FAILED/\allowbreak{}EXCLUDED & company\_\allowbreak{}overseas\_\allowbreak{}exposure | country\_\allowbreak{}event | asset\_\allowbreak{}monthly\_\allowbreak{}return | portfol… & 补充权威替代来源；无法取得时删除相应字段和结论承诺。\\
WEB-0346 & D/\allowbreak{}08\_\allowbreak{}企业案例与公开披露 & NOT\_\allowbreak{}APPLICABLE/\allowbreak{}NOT\_\allowbreak{}ACQUIRED & FAILED/\allowbreak{}EXCLUDED & company\_\allowbreak{}overseas\_\allowbreak{}exposure | country\_\allowbreak{}event | portfolio\_\allowbreak{}scenario | case\_\allowbreak{}evid… & 仅保留元数据；正式分析改用许可允许的权威替代源。\\
WEB-0347 & D/\allowbreak{}08\_\allowbreak{}企业案例与公开披露 & NOT\_\allowbreak{}APPLICABLE/\allowbreak{}NOT\_\allowbreak{}ACQUIRED & FAILED/\allowbreak{}EXCLUDED & company\_\allowbreak{}overseas\_\allowbreak{}exposure | country\_\allowbreak{}monthly\_\allowbreak{}risk | country\_\allowbreak{}event | portfol… & 仅保留元数据；正式分析改用许可允许的权威替代源。\\
WEB-0348 & B/\allowbreak{}05\_\allowbreak{}全球周期与历史危机 & ACQUIRED/\allowbreak{}NARRATIVE\_\allowbreak{}CONTENT & EXTRACTED/\allowbreak{}SUPPLEMENT & asset\_\allowbreak{}monthly\_\allowbreak{}return | historical\_\allowbreak{}crisis\_\allowbreak{}event | global\_\allowbreak{}cycle\_\allowbreak{}month & 按项目用途复核；不足以单独支持核心结论。\\
WEB-0349 & B/\allowbreak{}05\_\allowbreak{}全球周期与历史危机 & ACQUIRED/\allowbreak{}NARRATIVE\_\allowbreak{}CONTENT & EXTRACTED/\allowbreak{}SUPPLEMENT & source\_\allowbreak{}registry | historical\_\allowbreak{}crisis\_\allowbreak{}event | global\_\allowbreak{}cycle\_\allowbreak{}month & 按项目用途复核；不足以单独支持核心结论。\\
WEB-0350 & B/\allowbreak{}05\_\allowbreak{}全球周期与历史危机 & ACQUIRED/\allowbreak{}NARRATIVE\_\allowbreak{}CONTENT & EXTRACTED/\allowbreak{}SUPPLEMENT & source\_\allowbreak{}registry | historical\_\allowbreak{}crisis\_\allowbreak{}event | global\_\allowbreak{}cycle\_\allowbreak{}month & 按项目用途复核；不足以单独支持核心结论。\\
WEB-0351 & B/\allowbreak{}05\_\allowbreak{}全球周期与历史危机 & ACQUIRED/\allowbreak{}MIXED\_\allowbreak{}CONTENT & EXTRACTED/\allowbreak{}SUPPLEMENT & source\_\allowbreak{}registry | historical\_\allowbreak{}crisis\_\allowbreak{}event | global\_\allowbreak{}cycle\_\allowbreak{}month & 按项目用途复核；不足以单独支持核心结论。\\
WEB-0352 & B/\allowbreak{}08\_\allowbreak{}企业案例与公开披露 & ACQUIRED/\allowbreak{}DOCUMENT\_\allowbreak{}LOCAL & NOT\_\allowbreak{}REQUIRED/\allowbreak{}SUPPLEMENT & company\_\allowbreak{}overseas\_\allowbreak{}exposure | country\_\allowbreak{}event | asset\_\allowbreak{}monthly\_\allowbreak{}return | portfol… & 按项目用途复核；不足以单独支持核心结论。\\
WEB-0353 & B/\allowbreak{}05\_\allowbreak{}全球周期与历史危机 & NOT\_\allowbreak{}APPLICABLE/\allowbreak{}NOT\_\allowbreak{}ACQUIRED & FAILED/\allowbreak{}EXCLUDED & source\_\allowbreak{}registry | historical\_\allowbreak{}crisis\_\allowbreak{}event | global\_\allowbreak{}cycle\_\allowbreak{}month & 仅保留元数据；正式分析改用许可允许的权威替代源。\\
WEB-0354 & D/\allowbreak{}08\_\allowbreak{}企业案例与公开披露 & NOT\_\allowbreak{}APPLICABLE/\allowbreak{}NOT\_\allowbreak{}ACQUIRED & FAILED/\allowbreak{}EXCLUDED & company\_\allowbreak{}overseas\_\allowbreak{}exposure | country\_\allowbreak{}event | asset\_\allowbreak{}monthly\_\allowbreak{}return | portfol… & 由资料负责人合法登录、采购或申请授权；不得自动绕过。\\
WEB-0355 & D/\allowbreak{}08\_\allowbreak{}企业案例与公开披露 & NOT\_\allowbreak{}APPLICABLE/\allowbreak{}NOT\_\allowbreak{}ACQUIRED & FAILED/\allowbreak{}EXCLUDED & company\_\allowbreak{}overseas\_\allowbreak{}exposure | country\_\allowbreak{}monthly\_\allowbreak{}risk | country\_\allowbreak{}event | portfol… & 仅保留元数据；正式分析改用许可允许的权威替代源。\\
WEB-0356 & D/\allowbreak{}08\_\allowbreak{}企业案例与公开披露 & NOT\_\allowbreak{}APPLICABLE/\allowbreak{}NOT\_\allowbreak{}ACQUIRED & FAILED/\allowbreak{}EXCLUDED & company\_\allowbreak{}overseas\_\allowbreak{}exposure | country\_\allowbreak{}event | portfolio\_\allowbreak{}scenario | case\_\allowbreak{}evid… & 仅保留元数据；正式分析改用许可允许的权威替代源。\\
WEB-0357 & A/\allowbreak{}08\_\allowbreak{}企业案例与公开披露 & ACQUIRED/\allowbreak{}STRUCTURED\_\allowbreak{}DATA & EXTRACTED/\allowbreak{}CORE & company\_\allowbreak{}overseas\_\allowbreak{}exposure | country\_\allowbreak{}event | asset\_\allowbreak{}monthly\_\allowbreak{}return | portfol… & 人工复核关键字段与source\_\allowbreak{}locator；通过后映射至正式source\_\allowbreak{}registry。\\
WEB-0358 & A/\allowbreak{}08\_\allowbreak{}企业案例与公开披露 & ACQUIRED/\allowbreak{}STRUCTURED\_\allowbreak{}DATA & EXTRACTED/\allowbreak{}REPLACED & company\_\allowbreak{}overseas\_\allowbreak{}exposure | country\_\allowbreak{}event | asset\_\allowbreak{}monthly\_\allowbreak{}return | portfol… & 使用替代来源：WEB-0357\\
WEB-0359 & A/\allowbreak{}08\_\allowbreak{}企业案例与公开披露 & ACQUIRED/\allowbreak{}LOW\_\allowbreak{}CONTENT & PARTIAL/\allowbreak{}EXCLUDED & company\_\allowbreak{}overseas\_\allowbreak{}exposure | country\_\allowbreak{}event | asset\_\allowbreak{}monthly\_\allowbreak{}return | portfol… & 补充权威替代来源；无法取得时删除相应字段和结论承诺。\\
WEB-0360 & A/\allowbreak{}08\_\allowbreak{}企业案例与公开披露 & NOT\_\allowbreak{}APPLICABLE/\allowbreak{}NOT\_\allowbreak{}ACQUIRED & FAILED/\allowbreak{}EXCLUDED & company\_\allowbreak{}overseas\_\allowbreak{}exposure | country\_\allowbreak{}monthly\_\allowbreak{}risk | country\_\allowbreak{}event | portfol… & 仅保留元数据；正式分析改用许可允许的权威替代源。\\
WEB-0361 & D/\allowbreak{}08\_\allowbreak{}企业案例与公开披露 & NOT\_\allowbreak{}APPLICABLE/\allowbreak{}NOT\_\allowbreak{}ACQUIRED & FAILED/\allowbreak{}EXCLUDED & company\_\allowbreak{}overseas\_\allowbreak{}exposure | country\_\allowbreak{}event | asset\_\allowbreak{}monthly\_\allowbreak{}return | portfol… & 仅保留元数据；正式分析改用许可允许的权威替代源。\\
WEB-0362 & D/\allowbreak{}08\_\allowbreak{}企业案例与公开披露 & NOT\_\allowbreak{}APPLICABLE/\allowbreak{}NOT\_\allowbreak{}ACQUIRED & FAILED/\allowbreak{}EXCLUDED & company\_\allowbreak{}overseas\_\allowbreak{}exposure | country\_\allowbreak{}event | asset\_\allowbreak{}monthly\_\allowbreak{}return | portfol… & 仅保留元数据；正式分析改用许可允许的权威替代源。\\
WEB-0363 & D/\allowbreak{}06\_\allowbreak{}黄金美元与储备配置 & ACQUIRED/\allowbreak{}STRUCTURED\_\allowbreak{}DATA & EXTRACTED/\allowbreak{}EXCLUDED & asset\_\allowbreak{}monthly\_\allowbreak{}return | portfolio\_\allowbreak{}scenario & 补充权威替代来源；无法取得时删除相应字段和结论承诺。\\
WEB-0364 & D/\allowbreak{}06\_\allowbreak{}黄金美元与储备配置 & ACQUIRED/\allowbreak{}MIXED\_\allowbreak{}CONTENT & EXTRACTED/\allowbreak{}EXCLUDED & asset\_\allowbreak{}monthly\_\allowbreak{}return | portfolio\_\allowbreak{}scenario & 补充权威替代来源；无法取得时删除相应字段和结论承诺。\\
WEB-0365 & D/\allowbreak{}05\_\allowbreak{}全球周期与历史危机 & NOT\_\allowbreak{}APPLICABLE/\allowbreak{}NOT\_\allowbreak{}ACQUIRED & FAILED/\allowbreak{}EXCLUDED & asset\_\allowbreak{}monthly\_\allowbreak{}return | case\_\allowbreak{}evidence | historical\_\allowbreak{}crisis\_\allowbreak{}event | global\_\allowbreak{}cy… & 仅保留元数据；正式分析改用许可允许的权威替代源。\\
WEB-0366 & D/\allowbreak{}05\_\allowbreak{}全球周期与历史危机 & NOT\_\allowbreak{}APPLICABLE/\allowbreak{}NOT\_\allowbreak{}ACQUIRED & FAILED/\allowbreak{}EXCLUDED & asset\_\allowbreak{}monthly\_\allowbreak{}return | case\_\allowbreak{}evidence | historical\_\allowbreak{}crisis\_\allowbreak{}event | global\_\allowbreak{}cy… & 仅保留元数据；正式分析改用许可允许的权威替代源。\\
WEB-0367 & D/\allowbreak{}08\_\allowbreak{}企业案例与公开披露 & NOT\_\allowbreak{}APPLICABLE/\allowbreak{}NOT\_\allowbreak{}ACQUIRED & FAILED/\allowbreak{}EXCLUDED & company\_\allowbreak{}overseas\_\allowbreak{}exposure | country\_\allowbreak{}event | asset\_\allowbreak{}monthly\_\allowbreak{}return | portfol… & 仅保留元数据；正式分析改用许可允许的权威替代源。\\
WEB-0368 & D/\allowbreak{}08\_\allowbreak{}企业案例与公开披露 & ACQUIRED/\allowbreak{}DOCUMENT\_\allowbreak{}LOCAL & NOT\_\allowbreak{}REQUIRED/\allowbreak{}EXCLUDED & company\_\allowbreak{}overseas\_\allowbreak{}exposure | country\_\allowbreak{}event | asset\_\allowbreak{}monthly\_\allowbreak{}return | portfol… & 补充权威替代来源；无法取得时删除相应字段和结论承诺。\\
WEB-0369 & D/\allowbreak{}07\_\allowbreak{}会计合规与制裁 & NOT\_\allowbreak{}APPLICABLE/\allowbreak{}NOT\_\allowbreak{}ACQUIRED & FAILED/\allowbreak{}EXCLUDED & country\_\allowbreak{}policy\_\allowbreak{}year | country\_\allowbreak{}event | asset\_\allowbreak{}monthly\_\allowbreak{}return | case\_\allowbreak{}evidence & 由资料负责人合法登录、采购或申请授权；不得自动绕过。\\
WEB-0370 & D/\allowbreak{}07\_\allowbreak{}会计合规与制裁 & NOT\_\allowbreak{}APPLICABLE/\allowbreak{}NOT\_\allowbreak{}ACQUIRED & FAILED/\allowbreak{}EXCLUDED & country\_\allowbreak{}policy\_\allowbreak{}year | asset\_\allowbreak{}monthly\_\allowbreak{}return & 由资料负责人合法登录、采购或申请授权；不得自动绕过。\\
WEB-0371 & D/\allowbreak{}05\_\allowbreak{}全球周期与历史危机 & NOT\_\allowbreak{}APPLICABLE/\allowbreak{}NOT\_\allowbreak{}ACQUIRED & FAILED/\allowbreak{}EXCLUDED & case\_\allowbreak{}evidence | historical\_\allowbreak{}crisis\_\allowbreak{}event | global\_\allowbreak{}cycle\_\allowbreak{}month & 仅保留元数据；正式分析改用许可允许的权威替代源。\\
WEB-0372 & D/\allowbreak{}08\_\allowbreak{}企业案例与公开披露 & NOT\_\allowbreak{}APPLICABLE/\allowbreak{}NOT\_\allowbreak{}ACQUIRED & FAILED/\allowbreak{}EXCLUDED & company\_\allowbreak{}overseas\_\allowbreak{}exposure | country\_\allowbreak{}event | portfolio\_\allowbreak{}scenario | case\_\allowbreak{}evid… & 仅保留元数据；正式分析改用许可允许的权威替代源。\\
WEB-0373 & C/\allowbreak{}06\_\allowbreak{}黄金美元与储备配置 & NOT\_\allowbreak{}APPLICABLE/\allowbreak{}NOT\_\allowbreak{}ACQUIRED & FAILED/\allowbreak{}EXCLUDED & country\_\allowbreak{}monthly\_\allowbreak{}risk | country\_\allowbreak{}event | asset\_\allowbreak{}monthly\_\allowbreak{}return | portfolio\_\allowbreak{}sc… & 仅保留元数据；正式分析改用许可允许的权威替代源。\\
WEB-0374 & C/\allowbreak{}07\_\allowbreak{}会计合规与制裁 & NOT\_\allowbreak{}APPLICABLE/\allowbreak{}NOT\_\allowbreak{}ACQUIRED & FAILED/\allowbreak{}EXCLUDED & company\_\allowbreak{}overseas\_\allowbreak{}exposure | country\_\allowbreak{}policy\_\allowbreak{}year | country\_\allowbreak{}event | asset\_\allowbreak{}mo… & 仅保留元数据；正式分析改用许可允许的权威替代源。\\
WEB-0375 & C/\allowbreak{}06\_\allowbreak{}黄金美元与储备配置 & NOT\_\allowbreak{}APPLICABLE/\allowbreak{}NOT\_\allowbreak{}ACQUIRED & FAILED/\allowbreak{}EXCLUDED & country\_\allowbreak{}monthly\_\allowbreak{}risk | country\_\allowbreak{}event | asset\_\allowbreak{}monthly\_\allowbreak{}return | portfolio\_\allowbreak{}sc… & 仅保留元数据；正式分析改用许可允许的权威替代源。\\
WEB-0376 & C/\allowbreak{}07\_\allowbreak{}会计合规与制裁 & NOT\_\allowbreak{}APPLICABLE/\allowbreak{}NOT\_\allowbreak{}ACQUIRED & FAILED/\allowbreak{}EXCLUDED & company\_\allowbreak{}overseas\_\allowbreak{}exposure | country\_\allowbreak{}monthly\_\allowbreak{}risk | country\_\allowbreak{}policy\_\allowbreak{}year | c… & 仅保留元数据；正式分析改用许可允许的权威替代源。\\
WEB-0377 & C/\allowbreak{}04\_\allowbreak{}国别货币与中国出海 & NOT\_\allowbreak{}APPLICABLE/\allowbreak{}NOT\_\allowbreak{}ACQUIRED & FAILED/\allowbreak{}EXCLUDED & country\_\allowbreak{}monthly\_\allowbreak{}risk | country\_\allowbreak{}policy\_\allowbreak{}year | country\_\allowbreak{}event | case\_\allowbreak{}evidence & 仅保留元数据；正式分析改用许可允许的权威替代源。\\
WEB-0378 & C/\allowbreak{}08\_\allowbreak{}企业案例与公开披露 & NOT\_\allowbreak{}APPLICABLE/\allowbreak{}NOT\_\allowbreak{}ACQUIRED & FAILED/\allowbreak{}EXCLUDED & company\_\allowbreak{}overseas\_\allowbreak{}exposure | country\_\allowbreak{}event | asset\_\allowbreak{}monthly\_\allowbreak{}return | portfol… & 仅保留元数据；正式分析改用许可允许的权威替代源。\\
WEB-0379 & C/\allowbreak{}07\_\allowbreak{}会计合规与制裁 & NOT\_\allowbreak{}APPLICABLE/\allowbreak{}NOT\_\allowbreak{}ACQUIRED & FAILED/\allowbreak{}EXCLUDED & country\_\allowbreak{}policy\_\allowbreak{}year | country\_\allowbreak{}event | asset\_\allowbreak{}monthly\_\allowbreak{}return | case\_\allowbreak{}evidence & 仅保留元数据；正式分析改用许可允许的权威替代源。\\
WEB-0380 & D/\allowbreak{}08\_\allowbreak{}企业案例与公开披露 & NOT\_\allowbreak{}APPLICABLE/\allowbreak{}NOT\_\allowbreak{}ACQUIRED & FAILED/\allowbreak{}EXCLUDED & company\_\allowbreak{}overseas\_\allowbreak{}exposure | country\_\allowbreak{}event | asset\_\allowbreak{}monthly\_\allowbreak{}return | portfol… & 仅保留元数据；正式分析改用许可允许的权威替代源。\\
WEB-0381 & D/\allowbreak{}08\_\allowbreak{}企业案例与公开披露 & NOT\_\allowbreak{}APPLICABLE/\allowbreak{}NOT\_\allowbreak{}ACQUIRED & FAILED/\allowbreak{}EXCLUDED & company\_\allowbreak{}overseas\_\allowbreak{}exposure | country\_\allowbreak{}event | asset\_\allowbreak{}monthly\_\allowbreak{}return | portfol… & 仅保留元数据；正式分析改用许可允许的权威替代源。\\
WEB-0382 & A/\allowbreak{}05\_\allowbreak{}全球周期与历史危机 & ACQUIRED/\allowbreak{}NARRATIVE\_\allowbreak{}CONTENT & EXTRACTED/\allowbreak{}CORE & case\_\allowbreak{}evidence | historical\_\allowbreak{}crisis\_\allowbreak{}event | global\_\allowbreak{}cycle\_\allowbreak{}month & 人工复核关键字段与source\_\allowbreak{}locator；通过后映射至正式source\_\allowbreak{}registry。\\
WEB-0383 & D/\allowbreak{}05\_\allowbreak{}全球周期与历史危机 & NOT\_\allowbreak{}APPLICABLE/\allowbreak{}NOT\_\allowbreak{}ACQUIRED & FAILED/\allowbreak{}EXCLUDED & case\_\allowbreak{}evidence | historical\_\allowbreak{}crisis\_\allowbreak{}event | global\_\allowbreak{}cycle\_\allowbreak{}month & 仅保留元数据；正式分析改用许可允许的权威替代源。\\
WEB-0384 & D/\allowbreak{}08\_\allowbreak{}企业案例与公开披露 & ACQUIRED/\allowbreak{}DOCUMENT\_\allowbreak{}LOCAL & NOT\_\allowbreak{}REQUIRED/\allowbreak{}EXCLUDED & company\_\allowbreak{}overseas\_\allowbreak{}exposure | country\_\allowbreak{}event | asset\_\allowbreak{}monthly\_\allowbreak{}return | portfol… & 补充权威替代来源；无法取得时删除相应字段和结论承诺。\\
WEB-0385 & A/\allowbreak{}08\_\allowbreak{}企业案例与公开披露 & ACQUIRED/\allowbreak{}MIXED\_\allowbreak{}CONTENT & EXTRACTED/\allowbreak{}CORE & company\_\allowbreak{}overseas\_\allowbreak{}exposure | country\_\allowbreak{}event | portfolio\_\allowbreak{}scenario | case\_\allowbreak{}evid… & 人工复核关键字段与source\_\allowbreak{}locator；通过后映射至正式source\_\allowbreak{}registry。\\
WEB-0386 & A/\allowbreak{}08\_\allowbreak{}企业案例与公开披露 & ACQUIRED/\allowbreak{}MIXED\_\allowbreak{}CONTENT & EXTRACTED/\allowbreak{}CORE & company\_\allowbreak{}overseas\_\allowbreak{}exposure | country\_\allowbreak{}event | portfolio\_\allowbreak{}scenario | case\_\allowbreak{}evid… & 人工复核关键字段与source\_\allowbreak{}locator；通过后映射至正式source\_\allowbreak{}registry。\\
WEB-0387 & A/\allowbreak{}08\_\allowbreak{}企业案例与公开披露 & ACQUIRED/\allowbreak{}MIXED\_\allowbreak{}CONTENT & EXTRACTED/\allowbreak{}CORE & company\_\allowbreak{}overseas\_\allowbreak{}exposure | country\_\allowbreak{}monthly\_\allowbreak{}risk | country\_\allowbreak{}event | portfol… & 人工复核关键字段与source\_\allowbreak{}locator；通过后映射至正式source\_\allowbreak{}registry。\\
WEB-0388 & A/\allowbreak{}08\_\allowbreak{}企业案例与公开披露 & ACQUIRED/\allowbreak{}MIXED\_\allowbreak{}CONTENT & EXTRACTED/\allowbreak{}CORE & company\_\allowbreak{}overseas\_\allowbreak{}exposure | country\_\allowbreak{}event | portfolio\_\allowbreak{}scenario | case\_\allowbreak{}evid… & 人工复核关键字段与source\_\allowbreak{}locator；通过后映射至正式source\_\allowbreak{}registry。\\
WEB-0389 & A/\allowbreak{}08\_\allowbreak{}企业案例与公开披露 & ACQUIRED/\allowbreak{}NARRATIVE\_\allowbreak{}CONTENT & EXTRACTED/\allowbreak{}CORE & company\_\allowbreak{}overseas\_\allowbreak{}exposure | country\_\allowbreak{}event | asset\_\allowbreak{}monthly\_\allowbreak{}return | portfol… & 人工复核关键字段与source\_\allowbreak{}locator；通过后映射至正式source\_\allowbreak{}registry。\\
WEB-0390 & A/\allowbreak{}08\_\allowbreak{}企业案例与公开披露 & ACQUIRED/\allowbreak{}MIXED\_\allowbreak{}CONTENT & EXTRACTED/\allowbreak{}CORE & company\_\allowbreak{}overseas\_\allowbreak{}exposure | country\_\allowbreak{}event | asset\_\allowbreak{}monthly\_\allowbreak{}return | portfol… & 人工复核关键字段与source\_\allowbreak{}locator；通过后映射至正式source\_\allowbreak{}registry。\\
WEB-0391 & A/\allowbreak{}08\_\allowbreak{}企业案例与公开披露 & ACQUIRED/\allowbreak{}MIXED\_\allowbreak{}CONTENT & EXTRACTED/\allowbreak{}CORE & company\_\allowbreak{}overseas\_\allowbreak{}exposure | country\_\allowbreak{}event | asset\_\allowbreak{}monthly\_\allowbreak{}return | portfol… & 人工复核关键字段与source\_\allowbreak{}locator；通过后映射至正式source\_\allowbreak{}registry。\\
WEB-0392 & B/\allowbreak{}07\_\allowbreak{}会计合规与制裁 & ACQUIRED/\allowbreak{}STRUCTURED\_\allowbreak{}DATA & EXTRACTED/\allowbreak{}REPLACED & source\_\allowbreak{}registry | company\_\allowbreak{}overseas\_\allowbreak{}exposure | country\_\allowbreak{}policy\_\allowbreak{}year | asset\_\allowbreak{}… & 使用替代来源：WEB-0013\\
WEB-0393 & B/\allowbreak{}08\_\allowbreak{}企业案例与公开披露 & ACQUIRED/\allowbreak{}STRUCTURED\_\allowbreak{}DATA & EXTRACTED/\allowbreak{}SUPPLEMENT & company\_\allowbreak{}overseas\_\allowbreak{}exposure | country\_\allowbreak{}event | asset\_\allowbreak{}monthly\_\allowbreak{}return | portfol… & 按项目用途复核；不足以单独支持核心结论。\\
WEB-0394 & B/\allowbreak{}06\_\allowbreak{}黄金美元与储备配置 & ACQUIRED/\allowbreak{}DOWNLOAD\_\allowbreak{}LANDING & EXTRACTED/\allowbreak{}SUPPLEMENT & asset\_\allowbreak{}monthly\_\allowbreak{}return | portfolio\_\allowbreak{}scenario & 按项目用途复核；不足以单独支持核心结论。\\
WEB-0395 & B/\allowbreak{}06\_\allowbreak{}黄金美元与储备配置 & ACQUIRED/\allowbreak{}DOWNLOAD\_\allowbreak{}LANDING & EXTRACTED/\allowbreak{}SUPPLEMENT & asset\_\allowbreak{}monthly\_\allowbreak{}return | portfolio\_\allowbreak{}scenario & 按项目用途复核；不足以单独支持核心结论。\\
WEB-0396 & B/\allowbreak{}06\_\allowbreak{}黄金美元与储备配置 & ACQUIRED/\allowbreak{}MIXED\_\allowbreak{}CONTENT & EXTRACTED/\allowbreak{}SUPPLEMENT & asset\_\allowbreak{}monthly\_\allowbreak{}return | portfolio\_\allowbreak{}scenario & 按项目用途复核；不足以单独支持核心结论。\\
WEB-0397 & B/\allowbreak{}06\_\allowbreak{}黄金美元与储备配置 & ACQUIRED/\allowbreak{}MIXED\_\allowbreak{}CONTENT & EXTRACTED/\allowbreak{}SUPPLEMENT & asset\_\allowbreak{}monthly\_\allowbreak{}return | portfolio\_\allowbreak{}scenario & 按项目用途复核；不足以单独支持核心结论。\\
WEB-0398 & B/\allowbreak{}06\_\allowbreak{}黄金美元与储备配置 & ACQUIRED/\allowbreak{}MIXED\_\allowbreak{}CONTENT & EXTRACTED/\allowbreak{}SUPPLEMENT & asset\_\allowbreak{}monthly\_\allowbreak{}return | portfolio\_\allowbreak{}scenario & 按项目用途复核；不足以单独支持核心结论。\\
WEB-0399 & B/\allowbreak{}06\_\allowbreak{}黄金美元与储备配置 & ACQUIRED/\allowbreak{}MIXED\_\allowbreak{}CONTENT & EXTRACTED/\allowbreak{}SUPPLEMENT & asset\_\allowbreak{}monthly\_\allowbreak{}return | portfolio\_\allowbreak{}scenario & 按项目用途复核；不足以单独支持核心结论。\\
WEB-0400 & B/\allowbreak{}06\_\allowbreak{}黄金美元与储备配置 & ACQUIRED/\allowbreak{}MIXED\_\allowbreak{}CONTENT & EXTRACTED/\allowbreak{}SUPPLEMENT & asset\_\allowbreak{}monthly\_\allowbreak{}return | portfolio\_\allowbreak{}scenario & 按项目用途复核；不足以单独支持核心结论。\\
WEB-0401 & B/\allowbreak{}06\_\allowbreak{}黄金美元与储备配置 & ACQUIRED/\allowbreak{}MIXED\_\allowbreak{}CONTENT & EXTRACTED/\allowbreak{}SUPPLEMENT & asset\_\allowbreak{}monthly\_\allowbreak{}return | portfolio\_\allowbreak{}scenario & 按项目用途复核；不足以单独支持核心结论。\\
WEB-0402 & B/\allowbreak{}06\_\allowbreak{}黄金美元与储备配置 & ACQUIRED/\allowbreak{}MIXED\_\allowbreak{}CONTENT & EXTRACTED/\allowbreak{}SUPPLEMENT & asset\_\allowbreak{}monthly\_\allowbreak{}return | portfolio\_\allowbreak{}scenario & 按项目用途复核；不足以单独支持核心结论。\\
WEB-0403 & B/\allowbreak{}06\_\allowbreak{}黄金美元与储备配置 & ACQUIRED/\allowbreak{}MIXED\_\allowbreak{}CONTENT & EXTRACTED/\allowbreak{}SUPPLEMENT & asset\_\allowbreak{}monthly\_\allowbreak{}return | portfolio\_\allowbreak{}scenario & 按项目用途复核；不足以单独支持核心结论。\\
WEB-0404 & B/\allowbreak{}06\_\allowbreak{}黄金美元与储备配置 & ACQUIRED/\allowbreak{}MIXED\_\allowbreak{}CONTENT & EXTRACTED/\allowbreak{}SUPPLEMENT & asset\_\allowbreak{}monthly\_\allowbreak{}return | portfolio\_\allowbreak{}scenario & 按项目用途复核；不足以单独支持核心结论。\\
WEB-0405 & B/\allowbreak{}06\_\allowbreak{}黄金美元与储备配置 & ACQUIRED/\allowbreak{}MIXED\_\allowbreak{}CONTENT & EXTRACTED/\allowbreak{}SUPPLEMENT & asset\_\allowbreak{}monthly\_\allowbreak{}return | portfolio\_\allowbreak{}scenario & 按项目用途复核；不足以单独支持核心结论。\\
WEB-0406 & B/\allowbreak{}06\_\allowbreak{}黄金美元与储备配置 & ACQUIRED/\allowbreak{}MIXED\_\allowbreak{}CONTENT & EXTRACTED/\allowbreak{}SUPPLEMENT & asset\_\allowbreak{}monthly\_\allowbreak{}return | portfolio\_\allowbreak{}scenario & 按项目用途复核；不足以单独支持核心结论。\\
WEB-0407 & B/\allowbreak{}06\_\allowbreak{}黄金美元与储备配置 & ACQUIRED/\allowbreak{}MIXED\_\allowbreak{}CONTENT & EXTRACTED/\allowbreak{}SUPPLEMENT & asset\_\allowbreak{}monthly\_\allowbreak{}return | portfolio\_\allowbreak{}scenario & 按项目用途复核；不足以单独支持核心结论。\\
WEB-0408 & B/\allowbreak{}06\_\allowbreak{}黄金美元与储备配置 & ACQUIRED/\allowbreak{}MIXED\_\allowbreak{}CONTENT & EXTRACTED/\allowbreak{}SUPPLEMENT & asset\_\allowbreak{}monthly\_\allowbreak{}return | portfolio\_\allowbreak{}scenario & 按项目用途复核；不足以单独支持核心结论。\\
WEB-0409 & B/\allowbreak{}06\_\allowbreak{}黄金美元与储备配置 & ACQUIRED/\allowbreak{}MIXED\_\allowbreak{}CONTENT & EXTRACTED/\allowbreak{}SUPPLEMENT & asset\_\allowbreak{}monthly\_\allowbreak{}return | portfolio\_\allowbreak{}scenario & 按项目用途复核；不足以单独支持核心结论。\\
WEB-0410 & B/\allowbreak{}06\_\allowbreak{}黄金美元与储备配置 & ACQUIRED/\allowbreak{}MIXED\_\allowbreak{}CONTENT & EXTRACTED/\allowbreak{}SUPPLEMENT & asset\_\allowbreak{}monthly\_\allowbreak{}return | portfolio\_\allowbreak{}scenario & 按项目用途复核；不足以单独支持核心结论。\\
WEB-0411 & B/\allowbreak{}06\_\allowbreak{}黄金美元与储备配置 & ACQUIRED/\allowbreak{}MIXED\_\allowbreak{}CONTENT & EXTRACTED/\allowbreak{}SUPPLEMENT & asset\_\allowbreak{}monthly\_\allowbreak{}return | portfolio\_\allowbreak{}scenario & 按项目用途复核；不足以单独支持核心结论。\\
WEB-0412 & B/\allowbreak{}06\_\allowbreak{}黄金美元与储备配置 & ACQUIRED/\allowbreak{}MIXED\_\allowbreak{}CONTENT & EXTRACTED/\allowbreak{}SUPPLEMENT & asset\_\allowbreak{}monthly\_\allowbreak{}return | portfolio\_\allowbreak{}scenario & 按项目用途复核；不足以单独支持核心结论。\\
WEB-0413 & B/\allowbreak{}06\_\allowbreak{}黄金美元与储备配置 & ACQUIRED/\allowbreak{}MIXED\_\allowbreak{}CONTENT & EXTRACTED/\allowbreak{}SUPPLEMENT & asset\_\allowbreak{}monthly\_\allowbreak{}return | portfolio\_\allowbreak{}scenario & 按项目用途复核；不足以单独支持核心结论。\\
WEB-0414 & B/\allowbreak{}06\_\allowbreak{}黄金美元与储备配置 & ACQUIRED/\allowbreak{}MIXED\_\allowbreak{}CONTENT & EXTRACTED/\allowbreak{}SUPPLEMENT & asset\_\allowbreak{}monthly\_\allowbreak{}return | portfolio\_\allowbreak{}scenario & 按项目用途复核；不足以单独支持核心结论。\\
WEB-0415 & B/\allowbreak{}06\_\allowbreak{}黄金美元与储备配置 & ACQUIRED/\allowbreak{}MIXED\_\allowbreak{}CONTENT & EXTRACTED/\allowbreak{}SUPPLEMENT & asset\_\allowbreak{}monthly\_\allowbreak{}return | portfolio\_\allowbreak{}scenario & 按项目用途复核；不足以单独支持核心结论。\\
WEB-0416 & B/\allowbreak{}06\_\allowbreak{}黄金美元与储备配置 & ACQUIRED/\allowbreak{}MIXED\_\allowbreak{}CONTENT & EXTRACTED/\allowbreak{}SUPPLEMENT & asset\_\allowbreak{}monthly\_\allowbreak{}return | portfolio\_\allowbreak{}scenario & 按项目用途复核；不足以单独支持核心结论。\\
WEB-0417 & B/\allowbreak{}06\_\allowbreak{}黄金美元与储备配置 & ACQUIRED/\allowbreak{}MIXED\_\allowbreak{}CONTENT & EXTRACTED/\allowbreak{}SUPPLEMENT & asset\_\allowbreak{}monthly\_\allowbreak{}return | portfolio\_\allowbreak{}scenario & 按项目用途复核；不足以单独支持核心结论。\\
WEB-0418 & B/\allowbreak{}06\_\allowbreak{}黄金美元与储备配置 & ACQUIRED/\allowbreak{}MIXED\_\allowbreak{}CONTENT & EXTRACTED/\allowbreak{}SUPPLEMENT & asset\_\allowbreak{}monthly\_\allowbreak{}return | portfolio\_\allowbreak{}scenario & 按项目用途复核；不足以单独支持核心结论。\\
WEB-0419 & B/\allowbreak{}06\_\allowbreak{}黄金美元与储备配置 & ACQUIRED/\allowbreak{}MIXED\_\allowbreak{}CONTENT & EXTRACTED/\allowbreak{}SUPPLEMENT & asset\_\allowbreak{}monthly\_\allowbreak{}return | portfolio\_\allowbreak{}scenario & 按项目用途复核；不足以单独支持核心结论。\\
WEB-0420 & B/\allowbreak{}06\_\allowbreak{}黄金美元与储备配置 & ACQUIRED/\allowbreak{}MIXED\_\allowbreak{}CONTENT & EXTRACTED/\allowbreak{}SUPPLEMENT & asset\_\allowbreak{}monthly\_\allowbreak{}return | portfolio\_\allowbreak{}scenario & 按项目用途复核；不足以单独支持核心结论。\\
WEB-0421 & B/\allowbreak{}06\_\allowbreak{}黄金美元与储备配置 & ACQUIRED/\allowbreak{}MIXED\_\allowbreak{}CONTENT & EXTRACTED/\allowbreak{}SUPPLEMENT & asset\_\allowbreak{}monthly\_\allowbreak{}return | portfolio\_\allowbreak{}scenario & 按项目用途复核；不足以单独支持核心结论。\\
WEB-0422 & B/\allowbreak{}06\_\allowbreak{}黄金美元与储备配置 & ACQUIRED/\allowbreak{}MIXED\_\allowbreak{}CONTENT & EXTRACTED/\allowbreak{}SUPPLEMENT & asset\_\allowbreak{}monthly\_\allowbreak{}return | portfolio\_\allowbreak{}scenario & 按项目用途复核；不足以单独支持核心结论。\\
WEB-0423 & B/\allowbreak{}06\_\allowbreak{}黄金美元与储备配置 & ACQUIRED/\allowbreak{}MIXED\_\allowbreak{}CONTENT & EXTRACTED/\allowbreak{}SUPPLEMENT & asset\_\allowbreak{}monthly\_\allowbreak{}return | portfolio\_\allowbreak{}scenario & 按项目用途复核；不足以单独支持核心结论。\\
WEB-0424 & B/\allowbreak{}06\_\allowbreak{}黄金美元与储备配置 & ACQUIRED/\allowbreak{}MIXED\_\allowbreak{}CONTENT & EXTRACTED/\allowbreak{}SUPPLEMENT & asset\_\allowbreak{}monthly\_\allowbreak{}return | portfolio\_\allowbreak{}scenario & 按项目用途复核；不足以单独支持核心结论。\\
WEB-0425 & B/\allowbreak{}06\_\allowbreak{}黄金美元与储备配置 & ACQUIRED/\allowbreak{}MIXED\_\allowbreak{}CONTENT & EXTRACTED/\allowbreak{}SUPPLEMENT & asset\_\allowbreak{}monthly\_\allowbreak{}return | portfolio\_\allowbreak{}scenario & 按项目用途复核；不足以单独支持核心结论。\\
WEB-0426 & B/\allowbreak{}06\_\allowbreak{}黄金美元与储备配置 & ACQUIRED/\allowbreak{}MIXED\_\allowbreak{}CONTENT & EXTRACTED/\allowbreak{}SUPPLEMENT & asset\_\allowbreak{}monthly\_\allowbreak{}return | portfolio\_\allowbreak{}scenario & 按项目用途复核；不足以单独支持核心结论。\\
WEB-0427 & B/\allowbreak{}06\_\allowbreak{}黄金美元与储备配置 & ACQUIRED/\allowbreak{}MIXED\_\allowbreak{}CONTENT & EXTRACTED/\allowbreak{}SUPPLEMENT & asset\_\allowbreak{}monthly\_\allowbreak{}return | portfolio\_\allowbreak{}scenario & 按项目用途复核；不足以单独支持核心结论。\\
WEB-0428 & B/\allowbreak{}06\_\allowbreak{}黄金美元与储备配置 & ACQUIRED/\allowbreak{}MIXED\_\allowbreak{}CONTENT & EXTRACTED/\allowbreak{}SUPPLEMENT & asset\_\allowbreak{}monthly\_\allowbreak{}return | portfolio\_\allowbreak{}scenario & 按项目用途复核；不足以单独支持核心结论。\\
WEB-0429 & B/\allowbreak{}06\_\allowbreak{}黄金美元与储备配置 & ACQUIRED/\allowbreak{}MIXED\_\allowbreak{}CONTENT & EXTRACTED/\allowbreak{}SUPPLEMENT & asset\_\allowbreak{}monthly\_\allowbreak{}return | portfolio\_\allowbreak{}scenario & 按项目用途复核；不足以单独支持核心结论。\\
WEB-0430 & B/\allowbreak{}06\_\allowbreak{}黄金美元与储备配置 & ACQUIRED/\allowbreak{}MIXED\_\allowbreak{}CONTENT & EXTRACTED/\allowbreak{}SUPPLEMENT & asset\_\allowbreak{}monthly\_\allowbreak{}return | portfolio\_\allowbreak{}scenario & 按项目用途复核；不足以单独支持核心结论。\\
WEB-0431 & B/\allowbreak{}06\_\allowbreak{}黄金美元与储备配置 & ACQUIRED/\allowbreak{}MIXED\_\allowbreak{}CONTENT & EXTRACTED/\allowbreak{}SUPPLEMENT & asset\_\allowbreak{}monthly\_\allowbreak{}return | portfolio\_\allowbreak{}scenario & 按项目用途复核；不足以单独支持核心结论。\\
WEB-0432 & B/\allowbreak{}06\_\allowbreak{}黄金美元与储备配置 & ACQUIRED/\allowbreak{}MIXED\_\allowbreak{}CONTENT & EXTRACTED/\allowbreak{}SUPPLEMENT & asset\_\allowbreak{}monthly\_\allowbreak{}return | portfolio\_\allowbreak{}scenario & 按项目用途复核；不足以单独支持核心结论。\\
WEB-0433 & B/\allowbreak{}06\_\allowbreak{}黄金美元与储备配置 & ACQUIRED/\allowbreak{}MIXED\_\allowbreak{}CONTENT & EXTRACTED/\allowbreak{}SUPPLEMENT & asset\_\allowbreak{}monthly\_\allowbreak{}return | portfolio\_\allowbreak{}scenario & 按项目用途复核；不足以单独支持核心结论。\\
WEB-0434 & B/\allowbreak{}06\_\allowbreak{}黄金美元与储备配置 & ACQUIRED/\allowbreak{}MIXED\_\allowbreak{}CONTENT & EXTRACTED/\allowbreak{}SUPPLEMENT & asset\_\allowbreak{}monthly\_\allowbreak{}return | portfolio\_\allowbreak{}scenario & 按项目用途复核；不足以单独支持核心结论。\\
WEB-0435 & B/\allowbreak{}06\_\allowbreak{}黄金美元与储备配置 & ACQUIRED/\allowbreak{}MIXED\_\allowbreak{}CONTENT & EXTRACTED/\allowbreak{}SUPPLEMENT & asset\_\allowbreak{}monthly\_\allowbreak{}return | portfolio\_\allowbreak{}scenario & 按项目用途复核；不足以单独支持核心结论。\\
WEB-0436 & B/\allowbreak{}06\_\allowbreak{}黄金美元与储备配置 & ACQUIRED/\allowbreak{}MIXED\_\allowbreak{}CONTENT & EXTRACTED/\allowbreak{}SUPPLEMENT & asset\_\allowbreak{}monthly\_\allowbreak{}return | portfolio\_\allowbreak{}scenario & 按项目用途复核；不足以单独支持核心结论。\\
WEB-0437 & B/\allowbreak{}06\_\allowbreak{}黄金美元与储备配置 & ACQUIRED/\allowbreak{}MIXED\_\allowbreak{}CONTENT & EXTRACTED/\allowbreak{}SUPPLEMENT & asset\_\allowbreak{}monthly\_\allowbreak{}return | portfolio\_\allowbreak{}scenario & 按项目用途复核；不足以单独支持核心结论。\\
WEB-0438 & B/\allowbreak{}06\_\allowbreak{}黄金美元与储备配置 & ACQUIRED/\allowbreak{}MIXED\_\allowbreak{}CONTENT & EXTRACTED/\allowbreak{}SUPPLEMENT & asset\_\allowbreak{}monthly\_\allowbreak{}return | portfolio\_\allowbreak{}scenario & 按项目用途复核；不足以单独支持核心结论。\\
WEB-0439 & B/\allowbreak{}06\_\allowbreak{}黄金美元与储备配置 & ACQUIRED/\allowbreak{}MIXED\_\allowbreak{}CONTENT & EXTRACTED/\allowbreak{}SUPPLEMENT & asset\_\allowbreak{}monthly\_\allowbreak{}return | portfolio\_\allowbreak{}scenario & 按项目用途复核；不足以单独支持核心结论。\\
WEB-0440 & B/\allowbreak{}06\_\allowbreak{}黄金美元与储备配置 & ACQUIRED/\allowbreak{}MIXED\_\allowbreak{}CONTENT & EXTRACTED/\allowbreak{}SUPPLEMENT & asset\_\allowbreak{}monthly\_\allowbreak{}return | portfolio\_\allowbreak{}scenario & 按项目用途复核；不足以单独支持核心结论。\\
WEB-0441 & B/\allowbreak{}06\_\allowbreak{}黄金美元与储备配置 & ACQUIRED/\allowbreak{}MIXED\_\allowbreak{}CONTENT & EXTRACTED/\allowbreak{}SUPPLEMENT & asset\_\allowbreak{}monthly\_\allowbreak{}return | portfolio\_\allowbreak{}scenario & 按项目用途复核；不足以单独支持核心结论。\\
WEB-0442 & B/\allowbreak{}06\_\allowbreak{}黄金美元与储备配置 & ACQUIRED/\allowbreak{}MIXED\_\allowbreak{}CONTENT & EXTRACTED/\allowbreak{}SUPPLEMENT & asset\_\allowbreak{}monthly\_\allowbreak{}return | portfolio\_\allowbreak{}scenario & 按项目用途复核；不足以单独支持核心结论。\\
WEB-0443 & B/\allowbreak{}06\_\allowbreak{}黄金美元与储备配置 & ACQUIRED/\allowbreak{}MIXED\_\allowbreak{}CONTENT & EXTRACTED/\allowbreak{}SUPPLEMENT & asset\_\allowbreak{}monthly\_\allowbreak{}return | portfolio\_\allowbreak{}scenario & 按项目用途复核；不足以单独支持核心结论。\\
WEB-0444 & B/\allowbreak{}06\_\allowbreak{}黄金美元与储备配置 & ACQUIRED/\allowbreak{}MIXED\_\allowbreak{}CONTENT & EXTRACTED/\allowbreak{}SUPPLEMENT & asset\_\allowbreak{}monthly\_\allowbreak{}return | portfolio\_\allowbreak{}scenario & 按项目用途复核；不足以单独支持核心结论。\\
WEB-0445 & B/\allowbreak{}06\_\allowbreak{}黄金美元与储备配置 & ACQUIRED/\allowbreak{}MIXED\_\allowbreak{}CONTENT & EXTRACTED/\allowbreak{}SUPPLEMENT & asset\_\allowbreak{}monthly\_\allowbreak{}return | portfolio\_\allowbreak{}scenario & 按项目用途复核；不足以单独支持核心结论。\\
WEB-0446 & B/\allowbreak{}06\_\allowbreak{}黄金美元与储备配置 & ACQUIRED/\allowbreak{}MIXED\_\allowbreak{}CONTENT & EXTRACTED/\allowbreak{}SUPPLEMENT & asset\_\allowbreak{}monthly\_\allowbreak{}return | portfolio\_\allowbreak{}scenario & 按项目用途复核；不足以单独支持核心结论。\\
WEB-0447 & B/\allowbreak{}06\_\allowbreak{}黄金美元与储备配置 & ACQUIRED/\allowbreak{}MIXED\_\allowbreak{}CONTENT & EXTRACTED/\allowbreak{}SUPPLEMENT & asset\_\allowbreak{}monthly\_\allowbreak{}return | portfolio\_\allowbreak{}scenario & 按项目用途复核；不足以单独支持核心结论。\\
WEB-0448 & B/\allowbreak{}06\_\allowbreak{}黄金美元与储备配置 & ACQUIRED/\allowbreak{}MIXED\_\allowbreak{}CONTENT & EXTRACTED/\allowbreak{}SUPPLEMENT & asset\_\allowbreak{}monthly\_\allowbreak{}return | portfolio\_\allowbreak{}scenario & 按项目用途复核；不足以单独支持核心结论。\\
WEB-0449 & B/\allowbreak{}06\_\allowbreak{}黄金美元与储备配置 & ACQUIRED/\allowbreak{}MIXED\_\allowbreak{}CONTENT & EXTRACTED/\allowbreak{}SUPPLEMENT & asset\_\allowbreak{}monthly\_\allowbreak{}return | portfolio\_\allowbreak{}scenario & 按项目用途复核；不足以单独支持核心结论。\\
WEB-0450 & B/\allowbreak{}06\_\allowbreak{}黄金美元与储备配置 & ACQUIRED/\allowbreak{}MIXED\_\allowbreak{}CONTENT & EXTRACTED/\allowbreak{}SUPPLEMENT & asset\_\allowbreak{}monthly\_\allowbreak{}return | portfolio\_\allowbreak{}scenario & 按项目用途复核；不足以单独支持核心结论。\\
WEB-0451 & B/\allowbreak{}06\_\allowbreak{}黄金美元与储备配置 & ACQUIRED/\allowbreak{}MIXED\_\allowbreak{}CONTENT & EXTRACTED/\allowbreak{}SUPPLEMENT & asset\_\allowbreak{}monthly\_\allowbreak{}return | portfolio\_\allowbreak{}scenario & 按项目用途复核；不足以单独支持核心结论。\\
WEB-0452 & B/\allowbreak{}06\_\allowbreak{}黄金美元与储备配置 & ACQUIRED/\allowbreak{}MIXED\_\allowbreak{}CONTENT & EXTRACTED/\allowbreak{}SUPPLEMENT & asset\_\allowbreak{}monthly\_\allowbreak{}return | portfolio\_\allowbreak{}scenario & 按项目用途复核；不足以单独支持核心结论。\\
WEB-0453 & B/\allowbreak{}06\_\allowbreak{}黄金美元与储备配置 & ACQUIRED/\allowbreak{}MIXED\_\allowbreak{}CONTENT & EXTRACTED/\allowbreak{}SUPPLEMENT & asset\_\allowbreak{}monthly\_\allowbreak{}return | portfolio\_\allowbreak{}scenario & 按项目用途复核；不足以单独支持核心结论。\\
WEB-0454 & B/\allowbreak{}06\_\allowbreak{}黄金美元与储备配置 & ACQUIRED/\allowbreak{}MIXED\_\allowbreak{}CONTENT & EXTRACTED/\allowbreak{}SUPPLEMENT & asset\_\allowbreak{}monthly\_\allowbreak{}return | portfolio\_\allowbreak{}scenario & 按项目用途复核；不足以单独支持核心结论。\\
WEB-0455 & B/\allowbreak{}06\_\allowbreak{}黄金美元与储备配置 & ACQUIRED/\allowbreak{}MIXED\_\allowbreak{}CONTENT & EXTRACTED/\allowbreak{}SUPPLEMENT & asset\_\allowbreak{}monthly\_\allowbreak{}return | portfolio\_\allowbreak{}scenario & 按项目用途复核；不足以单独支持核心结论。\\
WEB-0456 & B/\allowbreak{}06\_\allowbreak{}黄金美元与储备配置 & ACQUIRED/\allowbreak{}MIXED\_\allowbreak{}CONTENT & EXTRACTED/\allowbreak{}SUPPLEMENT & asset\_\allowbreak{}monthly\_\allowbreak{}return | portfolio\_\allowbreak{}scenario & 按项目用途复核；不足以单独支持核心结论。\\
WEB-0457 & B/\allowbreak{}06\_\allowbreak{}黄金美元与储备配置 & ACQUIRED/\allowbreak{}MIXED\_\allowbreak{}CONTENT & EXTRACTED/\allowbreak{}SUPPLEMENT & asset\_\allowbreak{}monthly\_\allowbreak{}return | portfolio\_\allowbreak{}scenario & 按项目用途复核；不足以单独支持核心结论。\\
WEB-0458 & D/\allowbreak{}05\_\allowbreak{}全球周期与历史危机 & ACQUIRED/\allowbreak{}DOCUMENT\_\allowbreak{}LOCAL & NOT\_\allowbreak{}REQUIRED/\allowbreak{}EXCLUDED & asset\_\allowbreak{}monthly\_\allowbreak{}return | case\_\allowbreak{}evidence | historical\_\allowbreak{}crisis\_\allowbreak{}event | global\_\allowbreak{}cy… & 补充权威替代来源；无法取得时删除相应字段和结论承诺。\\
WEB-0459 & D/\allowbreak{}08\_\allowbreak{}企业案例与公开披露 & ACQUIRED/\allowbreak{}DOCUMENT\_\allowbreak{}LOCAL & NOT\_\allowbreak{}REQUIRED/\allowbreak{}EXCLUDED & company\_\allowbreak{}overseas\_\allowbreak{}exposure | country\_\allowbreak{}event | asset\_\allowbreak{}monthly\_\allowbreak{}return | portfol… & 补充权威替代来源；无法取得时删除相应字段和结论承诺。\\
WEB-0460 & C/\allowbreak{}08\_\allowbreak{}企业案例与公开披露 & NOT\_\allowbreak{}APPLICABLE/\allowbreak{}NOT\_\allowbreak{}ACQUIRED & FAILED/\allowbreak{}EXCLUDED & company\_\allowbreak{}overseas\_\allowbreak{}exposure | country\_\allowbreak{}event | asset\_\allowbreak{}monthly\_\allowbreak{}return | portfol… & 仅保留元数据；正式分析改用许可允许的权威替代源。\\
WEB-0461 & D/\allowbreak{}05\_\allowbreak{}全球周期与历史危机 & NOT\_\allowbreak{}APPLICABLE/\allowbreak{}NOT\_\allowbreak{}ACQUIRED & FAILED/\allowbreak{}EXCLUDED & case\_\allowbreak{}evidence | historical\_\allowbreak{}crisis\_\allowbreak{}event | global\_\allowbreak{}cycle\_\allowbreak{}month & 仅保留元数据；正式分析改用许可允许的权威替代源。\\
WEB-0462 & D/\allowbreak{}09\_\allowbreak{}Smartbi与比赛平台 & NOT\_\allowbreak{}APPLICABLE/\allowbreak{}NOT\_\allowbreak{}ACQUIRED & FAILED/\allowbreak{}EXCLUDED & source\_\allowbreak{}registry & 仅保留元数据；正式分析改用许可允许的权威替代源。\\
WEB-0463 & D/\allowbreak{}09\_\allowbreak{}Smartbi与比赛平台 & NOT\_\allowbreak{}APPLICABLE/\allowbreak{}NOT\_\allowbreak{}ACQUIRED & FAILED/\allowbreak{}EXCLUDED & source\_\allowbreak{}registry & 仅保留元数据；正式分析改用许可允许的权威替代源。\\
WEB-0464 & D/\allowbreak{}09\_\allowbreak{}Smartbi与比赛平台 & NOT\_\allowbreak{}APPLICABLE/\allowbreak{}NOT\_\allowbreak{}ACQUIRED & FAILED/\allowbreak{}EXCLUDED & source\_\allowbreak{}registry & 仅保留元数据；正式分析改用许可允许的权威替代源。\\
WEB-0465 & D/\allowbreak{}08\_\allowbreak{}企业案例与公开披露 & NOT\_\allowbreak{}APPLICABLE/\allowbreak{}NOT\_\allowbreak{}ACQUIRED & FAILED/\allowbreak{}EXCLUDED & company\_\allowbreak{}overseas\_\allowbreak{}exposure | country\_\allowbreak{}event | portfolio\_\allowbreak{}scenario | case\_\allowbreak{}evid… & 仅保留元数据；正式分析改用许可允许的权威替代源。\\
WEB-0466 & D/\allowbreak{}08\_\allowbreak{}企业案例与公开披露 & NOT\_\allowbreak{}APPLICABLE/\allowbreak{}NOT\_\allowbreak{}ACQUIRED & FAILED/\allowbreak{}EXCLUDED & company\_\allowbreak{}overseas\_\allowbreak{}exposure | country\_\allowbreak{}event | asset\_\allowbreak{}monthly\_\allowbreak{}return | portfol… & 仅保留元数据；正式分析改用许可允许的权威替代源。\\
WEB-0467 & D/\allowbreak{}08\_\allowbreak{}企业案例与公开披露 & NOT\_\allowbreak{}APPLICABLE/\allowbreak{}NOT\_\allowbreak{}ACQUIRED & FAILED/\allowbreak{}EXCLUDED & company\_\allowbreak{}overseas\_\allowbreak{}exposure | country\_\allowbreak{}event | asset\_\allowbreak{}monthly\_\allowbreak{}return | portfol… & 仅保留元数据；正式分析改用许可允许的权威替代源。\\
WEB-0468 & D/\allowbreak{}05\_\allowbreak{}全球周期与历史危机 & NOT\_\allowbreak{}APPLICABLE/\allowbreak{}NOT\_\allowbreak{}ACQUIRED & FAILED/\allowbreak{}EXCLUDED & asset\_\allowbreak{}monthly\_\allowbreak{}return | case\_\allowbreak{}evidence | historical\_\allowbreak{}crisis\_\allowbreak{}event | global\_\allowbreak{}cy… & 仅保留元数据；正式分析改用许可允许的权威替代源。\\
WEB-0469 & D/\allowbreak{}07\_\allowbreak{}会计合规与制裁 & ACQUIRED/\allowbreak{}MIXED\_\allowbreak{}CONTENT & EXTRACTED/\allowbreak{}EXCLUDED & company\_\allowbreak{}overseas\_\allowbreak{}exposure | country\_\allowbreak{}policy\_\allowbreak{}year | asset\_\allowbreak{}monthly\_\allowbreak{}return | c… & 补充权威替代来源；无法取得时删除相应字段和结论承诺。\\
WEB-0470 & D/\allowbreak{}05\_\allowbreak{}全球周期与历史危机 & NOT\_\allowbreak{}APPLICABLE/\allowbreak{}NOT\_\allowbreak{}ACQUIRED & FAILED/\allowbreak{}EXCLUDED & case\_\allowbreak{}evidence | historical\_\allowbreak{}crisis\_\allowbreak{}event | global\_\allowbreak{}cycle\_\allowbreak{}month & 仅保留元数据；正式分析改用许可允许的权威替代源。\\
WEB-0471 & A/\allowbreak{}08\_\allowbreak{}企业案例与公开披露 & ACQUIRED/\allowbreak{}MIXED\_\allowbreak{}CONTENT & EXTRACTED/\allowbreak{}CORE & company\_\allowbreak{}overseas\_\allowbreak{}exposure | country\_\allowbreak{}event | asset\_\allowbreak{}monthly\_\allowbreak{}return | portfol… & 人工复核关键字段与source\_\allowbreak{}locator；通过后映射至正式source\_\allowbreak{}registry。\\
WEB-0472 & A/\allowbreak{}08\_\allowbreak{}企业案例与公开披露 & ACQUIRED/\allowbreak{}MIXED\_\allowbreak{}CONTENT & EXTRACTED/\allowbreak{}CORE & company\_\allowbreak{}overseas\_\allowbreak{}exposure | country\_\allowbreak{}event | asset\_\allowbreak{}monthly\_\allowbreak{}return | portfol… & 人工复核关键字段与source\_\allowbreak{}locator；通过后映射至正式source\_\allowbreak{}registry。\\
WEB-0473 & D/\allowbreak{}05\_\allowbreak{}全球周期与历史危机 & NOT\_\allowbreak{}APPLICABLE/\allowbreak{}NOT\_\allowbreak{}ACQUIRED & FAILED/\allowbreak{}EXCLUDED & case\_\allowbreak{}evidence | historical\_\allowbreak{}crisis\_\allowbreak{}event | global\_\allowbreak{}cycle\_\allowbreak{}month & 仅保留元数据；正式分析改用许可允许的权威替代源。\\
WEB-0474 & A/\allowbreak{}04\_\allowbreak{}国别货币与中国出海 & NOT\_\allowbreak{}APPLICABLE/\allowbreak{}NOT\_\allowbreak{}ACQUIRED & FAILED/\allowbreak{}EXCLUDED & company\_\allowbreak{}overseas\_\allowbreak{}exposure | country\_\allowbreak{}monthly\_\allowbreak{}risk | country\_\allowbreak{}policy\_\allowbreak{}year | c… & 仅保留元数据；正式分析改用许可允许的权威替代源。\\
WEB-0475 & D/\allowbreak{}07\_\allowbreak{}会计合规与制裁 & NOT\_\allowbreak{}APPLICABLE/\allowbreak{}NOT\_\allowbreak{}ACQUIRED & FAILED/\allowbreak{}EXCLUDED & country\_\allowbreak{}monthly\_\allowbreak{}risk | country\_\allowbreak{}policy\_\allowbreak{}year | country\_\allowbreak{}event | asset\_\allowbreak{}monthly… & 仅保留元数据；正式分析改用许可允许的权威替代源。\\
WEB-0476 & A/\allowbreak{}04\_\allowbreak{}国别货币与中国出海 & NOT\_\allowbreak{}APPLICABLE/\allowbreak{}NOT\_\allowbreak{}ACQUIRED & FAILED/\allowbreak{}EXCLUDED & country\_\allowbreak{}monthly\_\allowbreak{}risk | country\_\allowbreak{}policy\_\allowbreak{}year | country\_\allowbreak{}event | asset\_\allowbreak{}monthly… & 仅保留元数据；正式分析改用许可允许的权威替代源。\\
WEB-0477 & D/\allowbreak{}08\_\allowbreak{}企业案例与公开披露 & NOT\_\allowbreak{}APPLICABLE/\allowbreak{}NOT\_\allowbreak{}ACQUIRED & FAILED/\allowbreak{}EXCLUDED & company\_\allowbreak{}overseas\_\allowbreak{}exposure | country\_\allowbreak{}event | asset\_\allowbreak{}monthly\_\allowbreak{}return | portfol… & 仅保留元数据；正式分析改用许可允许的权威替代源。\\
WEB-0478 & D/\allowbreak{}08\_\allowbreak{}企业案例与公开披露 & NOT\_\allowbreak{}APPLICABLE/\allowbreak{}NOT\_\allowbreak{}ACQUIRED & FAILED/\allowbreak{}EXCLUDED & company\_\allowbreak{}overseas\_\allowbreak{}exposure | country\_\allowbreak{}event | asset\_\allowbreak{}monthly\_\allowbreak{}return | portfol… & 仅保留元数据；正式分析改用许可允许的权威替代源。\\
WEB-0479 & C/\allowbreak{}08\_\allowbreak{}企业案例与公开披露 & NOT\_\allowbreak{}APPLICABLE/\allowbreak{}NOT\_\allowbreak{}ACQUIRED & FAILED/\allowbreak{}EXCLUDED & company\_\allowbreak{}overseas\_\allowbreak{}exposure | country\_\allowbreak{}monthly\_\allowbreak{}risk | country\_\allowbreak{}event | portfol… & 仅保留元数据；正式分析改用许可允许的权威替代源。\\
WEB-0480 & C/\allowbreak{}08\_\allowbreak{}企业案例与公开披露 & NOT\_\allowbreak{}APPLICABLE/\allowbreak{}NOT\_\allowbreak{}ACQUIRED & FAILED/\allowbreak{}EXCLUDED & company\_\allowbreak{}overseas\_\allowbreak{}exposure | country\_\allowbreak{}event | asset\_\allowbreak{}monthly\_\allowbreak{}return | portfol… & 仅保留元数据；正式分析改用许可允许的权威替代源。\\
WEB-0481 & A/\allowbreak{}04\_\allowbreak{}国别货币与中国出海 & NOT\_\allowbreak{}APPLICABLE/\allowbreak{}NOT\_\allowbreak{}ACQUIRED & FAILED/\allowbreak{}EXCLUDED & country\_\allowbreak{}monthly\_\allowbreak{}risk | country\_\allowbreak{}policy\_\allowbreak{}year | country\_\allowbreak{}event | case\_\allowbreak{}evidence & 仅保留元数据；正式分析改用许可允许的权威替代源。\\
WEB-0482 & A/\allowbreak{}07\_\allowbreak{}会计合规与制裁 & NOT\_\allowbreak{}APPLICABLE/\allowbreak{}NOT\_\allowbreak{}ACQUIRED & FAILED/\allowbreak{}EXCLUDED & country\_\allowbreak{}policy\_\allowbreak{}year | asset\_\allowbreak{}monthly\_\allowbreak{}return | case\_\allowbreak{}evidence & 仅保留元数据；正式分析改用许可允许的权威替代源。\\
WEB-0483 & D/\allowbreak{}05\_\allowbreak{}全球周期与历史危机 & NOT\_\allowbreak{}APPLICABLE/\allowbreak{}NOT\_\allowbreak{}ACQUIRED & FAILED/\allowbreak{}EXCLUDED & asset\_\allowbreak{}monthly\_\allowbreak{}return | case\_\allowbreak{}evidence | historical\_\allowbreak{}crisis\_\allowbreak{}event | global\_\allowbreak{}cy… & 仅保留元数据；正式分析改用许可允许的权威替代源。\\
WEB-0484 & D/\allowbreak{}07\_\allowbreak{}会计合规与制裁 & NOT\_\allowbreak{}APPLICABLE/\allowbreak{}NOT\_\allowbreak{}ACQUIRED & FAILED/\allowbreak{}EXCLUDED & country\_\allowbreak{}policy\_\allowbreak{}year | asset\_\allowbreak{}monthly\_\allowbreak{}return & 仅保留元数据；正式分析改用许可允许的权威替代源。\\
WEB-0485 & A/\allowbreak{}06\_\allowbreak{}黄金美元与储备配置 & NOT\_\allowbreak{}APPLICABLE/\allowbreak{}NOT\_\allowbreak{}ACQUIRED & FAILED/\allowbreak{}EXCLUDED & asset\_\allowbreak{}monthly\_\allowbreak{}return | portfolio\_\allowbreak{}scenario & 仅保留元数据；正式分析改用许可允许的权威替代源。\\
WEB-0486 & A/\allowbreak{}06\_\allowbreak{}黄金美元与储备配置 & NOT\_\allowbreak{}APPLICABLE/\allowbreak{}NOT\_\allowbreak{}ACQUIRED & FAILED/\allowbreak{}EXCLUDED & asset\_\allowbreak{}monthly\_\allowbreak{}return | portfolio\_\allowbreak{}scenario & 仅保留元数据；正式分析改用许可允许的权威替代源。\\
WEB-0487 & A/\allowbreak{}06\_\allowbreak{}黄金美元与储备配置 & NOT\_\allowbreak{}APPLICABLE/\allowbreak{}NOT\_\allowbreak{}ACQUIRED & FAILED/\allowbreak{}EXCLUDED & asset\_\allowbreak{}monthly\_\allowbreak{}return | portfolio\_\allowbreak{}scenario & 仅保留元数据；正式分析改用许可允许的权威替代源。\\
WEB-0488 & D/\allowbreak{}06\_\allowbreak{}黄金美元与储备配置 & NOT\_\allowbreak{}APPLICABLE/\allowbreak{}NOT\_\allowbreak{}ACQUIRED & FAILED/\allowbreak{}EXCLUDED & asset\_\allowbreak{}monthly\_\allowbreak{}return | portfolio\_\allowbreak{}scenario & 仅保留元数据；正式分析改用许可允许的权威替代源。\\
WEB-0489 & D/\allowbreak{}06\_\allowbreak{}黄金美元与储备配置 & NOT\_\allowbreak{}APPLICABLE/\allowbreak{}NOT\_\allowbreak{}ACQUIRED & FAILED/\allowbreak{}EXCLUDED & asset\_\allowbreak{}monthly\_\allowbreak{}return | portfolio\_\allowbreak{}scenario & 仅保留元数据；正式分析改用许可允许的权威替代源。\\
WEB-0490 & D/\allowbreak{}06\_\allowbreak{}黄金美元与储备配置 & NOT\_\allowbreak{}APPLICABLE/\allowbreak{}NOT\_\allowbreak{}ACQUIRED & FAILED/\allowbreak{}EXCLUDED & asset\_\allowbreak{}monthly\_\allowbreak{}return | portfolio\_\allowbreak{}scenario & 仅保留元数据；正式分析改用许可允许的权威替代源。\\
WEB-0491 & D/\allowbreak{}07\_\allowbreak{}会计合规与制裁 & NOT\_\allowbreak{}APPLICABLE/\allowbreak{}NOT\_\allowbreak{}ACQUIRED & FAILED/\allowbreak{}EXCLUDED & country\_\allowbreak{}monthly\_\allowbreak{}risk | country\_\allowbreak{}policy\_\allowbreak{}year | country\_\allowbreak{}event | asset\_\allowbreak{}monthly… & 仅保留元数据；正式分析改用许可允许的权威替代源。\\
WEB-0492 & B/\allowbreak{}06\_\allowbreak{}黄金美元与储备配置 & ACQUIRED/\allowbreak{}NARRATIVE\_\allowbreak{}CONTENT & EXTRACTED/\allowbreak{}SUPPLEMENT & asset\_\allowbreak{}monthly\_\allowbreak{}return | portfolio\_\allowbreak{}scenario & 按项目用途复核；不足以单独支持核心结论。\\
WEB-0493 & B/\allowbreak{}07\_\allowbreak{}会计合规与制裁 & ACQUIRED/\allowbreak{}NARRATIVE\_\allowbreak{}CONTENT & EXTRACTED/\allowbreak{}SUPPLEMENT & country\_\allowbreak{}policy\_\allowbreak{}year | asset\_\allowbreak{}monthly\_\allowbreak{}return | case\_\allowbreak{}evidence & 按项目用途复核；不足以单独支持核心结论。\\
WEB-0494 & B/\allowbreak{}06\_\allowbreak{}黄金美元与储备配置 & ACQUIRED/\allowbreak{}NARRATIVE\_\allowbreak{}CONTENT & EXTRACTED/\allowbreak{}SUPPLEMENT & country\_\allowbreak{}monthly\_\allowbreak{}risk | country\_\allowbreak{}event | asset\_\allowbreak{}monthly\_\allowbreak{}return | portfolio\_\allowbreak{}sc… & 按项目用途复核；不足以单独支持核心结论。\\
WEB-0495 & B/\allowbreak{}04\_\allowbreak{}国别货币与中国出海 & ACQUIRED/\allowbreak{}NARRATIVE\_\allowbreak{}CONTENT & EXTRACTED/\allowbreak{}SUPPLEMENT & country\_\allowbreak{}monthly\_\allowbreak{}risk | country\_\allowbreak{}policy\_\allowbreak{}year | country\_\allowbreak{}event | case\_\allowbreak{}evidence & 按项目用途复核；不足以单独支持核心结论。\\
WEB-0496 & B/\allowbreak{}06\_\allowbreak{}黄金美元与储备配置 & ACQUIRED/\allowbreak{}NARRATIVE\_\allowbreak{}CONTENT & EXTRACTED/\allowbreak{}SUPPLEMENT & country\_\allowbreak{}exposure | asset\_\allowbreak{}monthly\_\allowbreak{}return | portfolio\_\allowbreak{}scenario & 按项目用途复核；不足以单独支持核心结论。\\
WEB-0497 & B/\allowbreak{}05\_\allowbreak{}全球周期与历史危机 & ACQUIRED/\allowbreak{}MIXED\_\allowbreak{}CONTENT & EXTRACTED/\allowbreak{}SUPPLEMENT & case\_\allowbreak{}evidence | historical\_\allowbreak{}crisis\_\allowbreak{}event | global\_\allowbreak{}cycle\_\allowbreak{}month & 按项目用途复核；不足以单独支持核心结论。\\
WEB-0498 & B/\allowbreak{}07\_\allowbreak{}会计合规与制裁 & ACQUIRED/\allowbreak{}MIXED\_\allowbreak{}CONTENT & EXTRACTED/\allowbreak{}SUPPLEMENT & country\_\allowbreak{}exposure | country\_\allowbreak{}policy\_\allowbreak{}year | asset\_\allowbreak{}monthly\_\allowbreak{}return | case\_\allowbreak{}evide… & 按项目用途复核；不足以单独支持核心结论。\\
WEB-0499 & D/\allowbreak{}08\_\allowbreak{}企业案例与公开披露 & NOT\_\allowbreak{}APPLICABLE/\allowbreak{}NOT\_\allowbreak{}ACQUIRED & FAILED/\allowbreak{}EXCLUDED & company\_\allowbreak{}overseas\_\allowbreak{}exposure | country\_\allowbreak{}event | asset\_\allowbreak{}monthly\_\allowbreak{}return | portfol… & 仅保留元数据；正式分析改用许可允许的权威替代源。\\
WEB-0500 & D/\allowbreak{}08\_\allowbreak{}企业案例与公开披露 & NOT\_\allowbreak{}APPLICABLE/\allowbreak{}NOT\_\allowbreak{}ACQUIRED & FAILED/\allowbreak{}EXCLUDED & company\_\allowbreak{}overseas\_\allowbreak{}exposure | country\_\allowbreak{}event | asset\_\allowbreak{}monthly\_\allowbreak{}return | portfol… & 仅保留元数据；正式分析改用许可允许的权威替代源。\\
WEB-0501 & D/\allowbreak{}08\_\allowbreak{}企业案例与公开披露 & NOT\_\allowbreak{}APPLICABLE/\allowbreak{}NOT\_\allowbreak{}ACQUIRED & FAILED/\allowbreak{}EXCLUDED & company\_\allowbreak{}overseas\_\allowbreak{}exposure | country\_\allowbreak{}event | asset\_\allowbreak{}monthly\_\allowbreak{}return | portfol… & 仅保留元数据；正式分析改用许可允许的权威替代源。\\
WEB-0502 & D/\allowbreak{}08\_\allowbreak{}企业案例与公开披露 & NOT\_\allowbreak{}APPLICABLE/\allowbreak{}NOT\_\allowbreak{}ACQUIRED & FAILED/\allowbreak{}EXCLUDED & company\_\allowbreak{}overseas\_\allowbreak{}exposure | country\_\allowbreak{}event | asset\_\allowbreak{}monthly\_\allowbreak{}return | portfol… & 仅保留元数据；正式分析改用许可允许的权威替代源。\\
WEB-0503 & D/\allowbreak{}08\_\allowbreak{}企业案例与公开披露 & NOT\_\allowbreak{}APPLICABLE/\allowbreak{}NOT\_\allowbreak{}ACQUIRED & FAILED/\allowbreak{}EXCLUDED & company\_\allowbreak{}overseas\_\allowbreak{}exposure | country\_\allowbreak{}event | asset\_\allowbreak{}monthly\_\allowbreak{}return | portfol… & 仅保留元数据；正式分析改用许可允许的权威替代源。\\
WEB-0504 & D/\allowbreak{}08\_\allowbreak{}企业案例与公开披露 & NOT\_\allowbreak{}APPLICABLE/\allowbreak{}NOT\_\allowbreak{}ACQUIRED & FAILED/\allowbreak{}EXCLUDED & company\_\allowbreak{}overseas\_\allowbreak{}exposure | country\_\allowbreak{}event | asset\_\allowbreak{}monthly\_\allowbreak{}return | portfol… & 仅保留元数据；正式分析改用许可允许的权威替代源。\\
WEB-0505 & D/\allowbreak{}08\_\allowbreak{}企业案例与公开披露 & NOT\_\allowbreak{}APPLICABLE/\allowbreak{}NOT\_\allowbreak{}ACQUIRED & FAILED/\allowbreak{}EXCLUDED & company\_\allowbreak{}overseas\_\allowbreak{}exposure | country\_\allowbreak{}event | asset\_\allowbreak{}monthly\_\allowbreak{}return | portfol… & 仅保留元数据；正式分析改用许可允许的权威替代源。\\
WEB-0506 & D/\allowbreak{}08\_\allowbreak{}企业案例与公开披露 & NOT\_\allowbreak{}APPLICABLE/\allowbreak{}NOT\_\allowbreak{}ACQUIRED & FAILED/\allowbreak{}EXCLUDED & company\_\allowbreak{}overseas\_\allowbreak{}exposure | country\_\allowbreak{}event | asset\_\allowbreak{}monthly\_\allowbreak{}return | portfol… & 仅保留元数据；正式分析改用许可允许的权威替代源。\\
WEB-0507 & D/\allowbreak{}06\_\allowbreak{}黄金美元与储备配置 & NOT\_\allowbreak{}APPLICABLE/\allowbreak{}NOT\_\allowbreak{}ACQUIRED & FAILED/\allowbreak{}EXCLUDED & asset\_\allowbreak{}monthly\_\allowbreak{}return | portfolio\_\allowbreak{}scenario & 仅保留元数据；正式分析改用许可允许的权威替代源。\\
REF-0001 & D/\allowbreak{}07\_\allowbreak{}会计合规与制裁 & NOT\_\allowbreak{}APPLICABLE/\allowbreak{}NOT\_\allowbreak{}ACQUIRED & FAILED/\allowbreak{}EXCLUDED & country\_\allowbreak{}policy\_\allowbreak{}year & 按人工任务卡接受条款或完成机构授权后重新验收。\\
\bottomrule\end{longtable}\endgroup\end{landscape}
\chapter{159 个本地 HTML 文件逐项处理矩阵}
V1.0 清单另有37条扩展名为HTML但未取得本地文件的来源记录，故全部HTML型来源为196条；本章仅列用户关注的159个实际本地HTML。
\clearpage
\begin{landscape}\begingroup\setlength{\tabcolsep}{2pt}\scriptsize
\begin{longtable}{p{1.8cm}p{4.2cm}p{3.0cm}p{1.4cm}p{1.6cm}p{2.4cm}p{4.2cm}p{4.2cm}}
\toprule
ID & 标题 & 内容状态 & 可见字符 & 表格/段落 & 提取状态 & 目标表 & 输出或失败原因\\
\midrule\endfirsthead
\toprule ID & 标题 & 内容状态 & 可见字符 & 表格/段落 & 提取状态 & 目标表 & 输出或失败原因\\\midrule\endhead
LCL-0156 & heywhale\_\allowbreak{}dataset\_\allowbreak{}page.html & EMPTY\_\allowbreak{}SHELL & 0 & 0/\allowbreak{}0 & FAILED & case\_\allowbreak{}evidence & EMPTY\_\allowbreak{}SHELL\\
LCL-0157 & sge\_\allowbreak{}en\_\allowbreak{}daily.html & NARRATIVE\_\allowbreak{}CONTENT & 1988 & 0/\allowbreak{}5 & EXTRACTED & asset\_\allowbreak{}monthly\_\allowbreak{}return | case\_\allowbreak{}evidence & 02\_\allowbreak{}提取正文/\allowbreak{}LCL-0157/\allowbreak{}content.json\\
LCL-0160 & sge\_\allowbreak{}quotation\_\allowbreak{}daily\_\allowbreak{}new.html & NARRATIVE\_\allowbreak{}CONTENT & 726 & 0/\allowbreak{}1 & EXTRACTED & asset\_\allowbreak{}monthly\_\allowbreak{}return | case\_\allowbreak{}evidence & 02\_\allowbreak{}提取正文/\allowbreak{}LCL-0160/\allowbreak{}content.json\\
WEB-0002 & shanghai gold exchange & MIXED\_\allowbreak{}CONTENT & 2313 & 2/\allowbreak{}1 & EXTRACTED & asset\_\allowbreak{}monthly\_\allowbreak{}return | portfolio\_\allowbreak{}scenario & 02\_\allowbreak{}提取正文/\allowbreak{}WEB-0002/\allowbreak{}content.json\\
WEB-0011 & 中国人民银行办公厅关于印发《黄金租借业务管理暂行办法》的通知（银办发〔2022〕88号） & MIXED\_\allowbreak{}CONTENT & 2232 & 18/\allowbreak{}29 & EXTRACTED & company\_\allowbreak{}overseas\_\allowbreak{}exposure | country\_\allowbreak{}event | asset\_\allowbreak{}monthly\_\allowbreak{}return | po… & 02\_\allowbreak{}提取正文/\allowbreak{}WEB-0011/\allowbreak{}content.json\\
WEB-0013 & 首页-上海黄金交易所 & STRUCTURED\_\allowbreak{}DATA & 1602 & 1/\allowbreak{}3 & EXTRACTED & asset\_\allowbreak{}monthly\_\allowbreak{}return | portfolio\_\allowbreak{}scenario & 02\_\allowbreak{}提取正文/\allowbreak{}WEB-0013/\allowbreak{}content.json\\
WEB-0026 & Precious metals retreat from record highs & NARRATIVE\_\allowbreak{}CONTENT & 8310 & 0/\allowbreak{}5 & EXTRACTED & asset\_\allowbreak{}monthly\_\allowbreak{}return | portfolio\_\allowbreak{}scenario & 02\_\allowbreak{}提取正文/\allowbreak{}WEB-0026/\allowbreak{}content.json\\
WEB-0039 & Residential property prices - dashboards | BIS Data Portal & DOWNLOAD\_\allowbreak{}LANDING & 783 & 0/\allowbreak{}1 & EXTRACTED & source\_\allowbreak{}registry | country\_\allowbreak{}monthly\_\allowbreak{}risk | country\_\allowbreak{}policy\_\allowbreak{}year & 02\_\allowbreak{}提取正文/\allowbreak{}WEB-0039/\allowbreak{}content.json\\
WEB-0040 & WEO & NARRATIVE\_\allowbreak{}CONTENT & 3788 & 0/\allowbreak{}2 & EXTRACTED & source\_\allowbreak{}registry | asset\_\allowbreak{}monthly\_\allowbreak{}return | portfolio\_\allowbreak{}scenario & 02\_\allowbreak{}提取正文/\allowbreak{}WEB-0040/\allowbreak{}content.json\\
WEB-0042 & Inflation, consumer prices (annual \%) | Data & DOWNLOAD\_\allowbreak{}LANDING & 2364 & 0/\allowbreak{}1 & EXTRACTED & country\_\allowbreak{}monthly\_\allowbreak{}risk | country\_\allowbreak{}policy\_\allowbreak{}year & 02\_\allowbreak{}提取正文/\allowbreak{}WEB-0042/\allowbreak{}content.json\\
WEB-0043 & Official exchange rate (LCU per US\$, period average) | Data & DOWNLOAD\_\allowbreak{}LANDING & 1861 & 0/\allowbreak{}1 & EXTRACTED & country\_\allowbreak{}monthly\_\allowbreak{}risk | country\_\allowbreak{}policy\_\allowbreak{}year & 02\_\allowbreak{}提取正文/\allowbreak{}WEB-0043/\allowbreak{}content.json\\
WEB-0044 & Summary Terms of Use | Data & NARRATIVE\_\allowbreak{}CONTENT & 4866 & 0/\allowbreak{}60 & EXTRACTED & country\_\allowbreak{}monthly\_\allowbreak{}risk | country\_\allowbreak{}policy\_\allowbreak{}year & 02\_\allowbreak{}提取正文/\allowbreak{}WEB-0044/\allowbreak{}content.json\\
WEB-0045 & DataBank | The World Bank & MIXED\_\allowbreak{}CONTENT & 4959 & 1/\allowbreak{}6 & EXTRACTED & country\_\allowbreak{}monthly\_\allowbreak{}risk | country\_\allowbreak{}policy\_\allowbreak{}year | country\_\allowbreak{}event | case\_\allowbreak{}evi… & 02\_\allowbreak{}提取正文/\allowbreak{}WEB-0045/\allowbreak{}content.json\\
WEB-0046 & World Bank Data Catalog & NARRATIVE\_\allowbreak{}CONTENT & 579 & 0/\allowbreak{}1 & EXTRACTED & source\_\allowbreak{}registry | asset\_\allowbreak{}monthly\_\allowbreak{}return | portfolio\_\allowbreak{}scenario & 02\_\allowbreak{}提取正文/\allowbreak{}WEB-0046/\allowbreak{}content.json\\
WEB-0048 & Global Economic Monitor - World Bank Data Catalog & DOWNLOAD\_\allowbreak{}LANDING & 1889 & 0/\allowbreak{}1 & EXTRACTED & source\_\allowbreak{}registry | asset\_\allowbreak{}monthly\_\allowbreak{}return | portfolio\_\allowbreak{}scenario & 02\_\allowbreak{}提取正文/\allowbreak{}WEB-0048/\allowbreak{}content.json\\
WEB-0049 & World Bank Data Help Desk & NARRATIVE\_\allowbreak{}CONTENT & 13870 & 0/\allowbreak{}5 & EXTRACTED & source\_\allowbreak{}registry | country\_\allowbreak{}monthly\_\allowbreak{}risk | country\_\allowbreak{}policy\_\allowbreak{}year & 02\_\allowbreak{}提取正文/\allowbreak{}WEB-0049/\allowbreak{}content.json\\
WEB-0050 & About the Indicators API Documentation – World Bank Data Help Desk & NARRATIVE\_\allowbreak{}CONTENT & 9517 & 0/\allowbreak{}15 & EXTRACTED & source\_\allowbreak{}registry | asset\_\allowbreak{}monthly\_\allowbreak{}return | portfolio\_\allowbreak{}scenario & 02\_\allowbreak{}提取正文/\allowbreak{}WEB-0050/\allowbreak{}content.json\\
WEB-0060 & shanghai gold exchange & STRUCTURED\_\allowbreak{}DATA & 803 & 6/\allowbreak{}10 & EXTRACTED & asset\_\allowbreak{}monthly\_\allowbreak{}return | portfolio\_\allowbreak{}scenario & 02\_\allowbreak{}提取正文/\allowbreak{}WEB-0060/\allowbreak{}content.json\\
WEB-0098 & Treasury Prohibits Transactions with Central Bank of Russia and Imposes Sanctions on Key … & NARRATIVE\_\allowbreak{}CONTENT & 18142 & 0/\allowbreak{}25 & EXTRACTED & asset\_\allowbreak{}monthly\_\allowbreak{}return | portfolio\_\allowbreak{}scenario & 02\_\allowbreak{}提取正文/\allowbreak{}WEB-0098/\allowbreak{}content.json\\
WEB-0099 & Legislative Basis | U.S. Department of the Treasury & NARRATIVE\_\allowbreak{}CONTENT & 12582 & 0/\allowbreak{}5 & EXTRACTED & asset\_\allowbreak{}monthly\_\allowbreak{}return | case\_\allowbreak{}evidence | historical\_\allowbreak{}crisis\_\allowbreak{}event | glob… & 02\_\allowbreak{}提取正文/\allowbreak{}WEB-0099/\allowbreak{}content.json\\
WEB-0101 & Daily Treasury Rates | U.S. Department of the Treasury & MIXED\_\allowbreak{}CONTENT & 50414 & 1/\allowbreak{}5 & EXTRACTED & asset\_\allowbreak{}monthly\_\allowbreak{}return | portfolio\_\allowbreak{}scenario & 02\_\allowbreak{}提取正文/\allowbreak{}WEB-0101/\allowbreak{}content.json\\
WEB-0108 & 中国人寿完成国内保险机构首笔黄金询价交易\_\allowbreak{}中共上海市委金融委员会办公室、中共上海市金融工作委员会 & DOWNLOAD\_\allowbreak{}LANDING & 984 & 0/\allowbreak{}7 & EXTRACTED & company\_\allowbreak{}overseas\_\allowbreak{}exposure | country\_\allowbreak{}event | asset\_\allowbreak{}monthly\_\allowbreak{}return | po… & 02\_\allowbreak{}提取正文/\allowbreak{}WEB-0108/\allowbreak{}content.json\\
WEB-0109 & 上金所启用香港黄金交割仓库 中国黄金市场国际化迈出关键一步\_\allowbreak{}中共上海市委金融委员会办公室、中共上海市金融工作委员会 & NARRATIVE\_\allowbreak{}CONTENT & 1787 & 0/\allowbreak{}11 & EXTRACTED & company\_\allowbreak{}overseas\_\allowbreak{}exposure | country\_\allowbreak{}event | asset\_\allowbreak{}monthly\_\allowbreak{}return | po… & 02\_\allowbreak{}提取正文/\allowbreak{}WEB-0109/\allowbreak{}content.json\\
WEB-0128 & Home | Office of Foreign Assets Control & NARRATIVE\_\allowbreak{}CONTENT & 3940 & 0/\allowbreak{}1 & EXTRACTED & source\_\allowbreak{}registry | country\_\allowbreak{}policy\_\allowbreak{}year | case\_\allowbreak{}evidence & 02\_\allowbreak{}提取正文/\allowbreak{}WEB-0128/\allowbreak{}content.json\\
WEB-0129 & 1029 | Office of Foreign Assets Control & NARRATIVE\_\allowbreak{}CONTENT & 7506 & 0/\allowbreak{}13 & EXTRACTED & country\_\allowbreak{}policy\_\allowbreak{}year | asset\_\allowbreak{}monthly\_\allowbreak{}return & 02\_\allowbreak{}提取正文/\allowbreak{}WEB-0129/\allowbreak{}content.json\\
WEB-0139 & Forecasting the Price of Gold with Integrated Media Sentiment—A Prediction Framework Base… & MIXED\_\allowbreak{}CONTENT & 110344 & 55/\allowbreak{}397 & EXTRACTED & asset\_\allowbreak{}monthly\_\allowbreak{}return | portfolio\_\allowbreak{}scenario & 02\_\allowbreak{}提取正文/\allowbreak{}WEB-0139/\allowbreak{}content.json\\
WEB-0159 & World Bank Commodities Price Data (The Pink Sheet) & NARRATIVE\_\allowbreak{}CONTENT & 6909 & 0/\allowbreak{}1 & EXTRACTED & country\_\allowbreak{}exposure | asset\_\allowbreak{}monthly\_\allowbreak{}return | portfolio\_\allowbreak{}scenario & 02\_\allowbreak{}提取正文/\allowbreak{}WEB-0159/\allowbreak{}content.json\\
WEB-0162 & Front page | U.S. Department of the Treasury & NARRATIVE\_\allowbreak{}CONTENT & 22184 & 0/\allowbreak{}7 & EXTRACTED & asset\_\allowbreak{}monthly\_\allowbreak{}return | case\_\allowbreak{}evidence | historical\_\allowbreak{}crisis\_\allowbreak{}event | glob… & 02\_\allowbreak{}提取正文/\allowbreak{}WEB-0162/\allowbreak{}content.json\\
WEB-0187 & Bank for International Settlements & NARRATIVE\_\allowbreak{}CONTENT & 1489 & 0/\allowbreak{}19 & EXTRACTED & asset\_\allowbreak{}monthly\_\allowbreak{}return | portfolio\_\allowbreak{}scenario & 02\_\allowbreak{}提取正文/\allowbreak{}WEB-0187/\allowbreak{}content.json\\
WEB-0191 & What share for gold? On the interaction of gold and foreign exchange reserve returns & NARRATIVE\_\allowbreak{}CONTENT & 3587 & 0/\allowbreak{}6 & EXTRACTED & asset\_\allowbreak{}monthly\_\allowbreak{}return | portfolio\_\allowbreak{}scenario & 02\_\allowbreak{}提取正文/\allowbreak{}WEB-0191/\allowbreak{}content.json\\
WEB-0194 & CPI Home : U.S. Bureau of Labor Statistics & NARRATIVE\_\allowbreak{}CONTENT & 8182 & 0/\allowbreak{}8 & EXTRACTED & country\_\allowbreak{}monthly\_\allowbreak{}risk | asset\_\allowbreak{}monthly\_\allowbreak{}return | case\_\allowbreak{}evidence | histori… & 02\_\allowbreak{}提取正文/\allowbreak{}WEB-0194/\allowbreak{}content.json\\
WEB-0195 & Databases, Tables \& Calculators by Subject : U.S. Bureau of Labor Statistics & MIXED\_\allowbreak{}CONTENT & 12048 & 16/\allowbreak{}129 & EXTRACTED & asset\_\allowbreak{}monthly\_\allowbreak{}return | portfolio\_\allowbreak{}scenario & 02\_\allowbreak{}提取正文/\allowbreak{}WEB-0195/\allowbreak{}content.json\\
WEB-0197 & Home : 日本銀行 Bank of Japan & NARRATIVE\_\allowbreak{}CONTENT & 13469 & 0/\allowbreak{}12 & EXTRACTED & source\_\allowbreak{}registry | historical\_\allowbreak{}crisis\_\allowbreak{}event | global\_\allowbreak{}cycle\_\allowbreak{}month & 02\_\allowbreak{}提取正文/\allowbreak{}WEB-0197/\allowbreak{}content.json\\
WEB-0198 & Speech by Board Member Kamezaki in Mie (Recent Economic and Financial Developments in Jap… & NARRATIVE\_\allowbreak{}CONTENT & 39841 & 0/\allowbreak{}72 & EXTRACTED & source\_\allowbreak{}registry | historical\_\allowbreak{}crisis\_\allowbreak{}event | global\_\allowbreak{}cycle\_\allowbreak{}month & 02\_\allowbreak{}提取正文/\allowbreak{}WEB-0198/\allowbreak{}content.json\\
WEB-0199 & Speech by Board Member NAKAMURA in Miyazaki (Economic Activity, Prices, and Monetary Poli… & NARRATIVE\_\allowbreak{}CONTENT & 36885 & 0/\allowbreak{}59 & EXTRACTED & source\_\allowbreak{}registry | historical\_\allowbreak{}crisis\_\allowbreak{}event | global\_\allowbreak{}cycle\_\allowbreak{}month & 02\_\allowbreak{}提取正文/\allowbreak{}WEB-0199/\allowbreak{}content.json\\
WEB-0200 & Flow of Funds : 日本銀行 Bank of Japan & MIXED\_\allowbreak{}CONTENT & 15781 & 7/\allowbreak{}123 & EXTRACTED & case\_\allowbreak{}evidence | historical\_\allowbreak{}crisis\_\allowbreak{}event | global\_\allowbreak{}cycle\_\allowbreak{}month & 02\_\allowbreak{}提取正文/\allowbreak{}WEB-0200/\allowbreak{}content.json\\
WEB-0216 & National Financial Conditions Index: Current Data - Federal Reserve Bank of Chicago & NARRATIVE\_\allowbreak{}CONTENT & 8630 & 0/\allowbreak{}12 & EXTRACTED & source\_\allowbreak{}registry | country\_\allowbreak{}monthly\_\allowbreak{}risk | country\_\allowbreak{}policy\_\allowbreak{}year & 02\_\allowbreak{}提取正文/\allowbreak{}WEB-0216/\allowbreak{}content.json\\
WEB-0235 & 亚洲金融危机回顾与思考 - 国务院发展研究中心 & MIXED\_\allowbreak{}CONTENT & 7067 & 3/\allowbreak{}1 & EXTRACTED & case\_\allowbreak{}evidence | historical\_\allowbreak{}crisis\_\allowbreak{}event | global\_\allowbreak{}cycle\_\allowbreak{}month & 02\_\allowbreak{}提取正文/\allowbreak{}WEB-0235/\allowbreak{}content.json\\
WEB-0237 & European Central Bank & MIXED\_\allowbreak{}CONTENT & 54459 & 2/\allowbreak{}19 & EXTRACTED & asset\_\allowbreak{}monthly\_\allowbreak{}return | case\_\allowbreak{}evidence | historical\_\allowbreak{}crisis\_\allowbreak{}event | glob… & 02\_\allowbreak{}提取正文/\allowbreak{}WEB-0237/\allowbreak{}content.json\\
WEB-0238 & Publications by date & NARRATIVE\_\allowbreak{}CONTENT & 12315 & 0/\allowbreak{}6 & EXTRACTED & case\_\allowbreak{}evidence | historical\_\allowbreak{}crisis\_\allowbreak{}event | global\_\allowbreak{}cycle\_\allowbreak{}month & 02\_\allowbreak{}提取正文/\allowbreak{}WEB-0238/\allowbreak{}content.json\\
WEB-0253 & The Fed - The Federal Reserve’s responses to the post-Covid period of high inflation & MIXED\_\allowbreak{}CONTENT & 35610 & 1/\allowbreak{}68 & EXTRACTED & asset\_\allowbreak{}monthly\_\allowbreak{}return | portfolio\_\allowbreak{}scenario & 02\_\allowbreak{}提取正文/\allowbreak{}WEB-0253/\allowbreak{}content.json\\
WEB-0254 & The Fed - Meeting calendars and information & NARRATIVE\_\allowbreak{}CONTENT & 16657 & 0/\allowbreak{}49 & EXTRACTED & case\_\allowbreak{}evidence | historical\_\allowbreak{}crisis\_\allowbreak{}event | global\_\allowbreak{}cycle\_\allowbreak{}month & 02\_\allowbreak{}提取正文/\allowbreak{}WEB-0254/\allowbreak{}content.json\\
WEB-0256 & Federal Reserve Board - 2020 FOMC Press Releases & NARRATIVE\_\allowbreak{}CONTENT & 9459 & 0/\allowbreak{}3 & EXTRACTED & country\_\allowbreak{}monthly\_\allowbreak{}risk | country\_\allowbreak{}policy\_\allowbreak{}year & 02\_\allowbreak{}提取正文/\allowbreak{}WEB-0256/\allowbreak{}content.json\\
WEB-0257 & Federal Reserve Board - Federal Reserve issues FOMC statement & NARRATIVE\_\allowbreak{}CONTENT & 13084 & 0/\allowbreak{}7 & EXTRACTED & asset\_\allowbreak{}monthly\_\allowbreak{}return | portfolio\_\allowbreak{}scenario & 02\_\allowbreak{}提取正文/\allowbreak{}WEB-0257/\allowbreak{}content.json\\
WEB-0258 & Federal Reserve Board - Implementation Note issued March 15, 2020 & NARRATIVE\_\allowbreak{}CONTENT & 11993 & 0/\allowbreak{}12 & EXTRACTED & asset\_\allowbreak{}monthly\_\allowbreak{}return | portfolio\_\allowbreak{}scenario & 02\_\allowbreak{}提取正文/\allowbreak{}WEB-0258/\allowbreak{}content.json\\
WEB-0259 & Federal Reserve Board - Implementation Note issued December 15, 2021 & NARRATIVE\_\allowbreak{}CONTENT & 11535 & 0/\allowbreak{}18 & EXTRACTED & asset\_\allowbreak{}monthly\_\allowbreak{}return | portfolio\_\allowbreak{}scenario & 02\_\allowbreak{}提取正文/\allowbreak{}WEB-0259/\allowbreak{}content.json\\
WEB-0260 & Federal Reserve Board - Federal Reserve issues FOMC statement & NARRATIVE\_\allowbreak{}CONTENT & 10752 & 0/\allowbreak{}8 & EXTRACTED & asset\_\allowbreak{}monthly\_\allowbreak{}return | portfolio\_\allowbreak{}scenario & 02\_\allowbreak{}提取正文/\allowbreak{}WEB-0260/\allowbreak{}content.json\\
WEB-0261 & Federal Reserve Board - Federal Reserve Board announces it will make available additional… & NARRATIVE\_\allowbreak{}CONTENT & 11388 & 0/\allowbreak{}12 & EXTRACTED & asset\_\allowbreak{}monthly\_\allowbreak{}return | portfolio\_\allowbreak{}scenario & 02\_\allowbreak{}提取正文/\allowbreak{}WEB-0261/\allowbreak{}content.json\\
WEB-0262 & Federal Reserve Board - Federal Reserve issues FOMC statement & NARRATIVE\_\allowbreak{}CONTENT & 10860 & 0/\allowbreak{}9 & EXTRACTED & asset\_\allowbreak{}monthly\_\allowbreak{}return | portfolio\_\allowbreak{}scenario & 02\_\allowbreak{}提取正文/\allowbreak{}WEB-0262/\allowbreak{}content.json\\
WEB-0263 & Federal Reserve Board - Federal Reserve issues FOMC statement & NARRATIVE\_\allowbreak{}CONTENT & 11015 & 0/\allowbreak{}10 & EXTRACTED & asset\_\allowbreak{}monthly\_\allowbreak{}return | portfolio\_\allowbreak{}scenario & 02\_\allowbreak{}提取正文/\allowbreak{}WEB-0263/\allowbreak{}content.json\\
WEB-0264 & Federal Reserve Board - Industrial Production and Capacity Utilization - G.17 & MIXED\_\allowbreak{}CONTENT & 31322 & 30/\allowbreak{}170 & EXTRACTED & asset\_\allowbreak{}monthly\_\allowbreak{}return | portfolio\_\allowbreak{}scenario & 02\_\allowbreak{}提取正文/\allowbreak{}WEB-0264/\allowbreak{}content.json\\
WEB-0265 & Federal Reserve Board - Foreign Exchange Rates - H.10 - August 12, 2026 & MIXED\_\allowbreak{}CONTENT & 11132 & 1/\allowbreak{}17 & EXTRACTED & asset\_\allowbreak{}monthly\_\allowbreak{}return | portfolio\_\allowbreak{}scenario & 02\_\allowbreak{}提取正文/\allowbreak{}WEB-0265/\allowbreak{}content.json\\
WEB-0266 & Federal Reserve Board - Foreign Exchange Rates - H.10 - August 12, 2026 & MIXED\_\allowbreak{}CONTENT & 105588 & 1/\allowbreak{}6 & EXTRACTED & asset\_\allowbreak{}monthly\_\allowbreak{}return | portfolio\_\allowbreak{}scenario & 02\_\allowbreak{}提取正文/\allowbreak{}WEB-0266/\allowbreak{}content.json\\
WEB-0289 & 黄金及黄金制品进出口管理办法\_\allowbreak{}人民银行\_\allowbreak{}中国政府网 & NARRATIVE\_\allowbreak{}CONTENT & 5098 & 0/\allowbreak{}79 & EXTRACTED & company\_\allowbreak{}overseas\_\allowbreak{}exposure | country\_\allowbreak{}event | asset\_\allowbreak{}monthly\_\allowbreak{}return | po… & 02\_\allowbreak{}提取正文/\allowbreak{}WEB-0289/\allowbreak{}content.json\\
WEB-0298 & IATA - Home & NARRATIVE\_\allowbreak{}CONTENT & 7572 & 0/\allowbreak{}9 & EXTRACTED & company\_\allowbreak{}overseas\_\allowbreak{}exposure | country\_\allowbreak{}event | portfolio\_\allowbreak{}scenario | case… & 02\_\allowbreak{}提取正文/\allowbreak{}WEB-0298/\allowbreak{}content.json\\
WEB-0299 & IATA - Blocked Funds & MIXED\_\allowbreak{}CONTENT & 7901 & 1/\allowbreak{}22 & EXTRACTED & country\_\allowbreak{}monthly\_\allowbreak{}risk | country\_\allowbreak{}policy\_\allowbreak{}year & 02\_\allowbreak{}提取正文/\allowbreak{}WEB-0299/\allowbreak{}content.json\\
WEB-0300 & IATA - Pressroom & NARRATIVE\_\allowbreak{}CONTENT & 8210 & 0/\allowbreak{}14 & EXTRACTED & company\_\allowbreak{}overseas\_\allowbreak{}exposure | country\_\allowbreak{}event | portfolio\_\allowbreak{}scenario | case… & 02\_\allowbreak{}提取正文/\allowbreak{}WEB-0300/\allowbreak{}content.json\\
WEB-0301 & IATA - Blocked Funds Drop to \$1.8 billion with Major Clearance in Nigeria & MIXED\_\allowbreak{}CONTENT & 9061 & 1/\allowbreak{}15 & EXTRACTED & country\_\allowbreak{}monthly\_\allowbreak{}risk | country\_\allowbreak{}policy\_\allowbreak{}year & 02\_\allowbreak{}提取正文/\allowbreak{}WEB-0301/\allowbreak{}content.json\\
WEB-0303 & IFRS - IFRS Foundation & NARRATIVE\_\allowbreak{}CONTENT & 12900 & 0/\allowbreak{}8 & EXTRACTED & company\_\allowbreak{}overseas\_\allowbreak{}exposure | country\_\allowbreak{}policy\_\allowbreak{}year | asset\_\allowbreak{}monthly\_\allowbreak{}retur… & 02\_\allowbreak{}提取正文/\allowbreak{}WEB-0303/\allowbreak{}content.json\\
WEB-0309 & IFRS - IFRS Accounting Standards Navigator & NARRATIVE\_\allowbreak{}CONTENT & 10258 & 0/\allowbreak{}6 & EXTRACTED & country\_\allowbreak{}policy\_\allowbreak{}year | asset\_\allowbreak{}monthly\_\allowbreak{}return & 02\_\allowbreak{}提取正文/\allowbreak{}WEB-0309/\allowbreak{}content.json\\
WEB-0310 & IFRS - IAS 29 Financial Reporting in Hyperinflationary Economies & NARRATIVE\_\allowbreak{}CONTENT & 11114 & 0/\allowbreak{}7 & EXTRACTED & country\_\allowbreak{}policy\_\allowbreak{}year | asset\_\allowbreak{}monthly\_\allowbreak{}return & 02\_\allowbreak{}提取正文/\allowbreak{}WEB-0310/\allowbreak{}content.json\\
WEB-0318 & Money Matters, an IMF Exhibit -- The Importance of Global Cooperation, Reinventing the Sy… & STRUCTURED\_\allowbreak{}DATA & 1705 & 7/\allowbreak{}13 & EXTRACTED & case\_\allowbreak{}evidence | historical\_\allowbreak{}crisis\_\allowbreak{}event | global\_\allowbreak{}cycle\_\allowbreak{}month & 02\_\allowbreak{}提取正文/\allowbreak{}WEB-0318/\allowbreak{}content.json\\
WEB-0335 & \&nbsp; & BLOCK\_\allowbreak{}PAGE & 0 & 0/\allowbreak{}0 & FAILED & asset\_\allowbreak{}monthly\_\allowbreak{}return | portfolio\_\allowbreak{}scenario & BLOCK\_\allowbreak{}PAGE\\
WEB-0336 & \&nbsp; & BLOCK\_\allowbreak{}PAGE & 0 & 0/\allowbreak{}0 & FAILED & country\_\allowbreak{}monthly\_\allowbreak{}risk | country\_\allowbreak{}policy\_\allowbreak{}year & BLOCK\_\allowbreak{}PAGE\\
WEB-0344 & 商务部、国家统计局和国家外汇管理局联合发布《2024年度中国对外直接投资统计公报》 & NARRATIVE\_\allowbreak{}CONTENT & 2220 & 0/\allowbreak{}1 & EXTRACTED & source\_\allowbreak{}registry | country\_\allowbreak{}exposure | asset\_\allowbreak{}monthly\_\allowbreak{}return | portfolio… & 02\_\allowbreak{}提取正文/\allowbreak{}WEB-0344/\allowbreak{}content.json\\
WEB-0345 & 人社部，2020 & BLOCK\_\allowbreak{}PAGE & 0 & 0/\allowbreak{}0 & FAILED & company\_\allowbreak{}overseas\_\allowbreak{}exposure | country\_\allowbreak{}event | asset\_\allowbreak{}monthly\_\allowbreak{}return | po… & BLOCK\_\allowbreak{}PAGE\\
WEB-0348 & The Golden Dilemma | NBER & NARRATIVE\_\allowbreak{}CONTENT & 6991 & 0/\allowbreak{}7 & EXTRACTED & asset\_\allowbreak{}monthly\_\allowbreak{}return | historical\_\allowbreak{}crisis\_\allowbreak{}event | global\_\allowbreak{}cycle\_\allowbreak{}month & 02\_\allowbreak{}提取正文/\allowbreak{}WEB-0348/\allowbreak{}content.json\\
WEB-0349 & Business Cycle Dating | NBER & NARRATIVE\_\allowbreak{}CONTENT & 7140 & 0/\allowbreak{}1 & EXTRACTED & source\_\allowbreak{}registry | historical\_\allowbreak{}crisis\_\allowbreak{}event | global\_\allowbreak{}cycle\_\allowbreak{}month & 02\_\allowbreak{}提取正文/\allowbreak{}WEB-0349/\allowbreak{}content.json\\
WEB-0350 & Public Use Data Archive | NBER & NARRATIVE\_\allowbreak{}CONTENT & 4819 & 0/\allowbreak{}1 & EXTRACTED & source\_\allowbreak{}registry | historical\_\allowbreak{}crisis\_\allowbreak{}event | global\_\allowbreak{}cycle\_\allowbreak{}month & 02\_\allowbreak{}提取正文/\allowbreak{}WEB-0350/\allowbreak{}content.json\\
WEB-0351 & US Business Cycle Expansions and Contractions | NBER & MIXED\_\allowbreak{}CONTENT & 6807 & 1/\allowbreak{}1 & EXTRACTED & source\_\allowbreak{}registry | historical\_\allowbreak{}crisis\_\allowbreak{}event | global\_\allowbreak{}cycle\_\allowbreak{}month & 02\_\allowbreak{}提取正文/\allowbreak{}WEB-0351/\allowbreak{}content.json\\
WEB-0357 & 国家金融监督管理总局 & STRUCTURED\_\allowbreak{}DATA & 873 & 1/\allowbreak{}2 & EXTRACTED & company\_\allowbreak{}overseas\_\allowbreak{}exposure | country\_\allowbreak{}event | asset\_\allowbreak{}monthly\_\allowbreak{}return | po… & 02\_\allowbreak{}提取正文/\allowbreak{}WEB-0357/\allowbreak{}content.json\\
WEB-0358 & 金融监管总局原文，2025-02-07 & STRUCTURED\_\allowbreak{}DATA & 873 & 1/\allowbreak{}2 & EXTRACTED & company\_\allowbreak{}overseas\_\allowbreak{}exposure | country\_\allowbreak{}event | asset\_\allowbreak{}monthly\_\allowbreak{}return | po… & 02\_\allowbreak{}提取正文/\allowbreak{}WEB-0358/\allowbreak{}content.json\\
WEB-0359 & 国家金融监督管理总局 & LOW\_\allowbreak{}CONTENT & 110 & 0/\allowbreak{}0 & PARTIAL & company\_\allowbreak{}overseas\_\allowbreak{}exposure | country\_\allowbreak{}event | asset\_\allowbreak{}monthly\_\allowbreak{}return | po… & 02\_\allowbreak{}提取正文/\allowbreak{}WEB-0359/\allowbreak{}content.json\\
WEB-0363 & Economic Policy Uncertainty Index & STRUCTURED\_\allowbreak{}DATA & 1625 & 1/\allowbreak{}1 & EXTRACTED & asset\_\allowbreak{}monthly\_\allowbreak{}return | portfolio\_\allowbreak{}scenario & 02\_\allowbreak{}提取正文/\allowbreak{}WEB-0363/\allowbreak{}content.json\\
WEB-0364 & China Policy Uncertainty Indices Based on Mainland Papers & MIXED\_\allowbreak{}CONTENT & 3801 & 1/\allowbreak{}6 & EXTRACTED & asset\_\allowbreak{}monthly\_\allowbreak{}return | portfolio\_\allowbreak{}scenario & 02\_\allowbreak{}提取正文/\allowbreak{}WEB-0364/\allowbreak{}content.json\\
WEB-0382 & 国家外汇管理局门户网站 & NARRATIVE\_\allowbreak{}CONTENT & 6257 & 0/\allowbreak{}2 & EXTRACTED & case\_\allowbreak{}evidence | historical\_\allowbreak{}crisis\_\allowbreak{}event | global\_\allowbreak{}cycle\_\allowbreak{}month & 02\_\allowbreak{}提取正文/\allowbreak{}WEB-0382/\allowbreak{}content.json\\
WEB-0385 & Form 10-K & MIXED\_\allowbreak{}CONTENT & 610568 & 392/\allowbreak{}4069 & EXTRACTED & company\_\allowbreak{}overseas\_\allowbreak{}exposure | country\_\allowbreak{}event | portfolio\_\allowbreak{}scenario | case… & 02\_\allowbreak{}提取正文/\allowbreak{}WEB-0385/\allowbreak{}content.json\\
WEB-0386 & FWP & MIXED\_\allowbreak{}CONTENT & 21189 & 23/\allowbreak{}254 & EXTRACTED & company\_\allowbreak{}overseas\_\allowbreak{}exposure | country\_\allowbreak{}event | portfolio\_\allowbreak{}scenario | case… & 02\_\allowbreak{}提取正文/\allowbreak{}WEB-0386/\allowbreak{}content.json\\
WEB-0387 & Directory List of /\allowbreak{}Archives/\allowbreak{}edgar/\allowbreak{}data/\allowbreak{}1668717 & MIXED\_\allowbreak{}CONTENT & 14757 & 1/\allowbreak{}1 & EXTRACTED & company\_\allowbreak{}overseas\_\allowbreak{}exposure | country\_\allowbreak{}monthly\_\allowbreak{}risk | country\_\allowbreak{}event | po… & 02\_\allowbreak{}提取正文/\allowbreak{}WEB-0387/\allowbreak{}content.json\\
WEB-0388 & Directory List of /\allowbreak{}Archives/\allowbreak{}edgar/\allowbreak{}data/\allowbreak{}1668717/\allowbreak{}000119312524052002 & MIXED\_\allowbreak{}CONTENT & 12613 & 1/\allowbreak{}1 & EXTRACTED & company\_\allowbreak{}overseas\_\allowbreak{}exposure | country\_\allowbreak{}event | portfolio\_\allowbreak{}scenario | case… & 02\_\allowbreak{}提取正文/\allowbreak{}WEB-0388/\allowbreak{}content.json\\
WEB-0389 & EX-99.1 & NARRATIVE\_\allowbreak{}CONTENT & 768456 & 0/\allowbreak{}3434 & EXTRACTED & company\_\allowbreak{}overseas\_\allowbreak{}exposure | country\_\allowbreak{}event | asset\_\allowbreak{}monthly\_\allowbreak{}return | po… & 02\_\allowbreak{}提取正文/\allowbreak{}WEB-0389/\allowbreak{}content.json\\
WEB-0390 & Barrick 2009 年报，SEC & MIXED\_\allowbreak{}CONTENT & 369345 & 42/\allowbreak{}5195 & EXTRACTED & company\_\allowbreak{}overseas\_\allowbreak{}exposure | country\_\allowbreak{}event | asset\_\allowbreak{}monthly\_\allowbreak{}return | po… & 02\_\allowbreak{}提取正文/\allowbreak{}WEB-0390/\allowbreak{}content.json\\
WEB-0391 & EDGAR Entity Landing Page & MIXED\_\allowbreak{}CONTENT & 2482 & 1/\allowbreak{}9 & EXTRACTED & company\_\allowbreak{}overseas\_\allowbreak{}exposure | country\_\allowbreak{}event | asset\_\allowbreak{}monthly\_\allowbreak{}return | po… & 02\_\allowbreak{}提取正文/\allowbreak{}WEB-0391/\allowbreak{}content.json\\
WEB-0392 & 5. 2025 年人民币金价上涨（约 +20\%+）：上海黄金交易所（SGE）Au99.99 行情 & STRUCTURED\_\allowbreak{}DATA & 1602 & 1/\allowbreak{}3 & EXTRACTED & source\_\allowbreak{}registry | company\_\allowbreak{}overseas\_\allowbreak{}exposure | country\_\allowbreak{}policy\_\allowbreak{}year | a… & 02\_\allowbreak{}提取正文/\allowbreak{}WEB-0392/\allowbreak{}content.json\\
WEB-0393 & 现货实盘合约\_\allowbreak{}详情-上海黄金交易所 & STRUCTURED\_\allowbreak{}DATA & 1614 & 1/\allowbreak{}4 & EXTRACTED & company\_\allowbreak{}overseas\_\allowbreak{}exposure | country\_\allowbreak{}event | asset\_\allowbreak{}monthly\_\allowbreak{}return | po… & 02\_\allowbreak{}提取正文/\allowbreak{}WEB-0393/\allowbreak{}content.json\\
WEB-0394 & 每日行情-上海黄金交易所 & DOWNLOAD\_\allowbreak{}LANDING & 1319 & 0/\allowbreak{}3 & EXTRACTED & asset\_\allowbreak{}monthly\_\allowbreak{}return | portfolio\_\allowbreak{}scenario & 02\_\allowbreak{}提取正文/\allowbreak{}WEB-0394/\allowbreak{}content.json\\
WEB-0395 & 每日行情 & DOWNLOAD\_\allowbreak{}LANDING & 750 & 0/\allowbreak{}1 & EXTRACTED & asset\_\allowbreak{}monthly\_\allowbreak{}return | portfolio\_\allowbreak{}scenario & 02\_\allowbreak{}提取正文/\allowbreak{}WEB-0395/\allowbreak{}content.json\\
WEB-0396 & 每日行情 & MIXED\_\allowbreak{}CONTENT & 3431 & 1/\allowbreak{}3 & EXTRACTED & asset\_\allowbreak{}monthly\_\allowbreak{}return | portfolio\_\allowbreak{}scenario & 02\_\allowbreak{}提取正文/\allowbreak{}WEB-0396/\allowbreak{}content.json\\
WEB-0397 & 每日行情 & MIXED\_\allowbreak{}CONTENT & 3084 & 1/\allowbreak{}3 & EXTRACTED & asset\_\allowbreak{}monthly\_\allowbreak{}return | portfolio\_\allowbreak{}scenario & 02\_\allowbreak{}提取正文/\allowbreak{}WEB-0397/\allowbreak{}content.json\\
WEB-0398 & 每日行情 & MIXED\_\allowbreak{}CONTENT & 3000 & 1/\allowbreak{}3 & EXTRACTED & asset\_\allowbreak{}monthly\_\allowbreak{}return | portfolio\_\allowbreak{}scenario & 02\_\allowbreak{}提取正文/\allowbreak{}WEB-0398/\allowbreak{}content.json\\
WEB-0399 & 每日行情 & MIXED\_\allowbreak{}CONTENT & 2721 & 1/\allowbreak{}3 & EXTRACTED & asset\_\allowbreak{}monthly\_\allowbreak{}return | portfolio\_\allowbreak{}scenario & 02\_\allowbreak{}提取正文/\allowbreak{}WEB-0399/\allowbreak{}content.json\\
WEB-0400 & 每日行情 & MIXED\_\allowbreak{}CONTENT & 3434 & 1/\allowbreak{}3 & EXTRACTED & asset\_\allowbreak{}monthly\_\allowbreak{}return | portfolio\_\allowbreak{}scenario & 02\_\allowbreak{}提取正文/\allowbreak{}WEB-0400/\allowbreak{}content.json\\
WEB-0401 & 每日行情 & MIXED\_\allowbreak{}CONTENT & 3073 & 1/\allowbreak{}3 & EXTRACTED & asset\_\allowbreak{}monthly\_\allowbreak{}return | portfolio\_\allowbreak{}scenario & 02\_\allowbreak{}提取正文/\allowbreak{}WEB-0401/\allowbreak{}content.json\\
WEB-0402 & 每日行情 & MIXED\_\allowbreak{}CONTENT & 3447 & 1/\allowbreak{}3 & EXTRACTED & asset\_\allowbreak{}monthly\_\allowbreak{}return | portfolio\_\allowbreak{}scenario & 02\_\allowbreak{}提取正文/\allowbreak{}WEB-0402/\allowbreak{}content.json\\
WEB-0403 & 每日行情 & MIXED\_\allowbreak{}CONTENT & 3072 & 1/\allowbreak{}3 & EXTRACTED & asset\_\allowbreak{}monthly\_\allowbreak{}return | portfolio\_\allowbreak{}scenario & 02\_\allowbreak{}提取正文/\allowbreak{}WEB-0403/\allowbreak{}content.json\\
WEB-0404 & 每日行情 & MIXED\_\allowbreak{}CONTENT & 3432 & 1/\allowbreak{}3 & EXTRACTED & asset\_\allowbreak{}monthly\_\allowbreak{}return | portfolio\_\allowbreak{}scenario & 02\_\allowbreak{}提取正文/\allowbreak{}WEB-0404/\allowbreak{}content.json\\
WEB-0405 & 每日行情 & MIXED\_\allowbreak{}CONTENT & 3073 & 1/\allowbreak{}3 & EXTRACTED & asset\_\allowbreak{}monthly\_\allowbreak{}return | portfolio\_\allowbreak{}scenario & 02\_\allowbreak{}提取正文/\allowbreak{}WEB-0405/\allowbreak{}content.json\\
WEB-0406 & 每日行情 & MIXED\_\allowbreak{}CONTENT & 3333 & 1/\allowbreak{}3 & EXTRACTED & asset\_\allowbreak{}monthly\_\allowbreak{}return | portfolio\_\allowbreak{}scenario & 02\_\allowbreak{}提取正文/\allowbreak{}WEB-0406/\allowbreak{}content.json\\
WEB-0407 & 每日行情 & MIXED\_\allowbreak{}CONTENT & 2973 & 1/\allowbreak{}3 & EXTRACTED & asset\_\allowbreak{}monthly\_\allowbreak{}return | portfolio\_\allowbreak{}scenario & 02\_\allowbreak{}提取正文/\allowbreak{}WEB-0407/\allowbreak{}content.json\\
WEB-0408 & 每日行情 & MIXED\_\allowbreak{}CONTENT & 3429 & 1/\allowbreak{}3 & EXTRACTED & asset\_\allowbreak{}monthly\_\allowbreak{}return | portfolio\_\allowbreak{}scenario & 02\_\allowbreak{}提取正文/\allowbreak{}WEB-0408/\allowbreak{}content.json\\
WEB-0409 & 每日行情 & MIXED\_\allowbreak{}CONTENT & 3051 & 1/\allowbreak{}3 & EXTRACTED & asset\_\allowbreak{}monthly\_\allowbreak{}return | portfolio\_\allowbreak{}scenario & 02\_\allowbreak{}提取正文/\allowbreak{}WEB-0409/\allowbreak{}content.json\\
WEB-0410 & 每日行情 & MIXED\_\allowbreak{}CONTENT & 3442 & 1/\allowbreak{}3 & EXTRACTED & asset\_\allowbreak{}monthly\_\allowbreak{}return | portfolio\_\allowbreak{}scenario & 02\_\allowbreak{}提取正文/\allowbreak{}WEB-0410/\allowbreak{}content.json\\
WEB-0411 & 每日行情 & MIXED\_\allowbreak{}CONTENT & 3060 & 1/\allowbreak{}3 & EXTRACTED & asset\_\allowbreak{}monthly\_\allowbreak{}return | portfolio\_\allowbreak{}scenario & 02\_\allowbreak{}提取正文/\allowbreak{}WEB-0411/\allowbreak{}content.json\\
WEB-0412 & 每日行情 & MIXED\_\allowbreak{}CONTENT & 3323 & 1/\allowbreak{}3 & EXTRACTED & asset\_\allowbreak{}monthly\_\allowbreak{}return | portfolio\_\allowbreak{}scenario & 02\_\allowbreak{}提取正文/\allowbreak{}WEB-0412/\allowbreak{}content.json\\
WEB-0413 & 每日行情 & MIXED\_\allowbreak{}CONTENT & 2976 & 1/\allowbreak{}3 & EXTRACTED & asset\_\allowbreak{}monthly\_\allowbreak{}return | portfolio\_\allowbreak{}scenario & 02\_\allowbreak{}提取正文/\allowbreak{}WEB-0413/\allowbreak{}content.json\\
WEB-0414 & 每日行情 & MIXED\_\allowbreak{}CONTENT & 3212 & 1/\allowbreak{}3 & EXTRACTED & asset\_\allowbreak{}monthly\_\allowbreak{}return | portfolio\_\allowbreak{}scenario & 02\_\allowbreak{}提取正文/\allowbreak{}WEB-0414/\allowbreak{}content.json\\
WEB-0415 & 每日行情 & MIXED\_\allowbreak{}CONTENT & 2888 & 1/\allowbreak{}3 & EXTRACTED & asset\_\allowbreak{}monthly\_\allowbreak{}return | portfolio\_\allowbreak{}scenario & 02\_\allowbreak{}提取正文/\allowbreak{}WEB-0415/\allowbreak{}content.json\\
WEB-0416 & 每日行情 & MIXED\_\allowbreak{}CONTENT & 3437 & 1/\allowbreak{}3 & EXTRACTED & asset\_\allowbreak{}monthly\_\allowbreak{}return | portfolio\_\allowbreak{}scenario & 02\_\allowbreak{}提取正文/\allowbreak{}WEB-0416/\allowbreak{}content.json\\
WEB-0417 & 每日行情 & MIXED\_\allowbreak{}CONTENT & 3061 & 1/\allowbreak{}3 & EXTRACTED & asset\_\allowbreak{}monthly\_\allowbreak{}return | portfolio\_\allowbreak{}scenario & 02\_\allowbreak{}提取正文/\allowbreak{}WEB-0417/\allowbreak{}content.json\\
WEB-0418 & 每日行情 & MIXED\_\allowbreak{}CONTENT & 3436 & 1/\allowbreak{}3 & EXTRACTED & asset\_\allowbreak{}monthly\_\allowbreak{}return | portfolio\_\allowbreak{}scenario & 02\_\allowbreak{}提取正文/\allowbreak{}WEB-0418/\allowbreak{}content.json\\
WEB-0419 & 每日行情 & MIXED\_\allowbreak{}CONTENT & 3073 & 1/\allowbreak{}3 & EXTRACTED & asset\_\allowbreak{}monthly\_\allowbreak{}return | portfolio\_\allowbreak{}scenario & 02\_\allowbreak{}提取正文/\allowbreak{}WEB-0419/\allowbreak{}content.json\\
WEB-0420 & 每日行情 & MIXED\_\allowbreak{}CONTENT & 3219 & 1/\allowbreak{}3 & EXTRACTED & asset\_\allowbreak{}monthly\_\allowbreak{}return | portfolio\_\allowbreak{}scenario & 02\_\allowbreak{}提取正文/\allowbreak{}WEB-0420/\allowbreak{}content.json\\
WEB-0421 & 每日行情 & MIXED\_\allowbreak{}CONTENT & 2883 & 1/\allowbreak{}3 & EXTRACTED & asset\_\allowbreak{}monthly\_\allowbreak{}return | portfolio\_\allowbreak{}scenario & 02\_\allowbreak{}提取正文/\allowbreak{}WEB-0421/\allowbreak{}content.json\\
WEB-0422 & 每日行情 & MIXED\_\allowbreak{}CONTENT & 3226 & 1/\allowbreak{}3 & EXTRACTED & asset\_\allowbreak{}monthly\_\allowbreak{}return | portfolio\_\allowbreak{}scenario & 02\_\allowbreak{}提取正文/\allowbreak{}WEB-0422/\allowbreak{}content.json\\
WEB-0423 & 每日行情 & MIXED\_\allowbreak{}CONTENT & 2899 & 1/\allowbreak{}3 & EXTRACTED & asset\_\allowbreak{}monthly\_\allowbreak{}return | portfolio\_\allowbreak{}scenario & 02\_\allowbreak{}提取正文/\allowbreak{}WEB-0423/\allowbreak{}content.json\\
WEB-0424 & 每日行情 & MIXED\_\allowbreak{}CONTENT & 3438 & 1/\allowbreak{}3 & EXTRACTED & asset\_\allowbreak{}monthly\_\allowbreak{}return | portfolio\_\allowbreak{}scenario & 02\_\allowbreak{}提取正文/\allowbreak{}WEB-0424/\allowbreak{}content.json\\
WEB-0425 & 每日行情 & MIXED\_\allowbreak{}CONTENT & 3064 & 1/\allowbreak{}3 & EXTRACTED & asset\_\allowbreak{}monthly\_\allowbreak{}return | portfolio\_\allowbreak{}scenario & 02\_\allowbreak{}提取正文/\allowbreak{}WEB-0425/\allowbreak{}content.json\\
WEB-0426 & 每日行情 & MIXED\_\allowbreak{}CONTENT & 3447 & 1/\allowbreak{}3 & EXTRACTED & asset\_\allowbreak{}monthly\_\allowbreak{}return | portfolio\_\allowbreak{}scenario & 02\_\allowbreak{}提取正文/\allowbreak{}WEB-0426/\allowbreak{}content.json\\
WEB-0427 & 每日行情 & MIXED\_\allowbreak{}CONTENT & 3096 & 1/\allowbreak{}3 & EXTRACTED & asset\_\allowbreak{}monthly\_\allowbreak{}return | portfolio\_\allowbreak{}scenario & 02\_\allowbreak{}提取正文/\allowbreak{}WEB-0427/\allowbreak{}content.json\\
WEB-0428 & 每日行情 & MIXED\_\allowbreak{}CONTENT & 3328 & 1/\allowbreak{}3 & EXTRACTED & asset\_\allowbreak{}monthly\_\allowbreak{}return | portfolio\_\allowbreak{}scenario & 02\_\allowbreak{}提取正文/\allowbreak{}WEB-0428/\allowbreak{}content.json\\
WEB-0429 & 每日行情 & MIXED\_\allowbreak{}CONTENT & 2986 & 1/\allowbreak{}3 & EXTRACTED & asset\_\allowbreak{}monthly\_\allowbreak{}return | portfolio\_\allowbreak{}scenario & 02\_\allowbreak{}提取正文/\allowbreak{}WEB-0429/\allowbreak{}content.json\\
WEB-0430 & 每日行情 & MIXED\_\allowbreak{}CONTENT & 3443 & 1/\allowbreak{}3 & EXTRACTED & asset\_\allowbreak{}monthly\_\allowbreak{}return | portfolio\_\allowbreak{}scenario & 02\_\allowbreak{}提取正文/\allowbreak{}WEB-0430/\allowbreak{}content.json\\
WEB-0431 & 每日行情 & MIXED\_\allowbreak{}CONTENT & 3066 & 1/\allowbreak{}3 & EXTRACTED & asset\_\allowbreak{}monthly\_\allowbreak{}return | portfolio\_\allowbreak{}scenario & 02\_\allowbreak{}提取正文/\allowbreak{}WEB-0431/\allowbreak{}content.json\\
WEB-0432 & 每日行情 & MIXED\_\allowbreak{}CONTENT & 3434 & 1/\allowbreak{}3 & EXTRACTED & asset\_\allowbreak{}monthly\_\allowbreak{}return | portfolio\_\allowbreak{}scenario & 02\_\allowbreak{}提取正文/\allowbreak{}WEB-0432/\allowbreak{}content.json\\
WEB-0433 & 每日行情 & MIXED\_\allowbreak{}CONTENT & 3073 & 1/\allowbreak{}3 & EXTRACTED & asset\_\allowbreak{}monthly\_\allowbreak{}return | portfolio\_\allowbreak{}scenario & 02\_\allowbreak{}提取正文/\allowbreak{}WEB-0433/\allowbreak{}content.json\\
WEB-0434 & 每日行情 & MIXED\_\allowbreak{}CONTENT & 3428 & 1/\allowbreak{}3 & EXTRACTED & asset\_\allowbreak{}monthly\_\allowbreak{}return | portfolio\_\allowbreak{}scenario & 02\_\allowbreak{}提取正文/\allowbreak{}WEB-0434/\allowbreak{}content.json\\
WEB-0435 & 每日行情 & MIXED\_\allowbreak{}CONTENT & 3057 & 1/\allowbreak{}3 & EXTRACTED & asset\_\allowbreak{}monthly\_\allowbreak{}return | portfolio\_\allowbreak{}scenario & 02\_\allowbreak{}提取正文/\allowbreak{}WEB-0435/\allowbreak{}content.json\\
WEB-0436 & 每日行情 & MIXED\_\allowbreak{}CONTENT & 3473 & 1/\allowbreak{}3 & EXTRACTED & asset\_\allowbreak{}monthly\_\allowbreak{}return | portfolio\_\allowbreak{}scenario & 02\_\allowbreak{}提取正文/\allowbreak{}WEB-0436/\allowbreak{}content.json\\
WEB-0437 & 每日行情 & MIXED\_\allowbreak{}CONTENT & 3068 & 1/\allowbreak{}3 & EXTRACTED & asset\_\allowbreak{}monthly\_\allowbreak{}return | portfolio\_\allowbreak{}scenario & 02\_\allowbreak{}提取正文/\allowbreak{}WEB-0437/\allowbreak{}content.json\\
WEB-0438 & 每日行情 & MIXED\_\allowbreak{}CONTENT & 3198 & 1/\allowbreak{}3 & EXTRACTED & asset\_\allowbreak{}monthly\_\allowbreak{}return | portfolio\_\allowbreak{}scenario & 02\_\allowbreak{}提取正文/\allowbreak{}WEB-0438/\allowbreak{}content.json\\
WEB-0439 & 每日行情 & MIXED\_\allowbreak{}CONTENT & 2840 & 1/\allowbreak{}3 & EXTRACTED & asset\_\allowbreak{}monthly\_\allowbreak{}return | portfolio\_\allowbreak{}scenario & 02\_\allowbreak{}提取正文/\allowbreak{}WEB-0439/\allowbreak{}content.json\\
WEB-0440 & 每日行情 & MIXED\_\allowbreak{}CONTENT & 3497 & 1/\allowbreak{}3 & EXTRACTED & asset\_\allowbreak{}monthly\_\allowbreak{}return | portfolio\_\allowbreak{}scenario & 02\_\allowbreak{}提取正文/\allowbreak{}WEB-0440/\allowbreak{}content.json\\
WEB-0441 & 每日行情 & MIXED\_\allowbreak{}CONTENT & 3071 & 1/\allowbreak{}3 & EXTRACTED & asset\_\allowbreak{}monthly\_\allowbreak{}return | portfolio\_\allowbreak{}scenario & 02\_\allowbreak{}提取正文/\allowbreak{}WEB-0441/\allowbreak{}content.json\\
WEB-0442 & 每日行情 & MIXED\_\allowbreak{}CONTENT & 3557 & 1/\allowbreak{}3 & EXTRACTED & asset\_\allowbreak{}monthly\_\allowbreak{}return | portfolio\_\allowbreak{}scenario & 02\_\allowbreak{}提取正文/\allowbreak{}WEB-0442/\allowbreak{}content.json\\
WEB-0443 & 每日行情 & MIXED\_\allowbreak{}CONTENT & 3100 & 1/\allowbreak{}3 & EXTRACTED & asset\_\allowbreak{}monthly\_\allowbreak{}return | portfolio\_\allowbreak{}scenario & 02\_\allowbreak{}提取正文/\allowbreak{}WEB-0443/\allowbreak{}content.json\\
WEB-0444 & 每日行情 & MIXED\_\allowbreak{}CONTENT & 3508 & 1/\allowbreak{}3 & EXTRACTED & asset\_\allowbreak{}monthly\_\allowbreak{}return | portfolio\_\allowbreak{}scenario & 02\_\allowbreak{}提取正文/\allowbreak{}WEB-0444/\allowbreak{}content.json\\
WEB-0445 & 每日行情 & MIXED\_\allowbreak{}CONTENT & 3180 & 1/\allowbreak{}3 & EXTRACTED & asset\_\allowbreak{}monthly\_\allowbreak{}return | portfolio\_\allowbreak{}scenario & 02\_\allowbreak{}提取正文/\allowbreak{}WEB-0445/\allowbreak{}content.json\\
WEB-0446 & 每日行情 & MIXED\_\allowbreak{}CONTENT & 2809 & 1/\allowbreak{}3 & EXTRACTED & asset\_\allowbreak{}monthly\_\allowbreak{}return | portfolio\_\allowbreak{}scenario & 02\_\allowbreak{}提取正文/\allowbreak{}WEB-0446/\allowbreak{}content.json\\
WEB-0447 & 每日行情 & MIXED\_\allowbreak{}CONTENT & 2627 & 1/\allowbreak{}3 & EXTRACTED & asset\_\allowbreak{}monthly\_\allowbreak{}return | portfolio\_\allowbreak{}scenario & 02\_\allowbreak{}提取正文/\allowbreak{}WEB-0447/\allowbreak{}content.json\\
WEB-0448 & 每日行情 & MIXED\_\allowbreak{}CONTENT & 3451 & 1/\allowbreak{}3 & EXTRACTED & asset\_\allowbreak{}monthly\_\allowbreak{}return | portfolio\_\allowbreak{}scenario & 02\_\allowbreak{}提取正文/\allowbreak{}WEB-0448/\allowbreak{}content.json\\
WEB-0449 & 每日行情 & MIXED\_\allowbreak{}CONTENT & 3180 & 1/\allowbreak{}3 & EXTRACTED & asset\_\allowbreak{}monthly\_\allowbreak{}return | portfolio\_\allowbreak{}scenario & 02\_\allowbreak{}提取正文/\allowbreak{}WEB-0449/\allowbreak{}content.json\\
WEB-0450 & 每日行情 & MIXED\_\allowbreak{}CONTENT & 3534 & 1/\allowbreak{}3 & EXTRACTED & asset\_\allowbreak{}monthly\_\allowbreak{}return | portfolio\_\allowbreak{}scenario & 02\_\allowbreak{}提取正文/\allowbreak{}WEB-0450/\allowbreak{}content.json\\
WEB-0451 & 每日行情 & MIXED\_\allowbreak{}CONTENT & 3167 & 1/\allowbreak{}3 & EXTRACTED & asset\_\allowbreak{}monthly\_\allowbreak{}return | portfolio\_\allowbreak{}scenario & 02\_\allowbreak{}提取正文/\allowbreak{}WEB-0451/\allowbreak{}content.json\\
WEB-0452 & 每日行情 & MIXED\_\allowbreak{}CONTENT & 3329 & 1/\allowbreak{}3 & EXTRACTED & asset\_\allowbreak{}monthly\_\allowbreak{}return | portfolio\_\allowbreak{}scenario & 02\_\allowbreak{}提取正文/\allowbreak{}WEB-0452/\allowbreak{}content.json\\
WEB-0453 & 每日行情 & MIXED\_\allowbreak{}CONTENT & 2937 & 1/\allowbreak{}3 & EXTRACTED & asset\_\allowbreak{}monthly\_\allowbreak{}return | portfolio\_\allowbreak{}scenario & 02\_\allowbreak{}提取正文/\allowbreak{}WEB-0453/\allowbreak{}content.json\\
WEB-0454 & 每日行情 & MIXED\_\allowbreak{}CONTENT & 3555 & 1/\allowbreak{}3 & EXTRACTED & asset\_\allowbreak{}monthly\_\allowbreak{}return | portfolio\_\allowbreak{}scenario & 02\_\allowbreak{}提取正文/\allowbreak{}WEB-0454/\allowbreak{}content.json\\
WEB-0455 & 每日行情 & MIXED\_\allowbreak{}CONTENT & 3083 & 1/\allowbreak{}3 & EXTRACTED & asset\_\allowbreak{}monthly\_\allowbreak{}return | portfolio\_\allowbreak{}scenario & 02\_\allowbreak{}提取正文/\allowbreak{}WEB-0455/\allowbreak{}content.json\\
WEB-0456 & 每日行情 & MIXED\_\allowbreak{}CONTENT & 3321 & 1/\allowbreak{}3 & EXTRACTED & asset\_\allowbreak{}monthly\_\allowbreak{}return | portfolio\_\allowbreak{}scenario & 02\_\allowbreak{}提取正文/\allowbreak{}WEB-0456/\allowbreak{}content.json\\
WEB-0457 & 每日行情 & MIXED\_\allowbreak{}CONTENT & 2900 & 1/\allowbreak{}3 & EXTRACTED & asset\_\allowbreak{}monthly\_\allowbreak{}return | portfolio\_\allowbreak{}scenario & 02\_\allowbreak{}提取正文/\allowbreak{}WEB-0457/\allowbreak{}content.json\\
WEB-0469 & 首页 | 上海证券交易所 & MIXED\_\allowbreak{}CONTENT & 2310 & 1/\allowbreak{}1 & EXTRACTED & company\_\allowbreak{}overseas\_\allowbreak{}exposure | country\_\allowbreak{}policy\_\allowbreak{}year | asset\_\allowbreak{}monthly\_\allowbreak{}retur… & 02\_\allowbreak{}提取正文/\allowbreak{}WEB-0469/\allowbreak{}content.json\\
WEB-0471 & 2024年全国规模以上工业企业利润下降3.3\% - 国家统计局 & MIXED\_\allowbreak{}CONTENT & 13316 & 6/\allowbreak{}7 & EXTRACTED & company\_\allowbreak{}overseas\_\allowbreak{}exposure | country\_\allowbreak{}event | asset\_\allowbreak{}monthly\_\allowbreak{}return | po… & 02\_\allowbreak{}提取正文/\allowbreak{}WEB-0471/\allowbreak{}content.json\\
WEB-0472 & 国家统计局信息公开 & MIXED\_\allowbreak{}CONTENT & 5719 & 3/\allowbreak{}7 & EXTRACTED & company\_\allowbreak{}overseas\_\allowbreak{}exposure | country\_\allowbreak{}event | asset\_\allowbreak{}monthly\_\allowbreak{}return | po… & 02\_\allowbreak{}提取正文/\allowbreak{}WEB-0472/\allowbreak{}content.json\\
WEB-0492 & World Bank Group - International Development, Poverty and Sustainability & NARRATIVE\_\allowbreak{}CONTENT & 8658 & 0/\allowbreak{}78 & EXTRACTED & asset\_\allowbreak{}monthly\_\allowbreak{}return | portfolio\_\allowbreak{}scenario & 02\_\allowbreak{}提取正文/\allowbreak{}WEB-0492/\allowbreak{}content.json\\
WEB-0493 & Sudan | World Bank Group & NARRATIVE\_\allowbreak{}CONTENT & 17506 & 0/\allowbreak{}71 & EXTRACTED & country\_\allowbreak{}policy\_\allowbreak{}year | asset\_\allowbreak{}monthly\_\allowbreak{}return | case\_\allowbreak{}evidence & 02\_\allowbreak{}提取正文/\allowbreak{}WEB-0493/\allowbreak{}content.json\\
WEB-0494 & Lebanon’s Crisis: Great Denial in the Deliberate Depression & NARRATIVE\_\allowbreak{}CONTENT & 9059 & 0/\allowbreak{}11 & EXTRACTED & country\_\allowbreak{}monthly\_\allowbreak{}risk | country\_\allowbreak{}event | asset\_\allowbreak{}monthly\_\allowbreak{}return | portfol… & 02\_\allowbreak{}提取正文/\allowbreak{}WEB-0494/\allowbreak{}content.json\\
WEB-0495 & Lebanon: Normalization of Crisis is No Road to Stabilization & NARRATIVE\_\allowbreak{}CONTENT & 8610 & 0/\allowbreak{}8 & EXTRACTED & country\_\allowbreak{}monthly\_\allowbreak{}risk | country\_\allowbreak{}policy\_\allowbreak{}year | country\_\allowbreak{}event | case\_\allowbreak{}evi… & 02\_\allowbreak{}提取正文/\allowbreak{}WEB-0495/\allowbreak{}content.json\\
WEB-0496 & Conflict in Middle East Could Bring ‘Dual Shock’ to Global Commodity Markets & NARRATIVE\_\allowbreak{}CONTENT & 10023 & 0/\allowbreak{}18 & EXTRACTED & country\_\allowbreak{}exposure | asset\_\allowbreak{}monthly\_\allowbreak{}return | portfolio\_\allowbreak{}scenario & 02\_\allowbreak{}提取正文/\allowbreak{}WEB-0496/\allowbreak{}content.json\\
WEB-0497 & Global Economic Prospects & MIXED\_\allowbreak{}CONTENT & 21871 & 5/\allowbreak{}8 & EXTRACTED & case\_\allowbreak{}evidence | historical\_\allowbreak{}crisis\_\allowbreak{}event | global\_\allowbreak{}cycle\_\allowbreak{}month & 02\_\allowbreak{}提取正文/\allowbreak{}WEB-0497/\allowbreak{}content.json\\
WEB-0498 & Commodity Markets & MIXED\_\allowbreak{}CONTENT & 3009 & 1/\allowbreak{}3 & EXTRACTED & country\_\allowbreak{}exposure | country\_\allowbreak{}policy\_\allowbreak{}year | asset\_\allowbreak{}monthly\_\allowbreak{}return | case\_\allowbreak{}… & 02\_\allowbreak{}提取正文/\allowbreak{}WEB-0498/\allowbreak{}content.json\\
\bottomrule\end{longtable}\endgroup\end{landscape}
\chapter{公开底层资源补取结果}
本轮仅对公开、许可允许且无需认证的候选执行GET；服务器返回HTML而非所请求结构化文件时判为失败。
\clearpage
\begin{landscape}\begingroup\setlength{\tabcolsep}{2pt}\scriptsize
\begin{longtable}{p{1.8cm}p{1.8cm}p{2.6cm}p{3.5cm}p{2.0cm}p{5.0cm}p{6.5cm}}
\toprule
父资产 & 资源ID & 获取状态 & MIME & 大小 & 本地输出 & 失败原因/URL\\
\midrule\endfirsthead
\toprule 父资产 & 资源ID & 获取状态 & MIME & 大小 & 本地输出 & 失败原因/URL\\\midrule\endhead
WEB-0101 & WEB-0101-R01 & NOT\_\allowbreak{}ACQUIRED & text/\allowbreak{}html &  &  & HTTP\_\allowbreak{}403\\
WEB-0351 & WEB-0351-R01 & ACQUIRED & application/\allowbreak{}vnd.openxmlformats-office… & 20817 & 01\_\allowbreak{}新增原件/\allowbreak{}WEB-0351/\allowbreak{}WEB-0351-R01\_\allowbreak{}BCDC\_\allowbreak{}spreadsheet\_\allowbreak{}for\_\allowbreak{}website.xlsx & https:/\allowbreak{}/\allowbreak{}www.nber.org/\allowbreak{}sites/\allowbreak{}default/\allowbreak{}files/\allowbreak{}2023-03/\allowbreak{}BCDC\_\allowbreak{}spreadsheet\_\allowbreak{}for\_\allowbreak{}website.xlsx\\
WEB-0364 & WEB-0364-R01 & ACQUIRED & application/\allowbreak{}vnd.openxmlformats-office… & 48791 & 01\_\allowbreak{}新增原件/\allowbreak{}WEB-0364/\allowbreak{}WEB-0364-R01\_\allowbreak{}China\_\allowbreak{}Mainland\_\allowbreak{}Paper\_\allowbreak{}EPU.xlsx & https:/\allowbreak{}/\allowbreak{}www.policyuncertainty.com/\allowbreak{}media/\allowbreak{}China\_\allowbreak{}Mainland\_\allowbreak{}Paper\_\allowbreak{}EPU.xlsx\\
WEB-0042 & WEB-0042-R01 & ACQUIRED & application/\allowbreak{}zip & 113780 & 01\_\allowbreak{}新增原件/\allowbreak{}WEB-0042/\allowbreak{}WEB-0042-R01\_\allowbreak{}FP.CPI.TOTL.ZG.zip & https:/\allowbreak{}/\allowbreak{}api.worldbank.org/\allowbreak{}v2/\allowbreak{}en/\allowbreak{}indicator/\allowbreak{}FP.CPI.TOTL.ZG?downloadformat=csv\\
WEB-0043 & WEB-0043-R01 & ACQUIRED & application/\allowbreak{}zip & 68062 & 01\_\allowbreak{}新增原件/\allowbreak{}WEB-0043/\allowbreak{}WEB-0043-R01\_\allowbreak{}PA.NUS.FCRF.zip & https:/\allowbreak{}/\allowbreak{}api.worldbank.org/\allowbreak{}v2/\allowbreak{}en/\allowbreak{}indicator/\allowbreak{}PA.NUS.FCRF?downloadformat=csv\\
WEB-0040 & WEB-0040-R01 & ACQUIRED & application/\allowbreak{}vnd.openxmlformats-office… & 5585205 & 01\_\allowbreak{}新增原件/\allowbreak{}WEB-0040/\allowbreak{}WEB-0040-R01\_\allowbreak{}WEOApr2026all.xlsx & https:/\allowbreak{}/\allowbreak{}data.imf.org/\allowbreak{}-/\allowbreak{}media/\allowbreak{}iData/\allowbreak{}External-Storage/\allowbreak{}Documents/\allowbreak{}2F78EE59F79143A7921E5E203D3AAA80/\allowbreak{}en/\allowbreak{}WEOApr2026all.xlsx\\
WEB-0040 & WEB-0040-R02 & ACQUIRED & application/\allowbreak{}vnd.openxmlformats-office… & 10863437 & 01\_\allowbreak{}新增原件/\allowbreak{}WEB-0040/\allowbreak{}WEB-0040-R02\_\allowbreak{}WEOhistorical.xlsx & https:/\allowbreak{}/\allowbreak{}data.imf.org/\allowbreak{}-/\allowbreak{}media/\allowbreak{}iData/\allowbreak{}External-Storage/\allowbreak{}Documents/\allowbreak{}977574FA66914EAD900D3FEC55C7316A/\allowbreak{}en/\allowbreak{}WEOhistorical.xlsx\\
WEB-0048 & WEB-0048-R01 & ACQUIRED & application/\allowbreak{}octet-stream & 10568953 & 01\_\allowbreak{}新增原件/\allowbreak{}WEB-0048/\allowbreak{}WEB-0048-R01\_\allowbreak{}GemDataEXTR.zip & https:/\allowbreak{}/\allowbreak{}datacatalogfiles.worldbank.org/\allowbreak{}ddh-published/\allowbreak{}0037798/\allowbreak{}DR0092042/\allowbreak{}GemDataEXTR.zip\\
WEB-0497 & WEB-0497-R01 & NOT\_\allowbreak{}ACQUIRED & text/\allowbreak{}html; charset=UTF-8 &  &  & UNEXPECTED\_\allowbreak{}HTML\_\allowbreak{}INSTEAD\_\allowbreak{}OF\_\allowbreak{}STRUCTURED\_\allowbreak{}RESOURCE\\
WEB-0497 & WEB-0497-R02 & NOT\_\allowbreak{}ACQUIRED & text/\allowbreak{}html; charset=UTF-8 &  &  & UNEXPECTED\_\allowbreak{}HTML\_\allowbreak{}INSTEAD\_\allowbreak{}OF\_\allowbreak{}STRUCTURED\_\allowbreak{}RESOURCE\\
WEB-0497 & WEB-0497-R03 & NOT\_\allowbreak{}ACQUIRED & text/\allowbreak{}html; charset=UTF-8 &  &  & UNEXPECTED\_\allowbreak{}HTML\_\allowbreak{}INSTEAD\_\allowbreak{}OF\_\allowbreak{}STRUCTURED\_\allowbreak{}RESOURCE\\
WEB-0498 & WEB-0498-R01 & ACQUIRED & application/\allowbreak{}vnd.openxmlformats-office… & 399410 & 01\_\allowbreak{}新增原件/\allowbreak{}WEB-0498/\allowbreak{}WEB-0498-R01\_\allowbreak{}CMO-Historical-Data-Annual.xlsx & https:/\allowbreak{}/\allowbreak{}thedocs.worldbank.org/\allowbreak{}en/\allowbreak{}doc/\allowbreak{}74e8be41ceb20fa0da750cda2f6b9e4e-0050012026/\allowbreak{}related/\allowbreak{}CMO-Historical-Data-Annual.xlsx\\
WEB-0498 & WEB-0498-R02 & ACQUIRED & application/\allowbreak{}vnd.openxmlformats-office… & 577979 & 01\_\allowbreak{}新增原件/\allowbreak{}WEB-0498/\allowbreak{}WEB-0498-R02\_\allowbreak{}CMO-Historical-Data-Monthly.xlsx & https:/\allowbreak{}/\allowbreak{}thedocs.worldbank.org/\allowbreak{}en/\allowbreak{}doc/\allowbreak{}74e8be41ceb20fa0da750cda2f6b9e4e-0050012026/\allowbreak{}related/\allowbreak{}CMO-Historical-Data-Monthly.xl…\\
WEB-0197 & WEB-0197-R01 & ACQUIRED & application/\allowbreak{}xml;charset=UTF-8 & 16212 & 01\_\allowbreak{}新增原件/\allowbreak{}WEB-0197/\allowbreak{}WEB-0197-R01\_\allowbreak{}whatsnew.xml & https:/\allowbreak{}/\allowbreak{}www.boj.or.jp/\allowbreak{}en/\allowbreak{}rss/\allowbreak{}whatsnew.xml\\
WEB-0197 & WEB-0197-R02 & ACQUIRED & application/\allowbreak{}vnd.openxmlformats-office… & 28969 & 01\_\allowbreak{}新增原件/\allowbreak{}WEB-0197/\allowbreak{}WEB-0197-R02\_\allowbreak{}mei260810.xlsx & https:/\allowbreak{}/\allowbreak{}www.boj.or.jp/\allowbreak{}en/\allowbreak{}statistics/\allowbreak{}boj/\allowbreak{}other/\allowbreak{}mei/\allowbreak{}release/\allowbreak{}2026/\allowbreak{}mei260810.xlsx\\
WEB-0197 & WEB-0197-R03 & ACQUIRED & application/\allowbreak{}vnd.openxmlformats-office… & 33414 & 01\_\allowbreak{}新增原件/\allowbreak{}WEB-0197/\allowbreak{}WEB-0197-R03\_\allowbreak{}call.xlsx & https:/\allowbreak{}/\allowbreak{}www.boj.or.jp/\allowbreak{}en/\allowbreak{}statistics/\allowbreak{}market/\allowbreak{}short/\allowbreak{}call/\allowbreak{}call.xlsx\\
WEB-0200 & WEB-0200-R01 & ACQUIRED & application/\allowbreak{}vnd.openxmlformats-office… & 28679 & 01\_\allowbreak{}新增原件/\allowbreak{}WEB-0200/\allowbreak{}WEB-0200-R01\_\allowbreak{}sjds.xlsx & https:/\allowbreak{}/\allowbreak{}www.boj.or.jp/\allowbreak{}en/\allowbreak{}statistics/\allowbreak{}sj/\allowbreak{}sjds.xlsx\\
WEB-0200 & WEB-0200-R02 & ACQUIRED & application/\allowbreak{}vnd.openxmlformats-office… & 22877 & 01\_\allowbreak{}新增原件/\allowbreak{}WEB-0200/\allowbreak{}WEB-0200-R02\_\allowbreak{}sjfs.xlsx & https:/\allowbreak{}/\allowbreak{}www.boj.or.jp/\allowbreak{}en/\allowbreak{}statistics/\allowbreak{}sj/\allowbreak{}sjfs.xlsx\\
WEB-0200 & WEB-0200-R03 & ACQUIRED & application/\allowbreak{}vnd.openxmlformats-office… & 221195 & 01\_\allowbreak{}新增原件/\allowbreak{}WEB-0200/\allowbreak{}WEB-0200-R03\_\allowbreak{}sjfy.xlsx & https:/\allowbreak{}/\allowbreak{}www.boj.or.jp/\allowbreak{}en/\allowbreak{}statistics/\allowbreak{}sj/\allowbreak{}sjfy.xlsx\\
WEB-0253 & WEB-0253-R01 & ACQUIRED & text/\allowbreak{}xml & 19280 & 01\_\allowbreak{}新增原件/\allowbreak{}WEB-0253/\allowbreak{}WEB-0253-R01\_\allowbreak{}feds\_\allowbreak{}notes.xml & https:/\allowbreak{}/\allowbreak{}www.federalreserve.gov/\allowbreak{}feeds/\allowbreak{}feds\_\allowbreak{}notes.xml\\
WEB-0264 & WEB-0264-R01 & ACQUIRED & application/\allowbreak{}x-zip-compressed & 8569333 & 01\_\allowbreak{}新增原件/\allowbreak{}WEB-0264/\allowbreak{}WEB-0264-R01\_\allowbreak{}FRB\_\allowbreak{}g17\_\allowbreak{}xml.zip & https:/\allowbreak{}/\allowbreak{}www.federalreserve.gov/\allowbreak{}releases/\allowbreak{}g17/\allowbreak{}data/\allowbreak{}FRB\_\allowbreak{}g17\_\allowbreak{}xml.zip\\
WEB-0266 & WEB-0266-R01 & ACQUIRED & application/\allowbreak{}x-zip-compressed & 2075512 & 01\_\allowbreak{}新增原件/\allowbreak{}WEB-0266/\allowbreak{}WEB-0266-R01\_\allowbreak{}FRB\_\allowbreak{}h10\_\allowbreak{}xml.zip & https:/\allowbreak{}/\allowbreak{}www.federalreserve.gov/\allowbreak{}releases/\allowbreak{}h10/\allowbreak{}data/\allowbreak{}FRB\_\allowbreak{}h10\_\allowbreak{}xml.zip\\
WEB-0471 & WEB-0471-R01 & ACQUIRED & application/\allowbreak{}vnd.openxmlformats-office… & 31362 & 01\_\allowbreak{}新增原件/\allowbreak{}WEB-0471/\allowbreak{}WEB-0471-R01\_\allowbreak{}P020250127324149881531.xlsx & https:/\allowbreak{}/\allowbreak{}www.stats.gov.cn/\allowbreak{}sj/\allowbreak{}zxfb/\allowbreak{}202501/\allowbreak{}P020250127324149881531.xlsx\\
WEB-0472 & WEB-0472-R01 & ACQUIRED & application/\allowbreak{}vnd.openxmlformats-office… & 31356 & 01\_\allowbreak{}新增原件/\allowbreak{}WEB-0472/\allowbreak{}WEB-0472-R01\_\allowbreak{}P020241227323532210431.xlsx & https:/\allowbreak{}/\allowbreak{}www.stats.gov.cn/\allowbreak{}xxgk/\allowbreak{}sjfb/\allowbreak{}zxfb2020/\allowbreak{}202412/\allowbreak{}P020241227323532210431.xlsx\\
WEB-0050 & WEB-0050-R01 & NOT\_\allowbreak{}ACQUIRED & text/\allowbreak{}html; charset=utf-8 &  &  & UNEXPECTED\_\allowbreak{}HTML\_\allowbreak{}INSTEAD\_\allowbreak{}OF\_\allowbreak{}STRUCTURED\_\allowbreak{}RESOURCE\\
WEB-0042 & WEB-0042-R02 & ACQUIRED & application/\allowbreak{}vnd.ms-excel & 342016 & 01\_\allowbreak{}新增原件/\allowbreak{}WEB-0042/\allowbreak{}WEB-0042-R02\_\allowbreak{}FP.CPI.TOTL.ZG.xls & https:/\allowbreak{}/\allowbreak{}api.worldbank.org/\allowbreak{}v2/\allowbreak{}en/\allowbreak{}indicator/\allowbreak{}FP.CPI.TOTL.ZG?downloadformat=excel\\
WEB-0042 & WEB-0042-R03 & ACQUIRED & application/\allowbreak{}zip & 206787 & 01\_\allowbreak{}新增原件/\allowbreak{}WEB-0042/\allowbreak{}WEB-0042-R03\_\allowbreak{}FP.CPI.TOTL.ZG.zip & https:/\allowbreak{}/\allowbreak{}api.worldbank.org/\allowbreak{}v2/\allowbreak{}en/\allowbreak{}indicator/\allowbreak{}FP.CPI.TOTL.ZG?downloadformat=xml\\
WEB-0043 & WEB-0043-R02 & ACQUIRED & application/\allowbreak{}vnd.ms-excel & 339968 & 01\_\allowbreak{}新增原件/\allowbreak{}WEB-0043/\allowbreak{}WEB-0043-R02\_\allowbreak{}PA.NUS.FCRF.xls & https:/\allowbreak{}/\allowbreak{}api.worldbank.org/\allowbreak{}v2/\allowbreak{}en/\allowbreak{}indicator/\allowbreak{}PA.NUS.FCRF?downloadformat=excel\\
WEB-0043 & WEB-0043-R03 & ACQUIRED & application/\allowbreak{}zip & 156754 & 01\_\allowbreak{}新增原件/\allowbreak{}WEB-0043/\allowbreak{}WEB-0043-R03\_\allowbreak{}PA.NUS.FCRF.zip & https:/\allowbreak{}/\allowbreak{}api.worldbank.org/\allowbreak{}v2/\allowbreak{}en/\allowbreak{}indicator/\allowbreak{}PA.NUS.FCRF?downloadformat=xml\\
WEB-0040 & WEB-0040-R03 & NOT\_\allowbreak{}ACQUIRED & text/\allowbreak{}html; charset=utf-8 &  &  & UNEXPECTED\_\allowbreak{}HTML\_\allowbreak{}INSTEAD\_\allowbreak{}OF\_\allowbreak{}STRUCTURED\_\allowbreak{}RESOURCE\\
WEB-0351 & WEB-0351-R02 & NOT\_\allowbreak{}ACQUIRED & text/\allowbreak{}html;charset=ISO-8859-1 &  &  & UNEXPECTED\_\allowbreak{}HTML\_\allowbreak{}INSTEAD\_\allowbreak{}OF\_\allowbreak{}STRUCTURED\_\allowbreak{}RESOURCE\\
WEB-0101 & WEB-0101-R02 & NOT\_\allowbreak{}ACQUIRED & text/\allowbreak{}html; charset=UTF-8 &  &  & UNEXPECTED\_\allowbreak{}HTML\_\allowbreak{}INSTEAD\_\allowbreak{}OF\_\allowbreak{}STRUCTURED\_\allowbreak{}RESOURCE\\
WEB-0101 & WEB-0101-R03 & NOT\_\allowbreak{}ACQUIRED & text/\allowbreak{}html; charset=UTF-8 &  &  & UNEXPECTED\_\allowbreak{}HTML\_\allowbreak{}INSTEAD\_\allowbreak{}OF\_\allowbreak{}STRUCTURED\_\allowbreak{}RESOURCE\\
WEB-0264 & WEB-0264-R02 & ACQUIRED & text/\allowbreak{}plain & 176550 & 01\_\allowbreak{}新增原件/\allowbreak{}WEB-0264/\allowbreak{}WEB-0264-R02\_\allowbreak{}g17.txt & https:/\allowbreak{}/\allowbreak{}www.federalreserve.gov/\allowbreak{}releases/\allowbreak{}g17/\allowbreak{}19971215/\allowbreak{}g17.txt\\
WEB-0264 & WEB-0264-R03 & ACQUIRED & text/\allowbreak{}plain & 180665 & 01\_\allowbreak{}新增原件/\allowbreak{}WEB-0264/\allowbreak{}WEB-0264-R03\_\allowbreak{}g17.txt & https:/\allowbreak{}/\allowbreak{}www.federalreserve.gov/\allowbreak{}releases/\allowbreak{}g17/\allowbreak{}19980116/\allowbreak{}g17.txt\\
WEB-0388 & WEB-0388-R01 & NOT\_\allowbreak{}ACQUIRED & text/\allowbreak{}plain &  &  & DECLARED\_\allowbreak{}SIZE\_\allowbreak{}EXCEEDS\_\allowbreak{}50MB\\
\bottomrule\end{longtable}\endgroup\end{landscape}
\chapter{结构化资源展开与提取清单}
XLSX按工作表转为CSV；ZIP先检查路径后展开；XML、CSV和文本记录结构；旧XLS保留原件并进入人工转换。
\clearpage
\begin{landscape}\begingroup\setlength{\tabcolsep}{2pt}\scriptsize
\begin{longtable}{p{1.8cm}p{2.2cm}p{3.1cm}p{4.7cm}p{1.3cm}p{1.3cm}p{3.0cm}p{5.3cm}}
\toprule
资产 & 资源 & 类型 & 对象 & 行 & 列 & 状态 & 输出/限制\\
\midrule\endfirsthead
\toprule 资产 & 资源 & 类型 & 对象 & 行 & 列 & 状态 & 输出/限制\\\midrule\endhead
WEB-0351 & WEB-0351-R01 & xlsx\_\allowbreak{}sheet & NBER Chronology & 44 & 13 & EXTRACTED & 03\_\allowbreak{}提取表格/\allowbreak{}WEB-0351/\allowbreak{}WEB-0351-R01/\allowbreak{}NBER\_\allowbreak{}Chronology.csv\\
WEB-0364 & WEB-0364-R01 & xlsx\_\allowbreak{}sheet & EPU 1949-1978 & 352 & 6 & EXTRACTED & 03\_\allowbreak{}提取表格/\allowbreak{}WEB-0364/\allowbreak{}WEB-0364-R01/\allowbreak{}EPU\_\allowbreak{}1949-1978.csv\\
WEB-0364 & WEB-0364-R01 & xlsx\_\allowbreak{}sheet & EPU 1979-1999 & 253 & 6 & EXTRACTED & 03\_\allowbreak{}提取表格/\allowbreak{}WEB-0364/\allowbreak{}WEB-0364-R01/\allowbreak{}EPU\_\allowbreak{}1979-1999.csv\\
WEB-0364 & WEB-0364-R01 & xlsx\_\allowbreak{}sheet & EPU 2000 onwards & 320 & 6 & EXTRACTED & 03\_\allowbreak{}提取表格/\allowbreak{}WEB-0364/\allowbreak{}WEB-0364-R01/\allowbreak{}EPU\_\allowbreak{}2000\_\allowbreak{}onwards.csv\\
WEB-0364 & WEB-0364-R01 & xlsx\_\allowbreak{}sheet & TPU 2000 onwards & 320 & 6 & EXTRACTED & 03\_\allowbreak{}提取表格/\allowbreak{}WEB-0364/\allowbreak{}WEB-0364-R01/\allowbreak{}TPU\_\allowbreak{}2000\_\allowbreak{}onwards.csv\\
WEB-0042 & WEB-0042-R01 & zip\_\allowbreak{}archive & 3 entries & 3 &  & EXPANDED & 01\_\allowbreak{}新增原件/\allowbreak{}WEB-0042/\allowbreak{}expanded/\allowbreak{}WEB-0042-R01\\
WEB-0042 & WEB-0042-R01-E001 & csv & API\_\allowbreak{}FP.CPI.TOTL.ZG\_\allowbreak{}DS2\_\allowbreak{}en\_\allowbreak{}csv\_\allowbreak{}v2\_\allowbreak{}285.csv & 267 & 3 & INSPECTED & 01\_\allowbreak{}新增原件/\allowbreak{}WEB-0042/\allowbreak{}expanded/\allowbreak{}WEB-0042-R01/\allowbreak{}API\_\allowbreak{}FP.CPI.TOTL.ZG\_\allowbreak{}DS2\_\allowbreak{}en\_\allowbreak{}csv\_\allowbreak{}v2\_\allowbreak{}285.csv\\
WEB-0042 & WEB-0042-R01-E002 & csv & Metadata\_\allowbreak{}Country\_\allowbreak{}API\_\allowbreak{}FP.CPI.TOTL.ZG\_\allowbreak{}DS2\_\allowbreak{}en\_\allowbreak{}csv\_\allowbreak{}v2\_\allowbreak{}285.csv & 277 & 6 & INSPECTED & 01\_\allowbreak{}新增原件/\allowbreak{}WEB-0042/\allowbreak{}expanded/\allowbreak{}WEB-0042-R01/\allowbreak{}Metadata\_\allowbreak{}Country\_\allowbreak{}API\_\allowbreak{}FP.CPI.TOTL.ZG\_\allowbreak{}DS2\_\allowbreak{}en\_\allowbreak{}csv\_\allowbreak{}v2\_\allowbreak{}285.csv\\
WEB-0042 & WEB-0042-R01-E003 & csv & Metadata\_\allowbreak{}Indicator\_\allowbreak{}API\_\allowbreak{}FP.CPI.TOTL.ZG\_\allowbreak{}DS2\_\allowbreak{}en\_\allowbreak{}csv\_\allowbreak{}v2\_\allowbreak{}285.csv & 1 & 5 & INSPECTED & 01\_\allowbreak{}新增原件/\allowbreak{}WEB-0042/\allowbreak{}expanded/\allowbreak{}WEB-0042-R01/\allowbreak{}Metadata\_\allowbreak{}Indicator\_\allowbreak{}API\_\allowbreak{}FP.CPI.TOTL.ZG\_\allowbreak{}DS2\_\allowbreak{}en\_\allowbreak{}csv\_\allowbreak{}v2\_\allowbreak{}285.csv\\
WEB-0043 & WEB-0043-R01 & zip\_\allowbreak{}archive & 3 entries & 3 &  & EXPANDED & 01\_\allowbreak{}新增原件/\allowbreak{}WEB-0043/\allowbreak{}expanded/\allowbreak{}WEB-0043-R01\\
WEB-0043 & WEB-0043-R01-E001 & csv & API\_\allowbreak{}PA.NUS.FCRF\_\allowbreak{}DS2\_\allowbreak{}en\_\allowbreak{}csv\_\allowbreak{}v2\_\allowbreak{}32.csv & 267 & 3 & INSPECTED & 01\_\allowbreak{}新增原件/\allowbreak{}WEB-0043/\allowbreak{}expanded/\allowbreak{}WEB-0043-R01/\allowbreak{}API\_\allowbreak{}PA.NUS.FCRF\_\allowbreak{}DS2\_\allowbreak{}en\_\allowbreak{}csv\_\allowbreak{}v2\_\allowbreak{}32.csv\\
WEB-0043 & WEB-0043-R01-E002 & csv & Metadata\_\allowbreak{}Country\_\allowbreak{}API\_\allowbreak{}PA.NUS.FCRF\_\allowbreak{}DS2\_\allowbreak{}en\_\allowbreak{}csv\_\allowbreak{}v2\_\allowbreak{}32.csv & 277 & 6 & INSPECTED & 01\_\allowbreak{}新增原件/\allowbreak{}WEB-0043/\allowbreak{}expanded/\allowbreak{}WEB-0043-R01/\allowbreak{}Metadata\_\allowbreak{}Country\_\allowbreak{}API\_\allowbreak{}PA.NUS.FCRF\_\allowbreak{}DS2\_\allowbreak{}en\_\allowbreak{}csv\_\allowbreak{}v2\_\allowbreak{}32.csv\\
WEB-0043 & WEB-0043-R01-E003 & csv & Metadata\_\allowbreak{}Indicator\_\allowbreak{}API\_\allowbreak{}PA.NUS.FCRF\_\allowbreak{}DS2\_\allowbreak{}en\_\allowbreak{}csv\_\allowbreak{}v2\_\allowbreak{}32.csv & 1 & 5 & INSPECTED & 01\_\allowbreak{}新增原件/\allowbreak{}WEB-0043/\allowbreak{}expanded/\allowbreak{}WEB-0043-R01/\allowbreak{}Metadata\_\allowbreak{}Indicator\_\allowbreak{}API\_\allowbreak{}PA.NUS.FCRF\_\allowbreak{}DS2\_\allowbreak{}en\_\allowbreak{}csv\_\allowbreak{}v2\_\allowbreak{}32.csv\\
WEB-0040 & WEB-0040-R01 & xlsx\_\allowbreak{}sheet & April 2026 WEO & 19 & 10 & EXTRACTED & 03\_\allowbreak{}提取表格/\allowbreak{}WEB-0040/\allowbreak{}WEB-0040-R01/\allowbreak{}April\_\allowbreak{}2026\_\allowbreak{}WEO.csv\\
WEB-0040 & WEB-0040-R01 & xlsx\_\allowbreak{}sheet & Countries & 8669 & 81 & EXTRACTED & 03\_\allowbreak{}提取表格/\allowbreak{}WEB-0040/\allowbreak{}WEB-0040-R01/\allowbreak{}Countries.csv\\
WEB-0040 & WEB-0040-R01 & xlsx\_\allowbreak{}sheet & Country Groups & 765 & 81 & EXTRACTED & 03\_\allowbreak{}提取表格/\allowbreak{}WEB-0040/\allowbreak{}WEB-0040-R01/\allowbreak{}Country\_\allowbreak{}Groups.csv\\
WEB-0040 & WEB-0040-R01 & xlsx\_\allowbreak{}sheet & Commodity Prices & 76 & 81 & EXTRACTED & 03\_\allowbreak{}提取表格/\allowbreak{}WEB-0040/\allowbreak{}WEB-0040-R01/\allowbreak{}Commodity\_\allowbreak{}Prices.csv\\
WEB-0040 & WEB-0040-R01 & xlsx\_\allowbreak{}sheet & Country Group Composition & 627 & 6 & EXTRACTED & 03\_\allowbreak{}提取表格/\allowbreak{}WEB-0040/\allowbreak{}WEB-0040-R01/\allowbreak{}Country\_\allowbreak{}Group\_\allowbreak{}Composition.csv\\
WEB-0040 & WEB-0040-R02 & xlsx\_\allowbreak{}sheet & Info & 2 & 2 & EXTRACTED & 03\_\allowbreak{}提取表格/\allowbreak{}WEB-0040/\allowbreak{}WEB-0040-R02/\allowbreak{}Info.csv\\
WEB-0040 & WEB-0040-R02 & xlsx\_\allowbreak{}sheet & ngdp\_\allowbreak{}rpch & 8853 & 77 & EXTRACTED & 03\_\allowbreak{}提取表格/\allowbreak{}WEB-0040/\allowbreak{}WEB-0040-R02/\allowbreak{}ngdp\_\allowbreak{}rpch.csv\\
WEB-0040 & WEB-0040-R02 & xlsx\_\allowbreak{}sheet & pcpi\_\allowbreak{}pch & 8853 & 77 & EXTRACTED & 03\_\allowbreak{}提取表格/\allowbreak{}WEB-0040/\allowbreak{}WEB-0040-R02/\allowbreak{}pcpi\_\allowbreak{}pch.csv\\
WEB-0040 & WEB-0040-R02 & xlsx\_\allowbreak{}sheet & bca\_\allowbreak{}gdp\_\allowbreak{}bp6 & 8853 & 77 & EXTRACTED & 03\_\allowbreak{}提取表格/\allowbreak{}WEB-0040/\allowbreak{}WEB-0040-R02/\allowbreak{}bca\_\allowbreak{}gdp\_\allowbreak{}bp6.csv\\
WEB-0048 & WEB-0048-R01 & zip\_\allowbreak{}archive & 35 entries & 35 &  & EXPANDED & 01\_\allowbreak{}新增原件/\allowbreak{}WEB-0048/\allowbreak{}expanded/\allowbreak{}WEB-0048-R01\\
WEB-0048 & WEB-0048-R01-E001 & xlsx\_\allowbreak{}sheet & annual & 32 & 95 & EXTRACTED & 03\_\allowbreak{}提取表格/\allowbreak{}WEB-0048/\allowbreak{}WEB-0048-R01-E001/\allowbreak{}annual.csv\\
WEB-0048 & WEB-0048-R01-E001 & xlsx\_\allowbreak{}sheet & monthly & 362 & 95 & EXTRACTED & 03\_\allowbreak{}提取表格/\allowbreak{}WEB-0048/\allowbreak{}WEB-0048-R01-E001/\allowbreak{}monthly.csv\\
WEB-0048 & WEB-0048-R01-E002 & xlsx\_\allowbreak{}sheet & annual & 32 & 87 & EXTRACTED & 03\_\allowbreak{}提取表格/\allowbreak{}WEB-0048/\allowbreak{}WEB-0048-R01-E002/\allowbreak{}annual.csv\\
WEB-0048 & WEB-0048-R01-E002 & xlsx\_\allowbreak{}sheet & monthly & 362 & 87 & EXTRACTED & 03\_\allowbreak{}提取表格/\allowbreak{}WEB-0048/\allowbreak{}WEB-0048-R01-E002/\allowbreak{}monthly.csv\\
WEB-0048 & WEB-0048-R01-E003 & xlsx\_\allowbreak{}sheet & annual & 32 & 13 & EXTRACTED & 03\_\allowbreak{}提取表格/\allowbreak{}WEB-0048/\allowbreak{}WEB-0048-R01-E003/\allowbreak{}annual.csv\\
WEB-0048 & WEB-0048-R01-E003 & xlsx\_\allowbreak{}sheet & monthly & 362 & 13 & EXTRACTED & 03\_\allowbreak{}提取表格/\allowbreak{}WEB-0048/\allowbreak{}WEB-0048-R01-E003/\allowbreak{}monthly.csv\\
WEB-0048 & WEB-0048-R01-E004 & xlsx\_\allowbreak{}sheet & annual & 32 & 109 & EXTRACTED & 03\_\allowbreak{}提取表格/\allowbreak{}WEB-0048/\allowbreak{}WEB-0048-R01-E004/\allowbreak{}annual.csv\\
WEB-0048 & WEB-0048-R01-E004 & xlsx\_\allowbreak{}sheet & monthly & 362 & 109 & EXTRACTED & 03\_\allowbreak{}提取表格/\allowbreak{}WEB-0048/\allowbreak{}WEB-0048-R01-E004/\allowbreak{}monthly.csv\\
WEB-0048 & WEB-0048-R01-E005 & xlsx\_\allowbreak{}sheet & annual & 32 & 158 & EXTRACTED & 03\_\allowbreak{}提取表格/\allowbreak{}WEB-0048/\allowbreak{}WEB-0048-R01-E005/\allowbreak{}annual.csv\\
WEB-0048 & WEB-0048-R01-E005 & xlsx\_\allowbreak{}sheet & monthly & 362 & 158 & EXTRACTED & 03\_\allowbreak{}提取表格/\allowbreak{}WEB-0048/\allowbreak{}WEB-0048-R01-E005/\allowbreak{}monthly.csv\\
WEB-0048 & WEB-0048-R01-E006 & xlsx\_\allowbreak{}sheet & annual & 32 & 121 & EXTRACTED & 03\_\allowbreak{}提取表格/\allowbreak{}WEB-0048/\allowbreak{}WEB-0048-R01-E006/\allowbreak{}annual.csv\\
WEB-0048 & WEB-0048-R01-E006 & xlsx\_\allowbreak{}sheet & monthly & 362 & 121 & EXTRACTED & 03\_\allowbreak{}提取表格/\allowbreak{}WEB-0048/\allowbreak{}WEB-0048-R01-E006/\allowbreak{}monthly.csv\\
WEB-0048 & WEB-0048-R01-E007 & xlsx\_\allowbreak{}sheet & annual & 32 & 200 & EXTRACTED & 03\_\allowbreak{}提取表格/\allowbreak{}WEB-0048/\allowbreak{}WEB-0048-R01-E007/\allowbreak{}annual.csv\\
WEB-0048 & WEB-0048-R01-E007 & xlsx\_\allowbreak{}sheet & monthly & 362 & 200 & EXTRACTED & 03\_\allowbreak{}提取表格/\allowbreak{}WEB-0048/\allowbreak{}WEB-0048-R01-E007/\allowbreak{}monthly.csv\\
WEB-0048 & WEB-0048-R01-E008 & xlsx\_\allowbreak{}sheet & annual & 32 & 188 & EXTRACTED & 03\_\allowbreak{}提取表格/\allowbreak{}WEB-0048/\allowbreak{}WEB-0048-R01-E008/\allowbreak{}annual.csv\\
WEB-0048 & WEB-0048-R01-E008 & xlsx\_\allowbreak{}sheet & monthly & 362 & 188 & EXTRACTED & 03\_\allowbreak{}提取表格/\allowbreak{}WEB-0048/\allowbreak{}WEB-0048-R01-E008/\allowbreak{}monthly.csv\\
WEB-0048 & WEB-0048-R01-E009 & xlsx\_\allowbreak{}sheet & annual & 32 & 52 & EXTRACTED & 03\_\allowbreak{}提取表格/\allowbreak{}WEB-0048/\allowbreak{}WEB-0048-R01-E009/\allowbreak{}annual.csv\\
WEB-0048 & WEB-0048-R01-E009 & xlsx\_\allowbreak{}sheet & monthly & 362 & 52 & EXTRACTED & 03\_\allowbreak{}提取表格/\allowbreak{}WEB-0048/\allowbreak{}WEB-0048-R01-E009/\allowbreak{}monthly.csv\\
WEB-0048 & WEB-0048-R01-E009 & xlsx\_\allowbreak{}sheet & quarterly & 122 & 52 & EXTRACTED & 03\_\allowbreak{}提取表格/\allowbreak{}WEB-0048/\allowbreak{}WEB-0048-R01-E009/\allowbreak{}quarterly.csv\\
WEB-0048 & WEB-0048-R01-E010 & xlsx\_\allowbreak{}sheet & annual & 32 & 52 & EXTRACTED & 03\_\allowbreak{}提取表格/\allowbreak{}WEB-0048/\allowbreak{}WEB-0048-R01-E010/\allowbreak{}annual.csv\\
WEB-0048 & WEB-0048-R01-E010 & xlsx\_\allowbreak{}sheet & monthly & 362 & 52 & EXTRACTED & 03\_\allowbreak{}提取表格/\allowbreak{}WEB-0048/\allowbreak{}WEB-0048-R01-E010/\allowbreak{}monthly.csv\\
WEB-0048 & WEB-0048-R01-E010 & xlsx\_\allowbreak{}sheet & quarterly & 122 & 52 & EXTRACTED & 03\_\allowbreak{}提取表格/\allowbreak{}WEB-0048/\allowbreak{}WEB-0048-R01-E010/\allowbreak{}quarterly.csv\\
WEB-0048 & WEB-0048-R01-E011 & xlsx\_\allowbreak{}sheet & annual & 32 & 140 & EXTRACTED & 03\_\allowbreak{}提取表格/\allowbreak{}WEB-0048/\allowbreak{}WEB-0048-R01-E011/\allowbreak{}annual.csv\\
WEB-0048 & WEB-0048-R01-E011 & xlsx\_\allowbreak{}sheet & monthly & 362 & 140 & EXTRACTED & 03\_\allowbreak{}提取表格/\allowbreak{}WEB-0048/\allowbreak{}WEB-0048-R01-E011/\allowbreak{}monthly.csv\\
WEB-0048 & WEB-0048-R01-E011 & xlsx\_\allowbreak{}sheet & quarterly & 122 & 140 & EXTRACTED & 03\_\allowbreak{}提取表格/\allowbreak{}WEB-0048/\allowbreak{}WEB-0048-R01-E011/\allowbreak{}quarterly.csv\\
WEB-0048 & WEB-0048-R01-E012 & xlsx\_\allowbreak{}sheet & annual & 32 & 139 & EXTRACTED & 03\_\allowbreak{}提取表格/\allowbreak{}WEB-0048/\allowbreak{}WEB-0048-R01-E012/\allowbreak{}annual.csv\\
WEB-0048 & WEB-0048-R01-E012 & xlsx\_\allowbreak{}sheet & monthly & 362 & 139 & EXTRACTED & 03\_\allowbreak{}提取表格/\allowbreak{}WEB-0048/\allowbreak{}WEB-0048-R01-E012/\allowbreak{}monthly.csv\\
WEB-0048 & WEB-0048-R01-E012 & xlsx\_\allowbreak{}sheet & quarterly & 122 & 139 & EXTRACTED & 03\_\allowbreak{}提取表格/\allowbreak{}WEB-0048/\allowbreak{}WEB-0048-R01-E012/\allowbreak{}quarterly.csv\\
WEB-0048 & WEB-0048-R01-E013 & xlsx\_\allowbreak{}sheet & annual & 32 & 40 & EXTRACTED & 03\_\allowbreak{}提取表格/\allowbreak{}WEB-0048/\allowbreak{}WEB-0048-R01-E013/\allowbreak{}annual.csv\\
WEB-0048 & WEB-0048-R01-E013 & xlsx\_\allowbreak{}sheet & monthly & 362 & 52 & EXTRACTED & 03\_\allowbreak{}提取表格/\allowbreak{}WEB-0048/\allowbreak{}WEB-0048-R01-E013/\allowbreak{}monthly.csv\\
WEB-0048 & WEB-0048-R01-E013 & xlsx\_\allowbreak{}sheet & quarterly & 122 & 41 & EXTRACTED & 03\_\allowbreak{}提取表格/\allowbreak{}WEB-0048/\allowbreak{}WEB-0048-R01-E013/\allowbreak{}quarterly.csv\\
WEB-0048 & WEB-0048-R01-E014 & xlsx\_\allowbreak{}sheet & annual & 32 & 52 & EXTRACTED & 03\_\allowbreak{}提取表格/\allowbreak{}WEB-0048/\allowbreak{}WEB-0048-R01-E014/\allowbreak{}annual.csv\\
WEB-0048 & WEB-0048-R01-E014 & xlsx\_\allowbreak{}sheet & monthly & 362 & 52 & EXTRACTED & 03\_\allowbreak{}提取表格/\allowbreak{}WEB-0048/\allowbreak{}WEB-0048-R01-E014/\allowbreak{}monthly.csv\\
WEB-0048 & WEB-0048-R01-E014 & xlsx\_\allowbreak{}sheet & quarterly & 122 & 52 & EXTRACTED & 03\_\allowbreak{}提取表格/\allowbreak{}WEB-0048/\allowbreak{}WEB-0048-R01-E014/\allowbreak{}quarterly.csv\\
WEB-0048 & WEB-0048-R01-E015 & xlsx\_\allowbreak{}sheet & annual & 32 & 119 & EXTRACTED & 03\_\allowbreak{}提取表格/\allowbreak{}WEB-0048/\allowbreak{}WEB-0048-R01-E015/\allowbreak{}annual.csv\\
WEB-0048 & WEB-0048-R01-E015 & xlsx\_\allowbreak{}sheet & monthly & 362 & 119 & EXTRACTED & 03\_\allowbreak{}提取表格/\allowbreak{}WEB-0048/\allowbreak{}WEB-0048-R01-E015/\allowbreak{}monthly.csv\\
WEB-0048 & WEB-0048-R01-E016 & xlsx\_\allowbreak{}sheet & annual & 32 & 92 & EXTRACTED & 03\_\allowbreak{}提取表格/\allowbreak{}WEB-0048/\allowbreak{}WEB-0048-R01-E016/\allowbreak{}annual.csv\\
WEB-0048 & WEB-0048-R01-E016 & xlsx\_\allowbreak{}sheet & quarterly & 122 & 92 & EXTRACTED & 03\_\allowbreak{}提取表格/\allowbreak{}WEB-0048/\allowbreak{}WEB-0048-R01-E016/\allowbreak{}quarterly.csv\\
WEB-0048 & WEB-0048-R01-E017 & xlsx\_\allowbreak{}sheet & annual & 32 & 100 & EXTRACTED & 03\_\allowbreak{}提取表格/\allowbreak{}WEB-0048/\allowbreak{}WEB-0048-R01-E017/\allowbreak{}annual.csv\\
WEB-0048 & WEB-0048-R01-E017 & xlsx\_\allowbreak{}sheet & quarterly & 122 & 100 & EXTRACTED & 03\_\allowbreak{}提取表格/\allowbreak{}WEB-0048/\allowbreak{}WEB-0048-R01-E017/\allowbreak{}quarterly.csv\\
WEB-0048 & WEB-0048-R01-E018 & xlsx\_\allowbreak{}sheet & annual & 32 & 91 & EXTRACTED & 03\_\allowbreak{}提取表格/\allowbreak{}WEB-0048/\allowbreak{}WEB-0048-R01-E018/\allowbreak{}annual.csv\\
WEB-0048 & WEB-0048-R01-E018 & xlsx\_\allowbreak{}sheet & quarterly & 122 & 91 & EXTRACTED & 03\_\allowbreak{}提取表格/\allowbreak{}WEB-0048/\allowbreak{}WEB-0048-R01-E018/\allowbreak{}quarterly.csv\\
WEB-0048 & WEB-0048-R01-E019 & xlsx\_\allowbreak{}sheet & annual & 32 & 104 & EXTRACTED & 03\_\allowbreak{}提取表格/\allowbreak{}WEB-0048/\allowbreak{}WEB-0048-R01-E019/\allowbreak{}annual.csv\\
WEB-0048 & WEB-0048-R01-E019 & xlsx\_\allowbreak{}sheet & quarterly & 122 & 104 & EXTRACTED & 03\_\allowbreak{}提取表格/\allowbreak{}WEB-0048/\allowbreak{}WEB-0048-R01-E019/\allowbreak{}quarterly.csv\\
WEB-0048 & WEB-0048-R01-E020 & xlsx\_\allowbreak{}sheet & quarterly & 122 & 85 & EXTRACTED & 03\_\allowbreak{}提取表格/\allowbreak{}WEB-0048/\allowbreak{}WEB-0048-R01-E020/\allowbreak{}quarterly.csv\\
WEB-0048 & WEB-0048-R01-E021 & xlsx\_\allowbreak{}sheet & annual & 32 & 53 & EXTRACTED & 03\_\allowbreak{}提取表格/\allowbreak{}WEB-0048/\allowbreak{}WEB-0048-R01-E021/\allowbreak{}annual.csv\\
WEB-0048 & WEB-0048-R01-E021 & xlsx\_\allowbreak{}sheet & monthly & 362 & 53 & EXTRACTED & 03\_\allowbreak{}提取表格/\allowbreak{}WEB-0048/\allowbreak{}WEB-0048-R01-E021/\allowbreak{}monthly.csv\\
WEB-0048 & WEB-0048-R01-E021 & xlsx\_\allowbreak{}sheet & quarterly & 122 & 53 & EXTRACTED & 03\_\allowbreak{}提取表格/\allowbreak{}WEB-0048/\allowbreak{}WEB-0048-R01-E021/\allowbreak{}quarterly.csv\\
WEB-0048 & WEB-0048-R01-E022 & xlsx\_\allowbreak{}sheet & annual & 32 & 56 & EXTRACTED & 03\_\allowbreak{}提取表格/\allowbreak{}WEB-0048/\allowbreak{}WEB-0048-R01-E022/\allowbreak{}annual.csv\\
WEB-0048 & WEB-0048-R01-E022 & xlsx\_\allowbreak{}sheet & monthly & 362 & 56 & EXTRACTED & 03\_\allowbreak{}提取表格/\allowbreak{}WEB-0048/\allowbreak{}WEB-0048-R01-E022/\allowbreak{}monthly.csv\\
WEB-0048 & WEB-0048-R01-E022 & xlsx\_\allowbreak{}sheet & quarterly & 122 & 56 & EXTRACTED & 03\_\allowbreak{}提取表格/\allowbreak{}WEB-0048/\allowbreak{}WEB-0048-R01-E022/\allowbreak{}quarterly.csv\\
WEB-0048 & WEB-0048-R01-E023 & xlsx\_\allowbreak{}sheet & annual & 32 & 119 & EXTRACTED & 03\_\allowbreak{}提取表格/\allowbreak{}WEB-0048/\allowbreak{}WEB-0048-R01-E023/\allowbreak{}annual.csv\\
WEB-0048 & WEB-0048-R01-E023 & xlsx\_\allowbreak{}sheet & monthly & 362 & 119 & EXTRACTED & 03\_\allowbreak{}提取表格/\allowbreak{}WEB-0048/\allowbreak{}WEB-0048-R01-E023/\allowbreak{}monthly.csv\\
WEB-0048 & WEB-0048-R01-E023 & xlsx\_\allowbreak{}sheet & quarterly & 122 & 119 & EXTRACTED & 03\_\allowbreak{}提取表格/\allowbreak{}WEB-0048/\allowbreak{}WEB-0048-R01-E023/\allowbreak{}quarterly.csv\\
WEB-0048 & WEB-0048-R01-E024 & xlsx\_\allowbreak{}sheet & annual & 32 & 121 & EXTRACTED & 03\_\allowbreak{}提取表格/\allowbreak{}WEB-0048/\allowbreak{}WEB-0048-R01-E024/\allowbreak{}annual.csv\\
WEB-0048 & WEB-0048-R01-E024 & xlsx\_\allowbreak{}sheet & monthly & 362 & 121 & EXTRACTED & 03\_\allowbreak{}提取表格/\allowbreak{}WEB-0048/\allowbreak{}WEB-0048-R01-E024/\allowbreak{}monthly.csv\\
WEB-0048 & WEB-0048-R01-E024 & xlsx\_\allowbreak{}sheet & quarterly & 122 & 121 & EXTRACTED & 03\_\allowbreak{}提取表格/\allowbreak{}WEB-0048/\allowbreak{}WEB-0048-R01-E024/\allowbreak{}quarterly.csv\\
WEB-0048 & WEB-0048-R01-E025 & xlsx\_\allowbreak{}sheet & annual & 32 & 43 & EXTRACTED & 03\_\allowbreak{}提取表格/\allowbreak{}WEB-0048/\allowbreak{}WEB-0048-R01-E025/\allowbreak{}annual.csv\\
WEB-0048 & WEB-0048-R01-E025 & xlsx\_\allowbreak{}sheet & monthly & 362 & 55 & EXTRACTED & 03\_\allowbreak{}提取表格/\allowbreak{}WEB-0048/\allowbreak{}WEB-0048-R01-E025/\allowbreak{}monthly.csv\\
WEB-0048 & WEB-0048-R01-E025 & xlsx\_\allowbreak{}sheet & quarterly & 122 & 43 & EXTRACTED & 03\_\allowbreak{}提取表格/\allowbreak{}WEB-0048/\allowbreak{}WEB-0048-R01-E025/\allowbreak{}quarterly.csv\\
WEB-0048 & WEB-0048-R01-E026 & xlsx\_\allowbreak{}sheet & annual & 32 & 56 & EXTRACTED & 03\_\allowbreak{}提取表格/\allowbreak{}WEB-0048/\allowbreak{}WEB-0048-R01-E026/\allowbreak{}annual.csv\\
WEB-0048 & WEB-0048-R01-E026 & xlsx\_\allowbreak{}sheet & monthly & 362 & 56 & EXTRACTED & 03\_\allowbreak{}提取表格/\allowbreak{}WEB-0048/\allowbreak{}WEB-0048-R01-E026/\allowbreak{}monthly.csv\\
WEB-0048 & WEB-0048-R01-E026 & xlsx\_\allowbreak{}sheet & quarterly & 122 & 56 & EXTRACTED & 03\_\allowbreak{}提取表格/\allowbreak{}WEB-0048/\allowbreak{}WEB-0048-R01-E026/\allowbreak{}quarterly.csv\\
WEB-0048 & WEB-0048-R01-E027 & xlsx\_\allowbreak{}sheet & annual & 32 & 92 & EXTRACTED & 03\_\allowbreak{}提取表格/\allowbreak{}WEB-0048/\allowbreak{}WEB-0048-R01-E027/\allowbreak{}annual.csv\\
WEB-0048 & WEB-0048-R01-E027 & xlsx\_\allowbreak{}sheet & monthly & 362 & 92 & EXTRACTED & 03\_\allowbreak{}提取表格/\allowbreak{}WEB-0048/\allowbreak{}WEB-0048-R01-E027/\allowbreak{}monthly.csv\\
WEB-0048 & WEB-0048-R01-E028 & xlsx\_\allowbreak{}sheet & annual & 32 & 92 & EXTRACTED & 03\_\allowbreak{}提取表格/\allowbreak{}WEB-0048/\allowbreak{}WEB-0048-R01-E028/\allowbreak{}annual.csv\\
WEB-0048 & WEB-0048-R01-E028 & xlsx\_\allowbreak{}sheet & monthly & 362 & 92 & EXTRACTED & 03\_\allowbreak{}提取表格/\allowbreak{}WEB-0048/\allowbreak{}WEB-0048-R01-E028/\allowbreak{}monthly.csv\\
WEB-0048 & WEB-0048-R01-E029 & xlsx\_\allowbreak{}sheet & annual & 32 & 200 & EXTRACTED & 03\_\allowbreak{}提取表格/\allowbreak{}WEB-0048/\allowbreak{}WEB-0048-R01-E029/\allowbreak{}annual.csv\\
WEB-0048 & WEB-0048-R01-E029 & xlsx\_\allowbreak{}sheet & monthly & 362 & 200 & EXTRACTED & 03\_\allowbreak{}提取表格/\allowbreak{}WEB-0048/\allowbreak{}WEB-0048-R01-E029/\allowbreak{}monthly.csv\\
WEB-0048 & WEB-0048-R01-E030 & xlsx\_\allowbreak{}sheet & annual & 32 & 41 & EXTRACTED & 03\_\allowbreak{}提取表格/\allowbreak{}WEB-0048/\allowbreak{}WEB-0048-R01-E030/\allowbreak{}annual.csv\\
WEB-0048 & WEB-0048-R01-E030 & xlsx\_\allowbreak{}sheet & monthly & 362 & 41 & EXTRACTED & 03\_\allowbreak{}提取表格/\allowbreak{}WEB-0048/\allowbreak{}WEB-0048-R01-E030/\allowbreak{}monthly.csv\\
WEB-0048 & WEB-0048-R01-E031 & xlsx\_\allowbreak{}sheet & annual & 32 & 77 & EXTRACTED & 03\_\allowbreak{}提取表格/\allowbreak{}WEB-0048/\allowbreak{}WEB-0048-R01-E031/\allowbreak{}annual.csv\\
WEB-0048 & WEB-0048-R01-E031 & xlsx\_\allowbreak{}sheet & monthly & 362 & 77 & EXTRACTED & 03\_\allowbreak{}提取表格/\allowbreak{}WEB-0048/\allowbreak{}WEB-0048-R01-E031/\allowbreak{}monthly.csv\\
WEB-0048 & WEB-0048-R01-E032 & xlsx\_\allowbreak{}sheet & annual & 32 & 79 & EXTRACTED & 03\_\allowbreak{}提取表格/\allowbreak{}WEB-0048/\allowbreak{}WEB-0048-R01-E032/\allowbreak{}annual.csv\\
WEB-0048 & WEB-0048-R01-E032 & xlsx\_\allowbreak{}sheet & monthly & 362 & 79 & EXTRACTED & 03\_\allowbreak{}提取表格/\allowbreak{}WEB-0048/\allowbreak{}WEB-0048-R01-E032/\allowbreak{}monthly.csv\\
WEB-0048 & WEB-0048-R01-E033 & xlsx\_\allowbreak{}sheet & monthly & 362 & 42 & EXTRACTED & 03\_\allowbreak{}提取表格/\allowbreak{}WEB-0048/\allowbreak{}WEB-0048-R01-E033/\allowbreak{}monthly.csv\\
WEB-0048 & WEB-0048-R01-E034 & xlsx\_\allowbreak{}sheet & annual & 32 & 185 & EXTRACTED & 03\_\allowbreak{}提取表格/\allowbreak{}WEB-0048/\allowbreak{}WEB-0048-R01-E034/\allowbreak{}annual.csv\\
WEB-0048 & WEB-0048-R01-E034 & xlsx\_\allowbreak{}sheet & monthly & 362 & 185 & EXTRACTED & 03\_\allowbreak{}提取表格/\allowbreak{}WEB-0048/\allowbreak{}WEB-0048-R01-E034/\allowbreak{}monthly.csv\\
WEB-0048 & WEB-0048-R01-E035 & xlsx\_\allowbreak{}sheet & annual & 32 & 78 & EXTRACTED & 03\_\allowbreak{}提取表格/\allowbreak{}WEB-0048/\allowbreak{}WEB-0048-R01-E035/\allowbreak{}annual.csv\\
WEB-0048 & WEB-0048-R01-E035 & xlsx\_\allowbreak{}sheet & monthly & 361 & 78 & EXTRACTED & 03\_\allowbreak{}提取表格/\allowbreak{}WEB-0048/\allowbreak{}WEB-0048-R01-E035/\allowbreak{}monthly.csv\\
WEB-0498 & WEB-0498-R01 & xlsx\_\allowbreak{}sheet & AFOSHEET & 0 & 0 & EXTRACTED & 03\_\allowbreak{}提取表格/\allowbreak{}WEB-0498/\allowbreak{}WEB-0498-R01/\allowbreak{}AFOSHEET.csv\\
WEB-0498 & WEB-0498-R01 & xlsx\_\allowbreak{}sheet & Annual Prices (Nominal) & 247 & 87 & EXTRACTED & 03\_\allowbreak{}提取表格/\allowbreak{}WEB-0498/\allowbreak{}WEB-0498-R01/\allowbreak{}Annual\_\allowbreak{}Prices\_\allowbreak{}(Nominal).csv\\
WEB-0498 & WEB-0498-R01 & xlsx\_\allowbreak{}sheet & Annual Indices (Nominal) & 65219 & 81 & EXTRACTED & 03\_\allowbreak{}提取表格/\allowbreak{}WEB-0498/\allowbreak{}WEB-0498-R01/\allowbreak{}Annual\_\allowbreak{}Indices\_\allowbreak{}(Nominal).csv\\
WEB-0498 & WEB-0498-R01 & xlsx\_\allowbreak{}sheet & Annual Prices (Real) & 85 & 75 & EXTRACTED & 03\_\allowbreak{}提取表格/\allowbreak{}WEB-0498/\allowbreak{}WEB-0498-R01/\allowbreak{}Annual\_\allowbreak{}Prices\_\allowbreak{}(Real).csv\\
WEB-0498 & WEB-0498-R01 & xlsx\_\allowbreak{}sheet & Annual Indices (Real) & 328 & 82 & EXTRACTED & 03\_\allowbreak{}提取表格/\allowbreak{}WEB-0498/\allowbreak{}WEB-0498-R01/\allowbreak{}Annual\_\allowbreak{}Indices\_\allowbreak{}(Real).csv\\
WEB-0498 & WEB-0498-R01 & xlsx\_\allowbreak{}sheet & Description & 407 & 87 & EXTRACTED & 03\_\allowbreak{}提取表格/\allowbreak{}WEB-0498/\allowbreak{}WEB-0498-R01/\allowbreak{}Description.csv\\
WEB-0498 & WEB-0498-R01 & xlsx\_\allowbreak{}sheet & Definitions & 10 & 23 & EXTRACTED & 03\_\allowbreak{}提取表格/\allowbreak{}WEB-0498/\allowbreak{}WEB-0498-R01/\allowbreak{}Definitions.csv\\
WEB-0498 & WEB-0498-R01 & xlsx\_\allowbreak{}sheet & Index Weights & 92 & 251 & EXTRACTED & 03\_\allowbreak{}提取表格/\allowbreak{}WEB-0498/\allowbreak{}WEB-0498-R01/\allowbreak{}Index\_\allowbreak{}Weights.csv\\
WEB-0498 & WEB-0498-R02 & xlsx\_\allowbreak{}sheet & AFOSHEET & 0 & 0 & EXTRACTED & 03\_\allowbreak{}提取表格/\allowbreak{}WEB-0498/\allowbreak{}WEB-0498-R02/\allowbreak{}AFOSHEET.csv\\
WEB-0498 & WEB-0498-R02 & xlsx\_\allowbreak{}sheet & Monthly Prices & 805 & 89 & EXTRACTED & 03\_\allowbreak{}提取表格/\allowbreak{}WEB-0498/\allowbreak{}WEB-0498-R02/\allowbreak{}Monthly\_\allowbreak{}Prices.csv\\
WEB-0498 & WEB-0498-R02 & xlsx\_\allowbreak{}sheet & Monthly Indices & 1362 & 81 & EXTRACTED & 03\_\allowbreak{}提取表格/\allowbreak{}WEB-0498/\allowbreak{}WEB-0498-R02/\allowbreak{}Monthly\_\allowbreak{}Indices.csv\\
WEB-0498 & WEB-0498-R02 & xlsx\_\allowbreak{}sheet & Description & 409 & 87 & EXTRACTED & 03\_\allowbreak{}提取表格/\allowbreak{}WEB-0498/\allowbreak{}WEB-0498-R02/\allowbreak{}Description.csv\\
WEB-0498 & WEB-0498-R02 & xlsx\_\allowbreak{}sheet & Index Weights & 92 & 251 & EXTRACTED & 03\_\allowbreak{}提取表格/\allowbreak{}WEB-0498/\allowbreak{}WEB-0498-R02/\allowbreak{}Index\_\allowbreak{}Weights.csv\\
WEB-0197 & WEB-0197-R01 & xml & RSS & 294 &  & INSPECTED & 01\_\allowbreak{}新增原件/\allowbreak{}WEB-0197/\allowbreak{}WEB-0197-R01\_\allowbreak{}whatsnew.xml\\
WEB-0197 & WEB-0197-R02 & xlsx\_\allowbreak{}sheet & mei & 368 & 257 & EXTRACTED & 03\_\allowbreak{}提取表格/\allowbreak{}WEB-0197/\allowbreak{}WEB-0197-R02/\allowbreak{}mei.csv\\
WEB-0197 & WEB-0197-R03 & xlsx\_\allowbreak{}sheet & 業態別·By Sector (1) (2) & 86 & 13 & EXTRACTED & 03\_\allowbreak{}提取表格/\allowbreak{}WEB-0197/\allowbreak{}WEB-0197-R03/\allowbreak{}業態別·By\_\allowbreak{}Sector\_\allowbreak{}(1)\_\allowbreak{}(2).csv\\
WEB-0197 & WEB-0197-R03 & xlsx\_\allowbreak{}sheet & 業態別·By Sector (3) (4) & 78 & 13 & EXTRACTED & 03\_\allowbreak{}提取表格/\allowbreak{}WEB-0197/\allowbreak{}WEB-0197-R03/\allowbreak{}業態別·By\_\allowbreak{}Sector\_\allowbreak{}(3)\_\allowbreak{}(4).csv\\
WEB-0197 & WEB-0197-R03 & xlsx\_\allowbreak{}sheet & 期間別·By Period & 52 & 13 & EXTRACTED & 03\_\allowbreak{}提取表格/\allowbreak{}WEB-0197/\allowbreak{}WEB-0197-R03/\allowbreak{}期間別·By\_\allowbreak{}Period.csv\\
WEB-0200 & WEB-0200-R01 & xlsx\_\allowbreak{}sheet & SDDS & 55 & 52 & EXTRACTED & 03\_\allowbreak{}提取表格/\allowbreak{}WEB-0200/\allowbreak{}WEB-0200-R01/\allowbreak{}SDDS.csv\\
WEB-0200 & WEB-0200-R02 & xlsx\_\allowbreak{}sheet & OFCS & 30 & 92 & EXTRACTED & 03\_\allowbreak{}提取表格/\allowbreak{}WEB-0200/\allowbreak{}WEB-0200-R02/\allowbreak{}OFCS.csv\\
WEB-0200 & WEB-0200-R03 & xlsx\_\allowbreak{}sheet & 1 & 71 & 19 & EXTRACTED & 03\_\allowbreak{}提取表格/\allowbreak{}WEB-0200/\allowbreak{}WEB-0200-R03/\allowbreak{}1.csv\\
WEB-0200 & WEB-0200-R03 & xlsx\_\allowbreak{}sheet & 2 & 71 & 15 & EXTRACTED & 03\_\allowbreak{}提取表格/\allowbreak{}WEB-0200/\allowbreak{}WEB-0200-R03/\allowbreak{}2.csv\\
WEB-0200 & WEB-0200-R03 & xlsx\_\allowbreak{}sheet & 3 & 71 & 17 & EXTRACTED & 03\_\allowbreak{}提取表格/\allowbreak{}WEB-0200/\allowbreak{}WEB-0200-R03/\allowbreak{}3.csv\\
WEB-0200 & WEB-0200-R03 & xlsx\_\allowbreak{}sheet & 4 & 71 & 14 & EXTRACTED & 03\_\allowbreak{}提取表格/\allowbreak{}WEB-0200/\allowbreak{}WEB-0200-R03/\allowbreak{}4.csv\\
WEB-0200 & WEB-0200-R03 & xlsx\_\allowbreak{}sheet & 5 & 71 & 18 & EXTRACTED & 03\_\allowbreak{}提取表格/\allowbreak{}WEB-0200/\allowbreak{}WEB-0200-R03/\allowbreak{}5.csv\\
WEB-0200 & WEB-0200-R03 & xlsx\_\allowbreak{}sheet & 6 & 71 & 15 & EXTRACTED & 03\_\allowbreak{}提取表格/\allowbreak{}WEB-0200/\allowbreak{}WEB-0200-R03/\allowbreak{}6.csv\\
WEB-0200 & WEB-0200-R03 & xlsx\_\allowbreak{}sheet & 7 & 71 & 19 & EXTRACTED & 03\_\allowbreak{}提取表格/\allowbreak{}WEB-0200/\allowbreak{}WEB-0200-R03/\allowbreak{}7.csv\\
WEB-0200 & WEB-0200-R03 & xlsx\_\allowbreak{}sheet & 8 & 71 & 17 & EXTRACTED & 03\_\allowbreak{}提取表格/\allowbreak{}WEB-0200/\allowbreak{}WEB-0200-R03/\allowbreak{}8.csv\\
WEB-0200 & WEB-0200-R03 & xlsx\_\allowbreak{}sheet & 9 & 71 & 13 & EXTRACTED & 03\_\allowbreak{}提取表格/\allowbreak{}WEB-0200/\allowbreak{}WEB-0200-R03/\allowbreak{}9.csv\\
WEB-0200 & WEB-0200-R03 & xlsx\_\allowbreak{}sheet & 10 & 71 & 14 & EXTRACTED & 03\_\allowbreak{}提取表格/\allowbreak{}WEB-0200/\allowbreak{}WEB-0200-R03/\allowbreak{}10.csv\\
WEB-0200 & WEB-0200-R03 & xlsx\_\allowbreak{}sheet & 19 & 71 & 18 & EXTRACTED & 03\_\allowbreak{}提取表格/\allowbreak{}WEB-0200/\allowbreak{}WEB-0200-R03/\allowbreak{}19.csv\\
WEB-0200 & WEB-0200-R03 & xlsx\_\allowbreak{}sheet & 20 & 71 & 18 & EXTRACTED & 03\_\allowbreak{}提取表格/\allowbreak{}WEB-0200/\allowbreak{}WEB-0200-R03/\allowbreak{}20.csv\\
WEB-0200 & WEB-0200-R03 & xlsx\_\allowbreak{}sheet & 21 & 71 & 23 & EXTRACTED & 03\_\allowbreak{}提取表格/\allowbreak{}WEB-0200/\allowbreak{}WEB-0200-R03/\allowbreak{}21.csv\\
WEB-0200 & WEB-0200-R03 & xlsx\_\allowbreak{}sheet & 22 & 71 & 14 & EXTRACTED & 03\_\allowbreak{}提取表格/\allowbreak{}WEB-0200/\allowbreak{}WEB-0200-R03/\allowbreak{}22.csv\\
WEB-0200 & WEB-0200-R03 & xlsx\_\allowbreak{}sheet & 23 & 71 & 17 & EXTRACTED & 03\_\allowbreak{}提取表格/\allowbreak{}WEB-0200/\allowbreak{}WEB-0200-R03/\allowbreak{}23.csv\\
WEB-0200 & WEB-0200-R03 & xlsx\_\allowbreak{}sheet & 24 & 71 & 15 & EXTRACTED & 03\_\allowbreak{}提取表格/\allowbreak{}WEB-0200/\allowbreak{}WEB-0200-R03/\allowbreak{}24.csv\\
WEB-0200 & WEB-0200-R03 & xlsx\_\allowbreak{}sheet & 25 & 71 & 19 & EXTRACTED & 03\_\allowbreak{}提取表格/\allowbreak{}WEB-0200/\allowbreak{}WEB-0200-R03/\allowbreak{}25.csv\\
WEB-0200 & WEB-0200-R03 & xlsx\_\allowbreak{}sheet & 26 & 71 & 17 & EXTRACTED & 03\_\allowbreak{}提取表格/\allowbreak{}WEB-0200/\allowbreak{}WEB-0200-R03/\allowbreak{}26.csv\\
WEB-0200 & WEB-0200-R03 & xlsx\_\allowbreak{}sheet & 27 & 71 & 13 & EXTRACTED & 03\_\allowbreak{}提取表格/\allowbreak{}WEB-0200/\allowbreak{}WEB-0200-R03/\allowbreak{}27.csv\\
WEB-0200 & WEB-0200-R03 & xlsx\_\allowbreak{}sheet & 28 & 71 & 14 & EXTRACTED & 03\_\allowbreak{}提取表格/\allowbreak{}WEB-0200/\allowbreak{}WEB-0200-R03/\allowbreak{}28.csv\\
WEB-0200 & WEB-0200-R03 & xlsx\_\allowbreak{}sheet & 45 & 71 & 19 & EXTRACTED & 03\_\allowbreak{}提取表格/\allowbreak{}WEB-0200/\allowbreak{}WEB-0200-R03/\allowbreak{}45.csv\\
WEB-0200 & WEB-0200-R03 & xlsx\_\allowbreak{}sheet & 46 & 71 & 16 & EXTRACTED & 03\_\allowbreak{}提取表格/\allowbreak{}WEB-0200/\allowbreak{}WEB-0200-R03/\allowbreak{}46.csv\\
WEB-0200 & WEB-0200-R03 & xlsx\_\allowbreak{}sheet & 47 & 71 & 27 & EXTRACTED & 03\_\allowbreak{}提取表格/\allowbreak{}WEB-0200/\allowbreak{}WEB-0200-R03/\allowbreak{}47.csv\\
WEB-0200 & WEB-0200-R03 & xlsx\_\allowbreak{}sheet & 48 & 71 & 13 & EXTRACTED & 03\_\allowbreak{}提取表格/\allowbreak{}WEB-0200/\allowbreak{}WEB-0200-R03/\allowbreak{}48.csv\\
WEB-0200 & WEB-0200-R03 & xlsx\_\allowbreak{}sheet & 49 & 71 & 30 & EXTRACTED & 03\_\allowbreak{}提取表格/\allowbreak{}WEB-0200/\allowbreak{}WEB-0200-R03/\allowbreak{}49.csv\\
WEB-0200 & WEB-0200-R03 & xlsx\_\allowbreak{}sheet & 50 & 71 & 14 & EXTRACTED & 03\_\allowbreak{}提取表格/\allowbreak{}WEB-0200/\allowbreak{}WEB-0200-R03/\allowbreak{}50.csv\\
WEB-0200 & WEB-0200-R03 & xlsx\_\allowbreak{}sheet & 51 & 71 & 19 & EXTRACTED & 03\_\allowbreak{}提取表格/\allowbreak{}WEB-0200/\allowbreak{}WEB-0200-R03/\allowbreak{}51.csv\\
WEB-0200 & WEB-0200-R03 & xlsx\_\allowbreak{}sheet & 52 & 71 & 17 & EXTRACTED & 03\_\allowbreak{}提取表格/\allowbreak{}WEB-0200/\allowbreak{}WEB-0200-R03/\allowbreak{}52.csv\\
WEB-0200 & WEB-0200-R03 & xlsx\_\allowbreak{}sheet & 53 & 71 & 13 & EXTRACTED & 03\_\allowbreak{}提取表格/\allowbreak{}WEB-0200/\allowbreak{}WEB-0200-R03/\allowbreak{}53.csv\\
WEB-0200 & WEB-0200-R03 & xlsx\_\allowbreak{}sheet & 54 & 71 & 14 & EXTRACTED & 03\_\allowbreak{}提取表格/\allowbreak{}WEB-0200/\allowbreak{}WEB-0200-R03/\allowbreak{}54.csv\\
WEB-0253 & WEB-0253-R01 & xml & RSS & 111 &  & INSPECTED & 01\_\allowbreak{}新增原件/\allowbreak{}WEB-0253/\allowbreak{}WEB-0253-R01\_\allowbreak{}feds\_\allowbreak{}notes.xml\\
WEB-0264 & WEB-0264-R01 & zip\_\allowbreak{}archive & 17 entries & 17 &  & EXPANDED & 01\_\allowbreak{}新增原件/\allowbreak{}WEB-0264/\allowbreak{}expanded/\allowbreak{}WEB-0264-R01\\
WEB-0264 & WEB-0264-R01-E004 & xml & MESSAGE:MESSAGEGROUP & 1094475 &  & INSPECTED & 01\_\allowbreak{}新增原件/\allowbreak{}WEB-0264/\allowbreak{}expanded/\allowbreak{}WEB-0264-R01/\allowbreak{}G17\_\allowbreak{}data.xml\\
WEB-0264 & WEB-0264-R01-E017 & xml & MESSAGE:STRUCTURE & 3286 &  & INSPECTED & 01\_\allowbreak{}新增原件/\allowbreak{}WEB-0264/\allowbreak{}expanded/\allowbreak{}WEB-0264-R01/\allowbreak{}G17\_\allowbreak{}struct.xml\\
WEB-0266 & WEB-0266-R01 & zip\_\allowbreak{}archive & 4 entries & 4 &  & EXPANDED & 01\_\allowbreak{}新增原件/\allowbreak{}WEB-0266/\allowbreak{}expanded/\allowbreak{}WEB-0266-R01\\
WEB-0266 & WEB-0266-R01-E002 & xml & MESSAGE:MESSAGEGROUP & 347177 &  & INSPECTED & 01\_\allowbreak{}新增原件/\allowbreak{}WEB-0266/\allowbreak{}expanded/\allowbreak{}WEB-0266-R01/\allowbreak{}H10\_\allowbreak{}data.xml\\
WEB-0266 & WEB-0266-R01-E004 & xml & MESSAGE:STRUCTURE & 1803 &  & INSPECTED & 01\_\allowbreak{}新增原件/\allowbreak{}WEB-0266/\allowbreak{}expanded/\allowbreak{}WEB-0266-R01/\allowbreak{}H10\_\allowbreak{}struct.xml\\
WEB-0471 & WEB-0471-R01 & xlsx\_\allowbreak{}sheet & 附表1 & 16 & 7 & EXTRACTED & 03\_\allowbreak{}提取表格/\allowbreak{}WEB-0471/\allowbreak{}WEB-0471-R01/\allowbreak{}附表1.csv\\
WEB-0471 & WEB-0471-R01 & xlsx\_\allowbreak{}sheet & 附表2 & 13 & 9 & EXTRACTED & 03\_\allowbreak{}提取表格/\allowbreak{}WEB-0471/\allowbreak{}WEB-0471-R01/\allowbreak{}附表2.csv\\
WEB-0471 & WEB-0471-R01 & xlsx\_\allowbreak{}sheet & 附表3 & 53 & 7 & EXTRACTED & 03\_\allowbreak{}提取表格/\allowbreak{}WEB-0471/\allowbreak{}WEB-0471-R01/\allowbreak{}附表3.csv\\
WEB-0471 & WEB-0471-R01 & xlsx\_\allowbreak{}sheet & 12月用1 & 17 & 10 & EXTRACTED & 03\_\allowbreak{}提取表格/\allowbreak{}WEB-0471/\allowbreak{}WEB-0471-R01/\allowbreak{}12月用1.csv\\
WEB-0471 & WEB-0471-R01 & xlsx\_\allowbreak{}sheet & 12月用2 & 13 & 11 & EXTRACTED & 03\_\allowbreak{}提取表格/\allowbreak{}WEB-0471/\allowbreak{}WEB-0471-R01/\allowbreak{}12月用2.csv\\
WEB-0471 & WEB-0471-R01 & xlsx\_\allowbreak{}sheet & 12月用3 & 52 & 17 & EXTRACTED & 03\_\allowbreak{}提取表格/\allowbreak{}WEB-0471/\allowbreak{}WEB-0471-R01/\allowbreak{}12月用3.csv\\
WEB-0472 & WEB-0472-R01 & xlsx\_\allowbreak{}sheet & 附表1 & 16 & 7 & EXTRACTED & 03\_\allowbreak{}提取表格/\allowbreak{}WEB-0472/\allowbreak{}WEB-0472-R01/\allowbreak{}附表1.csv\\
WEB-0472 & WEB-0472-R01 & xlsx\_\allowbreak{}sheet & 附表2 & 13 & 9 & EXTRACTED & 03\_\allowbreak{}提取表格/\allowbreak{}WEB-0472/\allowbreak{}WEB-0472-R01/\allowbreak{}附表2.csv\\
WEB-0472 & WEB-0472-R01 & xlsx\_\allowbreak{}sheet & 附表3 & 53 & 7 & EXTRACTED & 03\_\allowbreak{}提取表格/\allowbreak{}WEB-0472/\allowbreak{}WEB-0472-R01/\allowbreak{}附表3.csv\\
WEB-0472 & WEB-0472-R01 & xlsx\_\allowbreak{}sheet & 12月用1 & 17 & 10 & EXTRACTED & 03\_\allowbreak{}提取表格/\allowbreak{}WEB-0472/\allowbreak{}WEB-0472-R01/\allowbreak{}12月用1.csv\\
WEB-0472 & WEB-0472-R01 & xlsx\_\allowbreak{}sheet & 12月用2 & 13 & 11 & EXTRACTED & 03\_\allowbreak{}提取表格/\allowbreak{}WEB-0472/\allowbreak{}WEB-0472-R01/\allowbreak{}12月用2.csv\\
WEB-0472 & WEB-0472-R01 & xlsx\_\allowbreak{}sheet & 12月用3 & 52 & 17 & EXTRACTED & 03\_\allowbreak{}提取表格/\allowbreak{}WEB-0472/\allowbreak{}WEB-0472-R01/\allowbreak{}12月用3.csv\\
WEB-0042 & WEB-0042-R02 & legacy\_\allowbreak{}xls & workbook &  &  & MANUAL\_\allowbreak{}CONVERSION\_\allowbreak{}REQUIRED & 旧二进制XLS已校验并保存；为避免高危依赖，本轮不自动解析。\\
WEB-0042 & WEB-0042-R03 & zip\_\allowbreak{}archive & 1 entries & 1 &  & EXPANDED & 01\_\allowbreak{}新增原件/\allowbreak{}WEB-0042/\allowbreak{}expanded/\allowbreak{}WEB-0042-R03\\
WEB-0042 & WEB-0042-R03-E001 & xml & ROOT & 87452 &  & INSPECTED & 01\_\allowbreak{}新增原件/\allowbreak{}WEB-0042/\allowbreak{}expanded/\allowbreak{}WEB-0042-R03/\allowbreak{}API\_\allowbreak{}FP.CPI.TOTL.ZG\_\allowbreak{}DS2\_\allowbreak{}en\_\allowbreak{}xml\_\allowbreak{}v2\_\allowbreak{}2812.xml\\
WEB-0043 & WEB-0043-R02 & legacy\_\allowbreak{}xls & workbook &  &  & MANUAL\_\allowbreak{}CONVERSION\_\allowbreak{}REQUIRED & 旧二进制XLS已校验并保存；为避免高危依赖，本轮不自动解析。\\
WEB-0043 & WEB-0043-R03 & zip\_\allowbreak{}archive & 1 entries & 1 &  & EXPANDED & 01\_\allowbreak{}新增原件/\allowbreak{}WEB-0043/\allowbreak{}expanded/\allowbreak{}WEB-0043-R03\\
WEB-0043 & WEB-0043-R03-E001 & xml & ROOT & 87452 &  & INSPECTED & 01\_\allowbreak{}新增原件/\allowbreak{}WEB-0043/\allowbreak{}expanded/\allowbreak{}WEB-0043-R03/\allowbreak{}API\_\allowbreak{}PA.NUS.FCRF\_\allowbreak{}DS2\_\allowbreak{}en\_\allowbreak{}xml\_\allowbreak{}v2\_\allowbreak{}753.xml\\
WEB-0264 & WEB-0264-R02 & text & WEB-0264-R02\_\allowbreak{}g17.txt & 1462 & 1 & INSPECTED & 01\_\allowbreak{}新增原件/\allowbreak{}WEB-0264/\allowbreak{}WEB-0264-R02\_\allowbreak{}g17.txt\\
WEB-0264 & WEB-0264-R03 & text & WEB-0264-R03\_\allowbreak{}g17.txt & 1505 & 1 & INSPECTED & 01\_\allowbreak{}新增原件/\allowbreak{}WEB-0264/\allowbreak{}WEB-0264-R03\_\allowbreak{}g17.txt\\
\bottomrule\end{longtable}\endgroup\end{landscape}

